% ============================================================
%  MODUL PEMBELAJARAN MATEMATIKA SMA
%  Kurikulum Merdeka — Fase E (Kelas 10)
%  Kompilasi: pdfLaTeX / XeLaTeX (disarankan Overleaf)
% ============================================================
\documentclass[11pt,a4paper]{report}

% ---------- PAKET ----------
\usepackage[utf8]{inputenc}
\usepackage[T1]{fontenc}
\usepackage[bahasa]{babel}
\usepackage{amsmath,amssymb,amsthm,mathtools}
\usepackage{geometry}
\usepackage{xcolor}
\usepackage{graphicx}
\usepackage{enumitem}
\usepackage{tikz}
\usetikzlibrary{calc,arrows.meta,angles,quotes,positioning}
\usepackage{pgfplots}
\pgfplotsset{compat=1.18}
\usepackage[most]{tcolorbox}
\usepackage{fancyhdr}
\usepackage{titlesec}
\usepackage{array,booktabs}
\usepackage[colorlinks=true,linkcolor=biru,urlcolor=biru]{hyperref}

\geometry{a4paper,margin=2.5cm,headheight=15pt}

% ---------- WARNA ----------
\definecolor{biru}{HTML}{1F4E79}
\definecolor{birumuda}{HTML}{D6E4F0}
\definecolor{hijau}{HTML}{2E7D32}
\definecolor{hijaumuda}{HTML}{E3F2E1}
\definecolor{oranye}{HTML}{C75B12}
\definecolor{oranyemuda}{HTML}{FBE8DA}
\definecolor{abu}{HTML}{F2F2F2}
\definecolor{ungu}{HTML}{6A1B9A}
\definecolor{ungumuda}{HTML}{EDE2F3}

% ---------- HEADER / FOOTER ----------
\pagestyle{fancy}
\fancyhf{}
\renewcommand{\headrulewidth}{0.8pt}
\renewcommand{\footrulewidth}{0pt}
\fancyhead[L]{\small\textcolor{biru}{\leftmark}}
\fancyhead[R]{\small\textcolor{biru}{Matematika SMA — Fase E}}
\fancyfoot[C]{\thepage}

% ---------- GAYA JUDUL BAB ----------
\titleformat{\chapter}[display]
  {\normalfont\huge\bfseries\color{biru}}
  {\filright\large BAB \thechapter}{8pt}{\Huge}
\titlespacing*{\chapter}{0pt}{0pt}{20pt}

% ---------- LINGKUNGAN TEOREMA ----------
\newtheoremstyle{gaya}{6pt}{6pt}{}{}{\bfseries\color{biru}}{.}{.5em}{}
\theoremstyle{gaya}
\newtheorem{definisi}{Definisi}[chapter]
\newtheorem{sifat}{Sifat}[chapter]
\newtheorem{teorema}{Teorema}[chapter]

% ---------- KOTAK KHUSUS ----------
% Tujuan Pembelajaran
\newtcolorbox{tujuan}{
  enhanced, colback=birumuda, colframe=biru, boxrule=0.6pt,
  arc=3pt, left=10pt, right=10pt, top=8pt, bottom=8pt,
  title=\textbf{Tujuan Pembelajaran}, fonttitle=\bfseries\color{white},
  coltitle=white, attach boxed title to top left={xshift=10pt,yshift=-10pt},
  boxed title style={colback=biru,arc=2pt}}

% Konsep / Catatan penting
\newtcolorbox{konsep}[1]{
  enhanced, breakable, colback=hijaumuda, colframe=hijau, boxrule=0.6pt,
  arc=3pt, left=10pt, right=10pt, top=6pt, bottom=6pt,
  title=#1, fonttitle=\bfseries\color{white}, coltitle=white,
  attach boxed title to top left={xshift=10pt,yshift=-10pt},
  boxed title style={colback=hijau,arc=2pt}}

% Contoh soal
\newtcolorbox{contoh}[1]{
  enhanced, breakable, colback=oranyemuda, colframe=oranye, boxrule=0.6pt,
  arc=3pt, left=10pt, right=10pt, top=6pt, bottom=6pt,
  title=\textbf{Contoh: #1}, fonttitle=\bfseries\color{white}, coltitle=white,
  attach boxed title to top left={xshift=10pt,yshift=-10pt},
  boxed title style={colback=oranye,arc=2pt}}

% Latihan
\newtcolorbox{latihan}{
  enhanced, breakable, colback=ungumuda, colframe=ungu, boxrule=0.6pt,
  arc=3pt, left=10pt, right=10pt, top=8pt, bottom=8pt,
  title=\textbf{Latihan Soal}, fonttitle=\bfseries\color{white}, coltitle=white,
  attach boxed title to top left={xshift=10pt,yshift=-10pt},
  boxed title style={colback=ungu,arc=2pt}}

\newcommand{\jawab}{\par\smallskip\noindent\textit{\textcolor{oranye}{Penyelesaian:}}\ }

% ---------- WARNA TAMBAHAN ----------
\definecolor{merah}{HTML}{C62828}
\definecolor{merahmuda}{HTML}{FCE4E4}
\definecolor{tosca}{HTML}{00838F}
\definecolor{toscamuda}{HTML}{E0F2F1}
\definecolor{emas}{HTML}{8D6E00}
\definecolor{emasmuda}{HTML}{FFF6D6}

\usepackage{multicol}

% ---------- KOTAK PEDAGOGIS TAMBAHAN ----------
% Pertanyaan Pemantik
\newtcolorbox{pemantik}{
  enhanced, breakable, colback=emasmuda, colframe=emas, boxrule=0.6pt,
  arc=3pt, left=10pt, right=10pt, top=8pt, bottom=8pt,
  title=\textbf{Pertanyaan Pemantik}, fonttitle=\bfseries\color{white}, coltitle=white,
  attach boxed title to top left={xshift=10pt,yshift=-10pt},
  boxed title style={colback=emas,arc=2pt}}

% Studi Kasus Kontekstual
\newtcolorbox{studikasus}[1]{
  enhanced, breakable, colback=abu, colframe=biru, boxrule=0.6pt,
  arc=3pt, left=10pt, right=10pt, top=8pt, bottom=8pt,
  title=\textbf{Studi Kasus: #1}, fonttitle=\bfseries\color{white}, coltitle=white,
  attach boxed title to top left={xshift=10pt,yshift=-10pt},
  boxed title style={colback=biru,arc=2pt}}

% Kekeliruan yang Sering Terjadi
\newtcolorbox{kesalahan}{
  enhanced, breakable, colback=merahmuda, colframe=merah, boxrule=0.6pt,
  arc=3pt, left=10pt, right=10pt, top=8pt, bottom=8pt,
  title=\textbf{Kekeliruan yang Sering Terjadi}, fonttitle=\bfseries\color{white}, coltitle=white,
  attach boxed title to top left={xshift=10pt,yshift=-10pt},
  boxed title style={colback=merah,arc=2pt}}

% Matematika dalam Dunia Nyata
\newtcolorbox{dunianyata}{
  enhanced, breakable, colback=toscamuda, colframe=tosca, boxrule=0.6pt,
  arc=3pt, left=10pt, right=10pt, top=8pt, bottom=8pt,
  title=\textbf{Matematika dalam Dunia Nyata}, fonttitle=\bfseries\color{white}, coltitle=white,
  attach boxed title to top left={xshift=10pt,yshift=-10pt},
  boxed title style={colback=tosca,arc=2pt}}

% Refleksi Diri
\newtcolorbox{refleksi}{
  enhanced, breakable, colback=toscamuda, colframe=tosca, boxrule=0.6pt,
  arc=3pt, left=10pt, right=10pt, top=8pt, bottom=8pt,
  title=\textbf{Refleksi Diri}, fonttitle=\bfseries\color{white}, coltitle=white,
  attach boxed title to top left={xshift=10pt,yshift=-10pt},
  boxed title style={colback=tosca,arc=2pt}}

% Challenge Corner
\newtcolorbox{challenge}{
  enhanced, breakable, colback=ungumuda, colframe=ungu, boxrule=1pt,
  arc=3pt, left=10pt, right=10pt, top=8pt, bottom=8pt,
  title=\textbf{\ding{72}\ Challenge Corner}, fonttitle=\bfseries\color{white}, coltitle=white,
  attach boxed title to top left={xshift=10pt,yshift=-10pt},
  boxed title style={colback=ungu,arc=2pt}}

% Catatan Penting (sidenote ringkas)
\newtcolorbox{catatan}{
  enhanced, breakable, colback=white, colframe=biru, boxrule=0.6pt,
  arc=2pt, left=8pt, right=8pt, top=5pt, bottom=5pt,
  borderline west={3pt}{0pt}{biru}}

\usepackage{pifont} % untuk \ding (checkbox & bintang)

% Penanda bagian soal akhir bab
\newcommand{\bagiansoal}[2]{%
  \par\medskip\noindent
  {\large\bfseries\color{biru}Bagian #1}\par\smallskip
  \addcontentsline{toc}{subsection}{Bagian #1}}

% Checklist refleksi
\newcommand{\cek}{\ding{113}\ } % kotak kosong

% ---------- FORMAT KUNCI JAWABAN ----------
% Nomor soal pada kunci jawaban
\newcommand{\knomor}[1]{\par\medskip\noindent{\bfseries\color{ungu}#1}\par\nopagebreak}
% Tiga bagian wajib: Soal -- Analisis -- Pembahasan
\newcommand{\ksoal}{\noindent\textbf{\color{biru}Soal.}\space}
\newcommand{\kanalisis}{\par\noindent\textbf{\color{oranye}Analisis.}\space}
\newcommand{\kbahas}{\par\noindent\textbf{\color{hijau}Pembahasan.}\space}
% Sub-judul bagian (A/B/C/D/Challenge) pada kunci
\newcommand{\kbagian}[1]{\par\medskip\noindent{\large\bfseries\color{biru}#1}\par\smallskip}

% ============================================================
\begin{document}

% ---------- HALAMAN JUDUL ----------
\begin{titlepage}
\centering
\vspace*{2cm}
{\color{biru}\rule{\textwidth}{2pt}}\\[0.6cm]
{\Huge\bfseries\color{biru} MODUL PEMBELAJARAN\\[0.3cm] MATEMATIKA SMA}\\[0.4cm]
{\Large Kurikulum Merdeka}\\[0.2cm]
{\color{biru}\rule{\textwidth}{2pt}}\\[1.2cm]

{\LARGE\bfseries FASE E}\\[0.3cm]
{\Large Kelas 10 — Matematika Umum}\\[2cm]

\vfill
{\large Disusun untuk pembelajaran mandiri dan tatap muka}\\[0.3cm]
{\large Tahun Pelajaran 2026/2027}
\end{titlepage}

% ---------- DAFTAR ISI ----------
\tableofcontents
\thispagestyle{empty}
\clearpage

% ============================================================
\part*{\centering FASE E — KELAS 10}
\addcontentsline{toc}{part}{FASE E — KELAS 10}

% ============================================================
\chapter{Eksponen dan Logaritma}
% ============================================================
\begin{tujuan}
Setelah mempelajari bab ini, peserta didik mampu:
\begin{itemize}[leftmargin=*,topsep=2pt]
  \item menjelaskan pengertian bilangan berpangkat (eksponen) dan memaknainya sebagai proses pelipatgandaan berulang;
  \item menerapkan sifat-sifat eksponen dalam menyederhanakan bentuk aljabar dan menyelesaikan persamaan;
  \item menjelaskan logaritma sebagai invers eksponen serta hubungan keduanya secara grafik;
  \item menerapkan sifat-sifat logaritma untuk menyelesaikan masalah kontekstual (pertumbuhan, peluruhan, skala);
  \item memodelkan fenomena pertumbuhan dan peluruhan eksponensial di dunia nyata.
\end{itemize}
\textbf{Capaian Pembelajaran (Fase E).} Peserta didik dapat menggeneralisasi sifat-sifat bilangan berpangkat (termasuk pangkat pecahan dan negatif), serta menggunakan eksponen dan logaritma untuk memodelkan dan menyelesaikan masalah.
\end{tujuan}

\begin{pemantik}
\begin{itemize}[leftmargin=*,topsep=2pt]
  \item Selembar kertas tebalnya $0{,}1$ mm. Jika dilipat 30 kali, mungkinkah tebalnya melampaui tinggi gunung? Mengapa intuisi kita sering keliru menebak angka seperti ini?
  \item Mengapa "menambah nol" pada sebuah bilangan terasa jauh lebih besar daripada "menambah satu"? Apa kaitannya dengan eksponen?
  \item Jika perkalian adalah penjumlahan yang diulang, maka eksponen adalah \dots\ yang diulang. Logaritma "membatalkan" proses apa?
\end{itemize}
\end{pemantik}

\begin{studikasus}{Pertumbuhan Subscriber Kanal YouTube}
Seorang kreator memulai kanal dengan $1000$ subscriber. Ia mengamati bahwa setiap bulan jumlah subscriber-nya menjadi $1{,}5$ kali lipat bulan sebelumnya.

\smallskip
\begin{center}
\begin{tabular}{c|c|c}
\toprule
Bulan ke-$n$ & Cara hitung & Subscriber \\
\midrule
0 & $1000$ & $1000$\\
1 & $1000 \times 1{,}5$ & $1500$\\
2 & $1000 \times 1{,}5 \times 1{,}5$ & $2250$\\
3 & $1000 \times 1{,}5^3$ & $3375$\\
$n$ & $1000 \times 1{,}5^{\,n}$ & ? \\
\bottomrule
\end{tabular}
\end{center}

\smallskip
\textbf{ Amati dan duga.} (a) Apa pola yang muncul di kolom "cara hitung"? (b) Bagaimana menuliskan jumlah subscriber pada bulan ke-$n$ secara ringkas? (c) Pada bulan ke berapa subscriber menembus $100.000$? Pertanyaan (c) tidak dapat dijawab dengan perkalian biasa --- kita memerlukan \emph{logaritma}. Inilah dua tokoh utama bab ini: \textbf{eksponen} (memproyeksikan pertumbuhan) dan \textbf{logaritma} (membongkar berapa lama pertumbuhan itu berlangsung).
\end{studikasus}

\section*{Peta Konsep Bab}
\addcontentsline{toc}{section}{Peta Konsep Bab}
\begin{center}
\begin{tikzpicture}[
  node distance=8mm,
  konsepbox/.style={draw=biru,fill=birumuda,rounded corners,align=center,inner sep=5pt,font=\small\bfseries},
  subbox/.style={draw=hijau,fill=hijaumuda,rounded corners,align=center,inner sep=4pt,font=\footnotesize},
  garis/.style={-Stealth,biru,thick}]
  \node[konsepbox] (akar) {EKSPONEN \&\ LOGARITMA};
  \node[subbox,below left=12mm and 6mm of akar] (eks) {Bilangan\\Berpangkat};
  \node[subbox,below=12mm of akar] (per) {Persamaan \&\\Pertidaksamaan};
  \node[subbox,below right=12mm and 6mm of akar] (log) {Logaritma\\(invers)};
  \node[subbox,below=10mm of eks] (sifat) {Sifat-sifat\\eksponen};
  \node[subbox,below=10mm of log] (siflog) {Sifat-sifat\\logaritma};
  \node[subbox,below=10mm of per] (apl) {Pertumbuhan \&\\peluruhan};
  \draw[garis] (akar)--(eks); \draw[garis] (akar)--(per); \draw[garis] (akar)--(log);
  \draw[garis] (eks)--(sifat); \draw[garis] (log)--(siflog); \draw[garis] (per)--(apl);
  \draw[garis,dashed,gray] (eks)--(log) node[midway,above,font=\tiny\itshape,black]{invers};
\end{tikzpicture}\\
\small \textbf{Peta konsep.} Eksponen dan logaritma adalah dua sisi mata uang yang sama; keduanya bertemu pada penerapan pertumbuhan dan peluruhan.
\end{center}

\section{Bilangan Berpangkat (Eksponen)}

\subsection*{Membangun konsep dari pola}
Perhatikan kembali studi kasus. Menulis $1000\times1{,}5\times1{,}5\times1{,}5$ itu melelahkan. Notasi pangkat lahir sebagai \emph{singkatan perkalian berulang}, persis seperti perkalian adalah singkatan penjumlahan berulang. Inilah inti "mengapa"-nya: eksponen bukan sekadar simbol, melainkan bahasa untuk menggambarkan sesuatu yang \emph{berlipat ganda}.

\begin{definisi}
Untuk bilangan real $a$ dan bilangan bulat positif $n$, \emph{bilangan berpangkat} didefinisikan sebagai
\[
a^n = \underbrace{a \times a \times \cdots \times a}_{n \text{ faktor}},
\]
dengan $a$ disebut \emph{basis} dan $n$ disebut \emph{pangkat} atau \emph{eksponen}.
\end{definisi}

\begin{konsep}{Sifat-Sifat Eksponen}
Untuk $a,b \neq 0$ dan $m,n$ bilangan rasional:
\[
\begin{array}{lll}
\text{(1) } a^m \cdot a^n = a^{m+n} & \quad &
\text{(2) } \dfrac{a^m}{a^n} = a^{m-n}\\[8pt]
\text{(3) } (a^m)^n = a^{mn} & &
\text{(4) } (ab)^n = a^n b^n\\[8pt]
\text{(5) } \left(\dfrac{a}{b}\right)^n = \dfrac{a^n}{b^n} & &
\text{(6) } a^0 = 1\\[8pt]
\text{(7) } a^{-n} = \dfrac{1}{a^n} & &
\text{(8) } a^{m/n} = \sqrt[n]{a^m}
\end{array}
\]
\end{konsep}

\textbf{Mengapa $a^0=1$ dan $a^{-n}=\tfrac{1}{a^n}$?} Sifat-sifat ini bukan aturan sembarangan, melainkan \emph{konsekuensi logis} agar sifat (1)--(2) tetap berlaku. Dari $\dfrac{a^m}{a^m}=a^{m-m}=a^0$ dan sekaligus $\dfrac{a^m}{a^m}=1$, maka haruslah $a^0=1$. Demikian pula $a^{-n}=a^{0-n}=\dfrac{a^0}{a^n}=\dfrac{1}{a^n}$. Memahami penurunan ini jauh lebih kuat daripada sekadar menghafal.

\begin{center}
\begin{tikzpicture}[
  katas/.style={draw=biru,fill=birumuda,rounded corners=3pt,
    minimum width=3cm,minimum height=0.62cm,align=center,font=\small},
  katbawah/.style={draw=oranye,fill=oranyemuda,rounded corners=3pt,
    minimum width=3cm,minimum height=0.62cm,align=center,font=\small},
  kat0/.style={draw=hijau,fill=hijaumuda,rounded corners=3pt,
    minimum width=3cm,minimum height=0.62cm,align=center,font=\small\bfseries},
  pd/.style={-Stealth,gray,semithick},
  pu/.style={-Stealth,oranye,semithick}
]
  \node[katas]    (an) {$2^3 = 8$};
  \node[katas]    (bn) [below=0.72cm of an] {$2^2 = 4$};
  \node[katas]    (cn) [below=0.72cm of bn] {$2^1 = 2$};
  \node[kat0]     (dn) [below=0.72cm of cn] {$2^0 = 1$};
  \node[katbawah] (en) [below=0.72cm of dn] {$2^{-1} = \tfrac{1}{2}$};
  \node[katbawah] (fn) [below=0.72cm of en] {$2^{-2} = \tfrac{1}{4}$};
  \node[font=\footnotesize\itshape,hijau] at ($(dn.east)+(1.7cm,0)$) {$\leftarrow$ patokan};
  % Panah turun (÷2)
  \draw[pd] ($(an.south)+(0.42cm,0)$)--($(bn.north)+(0.42cm,0)$) node[midway,right,font=\scriptsize]{$\div 2$};
  \draw[pd] ($(bn.south)+(0.42cm,0)$)--($(cn.north)+(0.42cm,0)$) node[midway,right,font=\scriptsize]{$\div 2$};
  \draw[pd] ($(cn.south)+(0.42cm,0)$)--($(dn.north)+(0.42cm,0)$) node[midway,right,font=\scriptsize]{$\div 2$};
  \draw[pd] ($(dn.south)+(0.42cm,0)$)--($(en.north)+(0.42cm,0)$) node[midway,right,font=\scriptsize]{$\div 2$};
  \draw[pd] ($(en.south)+(0.42cm,0)$)--($(fn.north)+(0.42cm,0)$) node[midway,right,font=\scriptsize]{$\div 2$};
  % Panah naik (×2)
  \draw[pu] ($(bn.north)+(-0.42cm,0)$)--($(an.south)+(-0.42cm,0)$) node[midway,left,font=\scriptsize,oranye]{$\times 2$};
  \draw[pu] ($(cn.north)+(-0.42cm,0)$)--($(bn.south)+(-0.42cm,0)$) node[midway,left,font=\scriptsize,oranye]{$\times 2$};
  \draw[pu] ($(dn.north)+(-0.42cm,0)$)--($(cn.south)+(-0.42cm,0)$) node[midway,left,font=\scriptsize,oranye]{$\times 2$};
  \draw[pu] ($(en.north)+(-0.42cm,0)$)--($(dn.south)+(-0.42cm,0)$) node[midway,left,font=\scriptsize,oranye]{$\times 2$};
  \draw[pu] ($(fn.north)+(-0.42cm,0)$)--($(en.south)+(-0.42cm,0)$) node[midway,left,font=\scriptsize,oranye]{$\times 2$};
\end{tikzpicture}\\[5pt]
\small\textbf{Rantai perkalian berulang (basis 2).} Naik satu tangga $= \times2$; turun satu tangga $= \div2$. Kotak hijau ($2^0=1$) adalah patokan. \textbf{Pangkat negatif bukan bilangan negatif} --- ia lahir dari melanjutkan pola $\div2$ melewati~$1$: satu langkah di bawah $1$ adalah $\tfrac12$, dua langkah adalah $\tfrac14$.
\end{center}

\begin{contoh}{Menyederhanakan Bentuk Eksponen}
Sederhanakan $\dfrac{(2a^3 b^{-2})^2}{4 a^{-1} b^3}$.

\jawab
\begin{align*}
\frac{(2a^3 b^{-2})^2}{4a^{-1}b^3}
&= \frac{4 a^6 b^{-4}}{4 a^{-1} b^3}
= a^{6-(-1)} \, b^{-4-3}
= a^{7} b^{-7} = \frac{a^7}{b^7}.
\end{align*}
\end{contoh}

\begin{contoh}{Pangkat Pecahan sebagai Akar}
Hitung $27^{2/3}$ dan jelaskan maknanya.

\jawab $27^{2/3}=\sqrt[3]{27^2}=\sqrt[3]{729}=9$, atau lebih cepat $\big(\sqrt[3]{27}\big)^2=3^2=9$. Pangkat pecahan $m/n$ berarti "akar ke-$n$ lalu dipangkatkan $m$".
\end{contoh}

\begin{contoh}{Menyelesaikan Studi Kasus}
Berapa subscriber kreator pada bulan ke-6?

\jawab $U_6=1000\times 1{,}5^6 = 1000 \times 11{,}39\ldots \approx 11.391$ subscriber. Pertumbuhan eksponensial membuat angka meledak dengan cepat.
\end{contoh}

\begin{center}
\begin{tikzpicture}
\begin{axis}[axis lines=middle,width=8.5cm,height=5.5cm,
  xlabel=$x$,ylabel=$y$,xmin=-3,xmax=3,ymin=-0.5,ymax=6,
  legend pos=north west,legend style={font=\small}]
\addplot[very thick,biru,domain=-3:2.5,samples=60]{2^x};
\addlegendentry{$y=2^x$ (naik)}
\addplot[very thick,oranye,domain=-3:2.5,samples=60]{0.5^x};
\addlegendentry{$y=(\tfrac12)^x$ (turun)}
\end{axis}
\end{tikzpicture}\\
\small \textbf{Grafik fungsi eksponen.} Jika basis $a>1$ grafik \emph{naik} (pertumbuhan);
jika $0<a<1$ grafik \emph{turun} (peluruhan). Keduanya selalu melalui titik $(0,1)$.
\end{center}

\begin{latihan}
\begin{enumerate}[leftmargin=*]
  \item Sederhanakan $\dfrac{a^4 b^{-3}}{a^{-2} b}$.
  \item Hitung $16^{3/4}$ dan $8^{-2/3}$.
  \item Tulis $0{,}000\,000\,7$ dalam notasi ilmiah.
  \item Sederhanakan $\dfrac{(3x^2y)^3}{9x^4y^{-1}}$.
  \item Jika $2^x=5$, tentukan nilai $2^{x+3}$.
  \item Bandingkan $2^{10}$ dan $10^3$; mana yang lebih besar?
  \item Bola jatuh, tiap pantulan $0{,}8$ kali tinggi sebelumnya. Tinggi pantulan ke-5 jika awalnya $2$ m?
  \item Sederhanakan $\sqrt{a}\cdot\sqrt[3]{a}$ dalam bentuk pangkat tunggal.
  \item Tentukan nilai $\left(\tfrac{1}{2}\right)^{-3}$.
  \item Tabungan menjadi $1{,}05$ kali tiap tahun. Tulis saldo setelah $t$ tahun jika awalnya Rp1.000.000.
\end{enumerate}
\end{latihan}

\section{Persamaan dan Pertidaksamaan Eksponen}
\begin{konsep}{Bentuk Dasar}
Jika $a^{f(x)} = a^{g(x)}$ dengan $a>0,\ a\neq 1$, maka $f(x)=g(x)$.\\[3pt]
Untuk pertidaksamaan: jika $a>1$ tanda dipertahankan ($a^{f(x)}>a^{g(x)}\Rightarrow f(x)>g(x)$); jika $0<a<1$ tanda \emph{dibalik}.
\end{konsep}

\begin{contoh}{Persamaan Eksponen}
Tentukan $x$ dari $2^{3x-1} = 32$.

\jawab Karena $32 = 2^5$, maka $3x-1 = 5 \Rightarrow x = 2$.
\end{contoh}

\begin{contoh}{Persamaan dengan Pemisalan}
Selesaikan $2^{2x}-5\cdot 2^x + 4 = 0$.

\jawab Misal $p=2^x$. Maka $p^2-5p+4=0\Rightarrow(p-1)(p-4)=0$, sehingga $p=1$ atau $p=4$. Jadi $2^x=1\Rightarrow x=0$ atau $2^x=4\Rightarrow x=2$.
\end{contoh}

\begin{contoh}{Pertidaksamaan Eksponen}
Tentukan himpunan penyelesaian $\left(\tfrac13\right)^{x} > \tfrac{1}{9}$.

\jawab $\left(\tfrac13\right)^x>\left(\tfrac13\right)^2$. Karena basis $0<\tfrac13<1$, tanda dibalik: $x<2$.
\end{contoh}

\begin{latihan}
\begin{enumerate}[leftmargin=*]
  \item Selesaikan $3^{2x+1}=81$.
  \item Selesaikan $5^{x^2-3x}=5^{-2}$.
  \item Selesaikan $4^x = 8^{x-1}$.
  \item Selesaikan $9^x-4\cdot3^x+3=0$.
  \item Tentukan HP dari $2^{x} \le 16$.
  \item Tentukan HP dari $\left(\tfrac12\right)^{2x-1} \ge 8$.
  \item Selesaikan $7^{x+2}=7^{3x-4}$.
  \item Jika $2^{x}+2^{-x}=\tfrac52$, tentukan $2^x$.
  \item Selesaikan $25^x = 5\cdot 5^{x}$.
  \item Berapa lama (bulan) subscriber pada studi kasus menembus $100.000$? (Sisakan dalam bentuk persamaan eksponen.)
\end{enumerate}
\end{latihan}

\section{Logaritma}
\begin{definisi}
\emph{Logaritma} adalah invers dari eksponen. Untuk $a>0,\ a\neq 1$ dan $x>0$:
\[
{}^{a}\!\log x = y \iff a^y = x.
\]
\end{definisi}

\textbf{Interpretasi makna.} ${}^{a}\!\log x$ menjawab pertanyaan: \emph{"basis $a$ harus dipangkatkan berapa agar menghasilkan $x$?"} Inilah alat yang kita cari pada studi kasus untuk menentukan "berapa lama". Logaritma mengubah pertanyaan tentang perkalian/pangkat menjadi pertanyaan tentang penjumlahan --- itulah mengapa dahulu logaritma merevolusi perhitungan astronomi dan navigasi.

\begin{center}
\begin{tikzpicture}[
  kotak/.style={draw=biru,fill=birumuda,rounded corners=4pt,
    minimum width=2.8cm,minimum height=0.85cm,align=center,font=\small\bfseries}
]
  \node[kotak] (exp) {$2^3 = 8$};
  \node[kotak,right=2.4cm of exp] (log) {${}^{2}\!\log 8 = 3$};
  \node[above=4pt,font=\footnotesize,biru] at (exp.north) {baca: "$2$ pangkat $3$"};
  \node[above=4pt,font=\footnotesize,biru] at (log.north) {baca: "log basis $2$ dari $8$"};
  \draw[-Stealth,oranye,very thick]
    ($(exp.east)+(0,0.18cm)$)--($(log.west)+(0,0.18cm)$)
    node[midway,above,font=\scriptsize\bfseries,oranye]{$\Rightarrow$\ tulis sebagai};
  \draw[-Stealth,hijau,very thick]
    ($(log.west)+(0,-0.18cm)$)--($(exp.east)+(0,-0.18cm)$)
    node[midway,below,font=\scriptsize\bfseries,hijau]{$\Leftarrow$\ invers};
  \node[below=8pt,font=\footnotesize,biru,align=center] at ($(exp)!0.5!(log)$) {};
\end{tikzpicture}
\end{center}

\begin{center}
\begin{tikzpicture}
\begin{axis}[axis lines=middle,width=8cm,height=6cm,
  xlabel=$x$,ylabel=$y$,xmin=-3,xmax=5,ymin=-3,ymax=5,
  legend pos=south east,legend style={font=\small}]
\addplot[very thick,biru,domain=-2.5:2.2,samples=60]{2^x};
\addlegendentry{$y=2^x$}
\addplot[very thick,hijau,domain=0.06:4.8,samples=60]{ln(x)/ln(2)};
\addlegendentry{$y={}^{2}\!\log x$}
\addplot[dashed,gray,domain=-2.5:4.8,samples=2]{x};
\addlegendentry{$y=x$}
\end{axis}
\end{tikzpicture}\\
\small \textbf{Logaritma adalah invers eksponen.} Grafik $y={}^{2}\!\log x$ merupakan
pencerminan grafik $y=2^x$ terhadap garis $y=x$.
\end{center}

\begin{konsep}{Sifat-Sifat Logaritma}
\[
\begin{array}{ll}
\text{(1) } {}^{a}\!\log(xy) = {}^{a}\!\log x + {}^{a}\!\log y &
\text{(2) } {}^{a}\!\log\dfrac{x}{y} = {}^{a}\!\log x - {}^{a}\!\log y\\[8pt]
\text{(3) } {}^{a}\!\log x^n = n\,{}^{a}\!\log x &
\text{(4) } {}^{a}\!\log a = 1,\quad {}^{a}\!\log 1 = 0\\[8pt]
\text{(5) } {}^{a}\!\log x = \dfrac{{}^{c}\!\log x}{{}^{c}\!\log a} &
\text{(6) } a^{{}^{a}\!\log x} = x
\end{array}
\]
\end{konsep}

\textbf{Mengapa sifat (1) berlaku?} Misal ${}^{a}\!\log x=m$ dan ${}^{a}\!\log y=n$, maka $x=a^m$, $y=a^n$. Akibatnya $xy=a^{m+n}$, sehingga ${}^{a}\!\log(xy)=m+n={}^{a}\!\log x+{}^{a}\!\log y$. Sifat logaritma hanyalah "terjemahan" sifat eksponen ke bahasa invers.

\begin{contoh}{Menghitung Nilai Logaritma}
Hitung ${}^{2}\!\log 8 + {}^{3}\!\log 27 - {}^{5}\!\log 1$.

\jawab
${}^{2}\!\log 2^3 + {}^{3}\!\log 3^3 - 0 = 3 + 3 - 0 = 6.$
\end{contoh}

\begin{contoh}{Menggunakan Sifat Penjumlahan}
Jika $\log 2 = 0{,}301$ dan $\log 3 = 0{,}477$ (basis 10), hitung $\log 12$.

\jawab $\log 12 = \log(2^2\cdot 3)=2\log 2+\log 3 = 2(0{,}301)+0{,}477 = 1{,}079.$
\end{contoh}

\begin{contoh}{Logaritma untuk "Berapa Lama" (HOTS)}
Pada studi kasus, kapan subscriber menembus $100.000$? Selesaikan $1000\cdot 1{,}5^{n}=100.000$.

\jawab $1{,}5^n=100 \Rightarrow n = {}^{1{,}5}\!\log 100 = \dfrac{\log 100}{\log 1{,}5}=\dfrac{2}{0{,}176}\approx 11{,}4$. Jadi pada bulan ke-12 subscriber pertama kali melampaui $100.000$. Perhatikan bagaimana logaritma "membongkar" eksponen untuk menemukan waktu.
\end{contoh}

\begin{latihan}
\begin{enumerate}[leftmargin=*]
  \item Hitung ${}^{2}\!\log 16 + {}^{4}\!\log 64$.
  \item Jika ${}^{2}\!\log 3 = p$, nyatakan ${}^{2}\!\log 18$ dalam $p$.
  \item Hitung ${}^{3}\!\log 81 - {}^{3}\!\log 9$.
  \item Sederhanakan ${}^{2}\!\log 6 + {}^{2}\!\log 8 - {}^{2}\!\log 3$.
  \item Tentukan nilai $x$ dari ${}^{2}\!\log x = 5$.
  \item Hitung ${}^{5}\!\log 10 \cdot {}^{10}\!\log 25$ (gunakan sifat rantai).
  \item Diketahui $\log 2 = 0{,}301$. Hitung $\log 5$.
  \item Selesaikan ${}^{3}\!\log(x-1)=2$.
  \item Nyatakan ${}^{a}\!\log b \cdot {}^{b}\!\log c \cdot {}^{c}\!\log a$ dalam bentuk paling sederhana.
  \item Tentukan ${}^{2}\!\log 5$ jika $\log 2=0{,}301$ dan $\log 5 = 0{,}699$.
\end{enumerate}
\end{latihan}

\begin{kesalahan}
\begin{enumerate}[leftmargin=*,topsep=2pt]
  \item \textbf{Menjumlah basis pada perkalian pangkat.} Banyak siswa menulis $2^3\cdot 2^4 = 4^7$. \emph{Salah!} Basis tetap, pangkat dijumlah: $2^3\cdot 2^4=2^7$. Penyebabnya: tertukar antara aturan perkalian basis dan perkalian pangkat.
  \item \textbf{Menganggap $(a+b)^2=a^2+b^2$.} Ini bukan sifat eksponen; yang benar $(a+b)^2=a^2+2ab+b^2$. Sifat $(ab)^n=a^nb^n$ hanya untuk \emph{perkalian}, bukan penjumlahan.
  \item \textbf{Menjumlahkan pangkat.} $a^m+a^n \neq a^{m+n}$. Sifat $a^m\cdot a^n = a^{m+n}$ hanya berlaku untuk \emph{perkalian}, bukan penjumlahan. Contoh: $2^2+2^3 = 4+8 = 12$, bukan $2^5=32$.
  \item \textbf{Logaritma jumlah.} ${}^{a}\!\log(x+y)\neq {}^{a}\!\log x+{}^{a}\!\log y$. Sifat penjumlahan berlaku untuk ${}^{a}\!\log(xy)$. Hindari dengan mengingat: logaritma mengubah \emph{kali} menjadi \emph{tambah}, bukan tambah menjadi tambah.
  \item \textbf{Lupa membalik tanda pertidaksamaan} saat basis $0<a<1$. Selalu cek besar basis sebelum menyamakan pangkat.
\end{enumerate}
\end{kesalahan}

\begin{dunianyata}
\begin{itemize}[leftmargin=*,topsep=2pt]
  \item \textbf{Keuangan:} bunga majemuk $M=M_0(1+i)^t$ adalah model eksponen murni.
  \item \textbf{Kesehatan/biologi:} pertumbuhan bakteri dan penyebaran virus (model awal epidemi) bersifat eksponensial; waktu paruh obat memakai peluruhan eksponen.
  \item \textbf{Teknologi \&\ data science:} kompleksitas algoritma $O(2^n)$, dan skala logaritmik pada grafik data yang rentangnya sangat lebar.
  \item \textbf{AI:} "log-likelihood" dan fungsi \emph{log-loss} pada pelatihan model menggunakan logaritma untuk mengubah perkalian peluang menjadi penjumlahan.
  \item \textbf{Sains:} skala Richter (gempa), pH (kimia), dan desibel (bunyi) semuanya berbasis logaritma --- tiap kenaikan 1 satuan berarti 10 kali lipat.
\end{itemize}
\end{dunianyata}

\section*{Ringkasan Bab}
\addcontentsline{toc}{section}{Ringkasan Bab}
\begin{center}
\begin{tabular}{>{\bfseries}p{3.2cm} p{10.5cm}}
\toprule
Konsep & Inti yang harus diingat \\
\midrule
Eksponen & $a^n$ = perkalian berulang; berlaku sifat (1)--(8). \\
Pangkat khusus & $a^0=1$, $a^{-n}=\tfrac{1}{a^n}$, $a^{m/n}=\sqrt[n]{a^m}$ (turunan dari sifat dasar). \\
Persamaan eksponen & Samakan basis $\Rightarrow$ samakan pangkat; gunakan pemisalan untuk bentuk kuadrat. \\
Pertidaksamaan & Basis $>1$ tanda tetap; basis $<1$ tanda dibalik. \\
Logaritma & Invers eksponen: ${}^{a}\!\log x=y\iff a^y=x$. \\
Sifat logaritma & Mengubah kali$\to$tambah, bagi$\to$kurang, pangkat$\to$kali. \\
Aplikasi & Pertumbuhan/peluruhan, bunga majemuk, skala (Richter, pH, dB). \\
\bottomrule
\end{tabular}
\end{center}

\begin{refleksi}
Beri tanda centang pada hal yang sudah kamu kuasai:
\begin{itemize}[leftmargin=*,label={},topsep=2pt]
  \item \cek Saya dapat menjelaskan eksponen sebagai perkalian berulang.
  \item \cek Saya dapat menurunkan mengapa $a^0=1$ dan $a^{-n}=\tfrac1{a^n}$.
  \item \cek Saya dapat menyelesaikan persamaan \&\ pertidaksamaan eksponen.
  \item \cek Saya memahami logaritma sebagai invers eksponen.
  \item \cek Saya dapat menerapkan eksponen/logaritma pada masalah pertumbuhan nyata.
\end{itemize}
\textit{Bila ada yang belum tercentang, tinjau kembali subbab terkait sebelum mengerjakan soal akhir bab.}
\end{refleksi}

% ---------- SOAL AKHIR BAB ----------
\section*{Soal Akhir Bab 1}
\addcontentsline{toc}{section}{Soal Akhir Bab 1}

\bagiansoal{A}{Fundamental (setara SNBT Dasar)}
\begin{enumerate}[leftmargin=*]
  \item Sederhanakan $\dfrac{a^5 b^2}{a^2 b^5}$.
  \item Hitung $32^{2/5}$.
  \item Tentukan nilai $x$ dari $3^{x}=27$.
  \item Hitung ${}^{2}\!\log 32$.
  \item Sederhanakan $\dfrac{x^{-2}y^3}{x^4 y^{-1}}$.
  \item Tentukan HP dari $2^{x+1}=64$.
  \item Hitung ${}^{3}\!\log 9 + {}^{2}\!\log 8$.
  \item Nyatakan $\sqrt[3]{a^2}$ dalam bentuk pangkat.
  \item Jika ${}^{2}\!\log 5 = q$, nyatakan ${}^{2}\!\log 20$ dalam $q$.
  \item Hitung nilai $\left(\tfrac{2}{3}\right)^{-2}$.
\end{enumerate}

\bagiansoal{B}{Intermediate (setara SNBT Sulit)}
\begin{enumerate}[leftmargin=*]
  \item Selesaikan $4^{x}-6\cdot 2^{x}+8=0$.
  \item Diketahui ${}^{a}\!\log 2 = m$ dan ${}^{a}\!\log 3 = n$. Nyatakan ${}^{a}\!\log 0{,}75$ dalam $m$ dan $n$.
  \item Tentukan himpunan penyelesaian $\left(\tfrac19\right)^{x-1} \ge 3^{x+2}$.
  \item Jika $2^x+2^{-x}=3$, tentukan nilai $4^x+4^{-x}$.
  \item Penduduk sebuah kota tumbuh $2\%$ per tahun. Jika kini $500.000$ jiwa, susun dan selesaikan persamaan untuk memperkirakan tahun saat penduduk mencapai $1.000.000$ jiwa.
\end{enumerate}

\bagiansoal{C}{Advanced (setara Olimpiade SMA --- gabungan dengan Barisan/Deret)}
\begin{enumerate}[leftmargin=*]
  \item Barisan $a, ar, ar^2, \dots$ adalah geometri. Tunjukkan bahwa barisan $\log a, \log(ar), \log(ar^2),\dots$ membentuk barisan \emph{aritmetika}, dan tentukan bedanya.
  \item Jika ${}^{2}\!\log x$, ${}^{2}\!\log(x+2)$, dan ${}^{2}\!\log(x+6)$ membentuk barisan aritmetika, tentukan $x$.
  \item Diketahui $x^{\log x}=100x$ (log basis 10). Tentukan semua nilai $x$.
  \item Hasil kali $n$ suku pertama barisan geometri positif adalah $P$. Jika ${}^{2}\!\log U_1 + {}^{2}\!\log U_n = 10$, tentukan ${}^{2}\!\log P$ dalam $n$.
  \item Tentukan jumlah ${}^{2}\!\log\tfrac{2}{1}+{}^{2}\!\log\tfrac{3}{2}+\cdots+{}^{2}\!\log\tfrac{64}{63}$.
\end{enumerate}

\bagiansoal{D}{Expert (setara tahun pertama universitas --- sintesis antarbab)}
\begin{enumerate}[leftmargin=*]
  \item \textbf{(Eksponen + Fungsi).} Diberikan $f(x)=\dfrac{a^x}{a^x+\sqrt a}$ dengan $a>0$. Buktikan $f(x)+f(1-x)=1$, lalu hitung $\sum_{k=1}^{2026} f\!\left(\tfrac{k}{2027}\right)$.
  \item \textbf{(Logaritma + Peluang).} Sebuah eksperimen menghasilkan $n$ keluaran bebas dengan peluang masing-masing $p_i$. \emph{Entropi} didefinisikan $H=-\sum p_i\,{}^{2}\!\log p_i$. Hitung $H$ untuk pelemparan dua koin adil, dan jelaskan maknanya sebagai "rata-rata informasi".
  \item \textbf{(Eksponen + Kalkulus dasar).} Fungsi peluruhan $N(t)=N_0 e^{-kt}$. Jika waktu paruh (saat $N=\tfrac12 N_0$) adalah $T$, nyatakan $k$ dalam $T$, lalu tentukan kapan tersisa $\tfrac18 N_0$.
  \item \textbf{(Logaritma + Statistika).} Pada data berskala lebar, sering dipakai transformasi $y=\log x$. Jika $x_1,\dots,x_n$ punya rata-rata geometrik $G$, buktikan bahwa rata-rata aritmetika dari $\log x_i$ sama dengan $\log G$.
  \item \textbf{(Eksponen + Pembuktian).} Buktikan dengan induksi bahwa $2^n > n^2$ untuk setiap bilangan asli $n\ge 5$.
\end{enumerate}

\begin{challenge}
\begin{enumerate}[leftmargin=*,topsep=2pt]
  \item Tentukan banyaknya digit dari $2^{100}$. (Petunjuk: gunakan $\log 2 \approx 0{,}30103$ dan ide $\lfloor \log N\rfloor + 1$.)
  \item Selesaikan persamaan $x^{x^{x^{\cdot^{\cdot^{\cdot}}}}}=2$ (menara pangkat tak hingga). Untuk nilai apa menara ini konvergen?
  \item Buktikan bahwa ${}^{2}\!\log 3$ adalah bilangan \emph{irasional}. (Petunjuk: andaikan rasional $\tfrac{p}{q}$, lalu tinjau $2^p=3^q$.)
\end{enumerate}
\end{challenge}

% ============================================================
\chapter{Barisan dan Deret}
% ============================================================
\begin{tujuan}
Setelah mempelajari bab ini, peserta didik mampu:
\begin{itemize}[leftmargin=*,topsep=2pt]
  \item mengenali pola bilangan dan membedakan barisan aritmetika serta geometri dari strukturnya;
  \item menentukan rumus suku ke-$n$ dan menggunakannya untuk memprediksi suku jauh;
  \item menghitung jumlah $n$ suku pertama (deret) aritmetika dan geometri;
  \item menjelaskan syarat dan makna deret geometri tak hingga;
  \item memodelkan masalah nyata (cicilan, pertumbuhan, lintasan) dengan barisan dan deret.
\end{itemize}
\textbf{Capaian Pembelajaran (Fase E).} Peserta didik dapat menggeneralisasi pola bilangan menjadi barisan dan deret aritmetika maupun geometri, serta memakainya untuk menyelesaikan masalah.
\end{tujuan}

\begin{pemantik}
\begin{itemize}[leftmargin=*,topsep=2pt]
  \item Jika kamu menabung Rp10.000 di hari pertama, lalu menambah Rp2.000 tiap hari, berapa total setelah sebulan? Bagaimana jika tabungannya \emph{dikali} 2 tiap hari?
  \item Mengapa dua pola di atas tumbuh sangat berbeda padahal awalnya sama?
  \item Mungkinkah menjumlahkan tak hingga banyak bilangan positif namun hasilnya tetap berhingga?
\end{itemize}
\end{pemantik}

\begin{studikasus}{Tantangan Menabung dan Bunga Cicilan}
Dua sahabat, Aisya dan Bima, menabung dengan cara berbeda selama 12 bulan.
\begin{center}
\begin{tabular}{c|c|c}
\toprule
Bulan & Aisya (tambah Rp50rb) & Bima (kali $1{,}2$) \\
\midrule
1 & 100.000 & 100.000 \\
2 & 150.000 & 120.000 \\
3 & 200.000 & 144.000 \\
4 & 250.000 & 172.800 \\
$\vdots$ & $\vdots$ & $\vdots$ \\
\bottomrule
\end{tabular}
\end{center}
\textbf{Amati dan duga.} (a) Pola mana yang "selisihnya tetap" dan mana yang "rasionya tetap"? (b) Kapan tabungan Bima menyusul Aisya? (c) Bila ini berlangsung bertahun-tahun, pola mana yang akhirnya jauh lebih besar? Pola Aisya adalah \textbf{barisan aritmetika} (beda tetap), pola Bima adalah \textbf{barisan geometri} (rasio tetap). Bab ini mempelajari keduanya beserta jumlahnya (deret).
\end{studikasus}

\section*{Peta Konsep Bab}
\addcontentsline{toc}{section}{Peta Konsep Bab}
\begin{center}
\begin{tikzpicture}[
  konsepbox/.style={draw=biru,fill=birumuda,rounded corners,align=center,inner sep=5pt,font=\small\bfseries},
  subbox/.style={draw=hijau,fill=hijaumuda,rounded corners,align=center,inner sep=4pt,font=\footnotesize},
  garis/.style={-Stealth,biru,thick}]
  \node[konsepbox] (akar) {BARISAN \&\ DERET};
  \node[subbox,below left=12mm and 14mm of akar] (ar) {Aritmetika\\(beda $b$)};
  \node[subbox,below right=12mm and 14mm of akar] (ge) {Geometri\\(rasio $r$)};
  \node[subbox,below=10mm of ar] (sar) {$U_n=a+(n-1)b$\\$S_n=\tfrac n2(a+U_n)$};
  \node[subbox,below=10mm of ge] (sge) {$U_n=ar^{n-1}$\\$S_n=\tfrac{a(r^n-1)}{r-1}$};
  \node[subbox,below=10mm of sge] (inf) {Deret tak hingga\\$S_\infty=\tfrac{a}{1-r}$};
  \draw[garis] (akar)--(ar); \draw[garis] (akar)--(ge);
  \draw[garis] (ar)--(sar); \draw[garis] (ge)--(sge); \draw[garis] (sge)--(inf);
\end{tikzpicture}\\
\small \textbf{Peta konsep.} Barisan menggambarkan polanya; deret menjumlahkannya. Kasus khusus $|r|<1$ memberi deret tak hingga yang konvergen.
\end{center}

\section{Barisan dan Deret Aritmetika}

\subsection*{Membangun konsep dari pola}
Pada pola Aisya, setiap suku diperoleh dengan \emph{menambah} bilangan tetap. Selisih tetap inilah jiwa barisan aritmetika. Karena penambahan dilakukan $(n-1)$ kali dari suku pertama, wajar jika $U_n = a+(n-1)b$. Inilah "mengapa" rumusnya berbentuk demikian, bukan sekadar hafalan.

\begin{definisi}
\emph{Barisan aritmetika} adalah barisan dengan selisih dua suku berurutan tetap, disebut \emph{beda} $b = U_n - U_{n-1}$.
\end{definisi}

\begin{konsep}{Rumus Aritmetika}
\[
U_n = a + (n-1)b, \qquad
S_n = \frac{n}{2}\big(2a + (n-1)b\big) = \frac{n}{2}(a + U_n),
\]
dengan $a = U_1$ suku pertama.
\end{konsep}

\textbf{Penurunan rumus jumlah (ide Gauss).} Tulis $S_n$ maju lalu mundur:
\[
\begin{aligned}
S_n &= a + (a+b) + \cdots + U_n,\\
S_n &= U_n + (U_n-b) + \cdots + a.
\end{aligned}
\]
Jumlahkan kolom demi kolom: tiap pasangan bernilai $a+U_n$, dan ada $n$ pasang, sehingga $2S_n=n(a+U_n)$, yakni $S_n=\tfrac n2(a+U_n)$. Inilah trik yang konon dipakai Gauss kecil menjumlah $1$ sampai $100$.

\begin{contoh}{Suku ke-$n$ Aritmetika}
Diketahui barisan $3, 7, 11, 15, \dots$. Tentukan $U_{20}$ dan $S_{20}$.

\jawab $a = 3,\ b = 4$. Maka $U_{20} = 3 + 19(4) = 79$ dan
$S_{20} = \tfrac{20}{2}(3 + 79) = 10 \cdot 82 = 820.$
\end{contoh}

\begin{contoh}{Menentukan Beda dari Dua Suku}
Diketahui $U_4=14$ dan $U_9=29$. Tentukan rumus $U_n$.

\jawab $U_9-U_4=5b=29-14=15\Rightarrow b=3$. Dari $U_4=a+3b=14$ diperoleh $a=14-9=5$. Jadi $U_n=5+(n-1)3=3n+2$.
\end{contoh}

\begin{contoh}{Sisipan (HOTS ringan)}
Di antara $4$ dan $40$ disisipkan $5$ bilangan sehingga membentuk barisan aritmetika. Tentukan beda barunya.

\jawab Setelah penyisipan ada $7$ suku, beda baru $b'=\dfrac{40-4}{7-1}=6$. Barisannya $4,10,16,22,28,34,40$.
\end{contoh}

\begin{latihan}
\begin{enumerate}[leftmargin=*]
  \item Tentukan suku ke-15 dari barisan $5, 9, 13, \dots$.
  \item Diketahui $U_3=11$, $U_7=27$. Tentukan $a$ dan $b$.
  \item Hitung $S_{30}$ dari $2+5+8+\cdots$.
  \item Tentukan jumlah semua bilangan asli kelipatan 3 antara 1 dan 100.
  \item Suku ke berapa dari $7,11,15,\dots$ yang bernilai $107$?
  \item Sisipkan 3 bilangan antara 2 dan 18 agar aritmetika; tuliskan barisannya.
  \item Sebuah gedung punya kursi: baris depan 20, tiap baris belakang bertambah 4. Berapa total kursi pada 10 baris?
  \item Jika $S_n=2n^2+3n$, tentukan rumus $U_n$.
  \item Tentukan banyak suku jika $a=3$, $b=5$, dan $S_n=255$.
  \item Tiga bilangan aritmetika berjumlah 21 dan hasil kali 231. Tentukan ketiganya.
\end{enumerate}
\end{latihan}

\section{Barisan dan Deret Geometri}

\subsection*{Membangun konsep dari pola}
Pada pola Bima, setiap suku diperoleh dengan \emph{mengalikan} rasio tetap. Karena perkalian dilakukan $(n-1)$ kali, maka $U_n=ar^{n-1}$ --- inilah bentuk eksponen yang sudah kita kenal di Bab 1. Tidak heran pertumbuhan geometri akhirnya melesat melampaui aritmetika.

\begin{definisi}
\emph{Barisan geometri} adalah barisan dengan perbandingan dua suku berurutan tetap, disebut \emph{rasio} $r = \dfrac{U_n}{U_{n-1}}$.
\end{definisi}

\begin{konsep}{Rumus Geometri}
\[
U_n = a\, r^{\,n-1}, \qquad
S_n = \frac{a(r^n - 1)}{r-1}\ (r \neq 1),\qquad
S_\infty = \frac{a}{1-r}\ \ (|r|<1).
\]
\end{konsep}

\textbf{Mengapa $S_\infty$ ada hanya jika $|r|<1$?} Jika $|r|<1$, maka $r^n\to 0$ saat $n$ membesar, sehingga $S_n=\dfrac{a(1-r^n)}{1-r}\to \dfrac{a}{1-r}$. Jika $|r|\ge 1$, suku tidak mengecil dan jumlahnya membesar tanpa batas --- tak ada nilai berhingga. Inilah jawaban pertanyaan pemantik: menjumlah tak hingga bilangan positif \emph{bisa} berhingga, asalkan suku-sukunya menyusut cukup cepat.

\begin{center}
\begin{tikzpicture}
\begin{axis}[ybar,bar width=10pt,width=9cm,height=5cm,
  xlabel=suku ke-$n$,ylabel=nilai,xtick=data,ymin=0,
  axis lines=left,legend style={font=\small}]
\addplot[fill=birumuda,draw=biru] coordinates {(1,6)(2,3)(3,1.5)(4,0.75)(5,0.375)(6,0.1875)};
\addlegendentry{$6,3,1{,}5,\dots\ (r=\tfrac12)$}
\end{axis}
\end{tikzpicture}\\
\small \textbf{Suku geometri yang menyusut.} Tinggi batang mengecil cepat; total luasnya menuju $S_\infty=12$.
\end{center}

\begin{contoh}{Deret Geometri Tak Hingga}
Hitung jumlah $6 + 3 + \tfrac{3}{2} + \cdots$.

\jawab $a = 6,\ r = \tfrac12$, sehingga $S_\infty = \dfrac{6}{1-\frac12} = 12.$
\end{contoh}

\begin{contoh}{Bunga Majemuk (kaitan dengan Bab 1)}
Modal Rp1.000.000 dibungakan majemuk $10\%$ per tahun. Tentukan saldo setelah 3 tahun.

\jawab Ini barisan geometri $a=1.000.000$, $r=1{,}1$. Saldo tahun ke-3 (suku ke-4): $1.000.000\cdot 1{,}1^3 = 1.331.000$ rupiah.
\end{contoh}

\begin{contoh}{Lintasan Bola Memantul (HOTS)}
Bola jatuh dari 8 m, tiap pantulan $\tfrac34$ tinggi sebelumnya. Tentukan total lintasan hingga berhenti.

\jawab Lintasan turun: $8+6+\tfrac92+\cdots$, deret geometri $a=8$, $r=\tfrac34$, jumlah $=\dfrac{8}{1-\frac34}=32$. Lintasan naik sama dengan turun dikurangi jatuhan pertama: $32-8=24$. Total $=32+24=56$ m.
\end{contoh}

\begin{latihan}
\begin{enumerate}[leftmargin=*]
  \item Tentukan suku ke-6 dari $3, 6, 12, \dots$.
  \item Jumlah 10 suku pertama dari $2 + 6 + 18 + \cdots$.
  \item Tentukan $S_\infty$ dari $9+3+1+\cdots$.
  \item Rasio sebuah barisan geometri $\tfrac13$, suku pertama 27. Tentukan $U_5$.
  \item Diketahui $U_2=6$, $U_5=48$. Tentukan $a$ dan $r$.
  \item Ubah desimal berulang $0{,}3333\ldots$ menjadi pecahan memakai deret tak hingga.
  \item Modal Rp2.000.000 bunga majemuk $5\%$/tahun. Saldo setelah 4 tahun?
  \item Bola jatuh dari 10 m, memantul $0{,}6$ kali. Total lintasan hingga berhenti?
  \item Tiga bilangan geometri berjumlah 21, hasil kali 216. Tentukan ketiganya.
  \item Jumlah tak hingga $\sum (\tfrac25)^n$ untuk $n\ge 1$.
\end{enumerate}
\end{latihan}

\begin{center}
\begin{tikzpicture}[scale=2.6]
  % Persegi penuh lebar 2 (konvergen ke 2)
  \draw[biru,dashed,thick] (0,0) rectangle (2,1);
  % Suku 1 = 1 (strip pertama)
  \fill[biru!65] (0,0) rectangle (1,1);
  % Suku 2 = 1/2
  \fill[biru!40] (1,0) rectangle (1.5,1);
  % Suku 3 = 1/4
  \fill[biru!22] (1.5,0) rectangle (1.75,1);
  % Suku 4 = 1/8
  \fill[biru!12] (1.75,0) rectangle (1.875,1);
  % Suku 5 = 1/16 dst.
  \fill[biru!6] (1.875,0) rectangle (1.9375,1);
  % Garis pembatas
  \draw[biru,thick] (1,0)--(1,1);
  \draw[biru,thick] (1.5,0)--(1.5,1);
  \draw[biru,thick] (1.75,0)--(1.75,1);
  \draw[biru,thick] (1.875,0)--(1.875,1);
  \draw[biru,thick] (1.9375,0)--(1.9375,1);
  % Label suku
  \node[white,font=\small\bfseries] at (0.5,0.5) {$1$};
  \node[white,font=\footnotesize\bfseries] at (1.25,0.5) {$\tfrac12$};
  \node[white,font=\scriptsize] at (1.625,0.5) {$\tfrac14$};
  \node[biru!70,font=\tiny] at (1.8125,0.5) {$\tfrac18$};
  \node[biru,font=\tiny] at (1.906,0.5) {$\cdots$};
  % Brace dan total
  \draw[biru,thick] (0,-0.1)--(0,-0.16)--(2,-0.16)--(2,-0.1);
  \node[below,font=\small,biru] at (1,-0.17) {Total $\to 2$ (garis putus-putus)};
\end{tikzpicture}\\[5pt]
\small\textbf{Konvergensi visual: $1+\tfrac{1}{2}+\tfrac{1}{4}+\tfrac{1}{8}+\cdots\to 2$.} Tiap irisan mewakili satu suku; strip semakin sempit mendekati tepi kanan tanpa pernah mencapainya. Luas total semua strip mengisi persegi panjang lebar $2$ --- itulah $S_\infty=\dfrac{1}{1-\frac12}=2$. Deret "tak berhingga banyak suku" bisa punya jumlah \emph{berhingga} jika sukunya menyusut cukup cepat.
\end{center}

\section{Notasi Sigma (Pengayaan)}

\subsection*{Mengapa perlu notasi sigma?}
Menuliskan $1+2+3+\cdots+100$ atau $U_1+U_2+\cdots+U_n$ secara panjang melelahkan dan rawan keliru. \emph{Notasi sigma} (lambang $\Sigma$, huruf Yunani "S" untuk \emph{Sum}/jumlah) adalah cara ringkas dan baku menuliskan penjumlahan beruntun. Ia menjadi bahasa wajib pada deret, statistika (Bab 7), peluang, hingga kalkulus.

\begin{definisi}
\emph{Notasi sigma} menyatakan jumlah suku $a_i$ untuk indeks $i$ dari batas bawah $m$ sampai batas atas $n$:
\[
\sum_{i=m}^{n} a_i = a_m + a_{m+1} + a_{m+2} + \cdots + a_n .
\]
Di sini $i$ disebut \emph{indeks penjumlahan} (boleh diganti $k,j,\dots$ tanpa mengubah nilai), $m$ batas bawah, $n$ batas atas, dan $a_i$ \emph{suku umum}.
\end{definisi}

\begin{contoh}{Membaca dan Menguraikan Sigma}
Uraikan lalu hitung $\displaystyle\sum_{k=1}^{4}(2k+1)$.

\jawab Substitusikan $k=1,2,3,4$ pada suku umum $2k+1$:
\[
\sum_{k=1}^{4}(2k+1)=3+5+7+9=24.
\]
\end{contoh}

\subsection*{Sifat-sifat sigma}
Karena sigma hanyalah penjumlahan yang ditulis ringkas, semua "kebebasan" penjumlahan biasa tetap berlaku.

\begin{konsep}{Sifat-Sifat Notasi Sigma}
Untuk konstanta $c$ serta suku $a_i,b_i$:
\[
\begin{aligned}
&\textbf{(1) Sigma konstanta:} && \sum_{i=1}^{n} c = nc, \\
&\textbf{(2) Penskalaan:} && \sum_{i=1}^{n} c\,a_i = c\sum_{i=1}^{n} a_i, \\
&\textbf{(3) Linear (jumlah/selisih):} && \sum_{i=1}^{n}(a_i\pm b_i)=\sum_{i=1}^{n}a_i\pm\sum_{i=1}^{n}b_i, \\
&\textbf{(4) Pemecahan batas:} && \sum_{i=1}^{n}a_i=\sum_{i=1}^{p}a_i+\sum_{i=p+1}^{n}a_i, \\
&\textbf{(5) Pergeseran indeks:} && \sum_{i=1}^{n}a_i=\sum_{i=0}^{n-1}a_{i+1}.
\end{aligned}
\]
\end{konsep}

\textbf{Mengapa $\sum_{i=1}^n c=nc$?} Suku konstan $c$ tertulis sebanyak $n$ kali, sehingga $\underbrace{c+c+\cdots+c}_{n\text{ suku}}=nc$. Sifat (2) dan (3) menyatakan bahwa "mengeluarkan faktor" dan "memisah jumlah" boleh dilakukan --- karena itu sigma disebut operator \emph{linear}. \textbf{Awas:} sigma \emph{tidak} linear terhadap perkalian: $\sum a_i b_i \neq \big(\sum a_i\big)\big(\sum b_i\big)$.

\begin{contoh}{Sigma Konstanta dan Linear}
Hitung $\displaystyle\sum_{i=1}^{20}(3i-2)$ menggunakan sifat sigma.

\jawab Pisahkan dengan sifat (3) dan (1)--(2):
\[
\sum_{i=1}^{20}(3i-2)=3\sum_{i=1}^{20} i-\sum_{i=1}^{20}2
=3\cdot\frac{20\cdot21}{2}-2\cdot20=3(210)-40=590.
\]
\end{contoh}

\subsection*{Sigma barisan baku}
\begin{konsep}{Rumus Jumlah yang Wajib Dihafal}
\[
\sum_{i=1}^{n} i = \frac{n(n+1)}{2},\qquad
\sum_{i=1}^{n} i^2 = \frac{n(n+1)(2n+1)}{6},\qquad
\sum_{i=1}^{n} i^3 = \left(\frac{n(n+1)}{2}\right)^{\!2}.
\]
Perhatikan fakta indah: $\displaystyle\sum_{i=1}^{n} i^3=\Big(\sum_{i=1}^{n} i\Big)^2$ --- jumlah pangkat tiga sama dengan kuadrat jumlah biasa.
\end{konsep}

\textbf{Sigma barisan aritmetika dan geometri.} Untuk suku aritmetika $a_i=a+(i-1)b$ dan geometri $a_i=ar^{\,i-1}$, notasi sigma hanyalah penulisan ulang rumus deret yang sudah kita kenal:
\[
\sum_{i=1}^{n}\big(a+(i-1)b\big)=\frac n2\big(2a+(n-1)b\big),\qquad
\sum_{i=1}^{n}ar^{\,i-1}=\frac{a(r^{n}-1)}{r-1}\ \ (r\neq1).
\]

\begin{contoh}{Sigma Kuadrat dan Kubik}
Hitung (a) $\displaystyle\sum_{k=1}^{10}k^2$ dan (b) $\displaystyle\sum_{k=1}^{5}k^3$.

\jawab (a) $\displaystyle\sum_{k=1}^{10}k^2=\frac{10\cdot11\cdot21}{6}=385$.\quad
(b) $\displaystyle\sum_{k=1}^{5}k^3=\left(\frac{5\cdot6}{2}\right)^2=15^2=225$.
\end{contoh}

\begin{contoh}{Sigma Barisan Geometri}
Nyatakan $2+4+8+\cdots+256$ dalam notasi sigma lalu hitung.

\jawab Suku $a_i=2^i$, dari $i=1$ sampai $i=8$ (sebab $2^8=256$):
\[
\sum_{i=1}^{8}2^i=\frac{2(2^8-1)}{2-1}=2(255)=510.
\]
\end{contoh}

\subsection*{Aplikasi sigma dalam deret}
Banyak penjumlahan rumit dapat "dipecah" menjadi sigma-sigma baku, lalu dijumlahkan dengan rumus di atas. Inilah kekuatan utama notasi sigma.

\begin{contoh}{Memecah Suku Umum (Gabungan Rumus)}
Hitung $\displaystyle\sum_{i=1}^{n} i(i+1)$.

\jawab Pecah $i(i+1)=i^2+i$, lalu gunakan dua rumus baku:
\[
\sum_{i=1}^{n} i(i+1)=\sum_{i=1}^{n} i^2+\sum_{i=1}^{n} i
=\frac{n(n+1)(2n+1)}{6}+\frac{n(n+1)}{2}=\frac{n(n+1)(n+2)}{3}.
\]
\end{contoh}

\begin{contoh}{Deret Teleskopik via Sigma (HOTS)}
Hitung $\displaystyle\sum_{k=1}^{n}\frac{1}{k(k+1)}$.

\jawab Gunakan pecahan parsial $\dfrac{1}{k(k+1)}=\dfrac1k-\dfrac{1}{k+1}$. Saat dijumlahkan, suku tengah saling meniadakan (teleskopik):
\[
\sum_{k=1}^{n}\Big(\frac1k-\frac{1}{k+1}\Big)=1-\frac{1}{n+1}=\frac{n}{n+1}.
\]
\end{contoh}

\begin{latihan}
\begin{enumerate}[leftmargin=*]
  \item Uraikan dan hitung $\displaystyle\sum_{k=1}^{5}(3k-1)$.
  \item Nyatakan $4+7+10+\cdots+31$ dalam notasi sigma.
  \item Hitung $\displaystyle\sum_{i=1}^{50} i$ dan $\displaystyle\sum_{i=1}^{50}5$.
  \item Hitung $\displaystyle\sum_{k=1}^{8}k^2$.
  \item Hitung $\displaystyle\sum_{k=1}^{4}k^3$ dan bandingkan dengan $\big(\sum_{k=1}^{4}k\big)^2$.
  \item Tentukan $\displaystyle\sum_{i=1}^{6}3\cdot2^{\,i-1}$ (sigma geometri).
  \item Sederhanakan $\displaystyle\sum_{i=1}^{n}(2i-1)$ menjadi bentuk tertutup (jumlah bilangan ganjil pertama).
  \item Hitung $\displaystyle\sum_{i=1}^{n} i(i-1)$ dalam $n$.
  \item Tunjukkan $\displaystyle\sum_{k=1}^{n}\frac{1}{k(k+1)}=\frac{n}{n+1}$ untuk $n=3$ dengan menjumlah langsung.
  \item Tuliskan kembali $\displaystyle\sum_{i=2}^{n}a_i$ sebagai sigma berbatas bawah $1$ memakai pergeseran indeks.
\end{enumerate}
\end{latihan}

\begin{kesalahan}
\begin{enumerate}[leftmargin=*,topsep=2pt]
  \item \textbf{Tertukar $n$ dan $n-1$ pada pangkat.} Banyak yang menulis $U_n=ar^n$. Yang benar $U_n=ar^{n-1}$, sebab suku pertama $U_1=ar^0=a$. Uji selalu dengan $n=1$.
  \item \textbf{Memakai $S_\infty$ saat $|r|\ge 1$.} Rumus $\tfrac{a}{1-r}$ hanya sah untuk deret konvergen. Jika $r=2$, jumlah tak hingga tidak ada.
  \item \textbf{Salah hitung lintasan bola}: lupa bahwa setelah jatuh pertama, bola \emph{naik dan turun} pada tiap pantulan. Pisahkan deret naik dan deret turun.
  \item \textbf{Menganggap semua pola "tambah" itu aritmetika.} Pola $1,2,4,7,11,\dots$ bedanya tidak tetap (ini barisan bertingkat), bukan aritmetika murni.
\end{enumerate}
\end{kesalahan}

\begin{dunianyata}
\begin{itemize}[leftmargin=*,topsep=2pt]
  \item \textbf{Ekonomi/keuangan:} bunga majemuk, anuitas, dan cicilan KPR adalah deret geometri.
  \item \textbf{Teknologi:} alamat memori, struktur pohon biner ($1,2,4,8,\dots$ simpul per level), dan kompresi data.
  \item \textbf{Kesehatan:} dosis obat berkala dan kadar sisa dalam darah membentuk deret geometri.
  \item \textbf{Fisika/konstruksi:} susunan batu bata bertingkat, peluruhan amplitudo getaran teredam.
  \item \textbf{Game development:} sistem "level up" dengan XP yang meningkat aritmetika/geometri.
\end{itemize}
\end{dunianyata}

\section*{Ringkasan Bab}
\addcontentsline{toc}{section}{Ringkasan Bab}
\begin{center}
\begin{tabular}{>{\bfseries}p{3cm} p{5.3cm} p{5.3cm}}
\toprule
Aspek & Aritmetika & Geometri \\
\midrule
Ciri & beda tetap $b=U_n-U_{n-1}$ & rasio tetap $r=U_n/U_{n-1}$ \\
Suku ke-$n$ & $U_n=a+(n-1)b$ & $U_n=ar^{n-1}$ \\
Jumlah $n$ suku & $S_n=\tfrac n2(a+U_n)$ & $S_n=\tfrac{a(r^n-1)}{r-1}$ \\
Tak hingga & tidak ada & $S_\infty=\tfrac{a}{1-r}$, syarat $|r|<1$ \\
Pertumbuhan & linear & eksponensial \\
\bottomrule
\end{tabular}
\end{center}

\begin{refleksi}
Beri tanda centang pada hal yang sudah kamu kuasai:
\begin{itemize}[leftmargin=*,label={},topsep=2pt]
  \item \cek Saya dapat membedakan barisan aritmetika dan geometri dari polanya.
  \item \cek Saya dapat menurunkan rumus $S_n$ aritmetika dengan ide Gauss.
  \item \cek Saya memahami syarat $|r|<1$ untuk deret tak hingga.
  \item \cek Saya dapat memodelkan masalah nyata (cicilan, lintasan bola) dengan deret.
\end{itemize}
\end{refleksi}

% ---------- SOAL AKHIR BAB ----------
\section*{Soal Akhir Bab 2}
\addcontentsline{toc}{section}{Soal Akhir Bab 2}

\bagiansoal{A}{Fundamental (setara SNBT Dasar)}
\begin{enumerate}[leftmargin=*]
  \item Tentukan $U_{10}$ dari $4,7,10,\dots$.
  \item Hitung $S_{12}$ dari $1+3+5+\cdots$.
  \item Tentukan rasio dan $U_5$ dari $2,6,18,\dots$.
  \item Hitung $S_\infty$ dari $12+4+\tfrac43+\cdots$.
  \item Diketahui $a=5$, $b=-2$. Tentukan $U_8$.
  \item Suku ke berapa dari $3,7,11,\dots$ yang bernilai $99$?
  \item Tentukan jumlah 6 suku pertama dari $1,2,4,8,\dots$.
  \item Tiga suku berurutan aritmetika: $x-2, x, x+2$. Jika jumlahnya 30, tentukan $x$.
  \item Ubah $0{,}777\ldots$ menjadi pecahan.
  \item Tentukan beda jika $U_5=20$ dan $U_{12}=48$.
\end{enumerate}

\bagiansoal{B}{Intermediate (setara SNBT Sulit)}
\begin{enumerate}[leftmargin=*]
  \item Jumlah $n$ suku pertama deret aritmetika adalah $S_n=3n^2-n$. Tentukan $U_{10}$ dan beda.
  \item Suku ke-3 barisan geometri adalah 18 dan suku ke-6 adalah 486. Tentukan $S_5$.
  \item Sebuah tali dipotong menjadi 6 bagian membentuk barisan geometri. Jika potongan terpendek 3 cm dan terpanjang 96 cm, tentukan panjang tali semula.
  \item Tiga bilangan membentuk barisan aritmetika berjumlah 36. Jika suku pertama dan ketiga ditambah 1 sedangkan tengah tetap, terbentuk barisan geometri. Tentukan ketiga bilangan semula.
  \item Diketahui deret geometri tak hingga berjumlah 12 dan suku pertamanya 8. Tentukan rasio dan suku ke-4.
\end{enumerate}

\bagiansoal{C}{Advanced (setara Olimpiade SMA --- gabungan dengan Eksponen/Logaritma)}
\begin{enumerate}[leftmargin=*]
  \item Jika $a,b,c$ membentuk barisan geometri, buktikan $\log a,\log b,\log c$ membentuk barisan aritmetika.
  \item Tentukan jumlah $\displaystyle\sum_{k=1}^{n} k\cdot 2^{k}$ (gabungan aritmetika $\times$ geometri).
  \item Jika $\log 2,\ \log(2^x-1),\ \log(2^x+3)$ membentuk barisan aritmetika, tentukan $x$.
  \item Suku-suku barisan geometri positif memenuhi $U_1 U_2 \cdots U_9 = 3^{18}$. Tentukan $U_5$.
  \item Sebuah deret geometri tak hingga memiliki jumlah $S$ dan jumlah kuadrat sukunya $Q$. Nyatakan rasio $r$ dalam $S$ dan $Q$.
\end{enumerate}

\bagiansoal{D}{Expert (setara tahun pertama universitas --- sintesis antarbab)}
\begin{enumerate}[leftmargin=*]
  \item \textbf{(Deret + Pembuktian induksi).} Buktikan $\displaystyle\sum_{k=1}^{n} k^2 = \frac{n(n+1)(2n+1)}{6}$ dengan induksi matematika.
  \item \textbf{(Deret geometri + limit).} Pada kurva Koch, tiap iterasi panjang menjadi $\tfrac43$ kali sebelumnya. Tunjukkan panjang total menuju tak hingga, sedangkan luas yang dilingkupi konvergen. Jelaskan paradoks ini.
  \item \textbf{(Deret + Peluang).} Sebuah permainan: lempar koin sampai muncul Angka. Peluang berhenti pada lemparan ke-$n$ adalah $(\tfrac12)^n$. Buktikan jumlah seluruh peluang $=1$ dan hitung ekspektasi banyak lemparan $\sum n(\tfrac12)^n$.
  \item \textbf{(Deret + Eksponen).} Tunjukkan bahwa $\displaystyle\sum_{n=0}^{\infty}\frac{1}{2^n}=2$ dan bandingkan dengan deret harmonik $\sum \tfrac1n$ yang divergen. Mengapa keduanya berbeda nasib?
  \item \textbf{(Anuitas).} Pinjaman Rp12.000.000 dilunasi 12 angsuran bulanan tetap dengan bunga majemuk $1\%$/bulan. Turunkan rumus besar angsuran $A$ menggunakan deret geometri, lalu hitung $A$.
\end{enumerate}

\begin{challenge}
\begin{enumerate}[leftmargin=*,topsep=2pt]
  \item (Deret teleskopik) Hitung $\displaystyle\sum_{k=1}^{100}\frac{1}{k(k+1)}$ tanpa menjumlah satu per satu.
  \item Bilangan $1$ ditulis, lalu tiap langkah setiap angka "1" diganti "10" dan "0" diganti "1". Selidiki banyak digit setelah $n$ langkah --- barisan apakah itu?
  \item (Fibonacci) Buktikan bahwa $\displaystyle\sum_{k=1}^{n} F_k = F_{n+2}-1$, dengan $F_1=F_2=1$.
\end{enumerate}
\end{challenge}

% ============================================================
\chapter{Vektor}
% ============================================================
\begin{tujuan}
Setelah mempelajari bab ini, peserta didik mampu:
\begin{itemize}[leftmargin=*,topsep=2pt]
  \item menjelaskan vektor sebagai besaran yang memiliki besar dan arah serta membedakannya dari skalar;
  \item melakukan operasi penjumlahan, pengurangan, dan perkalian skalar secara aljabar maupun geometri;
  \item menghitung panjang (norm) vektor dan vektor satuan;
  \item menghitung hasil kali titik (dot product) dan menafsirkannya sebagai ukuran "kesejajaran";
  \item menerapkan vektor pada masalah perpindahan, gaya, dan navigasi.
\end{itemize}
\textbf{Capaian Pembelajaran (Fase E).} Peserta didik dapat menyatakan dan mengoperasikan vektor pada bidang serta menggunakannya untuk menyelesaikan masalah geometri dan fisika sederhana.
\end{tujuan}

\begin{pemantik}
\begin{itemize}[leftmargin=*,topsep=2pt]
  \item Sebuah perahu mendayung lurus ke seberang sungai, tetapi tiba jauh di hilir. Mengapa? Besaran apa yang "ditambahkan" arus pada gerak perahu?
  \item Mengapa "5 km" saja tidak cukup untuk memberi tahu seseorang lokasi suatu tempat, tetapi "5 km ke arah timur laut" cukup?
  \item Dua gaya sama besar dapat menghasilkan efek nol atau efek dua kali lipat. Apa yang menentukannya?
\end{itemize}
\end{pemantik}

\begin{studikasus}{Navigasi Drone Pengantar Barang}
Sebuah drone berangkat dari gudang $O(0,0)$, terbang ke titik pelanggan $A(6,0)$ km (timur), lalu ke pelanggan kedua $B(6,4)$ km (utara). Angin meniupkan "dorongan" tetap $\vec w=(1,1)$ km terhadap setiap segmen.
\begin{center}
\begin{tikzpicture}[scale=0.8]
  \draw[->,gray] (-0.5,0)--(8,0) node[right]{$x$ (km)};
  \draw[->,gray] (0,-0.5)--(0,5.5) node[above]{$y$ (km)};
  \draw[->,very thick,biru] (0,0)--(6,0) node[midway,below]{$\vec{OA}$};
  \draw[->,very thick,hijau] (6,0)--(6,4) node[midway,right]{$\vec{AB}$};
  \draw[->,very thick,oranye,dashed] (0,0)--(6,4) node[midway,above left]{$\vec{OB}$};
  \fill (0,0) circle (2pt) node[below left]{$O$};
  \fill (6,0) circle (2pt) node[below right]{$A$};
  \fill (6,4) circle (2pt) node[above right]{$B$};
\end{tikzpicture}
\end{center}
\textbf{Amati dan duga.} (a) Bagaimana hubungan $\vec{OA}+\vec{AB}$ dengan $\vec{OB}$? (b) Berapa jarak langsung gudang ke $B$? (c) Bagaimana memodelkan pengaruh angin? Pertanyaan-pertanyaan ini menuntut besaran yang membawa \emph{arah} sekaligus \emph{besar} --- yaitu vektor. Penjumlahan vektor ternyata mengikuti aturan "ujung-ke-pangkal" yang akan kita formalkan.
\end{studikasus}

\section*{Peta Konsep Bab}
\addcontentsline{toc}{section}{Peta Konsep Bab}
\begin{center}
\begin{tikzpicture}[
  konsepbox/.style={draw=biru,fill=birumuda,rounded corners,align=center,inner sep=5pt,font=\small\bfseries},
  subbox/.style={draw=hijau,fill=hijaumuda,rounded corners,align=center,inner sep=4pt,font=\footnotesize},
  garis/.style={-Stealth,biru,thick}]
  \node[konsepbox] (akar) {VEKTOR};
  \node[subbox,below left=12mm and 16mm of akar] (rep) {Representasi\\\&\ notasi};
  \node[subbox,below=12mm of akar] (op) {Operasi\\($+,-,k\vec a$)};
  \node[subbox,below right=12mm and 16mm of akar] (uk) {Ukuran\\(panjang)};
  \node[subbox,below=10mm of op] (dot) {Hasil kali titik\\\&\ sudut};
  \draw[garis] (akar)--(rep); \draw[garis] (akar)--(op); \draw[garis] (akar)--(uk);
  \draw[garis] (op)--(dot); \draw[garis,dashed,gray] (uk)--(dot);
\end{tikzpicture}\\
\small \textbf{Peta konsep.} Dari representasi, kita membangun operasi; panjang dan sudut menyatukannya pada hasil kali titik.
\end{center}

\section{Konsep Dasar dan Operasi Vektor}

\subsection*{Membangun konsep}
Pada studi kasus, perpindahan drone dari $O$ ke $A$ lalu ke $B$ "setara" dengan satu perpindahan langsung $O$ ke $B$. Inilah inti penjumlahan vektor: \emph{menyambung} vektor ujung-ke-pangkal. Vektor berbeda dari bilangan biasa (skalar) karena membawa arah --- itulah sebabnya "5 km" saja tak menentukan posisi.

\begin{definisi}
\emph{Vektor} adalah besaran yang memiliki nilai (besar) dan arah. Vektor di $\mathbb{R}^2$ dituliskan $\vec{a} = \begin{pmatrix} a_1 \\ a_2 \end{pmatrix}$ atau $\vec{a} = a_1 \vec{i} + a_2 \vec{j}$.
\end{definisi}

\begin{center}
\begin{tikzpicture}[scale=1]
  \draw[->,gray] (-0.5,0)--(4,0) node[right]{$x$};
  \draw[->,gray] (0,-0.5)--(0,3) node[above]{$y$};
  \draw[->,very thick,biru] (0,0)--(3,2) node[midway,above left]{$\vec{a}$};
  \node[below right] at (3,2) {$(3,2)$};
\end{tikzpicture}
\end{center}

\begin{konsep}{Operasi Vektor}
Untuk $\vec{a}=(a_1,a_2)$, $\vec{b}=(b_1,b_2)$, dan skalar $k$:
\[
\vec{a}\pm\vec{b} = (a_1 \pm b_1,\ a_2 \pm b_2), \qquad
k\vec{a} = (k a_1,\ k a_2).
\]
\textbf{Panjang vektor:} $|\vec{a}| = \sqrt{a_1^2 + a_2^2}$.\\
\textbf{Vektor satuan:} $\hat{a}=\dfrac{\vec a}{|\vec a|}$ (panjang 1, arah sama).\\
\textbf{Hasil kali titik:} $\vec{a}\cdot\vec{b} = a_1 b_1 + a_2 b_2 = |\vec a||\vec b|\cos\theta$.
\end{konsep}

\textbf{Mengapa $|\vec a|=\sqrt{a_1^2+a_2^2}$?} Karena komponen $a_1$ (mendatar) dan $a_2$ (tegak) saling tegak lurus, panjang vektor adalah sisi miring segitiga siku-siku --- ini langsung dari Teorema Pythagoras. Jadi geometri vektor "menyembunyikan" Pythagoras di dalamnya.

\begin{center}
\begin{tikzpicture}[scale=1.0,>=Stealth,line join=round]
\coordinate (Ov) at (0,0); \coordinate (Xa) at (3,0); \coordinate (Pa) at (3,2);
\draw[->,gray](-0.3,0)--(4,0) node[font=\scriptsize,right]{$x$};
\draw[->,gray](0,-0.3)--(0,3) node[font=\scriptsize,above]{$y$};
\draw[->,very thick,biru](Ov)--(Pa);
\draw[oranye,very thick](Ov)--(Xa);
\draw[hijau!60!black,very thick](Xa)--(Pa);
\draw[gray] (2.74,0)--(2.74,0.26)--(3,0.26);
\node[oranye,font=\scriptsize,below] at (1.5,0){$a_1$ (komponen-$x$)};
\node[hijau!60!black,font=\scriptsize,right] at (3,1){$a_2$};
\node[biru,font=\scriptsize,above left] at (1.4,1.05){$|\vec a|=\sqrt{a_1^2+a_2^2}$};
\pic[draw,angle radius=0.55cm,"$\theta$"]{angle=Xa--Ov--Pa};
\end{tikzpicture}\\
\small Komponen mendatar $a_1$ dan tegak $a_2$ membentuk segitiga siku-siku; panjang vektor $|\vec a|$ adalah sisi miringnya.
\end{center}

\begin{center}
\begin{tikzpicture}[scale=0.9]
  \draw[->,gray] (-0.3,0)--(4.5,0) node[right]{$x$};
  \draw[->,gray] (0,-0.3)--(0,3.5) node[above]{$y$};
  % a
  \draw[->,very thick,biru] (0,0)--(3,1) node[above]{$\vec a$};
  % b from tip of a
  \draw[->,very thick,hijau] (3,1)--(4,3) node[right]{$\vec b$};
  % resultant
  \draw[->,very thick,oranye] (0,0)--(4,3) node[midway,above left]{$\vec a+\vec b$};
\end{tikzpicture}\\
\small \textbf{Aturan segitiga (ujung-ke-pangkal).} Pangkal $\vec b$ ditempelkan ke ujung $\vec a$; resultan menghubungkan pangkal awal ke ujung akhir.
\end{center}

\smallskip\noindent\textbf{Cara lain memvisualkan operasi.}
\begin{center}
\begin{tabular}{ccc}
% jajar genjang
\begin{tikzpicture}[scale=0.62,>=Stealth,line join=round]
\draw[->,very thick,biru](0,0)--(3,0.6) node[font=\scriptsize,right]{$\vec a$};
\draw[->,very thick,hijau](0,0)--(1,2) node[font=\scriptsize,left]{$\vec b$};
\draw[gray,dashed](3,0.6)--(4,2.6); \draw[gray,dashed](1,2)--(4,2.6);
\draw[->,very thick,oranye](0,0)--(4,2.6) node[font=\scriptsize,above]{$\vec a+\vec b$};
\fill (0,0) circle (1.5pt);
\end{tikzpicture}
&
% pengurangan
\begin{tikzpicture}[scale=0.62,>=Stealth,line join=round]
\draw[->,very thick,biru](0,0)--(3,0.6) node[font=\scriptsize,right]{$\vec a$};
\draw[->,very thick,hijau](0,0)--(1,2) node[font=\scriptsize,left]{$\vec b$};
\draw[->,very thick,merah](1,2)--(3,0.6) node[font=\scriptsize,above right]{$\vec a-\vec b$};
\fill (0,0) circle (1.5pt);
\end{tikzpicture}
&
% perkalian skalar
\begin{tikzpicture}[scale=0.62,>=Stealth,line join=round]
\draw[->,very thick,oranye](0,0)--(3,2) node[font=\scriptsize,right]{$2\vec a$};
\draw[->,very thick,biru](0,0)--(1.5,1) node[font=\scriptsize,above left]{$\vec a$};
\draw[->,very thick,merah](0,0)--(-1.5,-1) node[font=\scriptsize,left]{$-\vec a$};
\fill (0,0) circle (1.5pt);
\end{tikzpicture}
\\[-1pt]
\footnotesize \textbf{Jajar genjang}: $\vec a+\vec b$ diagonal & \footnotesize \textbf{Pengurangan}: $\vec a-\vec b$ dari ujung $\vec b$ ke ujung $\vec a$ & \footnotesize \textbf{Skalar}: $2\vec a$ memanjang, $-\vec a$ membalik\\
\end{tabular}
\end{center}

\begin{contoh}{Operasi dan Sudut Vektor}
Diketahui $\vec{a}=(3,4)$ dan $\vec{b}=(1,2)$. Tentukan $|\vec a|$ dan sudut antara keduanya.

\jawab $|\vec a| = \sqrt{9+16} = 5$, $|\vec b| = \sqrt5$, dan
$\vec a \cdot \vec b = 3+8 = 11$. Maka
$\cos\theta = \dfrac{11}{5\sqrt5} \approx 0{,}983$, sehingga $\theta \approx 10{,}5^\circ.$
\end{contoh}

\begin{contoh}{Menyelesaikan Studi Kasus}
Hitung $\vec{OA}+\vec{AB}$ dan jarak langsung $|\vec{OB}|$.

\jawab $\vec{OA}=(6,0)$, $\vec{AB}=(0,4)$, sehingga $\vec{OA}+\vec{AB}=(6,4)=\vec{OB}$. Jaraknya $|\vec{OB}|=\sqrt{6^2+4^2}=\sqrt{52}=2\sqrt{13}\approx 7{,}2$ km.
\end{contoh}

\begin{contoh}{Vektor Satuan dan Tegak Lurus (HOTS)}
Tentukan vektor satuan searah $\vec a=(3,4)$, lalu tentukan satu vektor yang tegak lurus $\vec a$.

\jawab $\hat a=\tfrac15(3,4)=(0{,}6,\,0{,}8)$. Vektor tegak lurus memenuhi $\vec a\cdot\vec b=0$; pilih $\vec b=(4,-3)$ sebab $3(4)+4(-3)=0$. (Trik: tukar komponen, balik satu tanda.)
\end{contoh}

\begin{latihan}
\begin{enumerate}[leftmargin=*]
  \item Diketahui $\vec a=(2,-1)$, $\vec b=(-3,4)$. Tentukan $2\vec a - \vec b$ dan $|\vec a|$.
  \item Tentukan hasil kali titik $\vec a \cdot \vec b$ jika $\vec a=(5,0)$, $\vec b=(0,5)$. Apa artinya?
  \item Tentukan vektor satuan searah $\vec u=(6,8)$.
  \item Tentukan $k$ agar $\vec a=(2,k)$ tegak lurus $\vec b=(3,-6)$.
  \item Hitung sudut antara $\vec p=(1,0)$ dan $\vec q=(1,1)$.
  \item Diketahui $\vec a=(1,2)$, $\vec b=(3,-1)$. Hitung $|\vec a+\vec b|$.
  \item Tentukan vektor $\vec{AB}$ jika $A(2,3)$ dan $B(5,-1)$.
  \item Jika $|\vec a|=5$, $|\vec b|=3$, dan sudutnya $60^\circ$, hitung $\vec a\cdot\vec b$.
  \item Titik $P(1,1)$, $Q(4,5)$. Tentukan jarak $PQ$.
  \item Tentukan proyeksi skalar $\vec a=(4,3)$ pada $\vec b=(1,0)$.
\end{enumerate}
\end{latihan}

\section{Hasil Kali Titik dan Sudut}

\subsection*{Makna dot product}
Mengapa $\vec a\cdot\vec b=|\vec a||\vec b|\cos\theta$ penting? Karena ia mengukur \emph{seberapa searah} dua vektor. Jika tegak lurus ($\theta=90^\circ$), $\cos\theta=0$ sehingga $\vec a\cdot\vec b=0$ --- inilah uji tegak lurus yang sangat berguna. Jika searah, dot product maksimum; jika berlawanan, negatif.

\begin{konsep}{Sifat Hasil Kali Titik}
\begin{itemize}[leftmargin=*,topsep=2pt]
  \item $\vec a\cdot\vec b=\vec b\cdot\vec a$ (komutatif).
  \item $\vec a\cdot\vec a=|\vec a|^2$.
  \item $\vec a\perp\vec b \iff \vec a\cdot\vec b=0$.
  \item Sudut: $\cos\theta=\dfrac{\vec a\cdot\vec b}{|\vec a|\,|\vec b|}$.
\end{itemize}
\end{konsep}

\begin{center}
\begin{tabular}{cc}
% proyeksi
\begin{tikzpicture}[scale=0.8,>=Stealth,line join=round]
\coordinate (Od) at (0,0); \coordinate (Bd) at (4,0); \coordinate (Ad) at (2.5,2);
\draw[->,very thick,hijau](Od)--(Bd) node[font=\scriptsize,right]{$\vec b$};
\draw[->,very thick,biru](Od)--(Ad) node[font=\scriptsize,above]{$\vec a$};
\draw[gray,dashed](Ad)--(2.5,0);
\draw[gray](2.5,0.24)--(2.26,0.24)--(2.26,0);
\draw[oranye,very thick](0,0)--(2.5,0);
\node[oranye,font=\scriptsize,below] at (1.25,-0.05){proyeksi $\vec a$ pada $\vec b$};
\pic[draw,angle radius=0.6cm,"$\theta$"]{angle=Bd--Od--Ad};
\end{tikzpicture}
&
% tegak lurus
\begin{tikzpicture}[scale=0.8,>=Stealth,line join=round]
\draw[->,very thick,hijau](0,0)--(3,0) node[font=\scriptsize,right]{$\vec b$};
\draw[->,very thick,biru](0,0)--(0,2.5) node[font=\scriptsize,above]{$\vec a$};
\draw[gray](0.26,0)--(0.26,0.26)--(0,0.26);
\node[font=\scriptsize] at (1.7,1.6){$\vec a\cdot\vec b=0$};
\fill (0,0) circle (1.3pt);
\end{tikzpicture}
\\[-1pt]
\footnotesize \textbf{Proyeksi}: $\vec a\cdot\vec b=|\vec a||\vec b|\cos\theta$ = (panjang $\vec b$)$\times$(bayangan $\vec a$ pada $\vec b$) & \footnotesize \textbf{Tegak lurus}: $\theta=90^\circ\Rightarrow\cos\theta=0\Rightarrow\vec a\cdot\vec b=0$\\
\end{tabular}
\end{center}

\begin{contoh}{Membuktikan Segitiga Siku-siku}
Titik $A(1,1)$, $B(4,2)$, $C(3,5)$. Tunjukkan segitiga $ABC$ siku-siku di $B$.

\jawab $\vec{BA}=(-3,-1)$, $\vec{BC}=(-1,3)$. Maka $\vec{BA}\cdot\vec{BC}=(-3)(-1)+(-1)(3)=3-3=0$, jadi $\vec{BA}\perp\vec{BC}$ dan segitiga siku-siku di $B$.
\end{contoh}

\begin{latihan}
\begin{enumerate}[leftmargin=*]
  \item Hitung $\vec a\cdot\vec b$ untuk $\vec a=(2,3)$, $\vec b=(4,-1)$.
  \item Tentukan sudut antara $\vec a=(1,\sqrt3)$ dan $\vec b=(1,0)$.
  \item Buktikan $\vec a=(2,4)$ dan $\vec b=(6,-3)$ saling tegak lurus.
  \item Diketahui $\vec a\cdot\vec b=10$, $|\vec a|=5$, $|\vec b|=4$. Tentukan $\cos\theta$.
  \item Cari $x$ agar sudut $(x,1)$ dan $(2,3)$ adalah $90^\circ$.
  \item Hitung $|\vec a|$ jika $\vec a\cdot\vec a=49$.
  \item Tentukan proyeksi vektor $\vec a=(3,4)$ pada $\vec b=(0,1)$.
  \item Tiga titik $A(0,0)$, $B(2,0)$, $C(2,2)$: hitung sudut $BAC$.
  \item Apakah $(1,2)$, $(2,4)$ sejajar? Jelaskan lewat dot product dan rasio komponen.
  \item Gaya $\vec F=(6,8)$ N memindahkan benda sejauh $\vec s=(3,0)$ m. Hitung usaha $W=\vec F\cdot\vec s$.
\end{enumerate}
\end{latihan}

\section{Vektor Lanjutan (Pengayaan)}

\subsection*{Vektor posisi}
Setiap titik $A(a_1,a_2)$ di bidang dapat "ditunjuk" oleh sebuah vektor dari titik asal $O$ ke titik itu. Vektor inilah jembatan antara \emph{geometri titik} dan \emph{aljabar vektor}.

\begin{definisi}
\emph{Vektor posisi} titik $A$ adalah vektor $\vec{a}=\vec{OA}=\begin{pmatrix}a_1\\a_2\end{pmatrix}$ yang berpangkal di titik asal $O(0,0)$ dan berujung di $A$. Akibatnya, vektor dari $A$ ke $B$ adalah selisih vektor posisinya: $\vec{AB}=\vec{b}-\vec{a}$.
\end{definisi}

\textbf{Mengapa $\vec{AB}=\vec b-\vec a$?} Untuk pergi dari $A$ ke $B$, kita boleh "mampir" ke titik asal: $\vec{AB}=\vec{AO}+\vec{OB}=-\vec a+\vec b$. Inilah alasan rumus "ujung dikurangi pangkal".

\begin{contoh}{Titik Tengah dan Pembagian Ruas Garis}
Diketahui $A(2,1)$ dan $B(8,5)$. Tentukan vektor posisi titik tengah $M$ ruas $AB$.

\jawab $\vec m=\tfrac12(\vec a+\vec b)=\tfrac12\big((2,1)+(8,5)\big)=\tfrac12(10,6)=(5,3)$. Jadi $M(5,3)$. Secara umum titik yang membagi $AB$ dengan rasio $m:n$ berposisi $\dfrac{n\vec a+m\vec b}{m+n}$.
\end{contoh}

\subsection*{Proyeksi vektor}
Sering kita perlu mengetahui "seberapa banyak" sebuah vektor mengarah \emph{searah} vektor lain --- misalnya komponen gaya yang searah perpindahan. Inilah proyeksi.

\begin{konsep}{Proyeksi Skalar dan Vektor}
Proyeksi $\vec a$ pada $\vec b$:
\[
\text{proyeksi skalar: } a_{\parallel}=\frac{\vec a\cdot\vec b}{|\vec b|},\qquad
\text{proyeksi vektor: } \vec a_{\parallel}=\left(\frac{\vec a\cdot\vec b}{|\vec b|^2}\right)\vec b .
\]
Proyeksi skalar adalah \emph{panjang bertanda}; proyeksi vektor adalah "bayangan" $\vec a$ pada garis arah $\vec b$.
\end{konsep}

\textbf{Dari mana rumus ini?} Karena $\vec a\cdot\vec b=|\vec a||\vec b|\cos\theta$, panjang bayangan $\vec a$ pada arah $\vec b$ adalah $|\vec a|\cos\theta=\dfrac{\vec a\cdot\vec b}{|\vec b|}$. Kalikan dengan vektor satuan $\hat b=\vec b/|\vec b|$ untuk memperoleh proyeksi vektornya.

\begin{contoh}{Menghitung Proyeksi}
Tentukan proyeksi vektor $\vec a=(4,3)$ pada $\vec b=(1,1)$.

\jawab $\vec a\cdot\vec b=4+3=7$, $|\vec b|^2=2$. Maka
\[
\vec a_{\parallel}=\frac{7}{2}(1,1)=\Big(\tfrac72,\tfrac72\Big),\qquad
a_{\parallel}=\frac{7}{\sqrt2}=\tfrac{7\sqrt2}{2}.
\]
\end{contoh}

\subsection*{Persamaan garis dalam bentuk vektor}
Selain bentuk $y=mx+c$, sebuah garis dapat dinyatakan memakai \emph{satu titik acuan} dan \emph{satu vektor arah}.

\begin{konsep}{Persamaan Garis Vektor}
Garis yang melalui titik berposisi $\vec a$ dan sejajar vektor arah $\vec u$ adalah himpunan titik $\vec r$ dengan
\[
\vec r = \vec a + t\,\vec u,\qquad t\in\mathbb{R}.
\]
Parameter $t$ "menggeser" titik di sepanjang garis. Bentuk ini mudah diperluas ke ruang $\mathbb R^3$, tempat $y=mx+c$ tak lagi cukup.
\end{konsep}

\begin{contoh}{Menyusun Persamaan Garis Vektor}
Tentukan persamaan vektor garis melalui $A(1,2)$ dan $B(4,6)$.

\jawab Vektor arah $\vec u=\vec{AB}=(3,4)$. Maka $\vec r=(1,2)+t(3,4)$, yakni $x=1+3t,\ y=2+4t$. Eliminasi $t$ memberi $\dfrac{x-1}{3}=\dfrac{y-2}{4}$, setara $y=\tfrac43 x+\tfrac23$ --- konsisten dengan bentuk biasa.
\end{contoh}

\subsection*{Aplikasi vektor dalam fisika}
Vektor lahir dari kebutuhan fisika: gaya, kecepatan, dan perpindahan semuanya berarah.

\begin{dunianyata}
\begin{itemize}[leftmargin=*,topsep=2pt]
  \item \textbf{Resultan gaya:} beberapa gaya dijumlahkan secara vektor; benda setimbang jika resultannya $\vec 0$.
  \item \textbf{Gerak relatif:} kecepatan perahu terhadap tanah $=$ kecepatan perahu terhadap air $+$ kecepatan arus (penjumlahan vektor).
  \item \textbf{Usaha:} $W=\vec F\cdot\vec s$ --- hanya komponen gaya searah perpindahan (proyeksi!) yang melakukan usaha.
  \item \textbf{Bidang miring:} gaya berat diuraikan menjadi komponen sejajar dan tegak lurus bidang dengan proyeksi.
\end{itemize}
\end{dunianyata}

\begin{contoh}{Usaha sebagai Proyeksi (Fisika)}
Gaya $\vec F=(10,6)$ N menarik kotak yang berpindah $\vec s=(5,0)$ m. Hitung usaha yang dilakukan.

\jawab $W=\vec F\cdot\vec s=10(5)+6(0)=50$ joule. Hanya komponen mendatar gaya (proyeksi pada arah perpindahan) yang berkontribusi; komponen tegak $6$ N tidak melakukan usaha.
\end{contoh}

\begin{latihan}
\begin{enumerate}[leftmargin=*]
  \item Diketahui $A(3,-1)$, $B(7,5)$. Tentukan vektor posisi titik tengah $AB$.
  \item Titik $P$ membagi $AB$ dengan rasio $2:1$, $A(0,0)$, $B(6,9)$. Tentukan posisi $P$.
  \item Tentukan proyeksi skalar $\vec a=(3,4)$ pada $\vec b=(1,0)$.
  \item Tentukan proyeksi vektor $\vec a=(2,6)$ pada $\vec b=(1,1)$.
  \item Susun persamaan vektor garis melalui $(2,1)$ dengan arah $\vec u=(1,3)$.
  \item Ubah $\vec r=(0,1)+t(2,2)$ menjadi bentuk $y=mx+c$.
  \item Gaya $\vec F=(8,6)$ N memindahkan benda sejauh $\vec s=(4,3)$ m searah. Hitung usaha $W=\vec F\cdot\vec s$.
  \item Tentukan vektor posisi titik berat segitiga $A(0,0)$, $B(6,0)$, $C(0,6)$ (rata-rata ketiga posisi).
  \item Apakah titik $(7,11)$ terletak pada garis $\vec r=(1,2)+t(3,4)$? Periksa.
  \item Perahu berkecepatan $(0,5)$ m/s menyeberang sungai berarus $(2,0)$ m/s. Tentukan vektor kecepatan resultan dan besarnya.
\end{enumerate}
\end{latihan}

\begin{kesalahan}
\begin{enumerate}[leftmargin=*,topsep=2pt]
  \item \textbf{Menjumlah vektor seperti skalar.} $|\vec a+\vec b|\neq|\vec a|+|\vec b|$ kecuali keduanya searah. Selalu jumlahkan per komponen dulu, baru hitung panjang.
  \item \textbf{Hasil kali titik dikira vektor.} $\vec a\cdot\vec b$ menghasilkan \emph{skalar} (bilangan), bukan vektor. Menulis "$\vec a\cdot\vec b=(\ldots,\ldots)$" adalah keliru.
  \item \textbf{Lupa arah pada pengurangan.} $\vec{AB}=B-A$, bukan $A-B$. Salah urutan membalik arah vektor.
  \item \textbf{Menyamakan sejajar dengan tegak lurus.} Sejajar: komponen proporsional; tegak lurus: dot product nol. Keduanya berlawanan makna.
\end{enumerate}
\end{kesalahan}

\begin{dunianyata}
\begin{itemize}[leftmargin=*,topsep=2pt]
  \item \textbf{GPS \&\ navigasi:} posisi dan perpindahan dihitung sebagai vektor; arah angin/arus ditambahkan secara vektor.
  \item \textbf{Game development \&\ grafika:} pergerakan karakter, deteksi tumbukan, pencahayaan (dot product menentukan terang permukaan).
  \item \textbf{Robotika:} arah dan kecepatan aktuator, perencanaan lintasan.
  \item \textbf{Fisika/teknik:} gaya, kecepatan, momentum, dan usaha ($W=\vec F\cdot\vec s$).
  \item \textbf{AI/data science:} representasi data sebagai vektor; "cosine similarity" mengukur kemiripan dokumen persis dengan rumus $\cos\theta$.
\end{itemize}
\end{dunianyata}

\begin{konsep}{Bahasa Alternatif yang Sering Mengecoh di Soal}
\begin{itemize}[leftmargin=*,topsep=2pt]
  \item \textbf{Vektor} (punya besar \emph{dan} arah) vs \textbf{skalar} (hanya besar, mis. jarak, suhu). "Kecepatan" itu vektor; "kelajuan" skalar.
  \item $|\vec a|$ (panjang/besar, sebuah \emph{bilangan}) berbeda dari $\vec a$ (vektornya). Hasil $\vec a\cdot\vec b$ juga \emph{skalar}, bukan vektor.
  \item Notasi setara: $(a_1,a_2)$, $\begin{psmallmatrix}a_1\\a_2\end{psmallmatrix}$, dan $a_1\vec i+a_2\vec j$ menyatakan vektor yang sama.
  \item \textbf{Vektor posisi} (berpangkal di $O$) vs \textbf{vektor bebas} (boleh digeser); $\vec{AB}=\vec b-\vec a$ (\emph{ujung} dikurangi \emph{pangkal}), bukan $\vec a-\vec b$.
  \item \textbf{Proyeksi skalar} (sebuah bilangan) berbeda dari \textbf{proyeksi vektor} (vektor searah $\vec b$).
  \item \textbf{"Searah"} berarti sudut $0^\circ$; \textbf{"sejajar"} bisa searah \emph{atau} berlawanan ($0^\circ$ atau $180^\circ$).
  \item $\vec a\cdot\vec b=0$ berarti \emph{tegak lurus}, \textbf{bukan} berarti salah satu vektor nol.
  \item \textbf{"Resultan"} = jumlah (penjumlahan) beberapa vektor gaya/perpindahan.
\end{itemize}
\end{konsep}

\section*{Ringkasan Bab}
\addcontentsline{toc}{section}{Ringkasan Bab}
\begin{center}
\begin{tabular}{>{\bfseries}p{3.5cm} p{10.2cm}}
\toprule
Konsep & Inti yang harus diingat \\
\midrule
Vektor & besaran berarah; $\vec a=(a_1,a_2)$. \\
Operasi & jumlah/selisih per komponen; $k\vec a$ menskala panjang. \\
Panjang & $|\vec a|=\sqrt{a_1^2+a_2^2}$ (Pythagoras). \\
Vektor satuan & $\hat a=\vec a/|\vec a|$. \\
Dot product & $\vec a\cdot\vec b=a_1b_1+a_2b_2=|\vec a||\vec b|\cos\theta$ (skalar). \\
Tegak lurus & $\vec a\cdot\vec b=0$. \\
Aplikasi & navigasi, gaya, usaha, kemiripan data. \\
\bottomrule
\end{tabular}
\end{center}

\begin{refleksi}
Beri tanda centang pada hal yang sudah kamu kuasai:
\begin{itemize}[leftmargin=*,label={},topsep=2pt]
  \item \cek Saya dapat menjumlah vektor secara geometri (ujung-ke-pangkal) dan aljabar.
  \item \cek Saya memahami mengapa $|\vec a|$ berasal dari Pythagoras.
  \item \cek Saya dapat menghitung sudut antar vektor dengan dot product.
  \item \cek Saya dapat menguji ketegaklurusan dua vektor.
\end{itemize}
\end{refleksi}

% ---------- SOAL AKHIR BAB ----------
\section*{Soal Akhir Bab 3}
\addcontentsline{toc}{section}{Soal Akhir Bab 3}

\bagiansoal{A}{Fundamental (setara SNBT Dasar)}
\begin{enumerate}[leftmargin=*]
  \item Diketahui $\vec a=(3,-2)$, $\vec b=(1,5)$. Tentukan $\vec a+\vec b$.
  \item Hitung $|\vec a|$ untuk $\vec a=(8,6)$.
  \item Tentukan $3\vec a$ jika $\vec a=(-1,2)$.
  \item Hitung $\vec a\cdot\vec b$ untuk $\vec a=(2,1)$, $\vec b=(3,4)$.
  \item Tentukan $\vec{AB}$ jika $A(1,2)$, $B(4,6)$.
  \item Tentukan vektor satuan searah $(0,5)$.
  \item Apakah $(2,3)$ dan $(4,6)$ sejajar?
  \item Hitung $\vec a-\vec b$ untuk $\vec a=(5,5)$, $\vec b=(2,1)$.
  \item Tentukan $k$ agar $(k,2)\perp(2,1)$.
  \item Hitung jarak $P(0,0)$ ke $Q(3,4)$.
\end{enumerate}

\bagiansoal{B}{Intermediate (setara SNBT Sulit)}
\begin{enumerate}[leftmargin=*]
  \item Titik $A(1,2)$, $B(4,3)$, $C(2,7)$. Tentukan jenis segitiga $ABC$ (lihat panjang sisi).
  \item Diketahui $\vec a=(2,3)$, $\vec b=(t,4)$ membentuk sudut $45^\circ$. Tentukan $t$.
  \item Tentukan titik $D$ agar $ABCD$ jajaran genjang, dengan $A(1,1)$, $B(4,2)$, $C(5,5)$.
  \item Vektor $\vec a=(3,4)$ diuraikan menjadi komponen sejajar dan tegak lurus terhadap $\vec b=(1,0)$. Tentukan keduanya.
  \item Perahu menyeberang dengan kecepatan $(0,4)$ m/s, arus sungai $(3,0)$ m/s. Tentukan kecepatan resultan, arah, dan besarnya.
\end{enumerate}

\bagiansoal{C}{Advanced (setara Olimpiade SMA --- gabungan dengan Trigonometri)}
\begin{enumerate}[leftmargin=*]
  \item Buktikan aturan kosinus $c^2=a^2+b^2-2ab\cos C$ menggunakan vektor sisi-sisi segitiga.
  \item Diketahui $|\vec a|=|\vec b|=1$ dan $|\vec a+\vec b|=\sqrt3$. Tentukan sudut antara $\vec a$ dan $\vec b$.
  \item Titik $A$, $B$, $C$ pada lingkaran satuan dengan sudut pusat $0^\circ, 120^\circ, 240^\circ$. Buktikan $\vec{OA}+\vec{OB}+\vec{OC}=\vec 0$.
  \item Diberikan segitiga dengan titik berat $G$. Buktikan $\vec{GA}+\vec{GB}+\vec{GC}=\vec0$.
  \item Dua vektor satuan membentuk sudut $\theta$. Nyatakan $|\vec a-\vec b|$ dalam $\theta$ dan tafsirkan hasilnya untuk $\theta=60^\circ$.
\end{enumerate}

\bagiansoal{D}{Expert (setara tahun pertama universitas --- sintesis antarbab)}
\begin{enumerate}[leftmargin=*]
  \item \textbf{(Vektor + Trigonometri).} Buktikan ketaksamaan Cauchy--Schwarz $|\vec a\cdot\vec b|\le|\vec a||\vec b|$ dan jelaskan kapan kesamaan tercapai.
  \item \textbf{(Vektor + Fungsi Kuadrat).} Tinjau $f(t)=|\vec a+t\vec b|^2$ sebagai fungsi kuadrat dalam $t$. Tentukan $t$ yang meminimumkan $f$, dan tunjukkan hubungannya dengan proyeksi.
  \item \textbf{(Vektor + Statistika).} Tunjukkan bahwa "cosine similarity" dua vektor data sama dengan koefisien korelasi data yang telah dipusatkan (dikurangi rata-rata).
  \item \textbf{(Vektor + Geometri).} Buktikan bahwa diagonal-diagonal belah ketupat saling tegak lurus menggunakan dot product.
  \item \textbf{(Vektor 3D).} Perkenalan singkat: untuk $\vec a,\vec b\in\mathbb R^3$, definisikan cross product. Hitung $\vec a\times\vec b$ untuk $\vec a=(1,0,0)$, $\vec b=(0,1,0)$ dan tafsirkan arah serta besarnya sebagai luas.
\end{enumerate}

\begin{challenge}
\begin{enumerate}[leftmargin=*,topsep=2pt]
  \item Empat titik di bidang membentuk segi empat. Buktikan bahwa titik-titik tengah sisinya selalu membentuk jajaran genjang (Teorema Varignon), gunakan vektor.
  \item Tiga gaya sama besar bekerja pada satu titik dan setimbang ($\vec0$). Buktikan sudut antar gaya pasti $120^\circ$.
  \item Tentukan vektor $\vec v$ di bidang yang memenuhi $\vec v\cdot(1,2)=5$ dan $|\vec v|$ minimum. (Petunjuk: geometri proyeksi.)
\end{enumerate}
\end{challenge}

% ============================================================
\chapter{Trigonometri Dasar}
% ============================================================
\begin{tujuan}
Setelah mempelajari bab ini, peserta didik mampu:
\begin{itemize}[leftmargin=*,topsep=2pt]
  \item menjelaskan perbandingan trigonometri (sin, cos, tan) sebagai \emph{rasio} sisi segitiga siku-siku yang bergantung hanya pada sudut;
  \item menentukan nilai trigonometri sudut istimewa dan memahami asal-usulnya;
  \item menggunakan trigonometri untuk menghitung panjang dan sudut yang sulit diukur langsung;
  \item menerapkan identitas dasar $\sin^2\alpha+\cos^2\alpha=1$;
  \item memodelkan masalah ketinggian, jarak, dan kemiringan dengan trigonometri.
\end{itemize}
\textbf{Capaian Pembelajaran (Fase E).} Peserta didik dapat menjelaskan dan menggunakan perbandingan trigonometri pada segitiga siku-siku untuk menyelesaikan masalah.
\end{tujuan}

\begin{pemantik}
\begin{itemize}[leftmargin=*,topsep=2pt]
  \item Bagaimana cara mengukur tinggi sebuah menara tanpa memanjatnya?
  \item Dua segitiga siku-siku memiliki sudut yang sama tetapi ukuran berbeda. Mengapa "perbandingan" sisinya tetap sama? Apa makna pentingnya hal ini?
  \item Mengapa fungsi sin dan cos muncul di mana-mana: gelombang suara, arus listrik, hingga animasi? Apa yang "berulang" pada semuanya?
\end{itemize}
\end{pemantik}

\begin{studikasus}{Mengukur Tinggi Menara}
Seorang siswa berdiri 50 m dari kaki sebuah menara. Dengan klinometer ia mengukur sudut elevasi puncak menara sebesar $40^\circ$. Tinggi matanya $1{,}5$ m.
\begin{center}
\begin{tikzpicture}[scale=0.9]
  \coordinate (A) at (0,0); \coordinate (B) at (5,0); \coordinate (C) at (5,4);
  \draw[very thick,biru] (A)--(B)--(C)--cycle;
  \draw (4.7,0) rectangle (5,0.3);
  \pic[draw,angle radius=0.9cm,"$40^\circ$"] {angle=B--A--C} ;
  \node[below] at (2.5,0) {50 m};
  \node[right] at (5,2) {tinggi $=?$};
  \fill (0,0) circle (2pt) node[left]{mata};
\end{tikzpicture}
\end{center}
\textbf{Amati dan duga.} (a) Besaran apa yang diketahui (sisi samping) dan apa yang dicari (sisi depan)? (b) Perbandingan trigonometri mana yang menghubungkan keduanya? (c) Bagaimana menambahkan tinggi mata? Karena yang kita kaitkan adalah sisi \emph{depan} dan \emph{samping}, kuncinya adalah $\tan 40^\circ$. Inilah kekuatan trigonometri: mengubah sudut menjadi jarak yang tak terjangkau.
\end{studikasus}

\section*{Peta Konsep Bab}
\addcontentsline{toc}{section}{Peta Konsep Bab}
\begin{center}
\begin{tikzpicture}[
  konsepbox/.style={draw=biru,fill=birumuda,rounded corners,align=center,inner sep=5pt,font=\small\bfseries},
  subbox/.style={draw=hijau,fill=hijaumuda,rounded corners,align=center,inner sep=4pt,font=\footnotesize},
  garis/.style={-Stealth,biru,thick}]
  \node[konsepbox] (akar) {TRIGONOMETRI};
  \node[subbox,below left=12mm and 16mm of akar] (rasio) {Perbandingan\\sin, cos, tan};
  \node[subbox,below=12mm of akar] (ist) {Sudut\\istimewa};
  \node[subbox,below right=12mm and 16mm of akar] (id) {Identitas\\$\sin^2+\cos^2=1$};
  \node[subbox,below=10mm of ist] (apl) {Aplikasi:\\tinggi \&\ jarak};
  \draw[garis] (akar)--(rasio); \draw[garis] (akar)--(ist); \draw[garis] (akar)--(id);
  \draw[garis] (ist)--(apl); \draw[garis,dashed,gray] (rasio)--(apl);
\end{tikzpicture}\\
\small \textbf{Peta konsep.} Dari rasio sisi lahir nilai sudut istimewa dan identitas; semuanya bermuara pada pengukuran tak langsung.
\end{center}

\section{Perbandingan Trigonometri}

\subsection*{Membangun konsep dari kesebangunan}
Mengapa $\sin 30^\circ$ selalu $\tfrac12$, tak peduli besar segitiganya? Karena semua segitiga siku-siku bersudut $30^\circ$ \emph{sebangun} --- perbandingan sisinya identik. Trigonometri adalah cara memberi "nama" pada perbandingan tetap itu. Inilah jawaban pemantik: rasio bergantung pada sudut saja, bukan ukuran.

\begin{center}
\begin{tikzpicture}[scale=1.3]
  \coordinate (A) at (0,0);
  \coordinate (B) at (3,0);
  \coordinate (C) at (3,2);
  \draw[very thick,biru] (A)--(B)--(C)--cycle;
  \draw (B) rectangle ++(-0.3,0.3);
  \node[below] at (1.5,0) {sisi samping};
  \node[right] at (3,1) {sisi depan};
  \node[above left] at (1.5,1) {sisi miring};
  \pic[draw,angle radius=0.7cm,"$\alpha$"] {angle=B--A--C};
\end{tikzpicture}
\end{center}

\begin{konsep}{Definisi pada Segitiga Siku-Siku}
\[
\sin\alpha = \frac{\text{depan}}{\text{miring}},\quad
\cos\alpha = \frac{\text{samping}}{\text{miring}},\quad
\tan\alpha = \frac{\text{depan}}{\text{samping}}.
\]
Jembatan keledai: \textbf{sin}-de-mi, \textbf{cos}-sa-mi, \textbf{tan}-de-sa. Selain itu $\csc,\sec,\cot$ adalah kebalikannya.
\end{konsep}

\begin{konsep}{Sudut Istimewa}
\[
\begin{array}{c|ccccc}
\alpha & 0^\circ & 30^\circ & 45^\circ & 60^\circ & 90^\circ\\\hline
\sin\alpha & 0 & \frac12 & \frac{\sqrt2}{2} & \frac{\sqrt3}{2} & 1\\[4pt]
\cos\alpha & 1 & \frac{\sqrt3}{2} & \frac{\sqrt2}{2} & \frac12 & 0\\[4pt]
\tan\alpha & 0 & \frac{\sqrt3}{3} & 1 & \sqrt3 & -\\
\end{array}
\]
\end{konsep}

\textbf{Dari mana nilai istimewa itu?} Sudut $45^\circ$ berasal dari segitiga siku-siku sama kaki (sisi $1,1,\sqrt2$), sehingga $\sin45^\circ=\tfrac{1}{\sqrt2}=\tfrac{\sqrt2}{2}$. Sudut $30^\circ$ dan $60^\circ$ berasal dari setengah segitiga sama sisi (sisi $1,\sqrt3,2$). Sekali memahami asalnya, tabel ini tak perlu dihafal mati.

\begin{center}
\begin{tikzpicture}[scale=1.1]
% ---- Segitiga 45-45-90 (kiri) ----
\coordinate (La) at (0,0);\coordinate (Lb) at (2,0);\coordinate (Lc) at (2,2);
\draw[very thick,biru](La)--(Lb)--(Lc)--cycle;
\draw[biru](1.8,0)--(1.8,0.2)--(2,0.2);
\node[below,font=\scriptsize] at (1,0) {$1$};
\node[right,font=\scriptsize] at (2,1) {$1$};
\node[above left,font=\scriptsize] at (1,1) {$\sqrt{2}$};
\pic[draw,biru,angle radius=0.55cm,"$45^\circ$"]{angle=Lb--La--Lc};
\node[font=\small\bfseries,biru] at (1,-0.55) {Segitiga $45^\circ$-$45^\circ$-$90^\circ$};
% ---- Segitiga 30-60-90 (kanan) ----
\coordinate (Ra) at (3.5,0);\coordinate (Rb) at (6.5,0);\coordinate (Rc) at (6.5,1.73);
\draw[very thick,oranye](Ra)--(Rb)--(Rc)--cycle;
\draw[oranye](6.3,0)--(6.3,0.2)--(6.5,0.2);
\node[below,font=\scriptsize] at (5,0) {$\sqrt{3}$};
\node[right,font=\scriptsize] at (6.5,0.87) {$1$};
\node[above left,font=\scriptsize] at (5,0.87) {$2$};
\pic[draw,oranye,angle radius=0.6cm,"$30^\circ$"]{angle=Rb--Ra--Rc};
\pic[draw,oranye,angle radius=0.4cm,"$60^\circ$"]{angle=Ra--Rc--Rb};
\node[font=\small\bfseries,oranye] at (5,-0.55) {Segitiga $30^\circ$-$60^\circ$-$90^\circ$};
\end{tikzpicture}\\[6pt]
\small\textbf{Asal-usul sudut istimewa.} Kiri: siku-siku sama kaki (sisi $1,1,\sqrt{2}$) → nilai $45^\circ$. Kanan: setengah segitiga sama sisi bersisi $2$ (sisi $\sqrt{3},1,2$) → nilai $30^\circ$ dan $60^\circ$. Memahami asal-usulnya berarti tabel tidak perlu dihafal.
\end{center}

\begin{contoh}{Menghitung Sisi}
Pada segitiga siku-siku, sudut $\alpha = 30^\circ$ dan sisi miring $= 10$ cm. Tentukan panjang sisi depan.

\jawab $\sin 30^\circ = \dfrac{\text{depan}}{10} \Rightarrow \text{depan} = 10 \cdot \tfrac12 = 5$ cm.
\end{contoh}

\begin{contoh}{Menyelesaikan Studi Kasus Menara}
Hitung tinggi menara pada studi kasus ($\tan 40^\circ\approx 0{,}839$).

\jawab Tinggi di atas mata: $h=50\tan40^\circ\approx 50(0{,}839)=41{,}95$ m. Tambah tinggi mata: $41{,}95+1{,}5\approx 43{,}5$ m.
\end{contoh}

\begin{contoh}{Identitas Dasar (HOTS)}
Diketahui $\sin\alpha=\tfrac35$ dan $\alpha$ lancip. Tentukan $\cos\alpha$ dan $\tan\alpha$ tanpa mencari $\alpha$.

\jawab Dari $\sin^2\alpha+\cos^2\alpha=1$: $\cos^2\alpha=1-\tfrac{9}{25}=\tfrac{16}{25}$, jadi $\cos\alpha=\tfrac45$ (positif karena lancip). Maka $\tan\alpha=\tfrac{\sin\alpha}{\cos\alpha}=\tfrac34$.
\end{contoh}

\begin{latihan}
\begin{enumerate}[leftmargin=*]
  \item Diketahui sisi depan $= 6$ dan samping $= 8$. Tentukan $\sin\alpha$, $\cos\alpha$, $\tan\alpha$.
  \item Sebuah tangga bersandar membentuk sudut $60^\circ$ dengan tanah. Jika panjang tangga 4 m, berapa tinggi dinding yang dicapai?
  \item Hitung $\sin45^\circ+\cos45^\circ$.
  \item Jika $\cos\alpha=\tfrac{12}{13}$ ($\alpha$ lancip), tentukan $\sin\alpha$.
  \item Tentukan nilai $\tan60^\circ\cdot\tan30^\circ$.
  \item Sebuah layang-layang terbang dengan benang 100 m membentuk sudut $30^\circ$. Berapa tingginya?
  \item Buktikan $\dfrac{\sin\alpha}{\cos\alpha}=\tan\alpha$ dari definisi.
  \item Hitung $\sin^2 30^\circ+\cos^2 30^\circ$.
  \item Sudut elevasi matahari $60^\circ$, bayangan tiang 3 m. Tinggi tiang?
  \item Jika $\tan\alpha=1$, tentukan $\alpha$ (lancip) dan $\sin\alpha$.
\end{enumerate}
\end{latihan}

\section{Sudut di Berbagai Kuadran dan Identitas}

\subsection*{Memperluas ke lingkaran satuan}
Sudut tidak harus lancip. Dengan lingkaran satuan, $\cos\alpha$ dan $\sin\alpha$ menjadi koordinat-$x$ dan $y$ titik pada lingkaran berjari-jari 1. Tanda sin/cos pun bergantung kuadran --- inilah yang membuat fungsi trigonometri bersifat \emph{periodik}, menjelaskan mengapa ia memodelkan gelombang.

\begin{center}
\begin{tikzpicture}[scale=1.4]
  \draw[->,gray](-1.4,0)--(1.4,0) node[right]{$x=\cos\alpha$};
  \draw[->,gray](0,-1.4)--(0,1.4) node[above]{$y=\sin\alpha$};
  \draw[biru,very thick] (0,0) circle (1);
  \coordinate (O) at (0,0);\coordinate (B) at (1,0);\coordinate (P) at (0.7,0.714);
  \draw[->,oranye,very thick] (O)--(P);
  \pic[draw,angle radius=0.45cm,"$\alpha$"]{angle=B--O--P};
  \draw[dashed] (0.7,0)--(0.7,0.714);
  \node[below right] at (0.7,0) {\small$\cos\alpha$};
  \node[above] at (0.35,0.36) {\small$1$};
\end{tikzpicture}\\
\small \textbf{Lingkaran satuan.} $(\cos\alpha,\sin\alpha)$ adalah koordinat titik; identitas $\sin^2\alpha+\cos^2\alpha=1$ hanyalah Pythagoras pada jari-jari 1.
\end{center}

\begin{konsep}{Tanda per Kuadran \&\ Identitas Dasar}
\begin{itemize}[leftmargin=*,topsep=2pt]
  \item Kuadran I: semua positif; II: sin$+$; III: tan$+$; IV: cos$+$ (\emph{All-Sin-Tan-Cos}).
  \item Identitas pokok: $\sin^2\alpha+\cos^2\alpha=1$, $\ 1+\tan^2\alpha=\sec^2\alpha$.
  \item Sudut berelasi: $\sin(90^\circ-\alpha)=\cos\alpha$, $\ \cos(90^\circ-\alpha)=\sin\alpha$.
\end{itemize}
\end{konsep}

\begin{center}
\begin{tikzpicture}[scale=1.5]
  % Warna per kuadran
  \fill[birumuda!70](0,0)--(1.3,0) arc(0:90:1.3)--cycle;
  \fill[hijaumuda!70](0,0)--(0,1.3) arc(90:180:1.3)--cycle;
  \fill[oranyemuda!80](0,0)--(-1.3,0) arc(180:270:1.3)--cycle;
  \fill[ungumuda!70](0,0)--(0,-1.3) arc(270:360:1.3)--cycle;
  \draw[biru,thick](0,0) circle(1.3);
  \draw[->,gray,thick](-1.5,0)--(1.5,0) node[right,font=\footnotesize]{$x$};
  \draw[->,gray,thick](0,-1.5)--(0,1.5) node[above,font=\footnotesize]{$y$};
  % Label kuadran (fill=white supaya terbaca jelas di atas background warna)
  \node[align=center,font=\small\bfseries,biru,fill=white,inner sep=2pt,rounded corners=2pt] at (0.6,0.6) {I\\Semua $+$};
  \node[align=center,font=\small\bfseries,hijau,fill=white,inner sep=2pt,rounded corners=2pt] at (-0.6,0.6) {II\\\textbf{Sin}\,$+$};
  \node[align=center,font=\small\bfseries,oranye,fill=white,inner sep=2pt,rounded corners=2pt] at (-0.6,-0.6) {III\\\textbf{Tan}\,$+$};
  \node[align=center,font=\small\bfseries,ungu,fill=white,inner sep=2pt,rounded corners=2pt] at (0.6,-0.6) {IV\\\textbf{Cos}\,$+$};
  \node[font=\footnotesize\bfseries] at (0,-1.85) {\textbf{A}ll\;---\;\textbf{S}in\;---\;\textbf{T}an\;---\;\textbf{C}os};
\end{tikzpicture}\\[4pt]
\small\textbf{Diagram ASTC.} Searah jarum jam dari Kuadran I: \textbf{A}ll (semua positif), \textbf{S}in, \textbf{T}an, \textbf{C}os. Di tiap kuadran hanya fungsi-fungsi tersebut (beserta kebalikannya csc/sec/cot) yang bernilai positif.
\end{center}

\begin{center}
\begin{tikzpicture}[scale=1.5]
  \draw[->,gray](-1.5,0)--(1.5,0) node[right,font=\footnotesize]{$x$};
  \draw[->,gray](0,-1.5)--(0,1.5) node[above,font=\footnotesize]{$y$};
  \draw[gray!40,thick](0,0) circle(1.2);
  \coordinate (O) at (0,0);
  % Empat titik simetris (sudut acuan α ≈ 35°)
  \coordinate (P1) at (0.98,0.69);
  \coordinate (P2) at (-0.98,0.69);
  \coordinate (P3) at (-0.98,-0.69);
  \coordinate (P4) at (0.98,-0.69);
  % Panah dari O
  \draw[->,oranye,very thick](O)--(P1);
  \draw[->,hijau,very thick](O)--(P2);
  \draw[->,merah,very thick](O)--(P3);
  \draw[->,ungu,very thick](O)--(P4);
  \fill[oranye](P1) circle(2pt); \fill[hijau](P2) circle(2pt);
  \fill[merah](P3) circle(2pt); \fill[ungu](P4) circle(2pt);
  % Label koordinat
  \node[above right,font=\scriptsize,oranye] at (P1) {$(\cos\alpha,\sin\alpha)$};
  \node[above left,font=\scriptsize,hijau] at (P2) {$(-\!\cos\alpha,\sin\alpha)$};
  \node[below left,font=\scriptsize,merah] at (P3) {$(-\!\cos\alpha,-\!\sin\alpha)$};
  \node[below right,font=\scriptsize,ungu] at (P4) {$(\cos\alpha,-\!\sin\alpha)$};
  % Label sudut (digeser menjauhi sumbu supaya tidak nabrak garis)
  \node[oranye,font=\scriptsize,fill=white,inner sep=1pt] at (0.42,0.26) {$\alpha$};
  \node[hijau,font=\scriptsize,fill=white,inner sep=1pt] at (-0.72,0.28) {$180^\circ\!-\!\alpha$};
  \node[merah,font=\scriptsize,fill=white,inner sep=1pt] at (-0.75,-0.32) {$180^\circ\!+\!\alpha$};
  \node[ungu,font=\scriptsize,fill=white,inner sep=1pt] at (0.72,-0.32) {$360^\circ\!-\!\alpha$};
\end{tikzpicture}\\[4pt]
\small\textbf{Sudut berelasi.} Empat sudut dengan \emph{sudut acuan} $\alpha$ yang sama memiliki nilai $|\sin|$ dan $|\cos|$ identik. Yang berbeda hanya tanda --- ditentukan oleh ASTC. Contoh: $\sin150^\circ=\sin(180^\circ-30^\circ)=+\sin30^\circ=\tfrac12$.
\end{center}

\begin{contoh}{Sudut Tumpul}
Tentukan $\sin 150^\circ$ dan $\cos 150^\circ$.

\jawab $150^\circ$ di kuadran II (sin$+$, cos$-$). Sudut relasinya $180^\circ-150^\circ=30^\circ$. Maka $\sin150^\circ=\sin30^\circ=\tfrac12$ dan $\cos150^\circ=-\cos30^\circ=-\tfrac{\sqrt3}{2}$.
\end{contoh}

\begin{latihan}
\begin{enumerate}[leftmargin=*]
  \item Tentukan $\sin120^\circ$, $\cos120^\circ$, $\tan120^\circ$.
  \item Hitung $\cos180^\circ+\sin90^\circ$.
  \item Jika $\sin\alpha=\tfrac45$ di kuadran II, tentukan $\cos\alpha$.
  \item Sederhanakan $\sin(90^\circ-\alpha)\cdot\sec\alpha$.
  \item Buktikan $1+\tan^2\alpha=\sec^2\alpha$ dari identitas pokok.
  \item Tentukan $\tan 225^\circ$.
  \item Hitung $\sin^2 60^\circ+\sin^2 30^\circ$.
  \item Tentukan kuadran $\alpha$ jika $\sin\alpha>0$ dan $\cos\alpha<0$.
  \item Sederhanakan $\dfrac{1-\cos^2\alpha}{\sin\alpha}$.
  \item Hitung $\cos 300^\circ$.
\end{enumerate}
\end{latihan}

\section{Trigonometri Lanjutan Dasar (Pengayaan)}

\begin{catatan}
\textbf{Kaitan dengan MAF (Kelas 11).} Seluruh rumus pada bagian ini --- jumlah/selisih sudut, sudut rangkap, sudut setengah, persamaan trigonometri --- dibahas \emph{lebih lengkap} dan disertai visualisasi penuh di \textbf{Bab Trigonometri Lanjutan (MAF Kelas 11)}. Bagian ini bersifat \emph{pengayaan opsional} sebagai pratinjau: pelajari jika sudah menguasai perbandingan dasar dan ingin mengeksplorasi lebih jauh. Latihan soal di akhir bagian ini tetap berguna untuk melatih keterampilan awal.
\end{catatan}

\subsection*{Mengapa perlu rumus sudut majemuk?}
Kita sudah mengetahui nilai sudut istimewa seperti $30^\circ$ dan $45^\circ$. Tetapi bagaimana dengan $75^\circ=45^\circ+30^\circ$ atau $15^\circ=45^\circ-30^\circ$? Rumus \emph{jumlah dan selisih sudut} memungkinkan kita menghitung nilai trigonometri sudut "campuran" tanpa kalkulator, dan menjadi pondasi identitas yang lebih lanjut.

\begin{konsep}{Rumus Jumlah dan Selisih Sudut}
\[
\begin{aligned}
\sin(A\pm B)&=\sin A\cos B\pm\cos A\sin B,\\
\cos(A\pm B)&=\cos A\cos B\mp\sin A\sin B,\\
\tan(A\pm B)&=\frac{\tan A\pm\tan B}{1\mp\tan A\tan B}.
\end{aligned}
\]
Perhatikan pola tanda: pada $\cos$, tanda \emph{berlawanan} dengan ruas kiri.
\end{konsep}

\begin{contoh}{Nilai Sudut Campuran}
Hitung $\sin 75^\circ$ tanpa kalkulator.

\jawab Tulis $75^\circ=45^\circ+30^\circ$:
\[
\sin75^\circ=\sin45^\circ\cos30^\circ+\cos45^\circ\sin30^\circ
=\frac{\sqrt2}{2}\cdot\frac{\sqrt3}{2}+\frac{\sqrt2}{2}\cdot\frac12=\frac{\sqrt6+\sqrt2}{4}.
\]
\end{contoh}

\begin{contoh}{Selisih Sudut}
Hitung $\cos15^\circ$.

\jawab $15^\circ=45^\circ-30^\circ$, maka
\[
\cos15^\circ=\cos45^\circ\cos30^\circ+\sin45^\circ\sin30^\circ
=\frac{\sqrt2}{2}\cdot\frac{\sqrt3}{2}+\frac{\sqrt2}{2}\cdot\frac12=\frac{\sqrt6+\sqrt2}{4}.
\]
\end{contoh}

\subsection*{Sudut rangkap}
Jika pada rumus jumlah kita ambil $B=A$, lahirlah rumus \emph{sudut rangkap} --- alat yang sangat sering muncul pada penyederhanaan dan persamaan.

\begin{konsep}{Rumus Sudut Rangkap}
\[
\sin2A=2\sin A\cos A,\qquad
\cos2A=\cos^2 A-\sin^2 A=1-2\sin^2 A=2\cos^2 A-1,
\]
\[
\tan2A=\frac{2\tan A}{1-\tan^2 A}.
\]
\end{konsep}

\begin{contoh}{Menggunakan Sudut Rangkap}
Diketahui $\sin A=\tfrac35$ dengan $A$ lancip. Tentukan $\sin2A$ dan $\cos2A$.

\jawab $\cos A=\tfrac45$ (segitiga $3$-$4$-$5$). Maka $\sin2A=2\cdot\tfrac35\cdot\tfrac45=\tfrac{24}{25}$ dan $\cos2A=1-2\sin^2A=1-2\cdot\tfrac{9}{25}=\tfrac{7}{25}$.
\end{contoh}

\subsection*{Sudut setengah}
Dari $\cos2A=1-2\sin^2A$ dan $\cos2A=2\cos^2A-1$, dengan mengganti $A\to\tfrac{\alpha}{2}$ diperoleh rumus \emph{sudut setengah}.

\begin{konsep}{Rumus Sudut Setengah}
\[
\sin\frac{\alpha}{2}=\pm\sqrt{\frac{1-\cos\alpha}{2}},\qquad
\cos\frac{\alpha}{2}=\pm\sqrt{\frac{1+\cos\alpha}{2}},\qquad
\tan\frac{\alpha}{2}=\frac{1-\cos\alpha}{\sin\alpha}.
\]
Tanda $\pm$ dipilih sesuai kuadran $\tfrac{\alpha}{2}$.
\end{konsep}

\begin{contoh}{Nilai Sudut Setengah}
Hitung $\cos22{,}5^\circ$.

\jawab $22{,}5^\circ=\tfrac12(45^\circ)$ dan $22{,}5^\circ$ di kuadran I (positif):
\[
\cos22{,}5^\circ=\sqrt{\frac{1+\cos45^\circ}{2}}=\sqrt{\frac{1+\frac{\sqrt2}{2}}{2}}=\frac12\sqrt{2+\sqrt2}.
\]
\end{contoh}

\subsection*{Persamaan trigonometri dasar}
Persamaan trigonometri mencari semua sudut yang memenuhi suatu nilai. Karena fungsi trigonometri \emph{periodik}, solusinya tak hingga banyak --- kita tuliskan dengan pola umum.

\begin{konsep}{Penyelesaian Persamaan Dasar}
Untuk $0^\circ\le x<360^\circ$ (satu putaran):
\[
\sin x=\sin\alpha \Rightarrow x=\alpha \ \text{atau}\ x=180^\circ-\alpha;
\]
\[
\cos x=\cos\alpha \Rightarrow x=\alpha \ \text{atau}\ x=360^\circ-\alpha;
\]
\[
\tan x=\tan\alpha \Rightarrow x=\alpha \ \text{atau}\ x=\alpha+180^\circ.
\]
Untuk semua solusi, tambahkan kelipatan periode ($360^\circ$ untuk sin/cos, $180^\circ$ untuk tan).
\end{konsep}

\begin{contoh}{Menyelesaikan Persamaan Trigonometri}
Tentukan semua $x\in[0^\circ,360^\circ)$ yang memenuhi $\sin x=\tfrac12$.

\jawab Sudut acuan $\alpha=30^\circ$ (karena $\sin30^\circ=\tfrac12$). Maka $x=30^\circ$ atau $x=180^\circ-30^\circ=150^\circ$. Jadi $x\in\{30^\circ,150^\circ\}$.
\end{contoh}

\begin{contoh}{Persamaan dengan Sudut Rangkap (HOTS)}
Selesaikan $\cos2x=\sin x$ pada $0^\circ\le x<360^\circ$.

\jawab Gunakan $\cos2x=1-2\sin^2x$. Misalkan $s=\sin x$:
\[
1-2s^2=s \Rightarrow 2s^2+s-1=0 \Rightarrow (2s-1)(s+1)=0.
\]
Jadi $s=\tfrac12$ atau $s=-1$. Dari $\sin x=\tfrac12$: $x=30^\circ,150^\circ$; dari $\sin x=-1$: $x=270^\circ$. Solusi: $\{30^\circ,150^\circ,270^\circ\}$.
\end{contoh}

\begin{latihan}
\begin{enumerate}[leftmargin=*]
  \item Hitung $\cos75^\circ$ dengan rumus jumlah sudut.
  \item Hitung $\tan75^\circ$ dengan rumus jumlah sudut.
  \item Diketahui $\cos A=\tfrac{12}{13}$ ($A$ lancip). Tentukan $\sin2A$.
  \item Buktikan $\sin2A=2\sin A\cos A$ dari rumus jumlah sudut.
  \item Hitung $\sin22{,}5^\circ$ dengan rumus sudut setengah.
  \item Sederhanakan $2\sin15^\circ\cos15^\circ$ menjadi nilai tunggal.
  \item Selesaikan $\cos x=\tfrac{\sqrt3}{2}$ untuk $0^\circ\le x<360^\circ$.
  \item Selesaikan $\tan x=1$ untuk $0^\circ\le x<360^\circ$.
  \item Selesaikan $2\sin x-1=0$ untuk $0^\circ\le x<360^\circ$.
  \item Selesaikan $\sin2x=\cos x$ untuk $0^\circ\le x<360^\circ$ (petunjuk: $\sin2x=2\sin x\cos x$).
\end{enumerate}
\end{latihan}

\begin{kesalahan}
\begin{enumerate}[leftmargin=*,topsep=2pt]
  \item \textbf{Tertukar sisi depan dan samping.} "Depan" adalah sisi di \emph{seberang} sudut, "samping" yang mengapit. Selalu tandai sudut acuan lebih dulu.
  \item \textbf{Menulis $\sin^2\alpha$ sebagai $\sin\alpha^2$.} $\sin^2\alpha=(\sin\alpha)^2$, bukan $\sin(\alpha^2)$.
  \item \textbf{Salah tanda di kuadran.} Lupa bahwa cos negatif di kuadran II/III. Gunakan aturan All-Sin-Tan-Cos.
  \item \textbf{Memakai $\sin$ padahal seharusnya $\tan$.} Pilih perbandingan berdasarkan sisi yang \emph{diketahui} dan \emph{dicari}, bukan kebiasaan.
  \item \textbf{Mendistribusikan fungsi ke penjumlahan.} $\sin(A+B)\neq\sin A+\sin B$. Cek: $\sin(30^\circ+60^\circ)=\sin90^\circ=1$, tetapi $\sin30^\circ+\sin60^\circ=\tfrac12+\tfrac{\sqrt3}{2}\approx1{,}37\neq 1$. Fungsi trigonometri \emph{bukan} fungsi linear; ia tidak "tembus" tanda kurung penjumlahan.
\end{enumerate}
\end{kesalahan}

\begin{dunianyata}
\begin{itemize}[leftmargin=*,topsep=2pt]
  \item \textbf{Konstruksi \&\ survei:} mengukur tinggi gedung, kemiringan atap, dan lereng jalan.
  \item \textbf{GPS \&\ navigasi:} triangulasi posisi dari sudut dan jarak.
  \item \textbf{Teknologi sinyal:} gelombang suara, cahaya, dan listrik AC dimodelkan dengan $\sin$/$\cos$.
  \item \textbf{Game \&\ animasi:} rotasi objek, gerak melingkar, dan osilasi memakai trigonometri.
  \item \textbf{Kesehatan:} sinyal EKG dan gelombang otak dianalisis lewat komponen sinusoidal.
\end{itemize}
\end{dunianyata}

\begin{tcolorbox}[colback=oranyemuda,colframe=oranye,title=\textbf{Bahasa Alternatif yang Sering Mengecoh di Soal}]
\begin{itemize}[leftmargin=*,topsep=2pt]
  \item \textbf{"Sudut elevasi" vs "sudut depresi":} elevasi = melihat ke \emph{atas}, depresi = melihat ke \emph{bawah}. Keduanya diukur dari \emph{garis horizontal}, bukan dari garis vertikal.
  \item \textbf{"Sisi depan" = "sisi hadapan":} keduanya merujuk sisi di seberang sudut $\alpha$ --- jangan tertukar dengan sisi miring.
  \item \textbf{$\sin\alpha$ negatif $\ne$ sudut negatif:} sudut $210^\circ$ adalah sudut positif, namun $\sin210^\circ<0$ karena berada di Kuadran III (only Tan$+$).
  \item \textbf{$\sin^{-1}$ (arcsin) $\ne$ $\dfrac{1}{\sin}$ (csc):} pangkat $-1$ pada fungsi trigonometri berarti \emph{fungsi invers}, bukan kebalikan nilai.
  \item \textbf{"Perbandingan trigonometri" adalah rasio:} $\sin30^\circ=\tfrac12$ bukan berarti ada sisi sependek $\tfrac12$; itu rasio \emph{depan:miring} yang tidak bersatuan.
  \item \textbf{Sudut $90^\circ$ bukan berarti $\tan$ ada:} $\tan90^\circ$ tidak terdefinisi (penyebut $\cos90^\circ=0$). Jangan asal substitusi.
\end{itemize}
\end{tcolorbox}

\section*{Ringkasan Bab}
\addcontentsline{toc}{section}{Ringkasan Bab}
\begin{center}
\begin{tabular}{>{\bfseries}p{3.5cm} p{10.2cm}}
\toprule
Konsep & Inti yang harus diingat \\
\midrule
Perbandingan & $\sin=\tfrac{de}{mi}$, $\cos=\tfrac{sa}{mi}$, $\tan=\tfrac{de}{sa}$ (rasio, bergantung sudut saja). \\
Sudut istimewa & $0,30,45,60,90^\circ$; berasal dari segitiga $1$-$1$-$\sqrt2$ dan $1$-$\sqrt3$-$2$. \\
Lingkaran satuan & $(\cos\alpha,\sin\alpha)$ koordinat titik; sumber periodisitas. \\
Tanda kuadran & All-Sin-Tan-Cos. \\
Identitas & $\sin^2+\cos^2=1$; $1+\tan^2=\sec^2$. \\
Aplikasi & pengukuran tinggi/jarak tak langsung, gelombang. \\
\bottomrule
\end{tabular}
\end{center}

\begin{refleksi}
Beri tanda centang pada hal yang sudah kamu kuasai:
\begin{itemize}[leftmargin=*,label={},topsep=2pt]
  \item \cek Saya memahami trigonometri sebagai rasio sisi yang bergantung sudut (kesebangunan).
  \item \cek Saya dapat menurunkan nilai sudut istimewa dari segitiga dasar.
  \item \cek Saya dapat menentukan tanda nilai trigonometri tiap kuadran.
  \item \cek Saya dapat menggunakan trigonometri untuk masalah tinggi/jarak.
\end{itemize}
\end{refleksi}

% ---------- SOAL AKHIR BAB ----------
\section*{Soal Akhir Bab 4}
\addcontentsline{toc}{section}{Soal Akhir Bab 4}

\bagiansoal{A}{Fundamental (setara SNBT Dasar)}
\begin{enumerate}[leftmargin=*]
  \item Hitung $\sin30^\circ+\cos60^\circ$.
  \item Diketahui depan $=5$, miring $=13$. Tentukan $\cos\alpha$.
  \item Hitung $\tan45^\circ$.
  \item Tentukan $\sin90^\circ-\cos0^\circ$.
  \item Jika $\tan\alpha=\tfrac34$ (lancip), tentukan $\sin\alpha$.
  \item Hitung $\sin^2 45^\circ+\cos^2 45^\circ$.
  \item Tinggi tiang 6 m, bayangan 6 m. Berapa sudut elevasi matahari?
  \item Hitung $\cos60^\circ\cdot\tan45^\circ$.
  \item Tentukan $\sin150^\circ$.
  \item Sederhanakan $\dfrac{\sin\alpha}{\tan\alpha}$.
\end{enumerate}

\bagiansoal{B}{Intermediate (setara SNBT Sulit)}
\begin{enumerate}[leftmargin=*]
  \item Diketahui $\sin\alpha+\cos\alpha=\tfrac75$. Tentukan $\sin\alpha\cos\alpha$.
  \item Sederhanakan $\dfrac{1}{1+\cos\alpha}+\dfrac{1}{1-\cos\alpha}$ menjadi bentuk $\sin$.
  \item Dari atas mercusuar tinggi 40 m, sudut depresi dua kapal $30^\circ$ dan $45^\circ$ segaris. Tentukan jarak antar kapal.
  \item Jika $\sin\alpha=\tfrac{2}{3}$ dan $\alpha$ tumpul, tentukan $\tan\alpha$.
  \item Buktikan $\dfrac{\cos\alpha}{1-\sin\alpha}=\dfrac{1+\sin\alpha}{\cos\alpha}$.
\end{enumerate}

\bagiansoal{C}{Advanced (setara Olimpiade SMA --- gabungan dengan Vektor)}
\begin{enumerate}[leftmargin=*]
  \item Gunakan dot product untuk membuktikan rumus $\cos(A-B)=\cos A\cos B+\sin A\sin B$ (sudut dua vektor satuan).
  \item Sebuah benda meluncur pada bidang miring $30^\circ$ ditarik gaya 50 N. Uraikan gaya menjadi komponen sejajar dan tegak lurus bidang (vektor + trigonometri).
  \item Buktikan identitas $\sin3\alpha=3\sin\alpha-4\sin^3\alpha$.
  \item Pada segitiga $ABC$, buktikan aturan sinus $\dfrac{a}{\sin A}=\dfrac{b}{\sin B}=\dfrac{c}{\sin C}$.
  \item Vektor $\vec a$ membentuk sudut $\alpha$ dengan sumbu-$x$. Nyatakan $\vec a$ dalam $|\vec a|$, $\cos\alpha$, $\sin\alpha$, lalu tentukan sudut jika $\vec a=(-1,\sqrt3)$.
\end{enumerate}

\bagiansoal{D}{Expert (setara tahun pertama universitas --- sintesis antarbab)}
\begin{enumerate}[leftmargin=*]
  \item \textbf{(Trigonometri + Fungsi Kuadrat).} Tentukan nilai maksimum dan minimum $f(\alpha)=3\sin\alpha+4\cos\alpha$ dengan menuliskannya sebagai $R\sin(\alpha+\varphi)$.
  \item \textbf{(Trigonometri + Vektor).} Buktikan rumus luas segitiga $L=\tfrac12 ab\sin C$ dan kaitkan dengan besar cross product dua vektor sisi.
  \item \textbf{(Trigonometri + Deret).} Tunjukkan bahwa untuk sudut kecil (dalam radian), $\sin\theta\approx\theta$. Gunakan untuk menaksir $\sin 1^\circ$.
  \item \textbf{(Trigonometri + Identitas).} Tentukan semua $x\in[0^\circ,360^\circ)$ yang memenuhi $\sin^4x+\cos^4x=\tfrac{7}{8}$, lalu hitung jumlah semua solusinya. (Petunjuk: $\sin^4x+\cos^4x=(\sin^2x+\cos^2x)^2-2\sin^2x\cos^2x$.)
  \item \textbf{(Pemodelan).} Tinggi air pasang dimodelkan $h(t)=2+1{,}5\sin\!\left(\tfrac{\pi}{6}t\right)$ meter. Tentukan tinggi maksimum, periode, dan kapan pertama kali $h=3{,}5$ m.
\end{enumerate}

\begin{challenge}
\begin{enumerate}[leftmargin=*,topsep=2pt]
  \item Tanpa kalkulator, buktikan $\tan 1^\circ\tan2^\circ\tan3^\circ\cdots\tan89^\circ=1$.
  \item Hitung $\cos36^\circ$ dalam bentuk akar (kaitannya dengan rasio emas).
  \item Sebuah tangga harus melewati sudut koridor berbentuk L (lebar $a$ dan $b$). Tentukan panjang maksimum tangga sebagai masalah optimasi bersudut $\theta$.
\end{enumerate}
\end{challenge}

% ============================================================
\chapter{Sistem Persamaan dan Pertidaksamaan Linear}
% ============================================================
\begin{tujuan}
Setelah mempelajari bab ini, peserta didik mampu:
\begin{itemize}[leftmargin=*,topsep=2pt]
  \item memodelkan masalah nyata menjadi sistem persamaan linear;
  \item menyelesaikan SPLTV dengan eliminasi, substitusi, dan gabungan;
  \item menafsirkan banyaknya penyelesaian (tunggal, tak hingga, atau tak ada) secara geometris;
  \item menggambar daerah penyelesaian sistem pertidaksamaan linear dua variabel;
  \item menentukan nilai optimum fungsi objektif (pengantar program linear).
\end{itemize}
\textbf{Capaian Pembelajaran (Fase E).} Peserta didik dapat menyelesaikan sistem persamaan dan pertidaksamaan linear serta menafsirkan penyelesaiannya dalam konteks.
\end{tujuan}

\begin{pemantik}
\begin{itemize}[leftmargin=*,topsep=2pt]
  \item Di kantin, 2 nasi + 1 es = Rp18.000 dan 1 nasi + 2 es = Rp15.000. Bisakah kamu menemukan harga satuannya tanpa menebak?
  \item Mengapa untuk mencari 3 hal yang belum diketahui, kita biasanya butuh 3 informasi (persamaan)?
  \item Jika sebuah pabrik ingin untung maksimum dengan bahan baku terbatas, bagaimana matematika membantu memutuskan jumlah produksi?
\end{itemize}
\end{pemantik}

\begin{studikasus}{Menentukan Harga di Kantin}
Tiga teman jajan di kantin dengan catatan berikut:
\begin{center}
\begin{tabular}{c|ccc|c}
\toprule
& Nasi ($x$) & Ayam ($y$) & Es ($z$) & Total (Rp) \\
\midrule
Andi & 2 & 1 & 1 & 32.000 \\
Budi & 1 & 2 & 1 & 34.000 \\
Cici & 1 & 1 & 2 & 30.000 \\
\bottomrule
\end{tabular}
\end{center}
\textbf{Amati dan duga.} (a) Bisakah harga satuan ditemukan hanya dari satu baris? Mengapa perlu ketiganya? (b) Bagaimana menuliskan situasi ini secara matematis? (c) Strategi apa untuk "menghilangkan" satu variabel? Situasi ini adalah \textbf{SPLTV} --- tiga persamaan, tiga variabel. Kuncinya adalah \emph{eliminasi}: menggabungkan persamaan agar satu variabel lenyap.
\end{studikasus}

\section*{Peta Konsep Bab}
\addcontentsline{toc}{section}{Peta Konsep Bab}
\begin{center}
\begin{tikzpicture}[
  konsepbox/.style={draw=biru,fill=birumuda,rounded corners,align=center,inner sep=5pt,font=\small\bfseries},
  subbox/.style={draw=hijau,fill=hijaumuda,rounded corners,align=center,inner sep=4pt,font=\footnotesize},
  garis/.style={-Stealth,biru,thick}]
  \node[konsepbox] (akar) {SISTEM LINEAR};
  \node[subbox,below left=12mm and 12mm of akar] (eq) {SPLTV\\(persamaan)};
  \node[subbox,below right=12mm and 12mm of akar] (in) {SPtLDV\\(pertidaksamaan)};
  \node[subbox,below=10mm of eq] (met) {Eliminasi,\\substitusi};
  \node[subbox,below=10mm of in] (dp) {Daerah\\penyelesaian};
  \node[subbox,below=10mm of dp] (pl) {Nilai optimum\\(program linear)};
  \draw[garis] (akar)--(eq); \draw[garis] (akar)--(in);
  \draw[garis] (eq)--(met); \draw[garis] (in)--(dp); \draw[garis] (dp)--(pl);
\end{tikzpicture}\\
\small \textbf{Peta konsep.} Persamaan mencari titik temu; pertidaksamaan mencari daerah; program linear memilih titik terbaik dalam daerah itu.
\end{center}

\section{Sistem Persamaan Linear Tiga Variabel (SPLTV)}

\subsection*{Membangun konsep}
Satu persamaan dengan tiga variabel punya tak hingga penyelesaian (sebuah bidang di ruang). Untuk "memenjarakan" satu titik, kita butuh tiga bidang yang berpotongan di satu titik --- itulah mengapa diperlukan tiga persamaan. Strategi eliminasi sesungguhnya "menjatuhkan" masalah 3 variabel menjadi 2 variabel, lalu 1.

\begin{konsep}{Bentuk Umum}
\[
\begin{cases}
a_1 x + b_1 y + c_1 z = d_1\\
a_2 x + b_2 y + c_2 z = d_2\\
a_3 x + b_3 y + c_3 z = d_3
\end{cases}
\]
Diselesaikan dengan metode \emph{eliminasi}, \emph{substitusi}, atau gabungan. Secara geometris tiap persamaan adalah \emph{bidang} di ruang; penyelesaian adalah titik potong ketiganya.
\end{konsep}

\begin{center}
\begin{tikzpicture}[scale=0.88]
% ---- Panel 1: Satu Solusi ----
\begin{scope}[xshift=0cm]
  \fill[abu] (-1.2,-1.2) rectangle (1.4,1.4);
  \draw[gray!40,thin] (-1.2,-1.2) rectangle (1.4,1.4);
  \draw[->,gray!60,thin](-1.3,0)--(1.5,0);
  \draw[->,gray!60,thin](0,-1.3)--(0,1.5);
  \draw[biru,thick](-1.1,0.9)--(1.3,-0.5);
  \draw[oranye,thick](-1.1,-0.2)--(1.3,0.9);
  \draw[hijau,thick](-0.4,-1.2)--(0.6,1.2);
  \fill[ungu](0.17,0.33) circle(2.2pt);
  \node[below,font=\small,biru,align=center] at (0.1,-1.45) {\textbf{Satu solusi}\\tiga bidang bertemu\\di satu titik};
\end{scope}
% ---- Panel 2: Tidak Ada Solusi ----
\begin{scope}[xshift=4cm]
  \fill[abu] (-1.2,-1.2) rectangle (1.4,1.4);
  \draw[gray!40,thin] (-1.2,-1.2) rectangle (1.4,1.4);
  \draw[->,gray!60,thin](-1.3,0)--(1.5,0);
  \draw[->,gray!60,thin](0,-1.3)--(0,1.5);
  \draw[biru,thick](-1.1,0.5)--(1.3,0.9);
  \draw[oranye,thick](-1.1,-0.3)--(1.3,0.1);
  \draw[hijau,thick](-0.4,-1.1)--(0.8,1.1);
  \node[below,font=\small,biru,align=center] at (0.1,-1.45) {\textbf{Tidak ada solusi}\\ada bidang sejajar\\tak ada titik potong};
\end{scope}
% ---- Panel 3: Tak Hingga Solusi ----
\begin{scope}[xshift=8cm]
  \fill[abu] (-1.2,-1.2) rectangle (1.4,1.4);
  \draw[gray!40,thin] (-1.2,-1.2) rectangle (1.4,1.4);
  \draw[->,gray!60,thin](-1.3,0)--(1.5,0);
  \draw[->,gray!60,thin](0,-1.3)--(0,1.5);
  \draw[biru,ultra thick](-1.1,0.3)--(1.3,0.7);
  \draw[oranye,thick,dashed](-1.1,0.3)--(1.3,0.7);
  \draw[hijau,thick,dotted,line width=1.4pt](-1.1,0.3)--(1.3,0.7);
  \node[below,font=\small,biru,align=center] at (0.1,-1.45) {\textbf{Tak hingga solusi}\\bidang berimpit\\atau berpotongan di garis};
\end{scope}
\end{tikzpicture}\\[4pt]
\small\textbf{Tiga kemungkinan penyelesaian SPLTV (analogi 2D: garis $\approx$ bidang).} Tiap persamaan mendefinisikan sebuah \emph{bidang} di ruang 3D. Apakah ketiga bidang bertemu di satu titik, sejajar, atau berimpit menentukan apakah sistemnya punya satu, nol, atau tak hingga solusi.
\end{center}

\begin{contoh}{Penyelesaian SPLTV}
Selesaikan $\begin{cases} x+y+z=6\\ x-y+z=2\\ x+y-z=0 \end{cases}$.

\jawab Jumlahkan (1) dan (3): $2x+2y = 6 \Rightarrow x+y=3$.
Dari (1): $z = 6-(x+y) = 3$. Kurangkan (1)$-$(2): $2y = 4 \Rightarrow y=2$, sehingga $x=1$.
Jadi $(x,y,z)=(1,2,3)$.
\end{contoh}

\begin{contoh}{Menyelesaikan Studi Kasus Kantin}
Selesaikan sistem harga: $2x+y+z=32$, $x+2y+z=34$, $x+y+2z=30$ (ribuan).

\jawab Jumlahkan ketiganya: $4x+4y+4z=96\Rightarrow x+y+z=24$. Kurangkan tiap persamaan dari ini: dari (1) $x=32-24=8$? Hati-hati: (1) $-$ ($x+y+z$): $(2x+y+z)-(x+y+z)=x=32-24=8$. Demikian $y=34-24=10$, $z=30-24=6$. Jadi nasi Rp8.000, ayam Rp10.000, es Rp6.000.
\end{contoh}

\begin{contoh}{Banyak Penyelesaian (HOTS)}
Tentukan apakah sistem $\begin{cases} x+y=2\\ 2x+2y=5\end{cases}$ punya penyelesaian.

\jawab Persamaan kedua $2x+2y=5$ berarti $x+y=2{,}5$, bertentangan dengan $x+y=2$. Kedua garis \emph{sejajar} dan tak berpotongan --- \textbf{tidak ada} penyelesaian. Ini mengajarkan bahwa banyaknya penyelesaian punya makna geometris.
\end{contoh}

\begin{latihan}
\begin{enumerate}[leftmargin=*]
  \item Selesaikan $\begin{cases} 2x+y-z=3\\ x-y+z=2\\ x+y+z=6\end{cases}$.
  \item Selesaikan $\begin{cases} x+y=5\\ y+z=7\\ x+z=6\end{cases}$.
  \item Tiga bilangan berjumlah 30; yang terbesar 2 kali terkecil; selisih terbesar dan tengah 4. Tentukan ketiganya.
  \item Selesaikan $\begin{cases} x+2y+z=8\\ 2x-y+z=3\\ x+y-z=0\end{cases}$.
  \item Tentukan apakah $\begin{cases} x-y=1\\ 2x-2y=2\end{cases}$ punya satu, tak hingga, atau tanpa solusi.
  \item Harga 2 buku + 1 pena = Rp17rb; 1 buku + 3 pena = Rp16rb. Tentukan harga satuan.
  \item Selesaikan $\begin{cases} x+y+z=9\\ 2x-y+z=6\\ x+2y-z=3\end{cases}$.
  \item Jumlah tiga sudut segitiga $180^\circ$; sudut A dua kali B; C $20^\circ$ lebih besar dari A. Tentukan ketiga sudut.
  \item Selesaikan $\begin{cases} 3x+2y=12\\ x-y=1\end{cases}$ (SPLDV).
  \item Sebuah pecahan: pembilang + penyebut = 11; bila pembilang ditambah 3 nilainya 1. Tentukan pecahan.
\end{enumerate}
\end{latihan}

\section{Sistem Pertidaksamaan Linear Dua Variabel (SPtLDV)}

\subsection*{Membangun konsep}
Berbeda dari persamaan yang memberi \emph{garis}, pertidaksamaan memberi \emph{daerah} (setengah bidang). Gabungan beberapa pertidaksamaan menghasilkan irisan daerah --- "wilayah aman" yang memenuhi semua syarat. Inilah fondasi program linear yang dipakai industri untuk optimasi.

\begin{konsep}{Daerah Penyelesaian}
Pertidaksamaan linear $ax+by \le c$ membagi bidang menjadi dua. Daerah penyelesaian ditentukan dengan uji titik (biasanya titik asal $(0,0)$): jika titik memenuhi, daerah yang memuatnya adalah penyelesaian.
\end{konsep}

\begin{center}
\begin{tikzpicture}[scale=0.7]
  \fill[hijaumuda] (0,0)--(4,0)--(0,4)--cycle;
  \draw[->] (-0.5,0)--(5,0) node[right]{$x$};
  \draw[->] (0,-0.5)--(0,5) node[above]{$y$};
  \draw[very thick,hijau] (4,0)--(0,4) node[above right]{$x+y=4$};
  \node at (1,1) {\small DP};
\end{tikzpicture}\\
\small \textbf{Daerah penyelesaian} $x+y\le4$, $x\ge0$, $y\ge0$: segitiga di pojok pertama.
\end{center}

\begin{contoh}{Nilai Optimum (Program Linear)}
Maksimumkan $f=3x+4y$ pada kendala $x+y\le 4$, $x\ge0$, $y\ge0$.

\jawab Titik-titik pojok daerah: $(0,0),(4,0),(0,4)$. Uji: $f(0,0)=0$, $f(4,0)=12$, $f(0,4)=16$. Maksimum $16$ di $(0,4)$. \emph{Teorema titik pojok}: optimum selalu di salah satu titik sudut daerah.
\end{contoh}

\begin{latihan}
\begin{enumerate}[leftmargin=*]
  \item Gambarkan daerah penyelesaian dari $x+2y\le 8$, $x\ge0$, $y\ge0$.
  \item Tentukan apakah $(1,1)$ memenuhi $2x+3y\le 6$.
  \item Gambarkan $y\ge x$ dan $y\le 4$.
  \item Tuliskan sistem pertidaksamaan untuk daerah segitiga bertitik $(0,0),(3,0),(0,2)$.
  \item Maksimumkan $f=x+y$ pada $x+2y\le6$, $2x+y\le6$, $x,y\ge0$.
  \item Minimumkan $f=2x+3y$ pada $x+y\ge4$, $x,y\ge0$.
  \item Gambarkan daerah $x+y\le5$, $x\le3$, $y\ge0$, $x\ge0$.
  \item Pabrik membuat 2 produk; tuliskan kendala jika bahan A maksimum 10 dan B maksimum 8 unit.
  \item Apakah titik $(2,2)$ ada dalam daerah $x+y\le3$?
  \item Tentukan titik pojok dari $x+y\le4$, $x\ge1$, $y\ge0$.
\end{enumerate}
\end{latihan}

\section{Fungsi, Persamaan, dan Pertidaksamaan Linear (Pengayaan)}

\subsection*{Konsep fungsi linear}
Sebelum berbicara tentang \emph{sistem}, kita perlu memahami satu garis lurus secara utuh. Fungsi linear adalah model paling sederhana sekaligus paling sering dipakai untuk menggambarkan hubungan "perubahan tetap".

\begin{definisi}
\emph{Fungsi linear} adalah fungsi berbentuk $f(x)=mx+c$ dengan $m,c$ konstanta dan $m\neq0$. Grafiknya berupa \emph{garis lurus}; $m$ adalah \emph{gradien} (kemiringan) dan $c$ adalah perpotongan dengan sumbu-$y$.
\end{definisi}

\subsection*{Domain dan range}
Untuk fungsi linear $f(x)=mx+c$, setiap bilangan real boleh dimasukkan dan setiap bilangan real dapat dihasilkan, sehingga \emph{domain} $D=\mathbb{R}$ dan \emph{range} $R=\mathbb{R}$. Namun pada pemodelan nyata, domain sering dibatasi konteks (misalnya jumlah barang $x\ge0$ dan bilangan cacah).

\begin{contoh}{Menentukan Domain dan Range Kontekstual}
Biaya sewa mobil $f(x)=200000+5000x$ rupiah, dengan $x$ jumlah kilometer ($0\le x\le 100$). Tentukan domain dan range.

\jawab Domain $=\{x\mid 0\le x\le100\}$. Saat $x=0$, $f=200{.}000$; saat $x=100$, $f=700{.}000$. Maka range $=\{f\mid 200000\le f\le700000\}$.
\end{contoh}

\subsection*{Gradien}
\begin{konsep}{Gradien (Kemiringan)}
Gradien garis melalui dua titik $(x_1,y_1)$ dan $(x_2,y_2)$:
\[
m=\frac{y_2-y_1}{x_2-x_1}=\frac{\Delta y}{\Delta x}\quad(\text{"naik dibagi geser"}).
\]
$m>0$: garis naik; $m<0$: turun; $m=0$: mendatar; garis tegak (vertikal) tak punya gradien.
\end{konsep}

\begin{contoh}{Menghitung Gradien}
Tentukan gradien garis melalui $A(1,2)$ dan $B(5,10)$.

\jawab $m=\dfrac{10-2}{5-1}=\dfrac{8}{4}=2$. Tiap pertambahan $1$ satuan $x$, nilai $y$ naik $2$.
\end{contoh}

\subsection*{Persamaan garis}
\begin{konsep}{Bentuk-Bentuk Persamaan Garis}
\[
\text{Gradien-intersep: } y=mx+c;\qquad
\text{Titik-gradien: } y-y_1=m(x-x_1);
\]
\[
\text{Melalui dua titik: } \frac{y-y_1}{y_2-y_1}=\frac{x-x_1}{x_2-x_1}.
\]
\end{konsep}

\begin{contoh}{Menyusun Persamaan Garis}
Tentukan persamaan garis bergradien $3$ yang melalui $(2,1)$.

\jawab $y-1=3(x-2)\Rightarrow y=3x-5$.
\end{contoh}

\subsection*{Garis sejajar dan tegak lurus}
\begin{konsep}{Hubungan Dua Garis}
Misalkan dua garis bergradien $m_1$ dan $m_2$:
\[
\text{Sejajar} \iff m_1=m_2;\qquad
\text{Tegak lurus} \iff m_1\cdot m_2=-1.
\]
\end{konsep}

\textbf{Mengapa hasil kali gradien $=-1$ untuk tegak lurus?} Memutar arah garis $90^\circ$ menukar peran $\Delta x$ dan $\Delta y$ sekaligus membalik satu tanda, sehingga gradien menjadi kebalikan negatifnya, $m_2=-\tfrac{1}{m_1}$. (Bandingkan uji tegak lurus vektor $\vec a\cdot\vec b=0$ pada Bab 3.)

\begin{contoh}{Garis Sejajar dan Tegak Lurus}
Tentukan persamaan garis melalui $(0,4)$ yang (a) sejajar dan (b) tegak lurus terhadap $y=2x+1$.

\jawab Gradien acuan $m=2$. (a) Sejajar: $m_1=2$, garis $y=2x+4$. (b) Tegak lurus: $m_2=-\tfrac12$, garis $y=-\tfrac12 x+4$.
\end{contoh}

\subsection*{Interpretasi grafik}
Pada grafik linear, $c$ adalah \emph{nilai awal} (saat $x=0$) dan $m$ adalah \emph{laju perubahan}. Membaca kedua angka ini langsung dari masalah nyata adalah inti pemodelan: "ongkos awal" vs "tarif per satuan", "saldo awal" vs "tabungan per bulan", dan seterusnya.

\begin{contoh}{Membaca Makna $m$ dan $c$}
Sebuah taksi memasang tarif $y=7000+4000x$ (rupiah), $x$ dalam km. Apa makna $7000$ dan $4000$?

\jawab $7000$ = tarif buka pintu (saat $x=0$); $4000$ = tarif per kilometer (gradien). Untuk $5$ km, ongkos $=7000+4000(5)=27{.}000$ rupiah.
\end{contoh}

\subsection*{Persamaan dan pertidaksamaan linear satu variabel}
\begin{konsep}{Aturan Pokok}
Persamaan/pertidaksamaan linear dapat dimanipulasi dengan menambah/mengurangi atau mengali/membagi kedua ruas. \textbf{Penting:} mengali atau membagi pertidaksamaan dengan bilangan \emph{negatif} \emph{membalik} arah tanda.
\end{konsep}

\begin{contoh}{Pertidaksamaan Linear}
Selesaikan $-3x+5\le 14$.

\jawab $-3x\le 9 \Rightarrow x\ge -3$ (tanda dibalik karena dibagi $-3$). Himpunan penyelesaian $\{x\mid x\ge-3\}$.
\end{contoh}

\subsection*{Pemodelan linear}
\begin{contoh}{Memodelkan Masalah Nyata}
Sebuah usaha keripik bermodal awal Rp500.000 dan menghabiskan Rp3.000 bahan per bungkus. Jika tiap bungkus dijual Rp8.000, tentukan banyak bungkus agar tidak rugi (titik impas).

\jawab Biaya $B(x)=500000+3000x$; pendapatan $P(x)=8000x$. Impas saat $P=B$: $8000x=500000+3000x\Rightarrow 5000x=500000\Rightarrow x=100$. Untung jika $x>100$ bungkus.
\end{contoh}

\begin{latihan}
\begin{enumerate}[leftmargin=*]
  \item Tentukan gradien garis melalui $(2,3)$ dan $(6,11)$.
  \item Susun persamaan garis bergradien $-2$ melalui $(1,4)$.
  \item Tentukan persamaan garis melalui $(0,0)$ dan $(3,6)$.
  \item Apakah $y=3x+1$ dan $y=3x-4$ sejajar? Jelaskan.
  \item Tentukan gradien garis yang tegak lurus $y=\tfrac14 x+2$.
  \item Tentukan domain dan range $f(x)=2x-1$ untuk $-1\le x\le3$.
  \item Selesaikan $4x-7=2x+5$.
  \item Selesaikan pertidaksamaan $-2x+1>7$.
  \item Tarif parkir $y=2000+1000t$ ($t$ jam). Berapa biaya $3$ jam dan apa makna $2000$?
  \item Sebuah ojek online menetapkan $y=8000+2500x$. Tentukan jarak agar ongkos tepat Rp23.000.
\end{enumerate}
\end{latihan}

\section{Program Linear Lanjutan (Pengayaan)}

\subsection*{Daerah feasible (himpunan penyelesaian)}
Pada masalah nyata, kita memaksimalkan/meminimalkan sesuatu di bawah sejumlah \emph{kendala}. Irisan semua kendala membentuk \emph{daerah feasible} (daerah yang layak) --- satu-satunya wilayah tempat solusi boleh berada.

\begin{definisi}
\emph{Daerah feasible} adalah himpunan semua titik $(x,y)$ yang memenuhi seluruh kendala sekaligus. \emph{Titik pojok} (titik sudut) adalah perpotongan dua garis batas yang menjadi sudut daerah feasible.
\end{definisi}

\begin{konsep}{Teorema Titik Pojok}
Jika daerah feasible berupa daerah tertutup (terbatas), maka nilai \emph{maksimum} dan \emph{minimum} dari fungsi objektif linear $f=ax+by$ selalu tercapai di salah satu \emph{titik pojok} daerah tersebut.
\end{konsep}

\textbf{Mengapa optimum di titik pojok?} Garis $ax+by=k$ (garis dengan nilai objektif sama) bergeser sejajar saat $k$ berubah. Saat digeser melintasi daerah feasible, ia menyentuh daerah terakhir kali tepat di sebuah sudut --- itulah \emph{garis selidik}.

\subsection*{Metode garis selidik}
\begin{konsep}{Garis Selidik}
Gambar garis $ax+by=k$ untuk satu nilai $k$ sembarang, lalu geser sejajar ke arah memperbesar $f$ (untuk maksimum) atau memperkecil $f$ (untuk minimum). Titik pojok terakhir yang disentuh garis adalah letak optimum.
\end{konsep}

\subsection*{Optimasi: maksimum dan minimum}
Langkah baku menyelesaikan program linear:
\begin{enumerate}[leftmargin=*,topsep=2pt]
  \item Terjemahkan masalah menjadi variabel, fungsi objektif, dan kendala.
  \item Gambar daerah feasible dan tentukan semua titik pojoknya.
  \item Uji nilai $f$ pada setiap titik pojok.
  \item Pilih nilai terbesar (maksimum) atau terkecil (minimum).
\end{enumerate}

\begin{contoh}{Optimasi Lengkap}
Maksimumkan $f=4x+5y$ dengan kendala $x+y\le 8$, $2x+y\le 12$, $x\ge0$, $y\ge0$.

\jawab Titik pojok: $(0,0),(6,0),(0,8)$, dan perpotongan $x+y=8$ dengan $2x+y=12$ yaitu $x=4,y=4$. Uji:
\[
f(0,0)=0,\ f(6,0)=24,\ f(0,8)=40,\ f(4,4)=16+20=36.
\]
Maksimum $40$ di titik $(0,8)$.
\end{contoh}

\subsection*{Studi kasus produksi}
\begin{studikasus}{Pabrik Mebel}
Sebuah pabrik membuat meja ($x$) dan kursi ($y$). Tiap meja butuh $4$ jam tukang dan $2$ kg kayu; tiap kursi $2$ jam tukang dan $3$ kg kayu. Tersedia $40$ jam tukang dan $42$ kg kayu. Laba per meja Rp80.000 dan per kursi Rp60.000. Tentukan produksi yang memaksimalkan laba.

\smallskip
\jawab Kendala: $4x+2y\le40$ (atau $2x+y\le20$), $2x+3y\le42$, $x,y\ge0$. Fungsi objektif $f=80x+60y$ (ribuan). Titik pojok: $(0,0)$, $(10,0)$, $(0,14)$, dan perpotongan $2x+y=20$ dengan $2x+3y=42$ yaitu $x=4{,}5,\,y=11$. Uji (ribuan):
\[
f(0,0)=0,\ f(10,0)=800,\ f(0,14)=840,\ f(4{,}5,11)=360+660=1020.
\]
Laba maksimum Rp1.020.000 dengan memproduksi $4{,}5$ meja dan $11$ kursi (dalam praktik dibulatkan sesuai kebijakan produksi).
\end{studikasus}

\subsection*{Studi kasus transportasi}
\begin{studikasus}{Pengiriman Barang}
Sebuah toko mengirim barang memakai dua jenis truk. Truk kecil ($x$) memuat $2$ ton dan ongkos Rp300.000/rit; truk besar ($y$) memuat $5$ ton dan ongkos Rp600.000/rit. Barang yang harus dikirim minimal $30$ ton, dan jumlah rit truk maksimal $10$. Minimumkan ongkos.

\smallskip
\jawab Kendala: $2x+5y\ge30$, $x+y\le10$, $x,y\ge0$. Objektif (minimum) $f=300x+600y$ (ribuan). Titik pojok yang feasible meliputi perpotongan $2x+5y=30$ dengan $x+y=10$ yaitu $x=\tfrac{20}{3},y=\tfrac{10}{3}$, serta $(0,6)$ dan $(0,10)$. Uji nilai $f$ pada titik-titik pojok: minimum tercapai di $(0,6)$ dengan $f=600\cdot6=3600$, yakni ongkos Rp3.600.000 (gunakan $6$ truk besar). Studi kasus ini memperlihatkan bagaimana program linear menekan biaya logistik.
\end{studikasus}

\begin{latihan}
\begin{enumerate}[leftmargin=*]
  \item Tentukan titik-titik pojok daerah $x+y\le6$, $x\ge0$, $y\ge0$.
  \item Maksimumkan $f=3x+2y$ pada $x+y\le5$, $x,y\ge0$.
  \item Minimumkan $f=x+2y$ pada $x+y\ge4$, $x\le5$, $x,y\ge0$.
  \item Tentukan perpotongan garis $x+y=8$ dan $2x+y=12$.
  \item Maksimumkan $f=5x+4y$ pada $x+2y\le10$, $3x+y\le15$, $x,y\ge0$.
  \item Gambarkan daerah feasible $2x+y\le8$, $x+3y\le9$, $x,y\ge0$ dan sebutkan titik pojoknya.
  \item Seorang penjahit membuat baju ($x$) dan celana ($y$); baju butuh $2$ m kain, celana $3$ m, tersedia $24$ m. Laba baju Rp40rb, celana Rp50rb, dan maksimal $10$ baju. Modelkan dan maksimalkan laba.
  \item Mengapa optimum fungsi linear tidak mungkin berada di bagian dalam daerah feasible? Jelaskan singkat.
  \item Pada $f=2x+3y$ di daerah segitiga $(0,0),(4,0),(0,3)$, tentukan maksimum dan minimumnya.
  \item Suatu kebun butuh pupuk minimal $12$ kg N dan $8$ kg P. Pupuk A (per karung) beri $2$ N, $1$ P seharga Rp50rb; pupuk B beri $1$ N, $2$ P seharga Rp40rb. Minimalkan biaya (rumuskan kendala dan objektifnya).
\end{enumerate}
\end{latihan}

\begin{kesalahan}
\begin{enumerate}[leftmargin=*,topsep=2pt]
  \item \textbf{Mengira selalu ada satu penyelesaian.} Sistem bisa tak punya solusi (garis sejajar) atau tak hingga (garis berimpit). Cek konsistensi, jangan paksakan satu jawaban.
  \item \textbf{Salah arah arsiran.} Lupa menguji titik untuk menentukan sisi mana yang memenuhi. Selalu uji $(0,0)$ bila garis tak melaluinya.
  \item \textbf{Lupa kendala $x\ge0,y\ge0$} pada masalah nyata (jumlah barang tak mungkin negatif).
  \item \textbf{Mencari optimum di dalam daerah}, padahal optimum fungsi linear selalu di \emph{titik pojok}.
\end{enumerate}
\end{kesalahan}

\begin{dunianyata}
\begin{itemize}[leftmargin=*,topsep=2pt]
  \item \textbf{Ekonomi/industri:} program linear untuk memaksimalkan laba atau meminimalkan biaya dengan sumber daya terbatas.
  \item \textbf{Logistik:} perencanaan rute dan alokasi armada.
  \item \textbf{Gizi/kesehatan:} menyusun menu termurah yang memenuhi kebutuhan nutrisi.
  \item \textbf{Data science/AI:} sistem linear adalah inti regresi dan banyak algoritma optimasi.
  \item \textbf{Teknik:} analisis rangkaian listrik (hukum Kirchhoff menghasilkan sistem linear).
\end{itemize}
\end{dunianyata}

\section*{Ringkasan Bab}
\addcontentsline{toc}{section}{Ringkasan Bab}
\begin{center}
\begin{tabular}{>{\bfseries}p{3.5cm} p{10.2cm}}
\toprule
Konsep & Inti yang harus diingat \\
\midrule
SPLTV & 3 persamaan, 3 variabel; eliminasi/substitusi. \\
Geometri & tiap persamaan = bidang; solusi = titik potong. \\
Banyak solusi & tunggal (berpotongan), tak hingga (berimpit), nihil (sejajar). \\
SPtLDV & tiap pertidaksamaan = setengah bidang; solusi = irisan. \\
Uji titik & tentukan sisi yang memenuhi dengan menguji $(0,0)$. \\
Program linear & optimum fungsi objektif ada di titik pojok. \\
\bottomrule
\end{tabular}
\end{center}

\begin{refleksi}
Beri tanda centang pada hal yang sudah kamu kuasai:
\begin{itemize}[leftmargin=*,label={},topsep=2pt]
  \item \cek Saya dapat memodelkan masalah nyata menjadi sistem linear.
  \item \cek Saya dapat menyelesaikan SPLTV dengan eliminasi/substitusi.
  \item \cek Saya memahami makna geometris banyaknya penyelesaian.
  \item \cek Saya dapat menggambar daerah penyelesaian dan menentukan nilai optimum.
\end{itemize}
\end{refleksi}

% ---------- SOAL AKHIR BAB ----------
\section*{Soal Akhir Bab 5}
\addcontentsline{toc}{section}{Soal Akhir Bab 5}

\bagiansoal{A}{Fundamental (setara SNBT Dasar)}
\begin{enumerate}[leftmargin=*]
  \item Selesaikan $\begin{cases} x+y=7\\ x-y=1\end{cases}$.
  \item Selesaikan $\begin{cases} 2x+y=8\\ x=3\end{cases}$.
  \item Tentukan $z$ jika $x=1,y=2$ dan $x+y+z=6$.
  \item Apakah $(0,0)$ memenuhi $x+y\le5$?
  \item Selesaikan $\begin{cases} x+y+z=6\\ x=y\\ z=2\end{cases}$.
  \item Tentukan titik potong $x+y=4$ dan $x-y=0$.
  \item Tuliskan pertidaksamaan untuk daerah di bawah garis $y=2x$.
  \item Selesaikan $3x=12$, lalu $x+y=7$.
  \item Hitung $f=2x+y$ di titik $(3,1)$.
  \item Apakah garis $x+y=2$ dan $x+y=5$ berpotongan?
\end{enumerate}

\bagiansoal{B}{Intermediate (setara SNBT Sulit)}
\begin{enumerate}[leftmargin=*]
  \item Selesaikan $\begin{cases} x+y+z=12\\ 2x-y+3z=14\\ x+3y-z=4\end{cases}$.
  \item Harga $3$ buku $+\,2$ pensil $+\,1$ penghapus $=$ Rp28rb; $1$ buku $+\,1$ pensil $+\,1$ penghapus $=$ Rp12rb; $2$ buku $+\,1$ pensil $+\,3$ penghapus $=$ Rp22rb. Tentukan harga tiap barang.
  \item Maksimumkan $f=4x+3y$ pada $2x+y\le10$, $x+3y\le15$, $x,y\ge0$.
  \item Tentukan nilai $a$ agar sistem $\begin{cases} x+y=2\\ ax+2y=4\end{cases}$ punya tak hingga penyelesaian.
  \item Sebuah perusahaan membuat meja dan kursi. Tiap meja butuh 4 jam kerja, kursi 2 jam; tersedia 40 jam. Laba meja Rp80rb, kursi Rp30rb. Maksimalkan laba bila kursi maksimal 15 buah.
\end{enumerate}

\bagiansoal{C}{Advanced (setara Olimpiade SMA --- gabungan dengan Vektor/Fungsi)}
\begin{enumerate}[leftmargin=*]
  \item Selesaikan sistem non-linear $\begin{cases} x+y+z=6\\ xy+yz+zx=11\\ xyz=6\end{cases}$ (kaitkan dengan akar polinomial).
  \item Tunjukkan bahwa sistem linear $A\vec x=\vec b$ punya solusi tunggal jika dan hanya jdan vektor-vektor kolomnya bebas linear (jelaskan untuk kasus $2\times2$).
  \item Tiga garis $x+y=2$, $2x-y=1$, $ax+y=b$ bertemu di satu titik. Tentukan hubungan $a$ dan $b$.
  \item Tentukan semua $(x,y,z)$ bilangan bulat positif dengan $x+y+z=10$ dan $x\le y\le z$.
  \item Pada daerah dengan kendala $x+y\le6$, $x\ge0$, $y\ge0$, fungsi objektif $f=2x+ky$ mencapai maksimum di dua titik pojok sekaligus. Tentukan $k$.
\end{enumerate}

\bagiansoal{D}{Expert (setara tahun pertama universitas --- sintesis antarbab)}
\begin{enumerate}[leftmargin=*]
  \item \textbf{(Sistem + Matriks).} Selesaikan SPLTV pada studi kasus kantin menggunakan aturan Cramer (determinan), lalu bandingkan dengan eliminasi.
  \item \textbf{(Program linear + Pemodelan).} Sebuah pabrik membuat 2 pupuk dengan 3 kendala bahan. Rumuskan masalah optimasi lengkap, gambar daerah fisibel, dan temukan laba maksimum beserta interpretasinya.
  \item \textbf{(Sistem + Geometri).} Diberikan tiga bidang di ruang. Klasifikasikan semua kemungkinan posisi relatif (berpotongan di titik, di garis, sejajar) dan kaitkan dengan banyaknya solusi.
  \item \textbf{(Sistem + Statistika).} Turunkan persamaan normal regresi linear $y=ax+b$ dengan metode kuadrat terkecil, yang berbentuk sistem 2 persamaan 2 variabel dalam $a,b$.
  \item \textbf{(Sistem + Eksponen).} Selesaikan $\begin{cases} 2^x\cdot 3^y=12\\ 2^x\cdot 3^{-y}=\tfrac43\end{cases}$ dengan mengubahnya menjadi sistem linear dalam $x$ dan $y$ lewat logaritma.
\end{enumerate}

\begin{challenge}
\begin{enumerate}[leftmargin=*,topsep=2pt]
  \item Sebuah kebun ingin memaksimalkan hasil panen dua tanaman dengan kendala lahan, air, dan pupuk. Rancang sendiri data yang masuk akal, modelkan, dan selesaikan sebagai program linear.
  \item Tentukan banyaknya solusi bilangan bulat non-negatif dari $x+y+z=20$ (kombinatorika "bintang dan batang").
  \item Sistem $\begin{cases} x+y+z=1\\ x^2+y^2+z^2=1\\ x^3+y^3+z^3=1\end{cases}$: tentukan semua penyelesaian (gunakan identitas simetri).
\end{enumerate}
\end{challenge}

% ============================================================
\chapter{Fungsi Kuadrat}
% ============================================================
\begin{tujuan}
Setelah mempelajari bab ini, peserta didik mampu:
\begin{itemize}[leftmargin=*,topsep=2pt]
  \item menjelaskan fungsi kuadrat dan mengenali grafiknya sebagai parabola;
  \item menentukan titik potong sumbu, sumbu simetri, dan titik puncak;
  \item menafsirkan diskriminan untuk menentukan jenis akar dan posisi grafik;
  \item menyusun fungsi kuadrat dari informasi grafik yang diberikan;
  \item memodelkan masalah optimasi (luas maksimum, lintasan proyektil) dengan fungsi kuadrat.
\end{itemize}
\textbf{Capaian Pembelajaran (Fase E).} Peserta didik dapat menganalisis sifat fungsi kuadrat dan menggunakannya untuk menyelesaikan masalah, termasuk optimasi.
\end{tujuan}

\begin{pemantik}
\begin{itemize}[leftmargin=*,topsep=2pt]
  \item Mengapa lintasan bola yang dilempar, pancuran air mancur, dan kabel jembatan gantung semuanya berbentuk melengkung serupa?
  \item Dengan pagar sepanjang 40 m, bentuk kandang persegi panjang seperti apa yang memberi luas terbesar? Apakah selalu berbentuk persegi?
  \item Mengapa sebuah persamaan kuadrat bisa punya 2, 1, atau bahkan 0 penyelesaian real?
\end{itemize}
\end{pemantik}

\begin{studikasus}{Air Mancur dan Lintasan Proyektil}
Sebuah air mancur menyemburkan air. Ketinggian air (m) terhadap jarak mendatar $x$ (m) tercatat:
\begin{center}
\begin{tabular}{c|ccccc}
\toprule
$x$ (m) & 0 & 1 & 2 & 3 & 4 \\
\midrule
tinggi (m) & 0 & 3 & 4 & 3 & 0 \\
\bottomrule
\end{tabular}
\end{center}
\textbf{Amati dan duga.} (a) Di mana titik tertinggi tercapai? (b) Apakah datanya simetris? Terhadap garis apa? (c) Dapatkah kamu menebak rumus tingginya? Pola "naik lalu turun secara simetris" adalah ciri khas \textbf{fungsi kuadrat}. Titik tertinggi disebut \emph{puncak}, dan garis simetrinya melewati puncak itu. Bab ini membongkar mengapa parabola muncul di mana-mana.
\end{studikasus}

\section*{Peta Konsep Bab}
\addcontentsline{toc}{section}{Peta Konsep Bab}
\begin{center}
\begin{tikzpicture}[
  konsepbox/.style={draw=biru,fill=birumuda,rounded corners,align=center,inner sep=5pt,font=\small\bfseries},
  subbox/.style={draw=hijau,fill=hijaumuda,rounded corners,align=center,inner sep=4pt,font=\footnotesize},
  garis/.style={-Stealth,biru,thick}]
  \node[konsepbox] (akar) {FUNGSI KUADRAT};
  \node[subbox,below left=12mm and 18mm of akar] (graf) {Grafik\\(parabola)};
  \node[subbox,below=12mm of akar] (disk) {Diskriminan\\\&\ akar};
  \node[subbox,below right=12mm and 18mm of akar] (apl) {Optimasi\\(maks/min)};
  \node[subbox,below=10mm of graf] (puncak) {Puncak \&\\sumbu simetri};
  \draw[garis] (akar)--(graf); \draw[garis] (akar)--(disk); \draw[garis] (akar)--(apl);
  \draw[garis] (graf)--(puncak); \draw[garis,dashed,gray] (puncak)--(apl);
\end{tikzpicture}\\
\small \textbf{Peta konsep.} Bentuk, akar, dan puncak parabola adalah tiga sudut pandang yang saling melengkapi; puncak menjadi kunci optimasi.
\end{center}

\section{Bentuk, Grafik, dan Karakteristik}

\subsection*{Membangun konsep}
Dari data air mancur, tinggi naik lalu turun simetris. Kuncinya: suku $x^2$ membuat perubahan tidak konstan (berbeda dari fungsi linear). Tanda koefisien $a$ menentukan parabola "tersenyum" ($a>0$) atau "cemberut" ($a<0$). Inilah "mengapa" bentuk grafiknya melengkung.

\begin{definisi}
\emph{Fungsi kuadrat} adalah $f(x) = ax^2 + bx + c$ dengan $a \neq 0$. Grafiknya berupa \emph{parabola}.
\end{definisi}

\begin{konsep}{Karakteristik Grafik}
\begin{itemize}[leftmargin=*,topsep=2pt]
  \item Jika $a>0$ parabola terbuka ke atas (punya \emph{minimum}); jika $a<0$ terbuka ke bawah (punya \emph{maksimum}).
  \item Sumbu simetri: $x = -\dfrac{b}{2a}$.
  \item Titik puncak: $\left(-\dfrac{b}{2a},\ -\dfrac{D}{4a}\right)$ dengan $D = b^2-4ac$.
  \item Titik potong sumbu-$y$: $(0,c)$; titik potong sumbu-$x$: akar dari $ax^2+bx+c=0$.
\end{itemize}
\end{konsep}

\begin{center}
\begin{tikzpicture}
%% PANEL KIRI: Pengaruh besar |a| (lebar vs sempit)
\begin{axis}[
  name=panelkiri,
  axis lines=middle,width=6.8cm,height=6.5cm,
  title={\small\textbf{Pengaruh $|a|$: lebar vs sempit}},
  xlabel=$x$,ylabel=$y$,
  xmin=-3.2,xmax=3.8,ymin=-0.4,ymax=5.5,
  xtick=\empty,ytick=\empty,samples=80,clip=false]
% a=2: sempit
\addplot[very thick,biru,domain=-1.65:1.65]{2*x^2};
\node[biru,font=\small\bfseries,fill=white,inner sep=2pt]
  at (axis cs:1.5,4.5){$a=2$};
\node[biru,font=\scriptsize,fill=white,inner sep=1pt]
  at (axis cs:1.5,3.85){(sempit)};
% a=1: acuan
\addplot[very thick,hijau,domain=-2.3:2.3]{x^2};
\node[hijau,font=\small\bfseries,fill=white,inner sep=2pt]
  at (axis cs:2.3,5.2){$a=1$};
\node[hijau,font=\scriptsize,fill=white,inner sep=1pt]
  at (axis cs:2.3,4.6){(acuan)};
% a=0.5: lebar
\addplot[very thick,ungu,domain=-3.2:3.2]{0.5*x^2};
\node[ungu,font=\small\bfseries,fill=white,inner sep=2pt]
  at (axis cs:-3.0,4.5){$a=\tfrac12$};
\node[ungu,font=\scriptsize,fill=white,inner sep=1pt]
  at (axis cs:-3.0,3.85){(lebar)};
% Arrow annotation: |a| besar = sempit
\draw[-Stealth,black,thin] (axis cs:0.5,5.0)--(axis cs:0.35,2.4);
\node[black,font=\scriptsize,fill=white,inner sep=1pt,align=center]
  at (axis cs:1.25,5.0){$|a|{\uparrow}$\\makin sempit};
\end{axis}
%% PANEL KANAN: Pengaruh tanda a (buka atas vs bawah)
\begin{axis}[
  name=panelkanan,
  at={(panelkiri.east)},anchor=west,xshift=8pt,
  axis lines=middle,width=6.8cm,height=6.5cm,
  title={\small\textbf{Pengaruh tanda $a$: arah bukaan}},
  xlabel=$x$,ylabel=$y$,
  xmin=-2.8,xmax=3.5,ymin=-3.2,ymax=3.5,
  xtick=\empty,ytick=\empty,samples=80,clip=false]
% a > 0: buka atas, ada minimum
\addplot[very thick,biru,domain=-2.4:2.4]{x^2-2};
\node[biru,font=\small\bfseries,fill=white,inner sep=2pt]
  at (axis cs:2.4,3.8){$a>0$};
\node[biru,font=\scriptsize,fill=white,inner sep=1pt,align=left]
  at (axis cs:2.4,3.0){buka atas\\ada \textbf{min}};
\addplot[only marks,mark=*,biru,mark size=3pt] coordinates {(0,-2)};
\node[biru,font=\scriptsize,fill=white,inner sep=1pt,right]
  at (axis cs:0.1,-2.0){min};
% a < 0: buka bawah, ada maksimum
\addplot[very thick,oranye,domain=-2.4:2.4]{-x^2+2};
\node[oranye,font=\small\bfseries,fill=white,inner sep=2pt]
  at (axis cs:-2.4,3.8){$a<0$};
\node[oranye,font=\scriptsize,fill=white,inner sep=1pt,align=left]
  at (axis cs:-2.4,3.0){buka bawah\\ada \textbf{maks}};
\addplot[only marks,mark=*,oranye,mark size=3pt] coordinates {(0,2)};
\node[oranye,font=\scriptsize,fill=white,inner sep=1pt,right]
  at (axis cs:0.1,2.0){maks};
% Sumbu simetri (garis putus-putus)
\addplot[black,dashed,thin,domain=-3:3]{0*x};
\node[black,font=\tiny,right] at (axis cs:2.1,0.2){sumbu simetri};
\end{axis}
\end{tikzpicture}\\
\small \textbf{Panel kiri:} Makin besar $|a|$, parabola makin sempit (curam). \textbf{Panel kanan:} $a>0$ buka ke atas (punya \emph{minimum}); $a<0$ buka ke bawah (punya \emph{maksimum}). Sumbu simetri selalu tegak lurus sumbu-$x$.
\end{center}

\textbf{Dari mana sumbu simetri $x=-\tfrac{b}{2a}$?} Dengan melengkapkan kuadrat, $f(x)=a\left(x+\tfrac{b}{2a}\right)^2 - \tfrac{D}{4a}$. Bentuk kuadrat $(\cdot)^2$ bernilai minimum (nol) saat $x=-\tfrac{b}{2a}$ --- di situlah puncaknya. Jadi rumus puncak bukan hafalan, melainkan hasil melengkapkan kuadrat.

\begin{konsep}{Diskriminan dan Jenis Akar}
\[
D = b^2 - 4ac:\quad
\begin{cases}
D>0 & \text{dua akar real berbeda (memotong sumbu-}x\text{ di 2 titik)}\\
D=0 & \text{satu akar real (menyinggung sumbu-}x\text{)}\\
D<0 & \text{tidak punya akar real (tidak memotong sumbu-}x\text{)}
\end{cases}
\]
\end{konsep}

\begin{center}
\begin{tikzpicture}
\begin{axis}[axis lines=middle,width=11cm,height=5.5cm,
  xlabel=$x$,ylabel=$y$,xmin=-3,xmax=5,ymin=-2,ymax=6,
  xtick=\empty,ytick=\empty,legend style={font=\footnotesize},legend pos=north west]
\addplot[very thick,biru,domain=-0.5:4.5,samples=60]{x^2-4*x+3};
\addlegendentry{$D>0$}
\addplot[very thick,hijau,domain=-2.2:0.2,samples=60]{(x+1)^2};
\addlegendentry{$D=0$}
\addplot[very thick,oranye,domain=-1.2:1.2,samples=60]{x^2+2};
\addlegendentry{$D<0$}
\end{axis}
\end{tikzpicture}\\
\small \textbf{Tiga kasus diskriminan.} Banyaknya titik potong dengan sumbu-$x$ persis sama dengan banyaknya akar real.
\end{center}

\begin{contoh}{Analisis Fungsi Kuadrat}
Diketahui $f(x) = x^2 - 4x + 3$. Tentukan akar, sumbu simetri, dan titik puncak.

\jawab $D = 16-12 = 4 > 0$ (dua akar real). Akar: $x = \dfrac{4\pm2}{2} = 1$ atau $3$.
Sumbu simetri $x = 2$, titik puncak $\left(2, -\tfrac{4}{4}\right) = (2,-1)$.
\end{contoh}

\begin{contoh}{Menyusun Fungsi dari Studi Kasus}
Tentukan rumus tinggi air mancur (puncak $(2,4)$, melalui $(0,0)$).

\jawab Bentuk puncak: $f(x)=a(x-2)^2+4$. Substitusi $(0,0)$: $0=4a+4\Rightarrow a=-1$. Jadi $f(x)=-(x-2)^2+4=-x^2+4x$. Cek $(1,3)$: $-1+4=3$. Benar.
\end{contoh}

\begin{contoh}{Optimasi Luas Kandang (HOTS)}
Dengan pagar 40 m membentuk persegi panjang, tentukan luas maksimum.

\jawab Misal panjang $x$, lebar $20-x$ (karena keliling $2(x+\text{lebar})=40$). Luas $L(x)=x(20-x)=-x^2+20x$. Puncak di $x=-\tfrac{20}{2(-1)}=10$. Maka lebar $=10$, $L_{\max}=100$ m$^2$. Ternyata bentuk optimum adalah \emph{persegi} $10\times10$ --- menjawab pertanyaan pemantik.
\end{contoh}

\begin{latihan}
\begin{enumerate}[leftmargin=*]
  \item Tentukan jenis akar dari $2x^2+3x+5=0$ tanpa menghitung akarnya.
  \item Tentukan titik puncak $f(x)=x^2-6x+5$.
  \item Tentukan akar $x^2-5x+6=0$.
  \item Gambarkan grafik $f(x)=-x^2+2x+3$ lengkap dengan titik puncak.
  \item Tentukan sumbu simetri $f(x)=2x^2-8x+1$.
  \item Tentukan nilai $c$ agar $x^2+4x+c=0$ punya akar kembar.
  \item Sebuah fungsi kuadrat memotong sumbu-$x$ di $2$ dan $5$. Tulis salah satu bentuknya.
  \item Tentukan nilai minimum $f(x)=x^2+2x+5$.
  \item Tentukan titik potong sumbu-$y$ dari $f(x)=3x^2-x+2$.
  \item Sebuah batu dilempar, tingginya $h(t)=20t-5t^2$. Kapan mencapai tinggi maksimum?
\end{enumerate}
\end{latihan}

\section{Persamaan, Pertidaksamaan, dan Penerapan}

\subsection*{Hubungan akar dan koefisien}
Selain rumus $abc$ ($x=\tfrac{-b\pm\sqrt D}{2a}$), ada hubungan elegan antara akar dan koefisien yang sering muncul di soal HOTS dan olimpiade.

\begin{konsep}{Rumus Akar dan Jumlah--Hasil Kali (Vieta)}
Akar persamaan $ax^2+bx+c=0$:
\[
x_{1,2}=\frac{-b\pm\sqrt{b^2-4ac}}{2a},\qquad
x_1+x_2=-\frac{b}{a},\qquad x_1 x_2=\frac{c}{a}.
\]
\end{konsep}

\begin{contoh}{Menggunakan Vieta}
Akar $x^2-5x+6=0$ adalah $x_1,x_2$. Tanpa mencari akarnya, hitung $x_1^2+x_2^2$.

\jawab $x_1+x_2=5$, $x_1x_2=6$. Maka $x_1^2+x_2^2=(x_1+x_2)^2-2x_1x_2=25-12=13$.
\end{contoh}

\begin{contoh}{Pertidaksamaan Kuadrat}
Tentukan HP dari $x^2-x-6>0$.

\jawab Akar $x^2-x-6=0$ adalah $-2$ dan $3$. Karena $a>0$ (terbuka ke atas), nilai positif di \emph{luar} akar: $x<-2$ atau $x>3$.
\end{contoh}

\begin{latihan}
\begin{enumerate}[leftmargin=*]
  \item Selesaikan $x^2-7x+10=0$ dengan rumus $abc$.
  \item Jika $x_1,x_2$ akar $x^2-3x+1=0$, hitung $x_1+x_2$ dan $x_1x_2$.
  \item Hitung $\tfrac1{x_1}+\tfrac1{x_2}$ untuk $2x^2-5x+3=0$.
  \item Tentukan HP dari $x^2-4\le0$.
  \item Tentukan HP dari $x^2+x-12\ge0$.
  \item Susun persamaan kuadrat berakar $2$ dan $-5$.
  \item Jika satu akar $x^2+px+12=0$ adalah $3$, tentukan $p$ dan akar lain.
  \item Tentukan $m$ agar $x^2-mx+9=0$ berakar kembar.
  \item Selesaikan $-x^2+3x>0$.
  \item Akar $x^2-6x+k=0$ berbanding $1:2$. Tentukan $k$.
\end{enumerate}
\end{latihan}

\begin{kesalahan}
\begin{enumerate}[leftmargin=*,topsep=2pt]
  \item \textbf{Salah tanda pada sumbu simetri.} Rumusnya $x=-\tfrac{b}{2a}$ (dengan tanda minus), sering ditulis $\tfrac{b}{2a}$.
  \item \textbf{Membalik arah pertidaksamaan kuadrat} tanpa memperhatikan tanda $a$. Gambarkan parabola dan lihat di mana grafik di atas/bawah sumbu.
  \item \textbf{Mengira $D<0$ berarti tak ada penyelesaian apa pun.} $D<0$ hanya berarti tak ada akar \emph{real}; grafiknya tetap ada (tidak memotong sumbu-$x$).
  \item \textbf{Lupa $a\neq0$.} Jika $a=0$ fungsi menjadi linear, bukan kuadrat --- rumus puncak tak berlaku.
  \item \textbf{Lupa akar negatif.} Dari $x^2=9$, banyak siswa menulis $x=3$ saja. Yang benar: $x=\pm3$ karena $(-3)^2=9$ juga. Ingat: persamaan kuadrat punya \emph{dua} akar (bisa sama, bisa berbeda tanda).
\end{enumerate}
\end{kesalahan}

\begin{dunianyata}
\begin{itemize}[leftmargin=*,topsep=2pt]
  \item \textbf{Fisika:} lintasan proyektil, gerak jatuh, dan bentuk pancuran air.
  \item \textbf{Ekonomi:} fungsi laba dan biaya sering kuadratik; titik puncak = laba maksimum.
  \item \textbf{Konstruksi:} lengkung parabola pada jembatan, antena parabola memusatkan sinyal di fokus.
  \item \textbf{Teknologi/AI:} fungsi "loss" kuadrat (mean squared error) diminimalkan saat melatih model.
  \item \textbf{Optimasi:} masalah luas/volume maksimum dengan sumber daya terbatas.
\end{itemize}
\end{dunianyata}

\section*{Ringkasan Bab}
\addcontentsline{toc}{section}{Ringkasan Bab}
\begin{center}
\begin{tabular}{>{\bfseries}p{3.5cm} p{10.2cm}}
\toprule
Konsep & Inti yang harus diingat \\
\midrule
Bentuk & $f(x)=ax^2+bx+c$, $a\neq0$; grafik parabola. \\
Arah & $a>0$ terbuka atas (min), $a<0$ terbuka bawah (maks). \\
Sumbu simetri & $x=-\tfrac{b}{2a}$ (dari melengkapkan kuadrat). \\
Puncak & $\left(-\tfrac{b}{2a},-\tfrac{D}{4a}\right)$. \\
Diskriminan & $D=b^2-4ac$ menentukan banyak akar real. \\
Vieta & $x_1+x_2=-\tfrac ba$, $x_1x_2=\tfrac ca$. \\
Aplikasi & proyektil, optimasi luas/laba, loss MSE. \\
\bottomrule
\end{tabular}
\end{center}

\begin{refleksi}
Beri tanda centang pada hal yang sudah kamu kuasai:
\begin{itemize}[leftmargin=*,label={},topsep=2pt]
  \item \cek Saya dapat menentukan puncak \&\ sumbu simetri dan menjelaskan asalnya.
  \item \cek Saya dapat menafsirkan diskriminan secara aljabar dan grafik.
  \item \cek Saya dapat menyusun fungsi kuadrat dari informasi grafik.
  \item \cek Saya dapat menyelesaikan masalah optimasi dengan fungsi kuadrat.
\end{itemize}
\end{refleksi}

% ---------- SOAL AKHIR BAB ----------
\section*{Soal Akhir Bab 6}
\addcontentsline{toc}{section}{Soal Akhir Bab 6}

\bagiansoal{A}{Fundamental (setara SNBT Dasar)}
\begin{enumerate}[leftmargin=*]
  \item Tentukan akar $x^2-9=0$.
  \item Tentukan titik potong sumbu-$y$ dari $f(x)=x^2+2x-3$.
  \item Hitung diskriminan $x^2+4x+4=0$.
  \item Tentukan sumbu simetri $f(x)=x^2-2x$.
  \item Faktorkan $x^2+5x+6$.
  \item Tentukan nilai minimum $f(x)=x^2-1$.
  \item Apakah parabola $f(x)=-2x^2+x$ terbuka ke atas atau bawah?
  \item Tentukan akar $x^2-x-2=0$.
  \item Jika $x_1+x_2=4$, $x_1x_2=3$, susun persamaannya.
  \item Tentukan titik puncak $f(x)=(x-3)^2+2$.
\end{enumerate}

\bagiansoal{B}{Intermediate (setara SNBT Sulit)}
\begin{enumerate}[leftmargin=*]
  \item Tentukan $m$ agar $f(x)=x^2-(m+1)x+4$ menyinggung sumbu-$x$.
  \item Akar $x^2-6x+k=0$ memenuhi $x_1=2x_2$. Tentukan $k$ dan kedua akarnya.
  \item Tentukan HP dari $\dfrac{x^2-4}{x-1}\ge0$.
  \item Fungsi kuadrat melalui $(0,3)$, $(1,0)$, $(3,0)$. Tentukan rumusnya.
  \item Sebuah peluru ditembakkan, tingginya $h(t)=40t-5t^2$. Tentukan tinggi maksimum dan waktu kembali ke tanah.
\end{enumerate}

\bagiansoal{C}{Advanced (setara Olimpiade SMA --- gabungan dengan Trigonometri/Vektor)}
\begin{enumerate}[leftmargin=*]
  \item Tentukan nilai maksimum $f(\theta)=\sin\theta\cos\theta+\cos^2\theta$ dengan substitusi dan analisis kuadrat (Trigonometri + Kuadrat).
  \item Persamaan $x^2-2px+(p^2-1)=0$ punya kedua akar di antara $-2$ dan $4$. Tentukan rentang $p$.
  \item Diberikan $\vec a=(t,1)$ dan $\vec b=(1,t)$. Tentukan $t$ yang meminimalkan $|\vec a+\vec b|^2$ (Vektor + Kuadrat).
  \item Jika $x_1,x_2$ akar $x^2+bx+c=0$, susun persamaan baru berakar $x_1^2$ dan $x_2^2$.
  \item Tentukan banyaknya pasangan bilangan bulat $(x,y)$ dengan $y=x^2$ dan $y\le 100$.
\end{enumerate}

\bagiansoal{D}{Expert (setara tahun pertama universitas --- sintesis antarbab)}
\begin{enumerate}[leftmargin=*]
  \item \textbf{(Kuadrat + Vektor + Trigonometri).} Tunjukkan bahwa nilai $f(t)=|\vec u+t\vec v|^2$ adalah kuadrat dalam $t$; tentukan minimumnya dan tafsirkan sebagai jarak titik ke garis.
  \item \textbf{(Kuadrat + Kalkulus).} Buktikan bahwa puncak parabola yang diperoleh dengan melengkapkan kuadrat sama dengan titik di mana turunan $f'(x)=0$.
  \item \textbf{(Kuadrat + Statistika).} Tunjukkan bahwa nilai $c$ yang meminimalkan $\sum (x_i-c)^2$ adalah rata-rata $\bar x$ (minimasi fungsi kuadrat dalam $c$).
  \item \textbf{(Kuadrat + Peluang).} Variabel acak hasil dadu adil. Nyatakan ragam (variansi) sebagai $E[X^2]-(E[X])^2$ dan hubungkan dengan minimum fungsi kuadrat $g(c)=E[(X-c)^2]$.
  \item \textbf{(Optimasi terkendala).} Dari semua persegi panjang dengan keliling tetap $K$, buktikan persegi memberi luas maksimum; lalu untuk luas tetap $A$, buktikan persegi memberi keliling minimum.
\end{enumerate}

\begin{challenge}
\begin{enumerate}[leftmargin=*,topsep=2pt]
  \item Tentukan semua nilai $a$ agar grafik $y=x^2$ dan garis $y=ax+1$ berpotongan di dua titik yang jaraknya tepat $\sqrt{20}$.
  \item Sebuah bilangan kuadrat sempurna empat digit memiliki dua digit pertama sama dan dua digit terakhir sama. Temukan bilangan itu.
  \item Buktikan bahwa di antara semua segitiga dengan keliling tetap, segitiga sama sisi memiliki luas terbesar (kaitkan dengan ide optimasi kuadratik).
\end{enumerate}
\end{challenge}

% ============================================================
\chapter{Statistika}
% ============================================================
\begin{tujuan}
Setelah mempelajari bab ini, peserta didik mampu:
\begin{itemize}[leftmargin=*,topsep=2pt]
  \item menyajikan data dalam histogram, box plot, dan dot plot serta memilih diagram yang tepat;
  \item menghitung ukuran pemusatan (mean, median, modus) dan menafsirkan maknanya;
  \item menghitung ukuran penyebaran (jangkauan, kuartil, simpangan baku) dan memahami "mengapa" data tersebar;
  \item membandingkan dua kelompok data secara kritis;
  \item menarik kesimpulan berdasarkan data (literasi statistik) dan mewaspadai penyajian yang menyesatkan.
\end{itemize}
\textbf{Capaian Pembelajaran (Fase E).} Peserta didik dapat menyajikan dan menganalisis data dengan ukuran pemusatan dan penyebaran untuk membuat keputusan.
\end{tujuan}

\begin{pemantik}
\begin{itemize}[leftmargin=*,topsep=2pt]
  \item Dua kelas punya rata-rata nilai sama, tetapi guru menyebut salah satunya "lebih merata". Apa yang membedakan keduanya jika rata-ratanya sama?
  \item Mengapa "rata-rata gaji" sebuah perusahaan bisa menyesatkan jika ada satu direktur bergaji sangat tinggi?
  \item Bagaimana satu grafik yang sama bisa dipakai untuk "menipu" pembaca? Apa yang harus diperiksa?
\end{itemize}
\end{pemantik}

\begin{studikasus}{Menganalisis Kebiasaan Belajar Siswa}
Seorang guru mencatat lama belajar (jam/hari) dari dua kelompok siswa selama seminggu:
\begin{center}
\begin{tabular}{l|l}
\toprule
Kelompok & Data (jam) \\
\midrule
A & $2,2,3,3,3,3,4,4$ \\
B & $0,1,2,3,3,4,5,6$ \\
\bottomrule
\end{tabular}
\end{center}
\textbf{Amati dan duga.} (a) Hitung rata-rata tiap kelompok --- apakah sama? (b) Kelompok mana yang "lebih konsisten"? (c) Ukuran apa yang menangkap "kekonsistenan" itu? Kedua kelompok punya rata-rata $3$ jam, tetapi B jauh lebih \emph{tersebar}. Inilah inti statistika: rata-rata saja tak cukup; kita perlu \textbf{ukuran penyebaran}.
\end{studikasus}

\section*{Peta Konsep Bab}
\addcontentsline{toc}{section}{Peta Konsep Bab}
\begin{center}
\begin{tikzpicture}[
  konsepbox/.style={draw=biru,fill=birumuda,rounded corners,align=center,inner sep=5pt,font=\small\bfseries},
  subbox/.style={draw=hijau,fill=hijaumuda,rounded corners,align=center,inner sep=4pt,font=\footnotesize},
  garis/.style={-Stealth,biru,thick}]
  \node[konsepbox] (akar) {STATISTIKA};
  \node[subbox,below left=12mm and 16mm of akar] (saji) {Penyajian\\(histogram,box,dot)};
  \node[subbox,below=12mm of akar] (pus) {Pemusatan\\(mean,median,modus)};
  \node[subbox,below right=12mm and 16mm of akar] (seb) {Penyebaran\\(kuartil,$s$)};
  \node[subbox,below=10mm of pus] (lit) {Literasi data\\\&\ kesimpulan};
  \draw[garis] (akar)--(saji); \draw[garis] (akar)--(pus); \draw[garis] (akar)--(seb);
  \draw[garis] (pus)--(lit); \draw[garis,dashed,gray] (seb)--(lit);
\end{tikzpicture}\\
\small \textbf{Peta konsep.} Data disajikan, diringkas (pusat \&\ sebaran), lalu ditafsirkan untuk mengambil keputusan.
\end{center}

\section{Penyajian Data}

\subsection*{Mengapa banyak jenis diagram?}
Setiap diagram menonjolkan hal berbeda: histogram menunjukkan bentuk distribusi, box plot menyoroti sebaran dan pencilan, dot plot cocok untuk data kecil. Memilih diagram yang tepat adalah keterampilan literasi yang penting --- diagram yang salah dapat menyembunyikan atau melebih-lebihkan pola.

\begin{konsep}{Jenis Diagram}
\begin{description}[leftmargin=1.5cm,style=nextline]
  \item[Histogram] menampilkan distribusi frekuensi data berkelompok dengan batang yang berdempetan.
  \item[Box plot] menampilkan lima ukuran: nilai minimum, $Q_1$, median ($Q_2$), $Q_3$, dan nilai maksimum.
  \item[Dot plot] menampilkan tiap data sebagai satu titik di atas garis bilangan.
\end{description}
\end{konsep}

\begin{center}
\begin{tikzpicture}[scale=0.9]
  \draw[->] (0,0)--(8,0) node[right]{nilai};
  \foreach \x/\q in {1/min,2.5/$Q_1$,4/Med,6/$Q_3$,7.5/maks}
    \draw (\x,-0.1)--(\x,0.1) node[below=2pt]{\small\q};
  \draw[very thick,biru] (2.5,0.4) rectangle (6,1.2);
  \draw[very thick,biru] (4,0.4)--(4,1.2);
  \draw[very thick,biru] (1,0.8)--(2.5,0.8);
  \draw[very thick,biru] (6,0.8)--(7.5,0.8);
\end{tikzpicture}\\
\small Contoh \emph{box plot}: kotak meliputi $50\%$ data tengah; "kumis" menjangkau ekstrem.
\end{center}

\begin{center}
\begin{tikzpicture}
\begin{axis}[ybar,bar width=14pt,width=10cm,height=5cm,
  xlabel=interval nilai,ylabel=frekuensi,ymin=0,
  symbolic x coords={50-59,60-69,70-79,80-89,90-99},
  xtick=data,axis lines=left,nodes near coords,every node near coord/.append style={font=\tiny}]
\addplot[fill=birumuda,draw=biru] coordinates {(50-59,2)(60-69,5)(70-79,9)(80-89,6)(90-99,3)};
\end{axis}
\end{tikzpicture}\\
\small \textbf{Histogram nilai ulangan.} Puncak di interval $70$--$79$ menunjukkan modus kelompok.
\end{center}

\subsection*{Waspada: Grafik yang Menyesatkan}
Visualisasi yang sama dapat memberi kesan yang \emph{sangat berbeda} hanya dengan mengubah skala sumbu. Ini adalah kemampuan literasi data paling penting di era modern.

\begin{center}
\begin{tikzpicture}
  \begin{axis}[ybar,bar width=0.55cm,width=6cm,height=5.5cm,
    title={\small\bfseries Grafik A (sumbu dipotong)},
    xlabel={\scriptsize Tahun},ylabel={\scriptsize Nilai},
    ymin=99,ymax=104,ytick={99,100,101,102,103,104},
    symbolic x coords={2021,2022,2023,2024},
    xtick=data,axis lines=left,
    ticklabel style={font=\scriptsize},
    ymajorgrids=true,grid style={dashed,gray!30},
    title style={font=\small\bfseries}]
  \addplot[fill=birumuda,draw=biru] coordinates
    {(2021,100)(2022,101)(2023,102)(2024,103)};
  \end{axis}
\end{tikzpicture}
\qquad
\begin{tikzpicture}
  \begin{axis}[ybar,bar width=0.55cm,width=6cm,height=5.5cm,
    title={\small\bfseries Grafik B (sumbu dari nol)},
    xlabel={\scriptsize Tahun},ylabel={\scriptsize Nilai},
    ymin=0,ymax=115,ytick={0,20,40,60,80,100},
    symbolic x coords={2021,2022,2023,2024},
    xtick=data,axis lines=left,
    ticklabel style={font=\scriptsize},
    ymajorgrids=true,grid style={dashed,gray!30},
    title style={font=\small\bfseries}]
  \addplot[fill=oranyemuda,draw=oranye] coordinates
    {(2021,100)(2022,101)(2023,102)(2024,103)};
  \end{axis}
\end{tikzpicture}\\[4pt]
\small\textbf{Data identik, kesan berbeda.} Kedua grafik menampilkan nilai yang naik dari $100$ ke $103$ selama 4 tahun (kenaikan $3\%$). Grafik A memotong sumbu-$y$ mulai $99$ sehingga kenaikan terlihat \emph{dramatis}; Grafik B memulai dari $0$ sehingga kenaikan hampir tak terlihat. \textbf{Selalu periksa titik awal sumbu-$y$} sebelum menarik kesimpulan!
\end{center}

\section{Ukuran Pemusatan dan Penyebaran}

\subsection*{Membangun konsep}
Mean adalah "titik keseimbangan" data; median adalah "nilai tengah" yang tahan terhadap pencilan; modus adalah yang paling sering. Mengapa kita butuh ketiganya? Karena masing-masing bisa menyesatkan sendirian --- seperti contoh gaji direktur (mean melonjak, median tetap wajar).

\begin{konsep}{Rumus Penting}
\[
\bar{x} = \frac{\sum x_i}{n},\qquad
\bar{x}_{\text{kelompok}} = \frac{\sum f_i x_i}{\sum f_i}.
\]
\textbf{Jangkauan} $= x_{\max} - x_{\min}$. \quad
\textbf{Jangkauan antarkuartil} $= Q_3 - Q_1$.\\[4pt]
\textbf{Simpangan baku:}\quad
$s = \sqrt{\dfrac{\sum (x_i - \bar{x})^2}{n}}.$
\end{konsep}

\textbf{Mengapa dikuadratkan pada simpangan baku?} Jika kita rata-ratakan selisih $(x_i-\bar x)$ begitu saja, hasilnya selalu nol (sifat mean). Mengkuadratkan menghilangkan tanda dan menghukum simpangan besar lebih berat; akar di akhir mengembalikan satuannya. Inilah "mengapa" rumusnya tampak rumit --- ia mengukur sebaran rata-rata di sekitar mean.

\begin{contoh}{Ukuran Pemusatan}
Data: $4, 6, 7, 7, 9, 10$. Tentukan mean, median, dan modus.

\jawab Mean $= \dfrac{4+6+7+7+9+10}{6} = \dfrac{43}{6} \approx 7{,}17$.
Median $= \dfrac{7+7}{2} = 7$. Modus $= 7$.
\end{contoh}

\begin{contoh}{Membandingkan Sebaran (Studi Kasus)}
Hitung simpangan baku kelompok A ($2,2,3,3,3,3,4,4$), $\bar x=3$.

\jawab Selisih kuadrat: $1,1,0,0,0,0,1,1$ jumlah $=4$. $s=\sqrt{4/8}=\sqrt{0{,}5}\approx0{,}71$. Bandingkan kelompok B (yang sebarannya lebih besar) akan menghasilkan $s$ lebih besar --- secara kuantitatif membenarkan dugaan bahwa A lebih konsisten.
\end{contoh}

\begin{contoh}{Pengaruh Pencilan (HOTS)}
Gaji (juta): $4,5,5,6,60$. Bandingkan mean dan median, mana yang lebih mewakili?

\jawab Mean $=\tfrac{80}{5}=16$ juta; median $=5$ juta. Mean tertarik oleh pencilan $60$. Untuk data dengan pencilan ekstrem, \emph{median} lebih mewakili "gaji tipikal". Inilah kunci literasi statistik.
\end{contoh}

\begin{latihan}
\begin{enumerate}[leftmargin=*]
  \item Tentukan $Q_1$, $Q_2$, $Q_3$ dari data $3,5,5,7,8,9,11,12$.
  \item Hitung simpangan baku dari data $2,4,4,6,9$.
  \item Tentukan mean, median, modus dari $5,7,7,8,10$.
  \item Hitung jangkauan dan jangkauan antarkuartil data soal no.1.
  \item Tentukan mean data berkelompok: nilai $5,6,7,8$ dengan frekuensi $2,4,3,1$.
  \item Data: $10,12,12,15,20$. Hitung mean dan median.
  \item Jika rata-rata 5 bilangan $=8$ dan satu bilangan dibuang menjadi $\bar x=7{,}5$ untuk 4 sisanya, tentukan bilangan yang dibuang.
  \item Tentukan modus dari $3,3,4,5,5,5,6$.
  \item Hitung $Q_2$ dari $1,3,3,6,7,8,9$.
  \item Sebuah data simetris: apa hubungan mean, median, dan modusnya?
\end{enumerate}
\end{latihan}

\section{Statistika Lanjutan (Pengayaan)}

\subsection*{Ukuran letak: kuartil, desil, persentil}
Median membagi data menjadi dua bagian sama banyak. Gagasan ini dapat diperhalus: membagi data terurut menjadi $4$, $10$, atau $100$ bagian sama banyak menghasilkan \emph{kuartil}, \emph{desil}, dan \emph{persentil}. Semuanya disebut \emph{ukuran letak} karena menunjuk "posisi" suatu nilai dalam data.

\begin{konsep}{{Kuartil, Desil, Persentil}}
Untuk data terurut sebanyak $n$ datum:
\begin{itemize}[leftmargin=*,topsep=2pt]
  \item \textbf{Kuartil} $Q_1,Q_2,Q_3$ membagi data menjadi $4$ bagian; letak $Q_i$ pada datum ke-$\dfrac{i(n+1)}{4}$.
  \item \textbf{Desil} $D_1,\dots,D_9$ membagi data menjadi $10$ bagian; letak $D_i$ pada datum ke-$\dfrac{i(n+1)}{10}$.
  \item \textbf{Persentil} $P_1,\dots,P_{99}$ membagi data menjadi $100$ bagian; letak $P_i$ pada datum ke-$\dfrac{i(n+1)}{100}$.
\end{itemize}
Jika letak jatuh di antara dua datum, lakukan interpolasi (rata-rata berbobot).
\end{konsep}

\textbf{Hubungan antar ukuran.} Karena semuanya membagi data yang sama, banyak nilai berimpit: $Q_2=D_5=P_{50}=\text{median}$, dan $Q_1=P_{25}$, $Q_3=P_{75}$.

\begin{contoh}{Menentukan Kuartil}
Tentukan $Q_1,Q_2,Q_3$ dari data $3,5,6,8,9,11,13,15,18$ ($n=9$).

\jawab Letak $Q_1$ pada datum ke-$\tfrac{1(10)}{4}=2{,}5$, yakni antara datum ke-2 ($5$) dan ke-3 ($6$): $Q_1=5{,}5$. $Q_2$ pada datum ke-$5$: $Q_2=9$. $Q_3$ pada datum ke-$7{,}5$, antara $13$ dan $15$: $Q_3=14$.
\end{contoh}

\begin{contoh}{Persentil dan Maknanya}
Pada ujian, nilaimu berada di persentil ke-$80$ ($P_{80}$). Apa artinya?

\jawab Artinya sekitar $80\%$ peserta memperoleh nilai \emph{di bawah} nilaimu, dan hanya $20\%$ yang lebih tinggi --- ukuran letak menjelaskan posisi relatif, bukan sekadar nilai mentah.
\end{contoh}

\subsection*{Box plot dan jangkauan antarkuartil}
\begin{konsep}{Lima Serangkai dan Box Plot}
\emph{Box plot} menggambar lima ringkasan: nilai minimum, $Q_1$, $Q_2$ (median), $Q_3$, dan maksimum. Kotak membentang dari $Q_1$ ke $Q_3$; di dalamnya garis median. \emph{Jangkauan antarkuartil} (JAK / IQR) $=Q_3-Q_1$ mengukur lebar $50\%$ data tengah --- ukuran sebaran yang tahan terhadap pencilan.
\end{konsep}

\begin{center}
\begin{tikzpicture}[scale=0.85]
  \draw[->] (0,0)--(10,0) node[right]{nilai};
  \foreach \x/\q in {1/3,3/5{,}5,5/9,7.5/14,9/18}
    \draw (\x,-0.1)--(\x,0.1) node[below=2pt]{\small\q};
  \draw[very thick,biru] (3,0.4) rectangle (7.5,1.2);
  \draw[very thick,biru] (5,0.4)--(5,1.2);
  \draw[very thick,biru] (1,0.8)--(3,0.8);
  \draw[very thick,biru] (7.5,0.8)--(9,0.8);
  \node[above] at (1,1) {\scriptsize min};
  \node[above] at (3,1.3) {\scriptsize $Q_1$};
  \node[above] at (5,1.3) {\scriptsize $Q_2$};
  \node[above] at (7.5,1.3) {\scriptsize $Q_3$};
  \node[above] at (9,1) {\scriptsize maks};
\end{tikzpicture}\\
\small \textbf{Box plot data contoh.} Kotak memuat $50\%$ data tengah ($Q_1$ ke $Q_3$); "kumis" menjangkau nilai ekstrem.
\end{center}

\subsection*{Pencilan (outlier)}
Nilai yang "jauh menyimpang" dari kumpulan data disebut \emph{pencilan}. Pencilan dapat menyesatkan mean, maka perlu dideteksi secara objektif.

\begin{konsep}{Aturan Pagar (Deteksi Outlier)}
Dengan $\text{JAK}=Q_3-Q_1$, sebuah datum disebut \emph{pencilan} jika berada di luar \emph{pagar}:
\[
\text{pagar bawah}=Q_1-1{,}5\cdot\text{JAK},\qquad
\text{pagar atas}=Q_3+1{,}5\cdot\text{JAK}.
\]
\end{konsep}

\begin{contoh}{Mendeteksi Outlier}
Data $4,5,6,6,7,8,9,25$. Apakah $25$ merupakan pencilan? (Diketahui $Q_1=5{,}5$, $Q_3=8{,}5$.)

\jawab $\text{JAK}=8{,}5-5{,}5=3$. Pagar atas $=8{,}5+1{,}5(3)=13$. Karena $25>13$, nilai $25$ adalah \textbf{pencilan}.
\end{contoh}

\subsection*{Z-score (pengantar)}
Bagaimana membandingkan nilai dari dua ujian dengan skala berbeda? \emph{Z-score} menyatakan "berapa simpangan baku" sebuah nilai dari rata-ratanya, sehingga dapat dibandingkan adil.

\begin{konsep}{Z-Score (Skor Baku)}
\[
z=\frac{x-\bar x}{s},
\]
dengan $\bar x$ rata-rata dan $s$ simpangan baku. $z>0$ di atas rata-rata; $z<0$ di bawah; $|z|$ besar berarti makin jauh dari pusat.
\end{konsep}

\begin{contoh}{Membandingkan dengan Z-Score}
Nilai Matematika $80$ ($\bar x=70$, $s=5$) dan Fisika $85$ ($\bar x=75$, $s=10$). Pada mapel mana posisimu relatif lebih baik?

\jawab Matematika: $z=\tfrac{80-70}{5}=2$. Fisika: $z=\tfrac{85-75}{10}=1$. Walau nilai Fisika lebih tinggi, posisi relatif di Matematika ($z=2$) lebih unggul.
\end{contoh}

\subsection*{Interpretasi data}
Angka statistik baru bermakna jika ditafsirkan dengan benar. Mean menjawab "tipikalnya berapa", JAK dan simpangan baku menjawab "seberapa beragam", kuartil/persentil menjawab "posisi relatif", dan box plot menampilkan semuanya sekaligus termasuk pencilan. Selalu kaitkan angka dengan \emph{konteks} dan waspadai pencilan serta penyajian yang menyesatkan.

\begin{latihan}
\begin{enumerate}[leftmargin=*]
  \item Tentukan $Q_1,Q_2,Q_3$ dari $2,4,5,7,8,10,12$.
  \item Hitung JAK (IQR) data soal no.1.
  \item Tentukan $D_4$ dari $1,3,4,6,7,9,10,12,14,15$.
  \item Tentukan $P_{25}$ dari data $5,8,10,12,15,18,20,24$.
  \item Sebutkan lima serangkai dari $3,6,7,9,12,15,20$.
  \item Data $10,11,12,13,40$ dengan $Q_1=11,Q_3=13$. Apakah $40$ pencilan?
  \item Hitung $z$-score nilai $90$ jika $\bar x=75$ dan $s=10$.
  \item Nilai $z=-1{,}5$ pada data $\bar x=60$, $s=8$. Tentukan nilai aslinya.
  \item Seseorang berada di persentil ke-$90$ tinggi badan. Jelaskan maknanya.
  \item Gambarkan sketsa box plot untuk data $2,3,5,5,6,8,9,14$ (tentukan dulu lima serangkainya).
\end{enumerate}
\end{latihan}

\begin{kesalahan}
\begin{enumerate}[leftmargin=*,topsep=2pt]
  \item \textbf{Lupa mengurutkan data sebelum mencari median/kuartil.} Median adalah nilai tengah \emph{setelah diurutkan}.
  \item \textbf{Mengira mean selalu mewakili data.} Pada data dengan pencilan, gunakan median.
  \item \textbf{Salah membagi pada simpangan baku.} Untuk populasi bagi $n$; untuk sampel bagi $n-1$. Perhatikan konteks soal.
  \item \textbf{Membaca grafik tanpa cek sumbu.} Sumbu-$y$ yang tidak mulai dari nol dapat melebih-lebihkan perbedaan --- waspadai penyajian menyesatkan.
\end{enumerate}
\end{kesalahan}

\begin{dunianyata}
\begin{itemize}[leftmargin=*,topsep=2pt]
  \item \textbf{Data science \&\ AI:} statistik deskriptif adalah langkah pertama setiap analisis data dan pelatihan model.
  \item \textbf{Kesehatan:} uji klinis membandingkan rata-rata dan sebaran kelompok perlakuan vs kontrol.
  \item \textbf{Ekonomi:} indeks harga, pendapatan median, dan ketimpangan (gini) berbasis statistik.
  \item \textbf{Kualitas industri:} kontrol mutu memantau simpangan baku produksi.
  \item \textbf{Olahraga:} analisis performa pemain menggunakan rata-rata dan variabilitas.
\end{itemize}
\end{dunianyata}

\section*{Ringkasan Bab}
\addcontentsline{toc}{section}{Ringkasan Bab}
\begin{center}
\begin{tabular}{>{\bfseries}p{3.5cm} p{10.2cm}}
\toprule
Konsep & Inti yang harus diingat \\
\midrule
Penyajian & histogram (bentuk), box plot (sebaran/pencilan), dot plot (data kecil). \\
Mean & $\bar x=\tfrac{\sum x_i}{n}$; titik keseimbangan, peka pencilan. \\
Median & nilai tengah data terurut; tahan pencilan. \\
Modus & nilai paling sering. \\
Kuartil & $Q_1,Q_2,Q_3$ membagi data jadi empat; JAK $=Q_3-Q_1$. \\
Simpangan baku & $s=\sqrt{\tfrac{\sum(x_i-\bar x)^2}{n}}$; ukuran sebaran rata-rata. \\
Literasi & rata-rata saja tak cukup; cek sebaran \&\ cara penyajian. \\
\bottomrule
\end{tabular}
\end{center}

\begin{refleksi}
Beri tanda centang pada hal yang sudah kamu kuasai:
\begin{itemize}[leftmargin=*,label={},topsep=2pt]
  \item \cek Saya dapat memilih dan membaca diagram yang tepat.
  \item \cek Saya dapat menghitung mean, median, modus dan tahu kapan masing-masing dipakai.
  \item \cek Saya memahami mengapa simpangan baku mengkuadratkan selisih.
  \item \cek Saya dapat membandingkan dua kelompok data secara kritis.
\end{itemize}
\end{refleksi}

% ---------- SOAL AKHIR BAB ----------
\section*{Soal Akhir Bab 7}
\addcontentsline{toc}{section}{Soal Akhir Bab 7}

\bagiansoal{A}{Fundamental (setara SNBT Dasar)}
\begin{enumerate}[leftmargin=*]
  \item Hitung mean dari $4,6,8,10$.
  \item Tentukan median dari $3,7,9,11,15$.
  \item Tentukan modus dari $2,2,3,4,4,4,5$.
  \item Hitung jangkauan dari $5,8,12,20$.
  \item Tentukan $Q_2$ dari $1,2,3,4,5,6$.
  \item Mean 4 bilangan $=10$; tentukan jumlahnya.
  \item Tentukan $Q_1$ dari $2,4,6,8,10,12$.
  \item Apakah data $5,5,5,5$ punya simpangan baku? Berapa?
  \item Hitung mean dari frekuensi: nilai $1,2,3$ frekuensi $2,2,1$.
  \item Tentukan jangkauan antarkuartil data soal no.7.
\end{enumerate}

\bagiansoal{B}{Intermediate (setara SNBT Sulit)}
\begin{enumerate}[leftmargin=*]
  \item Rata-rata 10 nilai $=70$. Setelah satu nilai diralat dari $50$ menjadi $80$, tentukan rata-rata baru.
  \item Data $x, 2x, 3x, 4x, 5x$ memiliki simpangan baku $s$. Bila tiap data dikali 2, bagaimana $s$ berubah? Jelaskan.
  \item Tentukan median data berkelompok: interval $[10,20),[20,30),[30,40)$ dengan frekuensi $4,10,6$.
  \item Gabungan dua kelompok: kelompok I ($n=20$, $\bar x=6$) dan II ($n=30$, $\bar x=8$). Tentukan rata-rata gabungan.
  \item Diketahui $\sum x_i=100$, $\sum x_i^2=1100$, $n=10$. Hitung simpangan baku.
\end{enumerate}

\bagiansoal{C}{Advanced (setara Olimpiade SMA --- gabungan dengan Fungsi Kuadrat)}
\begin{enumerate}[leftmargin=*]
  \item Buktikan bahwa nilai $c$ yang meminimalkan $\sum(x_i-c)^2$ adalah $c=\bar x$ (gunakan analisis fungsi kuadrat dalam $c$).
  \item Tunjukkan identitas $s^2=\overline{x^2}-\bar x^2$ (rata-rata kuadrat dikurangi kuadrat rata-rata).
  \item Lima bilangan asli berbeda memiliki mean $=$ median $=5$. Tentukan nilai maksimum bilangan terbesar.
  \item Data $a,b,c$ punya mean $m$. Buktikan $(a-m)+(b-m)+(c-m)=0$ dan jelaskan maknanya pada simpangan baku.
  \item Jika menambahkan konstanta $k$ ke semua data, buktikan mean naik $k$ tetapi simpangan baku tetap.
\end{enumerate}

\bagiansoal{D}{Expert (setara tahun pertama universitas --- sintesis antarbab)}
\begin{enumerate}[leftmargin=*]
  \item \textbf{(Statistika + Peluang + Logaritma).} Untuk distribusi peluang $p_1,\dots,p_n$, entropi $H=-\sum p_i\,{}^2\!\log p_i$. Tunjukkan $H$ maksimum saat semua $p_i$ sama, dan kaitkan dengan "ketidakpastian" data.
  \item \textbf{(Statistika + Vektor).} Tunjukkan bahwa koefisien korelasi dua variabel sama dengan kosinus sudut antara dua vektor data terpusat.
  \item \textbf{(Statistika + Fungsi Kuadrat).} Turunkan garis regresi $y=ax+b$ kuadrat terkecil dan tunjukkan ia melalui titik $(\bar x,\bar y)$.
  \item \textbf{(Statistika + Deret).} Jika data adalah barisan aritmetika $a, a+d,\dots,a+(n-1)d$, turunkan rumus tertutup untuk mean dan simpangan bakunya.
  \item \textbf{(Literasi).} Diberikan dua grafik dari data sama dengan skala sumbu berbeda. Analisis bagaimana penyajian dapat menyesatkan dan rumuskan kriteria penyajian yang jujur.
\end{enumerate}

\begin{challenge}
\begin{enumerate}[leftmargin=*,topsep=2pt]
  \item Mungkinkah menambah satu data baru ke kumpulan data \emph{menurunkan} mean tetapi \emph{menaikkan} median sekaligus? Beri contoh atau buktikan mustahil.
  \item Sepuluh bilangan bulat positif memiliki mean $10$ dan median $9$. Tentukan nilai maksimum yang mungkin untuk bilangan terbesar.
  \item Rancang sebuah survei mini di kelasmu (mis. waktu tidur), kumpulkan data, sajikan dengan box plot, dan tulis satu kesimpulan yang didukung data beserta keterbatasannya.
\end{enumerate}
\end{challenge}

% ============================================================
\chapter{Peluang}
% ============================================================
\begin{tujuan}
Setelah mempelajari bab ini, peserta didik mampu:
\begin{itemize}[leftmargin=*,topsep=2pt]
  \item menjelaskan ruang sampel, kejadian, dan menghitung peluang klasik;
  \item menggunakan aturan komplemen, penjumlahan, dan kejadian saling lepas;
  \item membedakan peluang teoretis dan empiris (frekuensi relatif);
  \item menerapkan aturan pencacahan dan peluang bersyarat sederhana;
  \item menafsirkan peluang dalam pengambilan keputusan nyata (cuaca, asuransi, kesehatan).
\end{itemize}
\textbf{Capaian Pembelajaran (Fase E).} Peserta didik dapat menjelaskan dan menentukan peluang kejadian majemuk serta menggunakannya untuk mengambil keputusan.
\end{tujuan}

\begin{pemantik}
\begin{itemize}[leftmargin=*,topsep=2pt]
  \item Ramalan cuaca menyebut "peluang hujan $70\%$". Apa arti angka itu, dan dari mana ia berasal?
  \item Jika sebuah koin sudah muncul Angka 5 kali berturut-turut, apakah lemparan berikutnya "lebih mungkin" Gambar? Mengapa banyak orang keliru di sini?
  \item Mengapa kasino dan asuransi hampir selalu untung dalam jangka panjang meskipun tiap kejadian tampak acak?
\end{itemize}
\end{pemantik}

\begin{studikasus}{Prediksi Cuaca dan Peluang Hujan}
Badan cuaca mencatat selama 100 hari dengan kondisi awan tertentu, ternyata hujan terjadi pada 70 hari.
\begin{center}
\begin{tabular}{l|c|c}
\toprule
Kondisi & Hujan & Tidak hujan \\
\midrule
Berawan tebal (100 hari) & 70 & 30 \\
\bottomrule
\end{tabular}
\end{center}
\textbf{Amati dan duga.} (a) Bagaimana cara menaksir "peluang hujan" dari data ini? (b) Apakah peluang ini pasti benar untuk besok? (c) Apa bedanya dengan peluang dadu yang bisa dihitung tanpa percobaan? "Peluang hujan $70\%$" adalah \emph{peluang empiris} --- frekuensi relatif $\tfrac{70}{100}$. Bab ini membedakannya dari \emph{peluang teoretis} dan membangun aturan untuk kejadian yang lebih rumit.
\end{studikasus}

\section*{Peta Konsep Bab}
\addcontentsline{toc}{section}{Peta Konsep Bab}
\begin{center}
\begin{tikzpicture}[
  konsepbox/.style={draw=biru,fill=birumuda,rounded corners,align=center,inner sep=5pt,font=\small\bfseries},
  subbox/.style={draw=hijau,fill=hijaumuda,rounded corners,align=center,inner sep=4pt,font=\footnotesize},
  garis/.style={-Stealth,biru,thick}]
  \node[konsepbox] (akar) {PELUANG};
  \node[subbox,below left=12mm and 16mm of akar] (rs) {Ruang sampel\\\&\ kejadian};
  \node[subbox,below=12mm of akar] (at) {Aturan peluang\\(komplemen, $\cup$)};
  \node[subbox,below right=12mm and 16mm of akar] (cac) {Pencacahan\\(permutasi/kombinasi)};
  \node[subbox,below=10mm of at] (apl) {Empiris vs teoretis\\\&\ keputusan};
  \draw[garis] (akar)--(rs); \draw[garis] (akar)--(at); \draw[garis] (akar)--(cac);
  \draw[garis] (at)--(apl); \draw[garis,dashed,gray] (cac)--(apl);
\end{tikzpicture}\\
\small \textbf{Peta konsep.} Dari ruang sampel lahir aturan peluang; pencacahan membantu menghitung kasus banyak; semuanya bermuara pada pengambilan keputusan.
\end{center}

\section{Konsep Dasar Peluang}

\subsection*{Membangun konsep}
Peluang teoretis mengandaikan semua hasil "sama mungkin" (seperti dadu adil), sehingga $P(A)=\tfrac{n(A)}{n(S)}$. Peluang empiris mengukur dari pengamatan: $\tfrac{\text{frekuensi}}{\text{total percobaan}}$. \emph{Hukum bilangan besar} menjamin keduanya makin dekat saat percobaan diperbanyak --- inilah mengapa kasino dan asuransi untung dalam jangka panjang (menjawab pemantik).

\begin{definisi}
\emph{Ruang sampel} $S$ adalah himpunan semua hasil yang mungkin. \emph{Kejadian} $A$ adalah himpunan bagian dari $S$. Peluang kejadian $A$:
\[
P(A) = \frac{n(A)}{n(S)}, \qquad 0 \le P(A) \le 1.
\]
\end{definisi}

\begin{konsep}{Aturan Peluang}
\begin{itemize}[leftmargin=*,topsep=2pt]
  \item Komplemen: $P(A^c) = 1 - P(A)$.
  \item Kejadian saling lepas: $P(A \cup B) = P(A) + P(B)$.
  \item Kejadian tidak saling lepas: $P(A\cup B) = P(A)+P(B)-P(A\cap B)$.
  \item Kejadian saling bebas: $P(A\cap B)=P(A)\cdot P(B)$.
\end{itemize}
\end{konsep}

\textbf{Mengapa komplemen sangat berguna?} Sering lebih mudah menghitung "peluang TIDAK terjadi" lalu menguranginya dari 1. Misalnya peluang "muncul minimal satu Angka pada beberapa lemparan" lebih cepat lewat $1-P(\text{tidak ada Angka})$.

\begin{contoh}{Pelemparan Dadu}
Sebuah dadu dilempar. Tentukan peluang muncul mata dadu genap.

\jawab $S = \{1,2,3,4,5,6\}$, $A = \{2,4,6\}$. Maka $P(A) = \dfrac{3}{6} = \dfrac12.$
\end{contoh}

\begin{contoh}{Dua Kejadian}
Dari kartu bernomor 1--10 diambil satu kartu. Peluang terambil kelipatan 2 atau kelipatan 3?

\jawab $A=\{2,4,6,8,10\}$, $B=\{3,6,9\}$, $A\cap B = \{6\}$.
$P(A\cup B) = \tfrac{5}{10}+\tfrac{3}{10}-\tfrac{1}{10} = \tfrac{7}{10}.$
\end{contoh}

\begin{contoh}{Strategi Komplemen (HOTS)}
Tiga koin dilempar. Berapa peluang muncul \emph{minimal} satu Gambar?

\jawab Lebih mudah lewat komplemen: $P(\text{tak ada Gambar})=P(\text{semua Angka})=\left(\tfrac12\right)^3=\tfrac18$. Maka $P(\ge1\text{ Gambar})=1-\tfrac18=\tfrac78$.
\end{contoh}

\begin{latihan}
\begin{enumerate}[leftmargin=*]
  \item Dua koin dilempar bersamaan. Tentukan peluang muncul tepat satu Angka.
  \item Dalam kotak terdapat 5 bola merah dan 3 bola putih. Satu bola diambil. Tentukan peluang terambil merah dan peluang \emph{bukan} merah.
  \item Sebuah dadu dilempar. Peluang muncul mata $>4$?
  \item Dari kartu 1--20 diambil satu. Peluang kelipatan 5?
  \item Peluang hujan $0{,}3$. Berapa peluang tidak hujan?
  \item Dua dadu dilempar. Peluang jumlah mata $=7$?
  \item Sebuah huruf dipilih acak dari "MATEMATIKA". Peluang huruf "A"?
  \item Kejadian $A$ dan $B$ saling lepas, $P(A)=0{,}4$, $P(B)=0{,}3$. Hitung $P(A\cup B)$.
  \item Dua koin: peluang keduanya sama (AA atau GG)?
  \item Dari 1--10, peluang bilangan prima?
\end{enumerate}
\end{latihan}

\section{Pencacahan dan Peluang Kejadian Majemuk}

\subsection*{Mengapa perlu mencacah?}
Ketika ruang sampel besar (mis. menyusun 5 orang dalam barisan), mendaftar satu per satu mustahil. Aturan perkalian, permutasi, dan kombinasi adalah "alat hitung cepat" banyaknya kemungkinan, yang menjadi penyebut peluang.

\begin{konsep}{Kaidah Pencacahan}
\begin{itemize}[leftmargin=*,topsep=2pt]
  \item \textbf{Aturan perkalian:} jika tahap 1 punya $m$ cara dan tahap 2 punya $n$ cara, total $m\times n$.
  \item \textbf{Permutasi} (urutan penting): $P(n,r)=\dfrac{n!}{(n-r)!}$.
  \item \textbf{Kombinasi} (urutan tidak penting): $C(n,r)=\dbinom{n}{r}=\dfrac{n!}{r!(n-r)!}$.
\end{itemize}
\end{konsep}

\begin{contoh}{Permutasi vs Kombinasi}
Dari 5 siswa dipilih 3. (a) Berapa cara menyusun mereka dalam barisan? (b) Berapa cara memilih tim tanpa urutan?

\jawab (a) Urutan penting: $P(5,3)=\tfrac{5!}{2!}=60$. (b) Urutan tak penting: $C(5,3)=\tfrac{5!}{3!2!}=10$.
\end{contoh}

\begin{contoh}{Peluang dengan Kombinasi}
Dari 7 bola (4 merah, 3 putih) diambil 2 sekaligus. Peluang keduanya merah?

\jawab $n(S)=C(7,2)=21$. $n(A)=C(4,2)=6$. $P(A)=\tfrac{6}{21}=\tfrac27$.
\end{contoh}

\begin{latihan}
\begin{enumerate}[leftmargin=*]
  \item Berapa banyak bilangan 3 digit berbeda dari angka $1,2,3,4$?
  \item Hitung $C(6,2)$ dan $P(6,2)$.
  \item Dari 8 orang dipilih ketua dan wakil. Berapa cara?
  \item Berapa cara menyusun huruf "BUKU"?
  \item Dari 5 merah, 4 biru diambil 3. Peluang ketiganya biru?
  \item Sebuah PIN 4 digit (boleh berulang). Berapa kemungkinan?
  \item Berapa banyak diagonal segi lima?
  \item Dari 10 soal dipilih 6. Berapa cara?
  \item Dua dadu: peluang muncul angka kembar?
  \item Tiga koin: peluang muncul tepat dua Gambar?
\end{enumerate}
\end{latihan}

\section{Peluang Lanjutan (Pengayaan)}

\subsection*{Peluang bersyarat}
Sering kita ingin tahu peluang suatu kejadian \emph{dengan syarat} kejadian lain sudah terjadi --- misalnya peluang seseorang sakit \emph{jika} hasil tesnya positif. Informasi tambahan ini memperbarui peluang.

\begin{definisi}
\emph{Peluang bersyarat} $A$ dengan syarat $B$ (ditulis $P(A\mid B)$) adalah peluang $A$ terjadi jika diketahui $B$ telah terjadi:
\[
P(A\mid B)=\frac{P(A\cap B)}{P(B)},\qquad P(B)\neq0.
\]
\end{definisi}

\textbf{Gagasan intinya:} mengetahui $B$ "mempersempit" ruang sampel menjadi hanya $B$, sehingga kita membandingkan $A\cap B$ terhadap $B$ saja, bukan terhadap seluruh ruang sampel.

\begin{contoh}{Peluang Bersyarat dari Tabel}
Dari $100$ siswa: $60$ suka Matematika (M), $40$ suka Fisika (F), dan $25$ suka keduanya. Jika seorang siswa diketahui suka Matematika, berapa peluang ia juga suka Fisika?

\jawab $P(F\mid M)=\dfrac{P(F\cap M)}{P(M)}=\dfrac{25/100}{60/100}=\dfrac{25}{60}=\dfrac{5}{12}.$
\end{contoh}

\subsection*{Diagram pohon}
\emph{Diagram pohon} menggambar kejadian bertahap: tiap cabang diberi peluang, dan peluang satu lintasan diperoleh dengan \emph{mengalikan} peluang sepanjang cabang.

\begin{center}
\begin{tikzpicture}[
  merahN/.style={circle,draw=oranye,fill=oranyemuda,font=\small\bfseries,
                 inner sep=2pt,minimum size=9mm},
  putihN/.style={circle,draw=biru,fill=birumuda,font=\small\bfseries,
                 inner sep=2pt,minimum size=9mm},
  mulaiN/.style={draw=hijau,fill=hijaumuda,rounded corners,
                 font=\small\bfseries,inner sep=5pt},
  panah/.style={-Stealth,black!55,thick}]
%--- Akar
\node[mulaiN] (R) at (0,0) {Mulai};
%--- Kolom header
\node[font=\footnotesize\itshape,black!60] at (3.2, 4.0) {Ambil ke-1};
\node[font=\footnotesize\itshape,black!60] at (6.5, 4.0) {Ambil ke-2};
%--- Level 1: pengambilan pertama
\node[merahN] (M1) at (3.2, 1.8) {$M_1$};
\node[putihN] (P1) at (3.2,-1.8) {$P_1$};
\draw[panah] (R)--(M1) node[midway,above left,font=\small]{$\tfrac{3}{5}$};
\draw[panah] (R)--(P1) node[midway,below left,font=\small]{$\tfrac{2}{5}$};
%--- Level 2: pengambilan kedua (jarak vertikal cukup supaya tidak bertabrakan)
\node[merahN] (MM) at (6.5, 3.1) {$M_2$};
\node[putihN] (MP) at (6.5, 0.8) {$P_2$};
\node[merahN] (PM) at (6.5,-0.8) {$M_2$};
\node[putihN] (PP) at (6.5,-3.1) {$P_2$};
\draw[panah] (M1)--(MM) node[midway,above left,font=\small]{$\tfrac{2}{4}$};
\draw[panah] (M1)--(MP) node[midway,below,font=\small]{$\tfrac{2}{4}$};
\draw[panah] (P1)--(PM) node[midway,above,font=\small]{$\tfrac{3}{4}$};
\draw[panah] (P1)--(PP) node[midway,below left,font=\small]{$\tfrac{1}{4}$};
%--- Peluang lintasan (hasil kali) di kanan tiap daun
\node[font=\small,anchor=west] at ($(MM.east)+(0.18cm,0)$)
  {$MM:\ \tfrac{3}{5}\!\cdot\!\tfrac{2}{4}=\boldsymbol{\tfrac{3}{10}}$};
\node[font=\small,anchor=west] at ($(MP.east)+(0.18cm,0)$)
  {$MP:\ \tfrac{3}{5}\!\cdot\!\tfrac{2}{4}=\boldsymbol{\tfrac{3}{10}}$};
\node[font=\small,anchor=west] at ($(PM.east)+(0.18cm,0)$)
  {$PM:\ \tfrac{2}{5}\!\cdot\!\tfrac{3}{4}=\boldsymbol{\tfrac{3}{10}}$};
\node[font=\small,anchor=west] at ($(PP.east)+(0.18cm,0)$)
  {$PP:\ \tfrac{2}{5}\!\cdot\!\tfrac{1}{4}=\boldsymbol{\tfrac{1}{10}}$};
\end{tikzpicture}\\
\small \textbf{Diagram pohon} pengambilan $2$ bola dari $3$ merah (\tikz\node[circle,draw=oranye,fill=oranyemuda,font=\tiny\bfseries,inner sep=1pt]{M};)
dan $2$ putih (\tikz\node[circle,draw=biru,fill=birumuda,font=\tiny\bfseries,inner sep=1pt]{P};)
tanpa pengembalian. Peluang tiap lintasan $=$ hasil kali peluang cabangnya. Jumlah seluruh peluang $= \tfrac{3}{10}+\tfrac{3}{10}+\tfrac{3}{10}+\tfrac{1}{10}=1$.
\end{center}

\subsection*{Kejadian independen dan dependen}
\begin{konsep}{Independen vs Dependen}
Dua kejadian $A$ dan $B$:
\begin{itemize}[leftmargin=*,topsep=2pt]
  \item \textbf{Independen (bebas):} terjadinya $B$ tidak mengubah peluang $A$. Ciri: $P(A\cap B)=P(A)\cdot P(B)$, atau setara $P(A\mid B)=P(A)$. Contoh: dua kali lempar koin.
  \item \textbf{Dependen (bergantung):} peluang berubah karena kejadian sebelumnya. Contoh: mengambil dua kartu \emph{tanpa pengembalian}.
\end{itemize}
\end{konsep}

\begin{contoh}{Dependen --- Tanpa Pengembalian}
Dari $3$ bola merah dan $2$ putih, diambil $2$ bola berurutan tanpa pengembalian. Tentukan peluang keduanya merah.

\jawab Bola pertama merah: $\tfrac35$. Setelah satu merah diambil, tersisa $2$ merah dari $4$ bola: $\tfrac24$. Maka $P(MM)=\tfrac35\cdot\tfrac24=\tfrac{6}{20}=\tfrac{3}{10}$. Peluang tahap kedua \emph{bergantung} pada hasil tahap pertama --- inilah kejadian dependen.
\end{contoh}

\begin{contoh}{Independen --- Dengan Pengembalian}
Pada kasus yang sama, bola diambil dengan \emph{pengembalian}. Tentukan peluang keduanya merah.

\jawab Karena dikembalikan, peluang tiap tahap tetap $\tfrac35$. Maka $P(MM)=\tfrac35\cdot\tfrac35=\tfrac{9}{25}$ --- kejadian independen.
\end{contoh}

\subsection*{Teorema Bayes dasar}
Teorema Bayes "membalik" arah peluang bersyarat: dari $P(B\mid A)$ menuju $P(A\mid B)$. Ia sangat penting pada uji medis, penyaringan spam, dan kecerdasan buatan.

\begin{konsep}{Teorema Bayes}
\[
P(A\mid B)=\frac{P(B\mid A)\,P(A)}{P(B)},\qquad
P(B)=P(B\mid A)P(A)+P(B\mid A^c)P(A^c).
\]
\end{konsep}

\begin{contoh}{Uji Penyakit (Bayes)}
Penyakit menjangkiti $1\%$ populasi. Sebuah tes positif pada $90\%$ orang sakit (sensitivitas) dan keliru positif pada $5\%$ orang sehat. Jika seseorang dites positif, berapa peluang ia benar-benar sakit?

\jawab Misalkan $S$ sakit, $+$ tes positif. $P(S)=0{,}01$, $P(+\mid S)=0{,}9$, $P(+\mid S^c)=0{,}05$.
\[
P(+)=0{,}9(0{,}01)+0{,}05(0{,}99)=0{,}009+0{,}0495=0{,}0585.
\]
\[
P(S\mid +)=\frac{0{,}9(0{,}01)}{0{,}0585}=\frac{0{,}009}{0{,}0585}\approx0{,}154.
\]
Hanya sekitar $15{,}4\%$! Karena penyakitnya langka, banyak hasil positif justru berasal dari orang sehat (positif palsu) --- hasil yang berlawanan dengan intuisi.
\end{contoh}

\subsection*{Studi kasus peluang dunia nyata}
\begin{dunianyata}
\begin{itemize}[leftmargin=*,topsep=2pt]
  \item \textbf{Kedokteran:} menafsirkan hasil tes (Bayes) agar tidak panik atas positif palsu.
  \item \textbf{Penyaring spam:} email diklasifikasi berdasarkan peluang bersyarat kemunculan kata.
  \item \textbf{Cuaca \&\ asuransi:} peluang kejadian bertahap dihitung lewat diagram pohon.
  \item \textbf{Permainan \&\ olahraga:} strategi optimal mempertimbangkan peluang bersyarat lawan.
\end{itemize}
\end{dunianyata}

\begin{latihan}
\begin{enumerate}[leftmargin=*]
  \item Diketahui $P(A\cap B)=0{,}2$ dan $P(B)=0{,}5$. Tentukan $P(A\mid B)$.
  \item Dari $50$ orang, $30$ berkacamata, $20$ pria, $12$ pria berkacamata. Jika terpilih pria, peluang ia berkacamata?
  \item Dua dadu dilempar. Jika diketahui jumlahnya genap, peluang jumlahnya $8$?
  \item Apakah "hujan hari ini" dan "lulus ujian" independen? Jelaskan.
  \item Dari $4$ merah, $3$ biru diambil $2$ tanpa pengembalian. Peluang keduanya biru?
  \item Sama seperti no.5 tetapi \emph{dengan} pengembalian. Bandingkan hasilnya.
  \item Sebuah koin dilempar dua kali. Buktikan kedua lemparan independen dengan menghitung $P(\text{AA})$.
  \item Gambarkan diagram pohon pelemparan koin $2$ kali dan tuliskan peluang tiap lintasan.
  \item Suatu pabrik: mesin A ($60\%$ produk, cacat $2\%$) dan mesin B ($40\%$, cacat $5\%$). Sebuah produk cacat; peluang ia dari mesin B? (Bayes)
  \item Tes narkoba akurat $98\%$, hanya $0{,}5\%$ populasi memakai. Tanpa hitung penuh, jelaskan mengapa hasil positif belum tentu berarti pemakai.
\end{enumerate}
\end{latihan}

\begin{kesalahan}
\begin{enumerate}[leftmargin=*,topsep=2pt]
  \item \textbf{Keliru judi (gambler's fallacy).} Koin tak punya "ingatan"; setelah 5 Angka, peluang Angka berikutnya tetap $\tfrac12$. Kejadian bebas tidak saling memengaruhi.
  \item \textbf{Menjumlah peluang yang tidak saling lepas.} Jika $A\cap B\neq\varnothing$, jangan lupa kurangi $P(A\cap B)$.
  \item \textbf{Tertukar permutasi dan kombinasi.} Tanya: "apakah urutan penting?" Barisan/jabatan = permutasi; tim/komite = kombinasi.
  \item \textbf{Peluang $>1$.} Jika hasil hitung melebihi 1, pasti ada kesalahan (mungkin mencacah ganda).
\end{enumerate}
\end{kesalahan}

\begin{dunianyata}
\begin{itemize}[leftmargin=*,topsep=2pt]
  \item \textbf{Cuaca:} prakiraan berbasis peluang empiris dari data historis.
  \item \textbf{Asuransi \&\ keuangan:} premi dihitung dari peluang klaim; manajemen risiko investasi.
  \item \textbf{Kesehatan:} sensitivitas/spesifisitas tes, peluang positif palsu (Teorema Bayes).
  \item \textbf{AI \&\ data science:} model probabilistik, klasifikasi, dan sistem rekomendasi.
  \item \textbf{Game development:} sistem "loot/drop rate" dan keseimbangan permainan.
\end{itemize}
\end{dunianyata}

\section*{Ringkasan Bab}
\addcontentsline{toc}{section}{Ringkasan Bab}
\begin{center}
\begin{tabular}{>{\bfseries}p{3.5cm} p{10.2cm}}
\toprule
Konsep & Inti yang harus diingat \\
\midrule
Peluang klasik & $P(A)=\tfrac{n(A)}{n(S)}$, $0\le P(A)\le1$. \\
Empiris & frekuensi relatif; mendekati teoretis (hukum bilangan besar). \\
Komplemen & $P(A^c)=1-P(A)$; ampuh untuk "minimal". \\
Gabungan & $P(A\cup B)=P(A)+P(B)-P(A\cap B)$. \\
Bebas & $P(A\cap B)=P(A)P(B)$. \\
Pencacahan & perkalian, permutasi (urutan), kombinasi (tanpa urutan). \\
Aplikasi & cuaca, asuransi, kesehatan, AI, game. \\
\bottomrule
\end{tabular}
\end{center}

\begin{refleksi}
Beri tanda centang pada hal yang sudah kamu kuasai:
\begin{itemize}[leftmargin=*,label={},topsep=2pt]
  \item \cek Saya dapat membedakan peluang teoretis dan empiris.
  \item \cek Saya dapat menggunakan aturan komplemen dan gabungan.
  \item \cek Saya dapat memilih antara permutasi dan kombinasi.
  \item \cek Saya tidak terjebak gambler's fallacy dan dapat menafsirkan peluang dalam keputusan.
\end{itemize}
\end{refleksi}

% ---------- SOAL AKHIR BAB ----------
\section*{Soal Akhir Bab 8}
\addcontentsline{toc}{section}{Soal Akhir Bab 8}

\bagiansoal{A}{Fundamental (setara SNBT Dasar)}
\begin{enumerate}[leftmargin=*]
  \item Sebuah dadu dilempar. Peluang muncul angka 3?
  \item Peluang muncul Gambar pada satu koin?
  \item Dari 1--10, peluang bilangan genap?
  \item Jika $P(A)=0{,}25$, tentukan $P(A^c)$.
  \item Hitung $C(5,2)$.
  \item Hitung $4!$.
  \item Dua koin: peluang dua-duanya Angka?
  \item Dari 6 merah, 4 putih, peluang terambil putih (satu bola)?
  \item Hitung $P(4,2)$.
  \item Peluang dadu muncul faktor dari 6?
\end{enumerate}

\bagiansoal{B}{Intermediate (setara SNBT Sulit)}
\begin{enumerate}[leftmargin=*]
  \item Dua dadu dilempar. Peluang jumlah mata $\ge 10$?
  \item Dari 9 bola (5 merah, 4 biru) diambil 3 sekaligus. Peluang tepat 2 merah?
  \item Sebuah kotak berisi 3 bola merah, 2 putih. Diambil 2 berurutan tanpa pengembalian. Peluang keduanya merah?
  \item Tiga koin dilempar. Peluang muncul jumlah Gambar genap (0 atau 2)?
  \item Dari kata "STATISTIKA", berapa susunan berbeda yang dapat dibentuk?
\end{enumerate}

\bagiansoal{C}{Advanced (setara Olimpiade SMA --- gabungan dengan Statistika)}
\begin{enumerate}[leftmargin=*]
  \item Sebuah dadu adil dilempar. Tentukan nilai harapan (mean) dan variansi mata yang muncul (Peluang + Statistika).
  \item Dalam permainan, peluang menang $0{,}4$. Dimainkan 5 kali bebas. Peluang menang tepat 3 kali (gunakan binomial)?
  \item Tiga orang menembak target dengan peluang kena $0{,}6$, $0{,}7$, $0{,}8$ secara bebas. Peluang target terkena oleh tepat satu orang?
  \item Tunjukkan bahwa $\binom{n}{0}+\binom{n}{1}+\cdots+\binom{n}{n}=2^n$ dan kaitkan dengan banyaknya himpunan bagian.
  \item Sebuah peubah acak bernilai $1$ dengan peluang $p$ dan $0$ dengan peluang $1-p$. Tentukan mean dan variansinya dalam $p$, lalu cari $p$ yang memaksimalkan variansi.
\end{enumerate}

\bagiansoal{D}{Expert (setara tahun pertama universitas --- sintesis antarbab)}
\begin{enumerate}[leftmargin=*]
  \item \textbf{(Peluang + Statistika + Logaritma).} Hitung entropi $H=-\sum p_i\,{}^2\!\log p_i$ untuk pelemparan dadu adil dan jelaskan maknanya sebagai rata-rata informasi.
  \item \textbf{(Peluang bersyarat --- Bayes).} Suatu tes penyakit punya sensitivitas $99\%$ dan spesifisitas $95\%$; prevalensi penyakit $1\%$. Hitung peluang seseorang benar-benar sakit jika hasil tesnya positif, dan jelaskan mengapa hasilnya mengejutkan.
  \item \textbf{(Peluang + Deret).} Sebuah permainan dilempar koin sampai muncul Angka pertama. Hitung ekspektasi banyak lemparan $\sum n(\tfrac12)^n$ dan tunjukkan jumlah seluruh peluang $=1$.
  \item \textbf{(Peluang + Fungsi Kuadrat).} Tunjukkan bahwa $g(c)=E[(X-c)^2]$ minimum saat $c=E[X]$, dan nilai minimumnya adalah variansi.
  \item \textbf{(Peluang geometris).} Dua orang berjanji bertemu antara pukul 12.00--13.00, masing-masing datang acak dan menunggu 15 menit. Tentukan peluang mereka bertemu (gunakan luas daerah pada bidang).
\end{enumerate}

\begin{challenge}
\begin{enumerate}[leftmargin=*,topsep=2pt]
  \item \textbf{(Masalah ulang tahun).} Tentukan peluang minimal dua dari 23 orang berulang tahun di hari yang sama, dan jelaskan mengapa hasilnya melebihi $50\%$.
  \item \textbf{(Masalah Monty Hall).} Dalam kuis dengan 3 pintu (1 hadiah), setelah memilih, pembawa acara membuka satu pintu kosong lalu menawarkan pindah. Apakah sebaiknya pindah? Buktikan peluangnya.
  \item Sebuah semut berjalan acak pada titik-titik segitiga $ABC$, tiap langkah pindah ke salah satu titik tetangga dengan peluang sama. Tentukan peluang ia kembali ke titik awal setelah 3 langkah.
\end{enumerate}
\end{challenge}

% ============================================================
\chapter{Bunga Tunggal dan Bunga Majemuk}
% ============================================================
\begin{tujuan}
Setelah mempelajari bab ini, peserta didik mampu:
\begin{itemize}[leftmargin=*,topsep=2pt]
  \item membedakan konsep bunga tunggal dan bunga majemuk serta mengenali perbedaan pertumbuhannya;
  \item menghitung nilai akhir tabungan atau pinjaman dengan rumus bunga tunggal $M_n = M_0(1+ni)$;
  \item menghitung nilai akhir tabungan atau pinjaman dengan rumus bunga majemuk $M_n = M_0(1+i)^n$;
  \item menghubungkan bunga tunggal dengan barisan aritmetika dan bunga majemuk dengan barisan geometri;
  \item menerapkan konsep bunga dalam pengambilan keputusan finansial sehari-hari.
\end{itemize}
\textbf{Capaian Pembelajaran (Fase E).} Peserta didik dapat memodelkan situasi finansial nyata menggunakan deret aritmetika dan geometri, serta membandingkan pertumbuhan linear dengan pertumbuhan eksponensial.
\end{tujuan}

\begin{pemantik}
\begin{itemize}[leftmargin=*,topsep=2pt]
  \item Jika kamu menyimpan Rp~1.000.000 di bank, setelah 5 tahun apakah bank akan menghitung bunga dari uang aslimu saja, atau dari uang yang sudah bertumbuh?
  \item Mana yang lebih menguntungkan: bunga $10\%$ per tahun dihitung dari modal awal, atau $10\%$ per tahun dihitung dari saldo yang terus bertumbuh?
  \item Apakah kamu menyadari bahwa pinjaman kartu kredit dan tabungan di bank menggunakan \emph{mekanisme yang berbeda}? Mekanisme mana yang lebih "mahal"?
\end{itemize}
\end{pemantik}

\begin{studikasus}{Memilih Produk Tabungan}
Andi memiliki uang Rp~10.000.000 dan ingin menabung selama 3 tahun. Bank A menawarkan bunga \textbf{tunggal} $12\%$ per tahun. Bank B menawarkan bunga \textbf{majemuk} $12\%$ per tahun.

\smallskip
\begin{center}
\begin{tabular}{c|r|r}
\toprule
Tahun ke- & Bank A (Tunggal) & Bank B (Majemuk) \\
\midrule
0 & Rp 10.000.000 & Rp 10.000.000 \\
1 & Rp 11.200.000 & Rp 11.200.000 \\
2 & Rp 12.400.000 & Rp 12.544.000 \\
3 & Rp 13.600.000 & Rp 14.049.280 \\
\bottomrule
\end{tabular}
\end{center}

\smallskip
\textbf{Amati dan duga.} (a) Mengapa nilai akhir Bank B lebih besar meski suku bunganya sama? (b) Apa pola pertumbuhannya --- konstan atau bertumbuh semakin cepat? (c) Jika menabung 10 tahun, selisihnya akan semakin besar atau kecil? Inilah dua sistem bunga yang menjadi fondasi seluruh produk keuangan: \textbf{bunga tunggal} (pertumbuhan linear) dan \textbf{bunga majemuk} (pertumbuhan eksponensial).
\end{studikasus}

\section*{Peta Konsep Bab}
\addcontentsline{toc}{section}{Peta Konsep Bab}
\begin{center}
\begin{tikzpicture}[
  node distance=8mm,
  konsepbox/.style={draw=biru,fill=birumuda,rounded corners,align=center,inner sep=5pt,font=\small\bfseries},
  subbox/.style={draw=hijau,fill=hijaumuda,rounded corners,align=center,inner sep=4pt,font=\footnotesize},
  garis/.style={-Stealth,biru,thick}]
  \node[konsepbox] (akar) {BUNGA};
  \node[subbox,below left=12mm and 18mm of akar] (tung) {Bunga\\Tunggal};
  \node[subbox,below right=12mm and 18mm of akar] (maj) {Bunga\\Majemuk};
  \node[subbox,below=10mm of tung] (bt1) {$M_n=M_0(1+ni)$\\linear};
  \node[subbox,below=10mm of maj] (bm1) {$M_n=M_0(1+i)^n$\\eksponensial};
  \node[subbox,below=10mm of bt1] (bar) {Barisan\\Aritmetika};
  \node[subbox,below=10mm of bm1] (geo) {Barisan\\Geometri};
  \draw[garis] (akar)--(tung); \draw[garis] (akar)--(maj);
  \draw[garis] (tung)--(bt1); \draw[garis] (maj)--(bm1);
  \draw[garis] (bt1)--(bar); \draw[garis] (bm1)--(geo);
\end{tikzpicture}\\
\small \textbf{Peta konsep.} Bunga tunggal tumbuh secara linear (barisan aritmetika); bunga majemuk tumbuh secara eksponensial (barisan geometri). Keduanya memiliki rumus berbeda dengan implikasi jangka panjang yang sangat berbeda.
\end{center}

\section{Konsep Dasar Bunga}

\begin{definisi}[Modal, Bunga, dan Suku Bunga]
\begin{itemize}[leftmargin=*,topsep=2pt]
  \item \textbf{Modal awal} ($M_0$): jumlah uang yang ditabung atau dipinjam pada awal periode.
  \item \textbf{Bunga} ($B$): imbalan yang diterima/dibayarkan atas penggunaan uang selama satu periode.
  \item \textbf{Suku bunga} ($i$): rasio bunga terhadap modal, dinyatakan dalam persen per periode (misalnya per tahun, per bulan).
  \item \textbf{Modal akhir} ($M_n$): jumlah uang setelah $n$ periode, yaitu $M_n = M_0 + \text{total bunga}$.
\end{itemize}
\end{definisi}

\begin{konsep}{Periode Bunga}
Suku bunga selalu dikaitkan dengan \emph{satuan waktu}. Jika suku bunga $12\%$ per tahun dan dihitung per bulan, maka suku bunga per bulan adalah $i = \frac{12\%}{12} = 1\%$ per bulan. Perhatikan selalu kecocokan satuan waktu antara $i$ dan $n$.
\end{konsep}

\section{Bunga Tunggal}

Pada sistem \textbf{bunga tunggal}, bunga dihitung \emph{hanya dari modal awal} setiap periode. Besar bunga setiap periode sama, sehingga nilai akhir tumbuh secara \textbf{linear}.

\begin{konsep}{Rumus Bunga Tunggal}
Jika modal awal $M_0$, suku bunga $i$ per periode, dan $n$ adalah jumlah periode, maka:
\[
  \text{Bunga per periode} = M_0 \cdot i
\]
\[
  \text{Total bunga setelah } n \text{ periode} = M_0 \cdot n \cdot i
\]
\[
  \boxed{M_n = M_0(1 + ni)}
\]
Nilai $M_0,\ M_0(1+i),\ M_0(1+2i),\ \ldots$ membentuk \textbf{barisan aritmetika} dengan beda $M_0 \cdot i$.
\end{konsep}

\begin{contoh}{Tabungan Bunga Tunggal}
Budi menabung Rp~5.000.000 dengan bunga tunggal $8\%$ per tahun selama 4 tahun. Berapa nilai akhir tabungannya?

\jawab $M_0 = 5.000.000$, $i = 0{,}08$, $n = 4$.
\[
  M_4 = 5.000.000 \times (1 + 4 \times 0{,}08) = 5.000.000 \times 1{,}32 = \text{Rp}\ 6.600.000.
\]
Total bunga yang diterima: $6.600.000 - 5.000.000 = \text{Rp}\ 1.600.000$.
\end{contoh}

\begin{contoh}{Mencari Modal Awal}
Setelah 2 tahun dengan bunga tunggal $10\%$ per tahun, nilai tabungan menjadi Rp~12.000.000. Berapa modal awalnya?

\jawab $M_n = M_0(1+ni) \Rightarrow 12.000.000 = M_0(1 + 2 \times 0{,}1) = 1{,}2\,M_0$.
\[
  M_0 = \frac{12.000.000}{1{,}2} = \text{Rp}\ 10.000.000.
\]
\end{contoh}

\begin{contoh}{Mencari Lama Waktu}
Modal Rp~8.000.000 ditabung dengan bunga tunggal $5\%$ per tahun. Berapa tahun hingga nilai tabungan mencapai Rp~10.000.000?

\jawab $10.000.000 = 8.000.000(1 + 0{,}05\,n) \Rightarrow 1{,}25 = 1 + 0{,}05n \Rightarrow n = \dfrac{0{,}25}{0{,}05} = 5$ tahun.
\end{contoh}

\begin{kesalahan}
Pada bunga tunggal, bunga \emph{tidak bertambah} ke modal di akhir setiap periode. Jadi pada tahun ke-2, bunga dihitung dari $M_0$, bukan dari $M_1 = M_0(1+i)$. Ini berbeda dengan bunga majemuk.
\end{kesalahan}

\begin{latihan}
\begin{enumerate}[leftmargin=*,topsep=2pt]
  \item Rp~6.000.000 ditabung dengan bunga tunggal $9\%$ per tahun selama 3 tahun. Hitung nilai akhirnya.
  \item Setelah 5 tahun dengan bunga tunggal $6\%$ per tahun, nilai tabungan adalah Rp~13.000.000. Berapa modal awalnya?
  \item Pada suku bunga tunggal $8\%$ per tahun, berapa lama hingga uang Rp~4.000.000 menjadi Rp~5.280.000?
\end{enumerate}
\end{latihan}

\section{Bunga Majemuk}

Pada sistem \textbf{bunga majemuk}, bunga setiap periode \emph{ditambahkan ke modal}, sehingga periode berikutnya bunga dihitung dari modal yang sudah bertumbuh. Pertumbuhannya bersifat \textbf{eksponensial}.

\begin{konsep}{Rumus Bunga Majemuk}
Jika modal awal $M_0$, suku bunga $i$ per periode (majemuk), dan $n$ adalah jumlah periode, maka:
\[
  M_1 = M_0(1+i),\quad M_2 = M_1(1+i) = M_0(1+i)^2,\quad \ldots
\]
\[
  \boxed{M_n = M_0(1+i)^n}
\]
Nilai $M_0,\,M_0(1+i),\,M_0(1+i)^2,\,\ldots$ membentuk \textbf{barisan geometri} dengan rasio $(1+i)$.
\end{konsep}

\begin{contoh}{Tabungan Bunga Majemuk}
Dina menabung Rp~5.000.000 dengan bunga majemuk $8\%$ per tahun selama 4 tahun. Hitung nilai akhir tabungannya.

\jawab $M_0 = 5.000.000$, $i = 0{,}08$, $n = 4$.
\[
  M_4 = 5.000.000 \times (1{,}08)^4 = 5.000.000 \times 1{,}36049 \approx \text{Rp}\ 6.802.450.
\]
Bandingkan: bunga tunggal memberikan Rp~6.600.000. Selisih $\approx$ Rp~202.450 karena efek \emph{bunga berbunga}.
\end{contoh}

\begin{contoh}{Bunga Majemuk Per Bulan}
Modal Rp~3.000.000 didepositokan dengan bunga majemuk $12\%$ per tahun, diperhitungkan per bulan, selama 6 bulan. Hitung nilai akhirnya.

\jawab Suku bunga per bulan: $i = \frac{12\%}{12} = 1\% = 0{,}01$. Jumlah periode: $n = 6$ bulan.
\[
  M_6 = 3.000.000 \times (1{,}01)^6 \approx 3.000.000 \times 1{,}06152 \approx \text{Rp}\ 3.184.560.
\]
\end{contoh}

\begin{contoh}{Waktu Penggandaan Modal}
Pada bunga majemuk $6\%$ per tahun, berapa tahun hingga modal berlipat dua?

\jawab Cari $n$ sehingga $M_n = 2M_0$, yaitu $(1{,}06)^n = 2$.
\[
  n = \frac{\log 2}{\log 1{,}06} = \frac{0{,}3010}{0{,}02531} \approx 11{,}9 \approx 12 \text{ tahun}.
\]
\textit{Aturan 72}: estimasi cepat $n \approx \frac{72}{i\%} = \frac{72}{6} = 12$ tahun (sangat tepat!).
\end{contoh}

\begin{konsep}{Aturan 72 (Rule of 72)}
Pada bunga majemuk, estimasi cepat \emph{waktu penggandaan} modal:
\[
  n \approx \frac{72}{i\%}
\]
Contoh: bunga $9\%$ per tahun $\Rightarrow$ modal berlipat dua dalam $\approx 8$ tahun.
\end{konsep}

\begin{kesalahan}
Jangan mencampur satuan: jika $i = 12\%$ per \emph{tahun} dan $n$ dalam \emph{bulan}, hasilnya salah. Selalu konversikan: $i_{\text{bulan}} = \frac{i_{\text{tahun}}}{12}$ lalu gunakan $n$ dalam bulan.
\end{kesalahan}

\begin{latihan}
\begin{enumerate}[leftmargin=*,topsep=2pt]
  \item Rp~4.000.000 ditabung dengan bunga majemuk $10\%$ per tahun selama 3 tahun. Hitung nilai akhirnya (gunakan $(1{,}1)^3 = 1{,}331$).
  \item Modal Rp~2.000.000 didepositokan dengan bunga majemuk $6\%$ per tahun (per bulan) selama 1 tahun. Hitung nilai akhirnya.
  \item Pada bunga majemuk $8\%$ per tahun, berapa tahun hingga modal berlipat dua? (Gunakan $\log 2 = 0{,}301$, $\log 1{,}08 = 0{,}0334$.)
\end{enumerate}
\end{latihan}

\section{Perbandingan Bunga Tunggal dan Bunga Majemuk}

\begin{center}
\begin{tabular}{>{\bfseries}p{3.8cm} p{5cm} p{5cm}}
\toprule
Aspek & Bunga Tunggal & Bunga Majemuk \\
\midrule
Rumus & $M_n = M_0(1+ni)$ & $M_n = M_0(1+i)^n$ \\
Pola pertumbuhan & Linear (barisan aritmetika) & Eksponensial (barisan geometri) \\
Dasar perhitungan bunga & Selalu dari $M_0$ & Dari modal yang bertumbuh \\
Jangka pendek ($n$ kecil) & Hampir sama & Hampir sama \\
Jangka panjang ($n$ besar) & Jauh lebih kecil & Jauh lebih besar \\
Penggunaan umum & Pinjaman jangka pendek & Deposito, investasi, KPR \\
\bottomrule
\end{tabular}
\end{center}

\begin{dunianyata}
\begin{itemize}[leftmargin=*,topsep=2pt]
  \item \textbf{Perbankan:} tabungan dan deposito umumnya menggunakan bunga majemuk; beberapa pinjaman jangka pendek memakai bunga tunggal.
  \item \textbf{Investasi:} konsep \emph{compound interest} adalah kunci kekayaan jangka panjang --- Einstein konon menyebutnya "keajaiban dunia kedelapan."
  \item \textbf{Inflasi:} kenaikan harga juga bersifat majemuk; Rp~100.000 di tahun 2000 tidak setara dengan Rp~100.000 sekarang.
  \item \textbf{Pinjaman:} kartu kredit dan KPR menggunakan bunga majemuk sehingga total yang dibayar jauh melebihi pokok pinjaman.
\end{itemize}
\end{dunianyata}

\section*{Ringkasan Bab}
\addcontentsline{toc}{section}{Ringkasan Bab}
\begin{center}
\begin{tabular}{>{\bfseries}p{3.5cm} p{10.2cm}}
\toprule
Konsep & Inti yang harus diingat \\
\midrule
Modal awal $M_0$ & Jumlah uang awal; $i$ = suku bunga per periode; $n$ = jumlah periode. \\
Bunga tunggal & $M_n = M_0(1+ni)$; tumbuh linear; barisan aritmetika, beda $M_0i$. \\
Bunga majemuk & $M_n = M_0(1+i)^n$; tumbuh eksponensial; barisan geometri, rasio $(1+i)$. \\
Konversi satuan & Jika $i$ tahunan, $n$ harus dalam tahun; jika per bulan, bagi $i$ dengan 12. \\
Aturan 72 & Waktu penggandaan $\approx 72/i\%$ (estimasi cepat bunga majemuk). \\
Waktu & Cari $n$ dengan logaritma: $n = \log(M_n/M_0)/\log(1+i)$ (majemuk). \\
\bottomrule
\end{tabular}
\end{center}

\begin{refleksi}
Beri tanda centang pada hal yang sudah kamu kuasai:
\begin{itemize}[leftmargin=*,label={},topsep=2pt]
  \item \cek Saya dapat membedakan bunga tunggal dan bunga majemuk secara konseptual.
  \item \cek Saya dapat menghitung nilai akhir tabungan/pinjaman dengan kedua rumus.
  \item \cek Saya dapat mencari modal awal, suku bunga, atau lama waktu dari informasi yang diberikan.
  \item \cek Saya memahami hubungan bunga tunggal dengan barisan aritmetika dan bunga majemuk dengan barisan geometri.
\end{itemize}
\end{refleksi}

% ---------- SOAL AKHIR BAB ----------
\section*{Soal Akhir Bab \thechapter}
\addcontentsline{toc}{section}{Soal Akhir Bab \thechapter}

\bagiansoal{A}{Fundamental (setara SNBT Dasar)}
\begin{enumerate}[leftmargin=*]
  \item Rp~2.000.000 ditabung dengan bunga tunggal $10\%$/tahun selama 3 tahun. Nilai akhirnya?
  \item Tentukan bunga total yang diperoleh dari modal Rp~5.000.000, bunga tunggal $6\%$/tahun, selama 4 tahun.
  \item Modal Rp~4.000.000 ditabung dengan bunga majemuk $5\%$/tahun selama 2 tahun ($(1{,}05)^2=1{,}1025$). Nilai akhirnya?
  \item Suku bunga tunggal berapa per tahun agar modal Rp~8.000.000 menjadi Rp~9.600.000 dalam 2 tahun?
  \item Modal Rp~3.000.000 ditabung dengan bunga tunggal $12\%$/tahun. Berapa tahun hingga menjadi Rp~4.080.000?
  \item Modal Rp~1.000.000 dengan bunga majemuk $10\%$/tahun. Nilai setelah 3 tahun? ($(1{,}1)^3=1{,}331$)
  \item Pada bunga tunggal, nilai akhir setelah $n$ tahun membentuk barisan jenis apa? Sebutkan bedanya.
  \item Modal Rp~6.000.000 didepositokan dengan bunga majemuk $12\%$/tahun (per bulan). Suku bunga per bulannya berapa?
  \item Dengan bunga tunggal $15\%$/tahun, berapa besar modal awal agar setelah 2 tahun nilai tabungan Rp~11.500.000?
  \item Bandingkan hasil bunga tunggal dan majemuk $10\%$/tahun untuk modal Rp~1.000.000 setelah 1 tahun.
\end{enumerate}

\bagiansoal{B}{Intermediate (setara SNBT Sulit)}
\begin{enumerate}[leftmargin=*]
  \item Modal Rp~10.000.000 ditabung dengan bunga majemuk $8\%$/tahun selama 5 tahun. Hitung nilai akhirnya dan selisihnya dengan bunga tunggal $8\%$. ($(1{,}08)^5 \approx 1{,}4693$)
  \item Suatu tabungan setelah 3 tahun dengan bunga majemuk $10\%$/tahun bernilai Rp~13.310.000. Berapa modal awalnya?
  \item Dua investasi: A memberikan bunga tunggal $15\%$/tahun, B memberikan bunga majemuk $12\%$/tahun. Modal yang sama diinvestasikan 5 tahun. Mana yang lebih menguntungkan? ($(1{,}12)^5 \approx 1{,}7623$)
  \item Pada bunga majemuk $6\%$/tahun, berapa lama hingga modal berlipat tiga? ($\log 3 = 0{,}4771$, $\log 1{,}06 = 0{,}0253$)
  \item Modal Rp~5.000.000 ditabung dengan bunga majemuk $12\%$/tahun diperhitungkan per triwulan (3 bulan). Hitung nilai akhir setelah 1 tahun. ($(1{,}03)^4 \approx 1{,}1255$)
\end{enumerate}

\bagiansoal{C}{Advanced (setara Olimpiade SMA --- gabungan dengan Barisan/Deret \& Eksponen)}
\begin{enumerate}[leftmargin=*]
  \item \textbf{(Bunga + Deret).} Setiap awal tahun Rina menabung Rp~1.000.000 dengan bunga majemuk $10\%$/tahun. Tulis ekspresi nilai total tabungannya setelah 3 tahun (hitung sebagai jumlah deret geometri) dan hitung nilainya. ($(1{,}1)^3=1{,}331$, $(1{,}1)^2=1{,}21$)
  \item \textbf{(Bunga + Eksponen).} Tunjukkan bahwa pada bunga majemuk, rasio $M_{n+1}/M_n$ konstan, dan hubungkan hal ini dengan definisi barisan geometri.
  \item \textbf{(Bunga + Logaritma).} Suku bunga majemuk bank A adalah $6\%$/tahun dan bank B adalah $8\%$/tahun. Selisih waktu yang dibutuhkan untuk menggandakan modal antara bank A dan B adalah berapa tahun? ($\log 2=0{,}301$, $\log 1{,}06=0{,}0253$, $\log 1{,}08=0{,}0334$)
  \item \textbf{(Perbandingan dua sistem).} Buktikan bahwa untuk $n \ge 1$ dan $i > 0$, selalu berlaku $(1+i)^n > 1+ni$, sehingga bunga majemuk selalu lebih besar dari bunga tunggal (untuk $n > 1$). (Gunakan induksi atau ekspansi binomial.)
  \item \textbf{(Bunga + Fungsi).} Modal $M_0$ ditabung $n$ tahun. Definisikan $f(n) = M_0(1+i)^n - M_0(1+ni)$ (selisih majemuk minus tunggal). Tunjukkan $f(0)=f(1)=0$ dan $f(n)>0$ untuk $n\ge2$, lalu hitung $f(3)$ untuk $M_0=1$ dan $i=0{,}1$.
\end{enumerate}

\bagiansoal{D}{Expert (setara tahun pertama universitas --- sintesis antarbab)}
\begin{enumerate}[leftmargin=*]
  \item \textbf{(Nilai Sekarang / Present Value).} Berapa modal yang harus ditabung sekarang (dengan bunga majemuk $8\%$/tahun) agar 5 tahun mendatang nilainya Rp~14.693.000? ($(1{,}08)^5 \approx 1{,}4693$). Konsep ini disebut \emph{nilai sekarang} (present value); jelaskan maknanya.
  \item \textbf{(Inflasi dan Daya Beli).} Harga barang naik $5\%$/tahun (inflasi majemuk). Jika harga sekarang Rp~100.000, tentukan harga 10 tahun mendatang ($(1{,}05)^{10}\approx1{,}6289$) dan berapa tahun hingga harga berlipat dua.
  \item \textbf{(Anuitas Sederhana).} Setiap akhir tahun Budi membayar cicilan Rp~2.000.000 selama 3 tahun. Berapa nilai total cicilan ini pada awal tahun pertama jika bunga majemuk $10\%$/tahun? (Nilai sekarang anuitas: $A = \sum_{k=1}^{n} \frac{C}{(1+i)^k}$)
  \item \textbf{(Optimasi Frekuensi Majemuk).} Bunga nominal $12\%$/tahun dapat diperhitungkan per tahun, per triwulan, atau per bulan. Hitung nilai akhir Rp~1.000.000 setelah 1 tahun untuk ketiganya dan bandingkan. (($(1{,}03)^4\approx1{,}1255$, $(1{,}01)^{12}\approx1{,}1268$))
  \item \textbf{(Penerapan Logaritma Natural).} Pada limit frekuensi majemuk tak terhingga, nilai akhir modal $M_0$ setelah $t$ tahun dengan suku bunga nominal $r$ adalah $M_0 e^{rt}$ (bunga kontinu). Bandingkan $e^{0{,}12}$ dengan hasil per-bulan dan per-tahun untuk $r=12\%$, $t=1$ tahun. ($e^{0{,}12}\approx1{,}1275$)
\end{enumerate}

\begin{challenge}
\begin{enumerate}[leftmargin=*,topsep=2pt]
  \item \textbf{(Paradoks Bunga Majemuk).} Suatu investasi dengan bunga majemuk $100\%$/tahun diperhitungkan $n$ kali per tahun menghasilkan faktor $\left(1+\frac{1}{n}\right)^n$. Tunjukkan (secara numerik untuk $n=1,2,4,12,365$) bahwa nilainya mendekati $e \approx 2{,}718$ dan tidak pernah melampaui $3$. Apa implikasi praktisnya?
  \item \textbf{(Pertumbuhan Penduduk vs Inflasi).} Penduduk kota tumbuh $3\%$/tahun dan inflasi $5\%$/tahun. Berapa tahun hingga pendapatan per kapita (asumsikan konstan) mengalami penurunan daya beli $20\%$ secara riil? Modelkan menggunakan rasio $\frac{(1{,}03)^n}{(1{,}05)^n}$ dan cari $n$ sehingga rasio ini $= 0{,}8$.
  \item \textbf{(Kredit dengan Bunga Berbunga).} Seseorang meminjam Rp~10.000.000 dengan bunga majemuk $2\%$/bulan dan tidak membayar cicilan selama 6 bulan. Setelah itu ia ingin melunasi sekaligus. Berapa yang harus dibayar? ($(1{,}02)^6\approx1{,}1262$). Bandingkan dengan jika bunga tunggal $2\%$/bulan.
\end{enumerate}
\end{challenge}

% ============================================================
\part*{\centering KUNCI JAWABAN DAN PEMBAHASAN}
\addcontentsline{toc}{part}{KUNCI JAWABAN DAN PEMBAHASAN}
% ============================================================
\noindent\textit{Setiap pembahasan disusun dengan tiga bagian: \textbf{Soal} (pernyataan lengkap), \textbf{Analisis} (strategi dan konsep kunci), dan \textbf{Pembahasan} (langkah penyelesaian). Gunakan kunci ini untuk mengoreksi diri, bukan sekadar mencocokkan hasil.}

% ============================================================
\chapter{Pembahasan Bab 1: Eksponen dan Logaritma}
% ============================================================

\kbagian{Bagian A}

\knomor{A.1}
\ksoal Sederhanakan $\dfrac{a^5 b^2}{a^2 b^5}$.
\kanalisis Gunakan sifat pembagian pangkat $\dfrac{a^m}{a^n}=a^{m-n}$ pada tiap basis.
\kbahas $\dfrac{a^5 b^2}{a^2 b^5}=a^{5-2}b^{2-5}=a^3 b^{-3}=\dfrac{a^3}{b^3}.$

\knomor{A.2}
\ksoal Hitung $32^{2/5}$.
\kanalisis Nyatakan basis sebagai pangkat prima, lalu pakai $(a^m)^n=a^{mn}$.
\kbahas $32=2^5$, sehingga $32^{2/5}=(2^5)^{2/5}=2^{2}=4.$

\knomor{A.3}
\ksoal Tentukan nilai $x$ dari $3^{x}=27$.
\kanalisis Samakan basis kedua ruas.
\kbahas $27=3^3$, maka $3^x=3^3\Rightarrow x=3.$

\knomor{A.4}
\ksoal Hitung ${}^{2}\!\log 32$.
\kanalisis Logaritma menanyakan "pangkat dari basis"; tulis $32$ sebagai pangkat $2$.
\kbahas ${}^{2}\!\log 32={}^{2}\!\log 2^5=5.$

\knomor{A.5}
\ksoal Sederhanakan $\dfrac{x^{-2}y^3}{x^4 y^{-1}}$.
\kanalisis Kurangkan pangkat tiap basis; perhatikan tanda pada pangkat negatif.
\kbahas $x^{-2-4}y^{3-(-1)}=x^{-6}y^{4}=\dfrac{y^4}{x^6}.$

\knomor{A.6}
\ksoal Tentukan HP dari $2^{x+1}=64$.
\kanalisis Samakan basis; $64$ adalah pangkat $2$.
\kbahas $64=2^6$, maka $x+1=6\Rightarrow x=5.$ HP $=\{5\}$.

\knomor{A.7}
\ksoal Hitung ${}^{3}\!\log 9 + {}^{2}\!\log 8$.
\kanalisis Hitung tiap logaritma secara terpisah dengan menyatakan argumen sebagai pangkat basisnya.
\kbahas ${}^{3}\!\log 3^2 + {}^{2}\!\log 2^3 = 2+3 = 5.$

\knomor{A.8}
\ksoal Nyatakan $\sqrt[3]{a^2}$ dalam bentuk pangkat.
\kanalisis Akar ke-$n$ setara pangkat $1/n$: $\sqrt[n]{a^m}=a^{m/n}$.
\kbahas $\sqrt[3]{a^2}=a^{2/3}.$

\knomor{A.9}
\ksoal Jika ${}^{2}\!\log 5 = q$, nyatakan ${}^{2}\!\log 20$ dalam $q$.
\kanalisis Faktorkan $20=4\cdot5=2^2\cdot5$, lalu pakai sifat penjumlahan logaritma.
\kbahas ${}^{2}\!\log 20={}^{2}\!\log(2^2\cdot5)={}^{2}\!\log 2^2+{}^{2}\!\log 5=2+q.$

\knomor{A.10}
\ksoal Hitung nilai $\left(\tfrac{2}{3}\right)^{-2}$.
\kanalisis Pangkat negatif membalik pecahan: $\left(\tfrac ab\right)^{-n}=\left(\tfrac ba\right)^{n}$.
\kbahas $\left(\tfrac23\right)^{-2}=\left(\tfrac32\right)^{2}=\dfrac94.$

\kbagian{Bagian B}

\knomor{B.1}
\ksoal Selesaikan $4^{x}-6\cdot 2^{x}+8=0$.
\kanalisis Bentuk $4^x=(2^x)^2$; misalkan $p=2^x>0$ agar menjadi persamaan kuadrat.
\kbahas $p^2-6p+8=0\Rightarrow(p-2)(p-4)=0\Rightarrow p=2$ atau $p=4$. Maka $2^x=2\Rightarrow x=1$ atau $2^x=4\Rightarrow x=2$. HP $=\{1,2\}$.

\knomor{B.2}
\ksoal Diketahui ${}^{a}\!\log 2 = m$ dan ${}^{a}\!\log 3 = n$. Nyatakan ${}^{a}\!\log 0{,}75$ dalam $m$ dan $n$.
\kanalisis Tulis $0{,}75=\tfrac34=\dfrac{3}{2^2}$, lalu pakai sifat selisih dan pangkat.
\kbahas ${}^{a}\!\log\tfrac{3}{2^2}={}^{a}\!\log 3-2\,{}^{a}\!\log 2=n-2m.$

\knomor{B.3}
\ksoal Tentukan himpunan penyelesaian $\left(\tfrac19\right)^{x-1} \ge 3^{x+2}$.
\kanalisis Samakan basis menjadi $3$. Karena basis $3>1$, arah tanda dipertahankan.
\kbahas $\left(3^{-2}\right)^{x-1}=3^{-2x+2}\ge 3^{x+2}\Rightarrow -2x+2\ge x+2\Rightarrow -3x\ge0\Rightarrow x\le0.$ HP $=\{x\mid x\le0\}$.

\knomor{B.4}
\ksoal Jika $2^x+2^{-x}=3$, tentukan nilai $4^x+4^{-x}$.
\kanalisis Kuadratkan kedua ruas; perhatikan $4^x=(2^x)^2$ dan suku silang $2^x\cdot2^{-x}=1$.
\kbahas $(2^x+2^{-x})^2=4^x+2+4^{-x}=9\Rightarrow 4^x+4^{-x}=7.$

\knomor{B.5}
\ksoal Penduduk sebuah kota tumbuh $2\%$ per tahun. Jika kini $500.000$ jiwa, susun dan selesaikan persamaan untuk memperkirakan tahun saat penduduk mencapai $1.000.000$ jiwa.
\kanalisis Model pertumbuhan eksponen $P(t)=P_0(1+i)^t$; selesaikan untuk $t$ dengan logaritma.
\kbahas $500000(1{,}02)^t=1000000\Rightarrow 1{,}02^t=2\Rightarrow t=\dfrac{\log 2}{\log 1{,}02}\approx\dfrac{0{,}30103}{0{,}0086}\approx 35.$ Sekitar \textbf{35 tahun}.

\kbagian{Bagian C}

\knomor{C.1}
\ksoal Barisan $a, ar, ar^2, \dots$ adalah geometri. Tunjukkan bahwa $\log a, \log(ar), \log(ar^2),\dots$ membentuk barisan aritmetika, dan tentukan bedanya.
\kanalisis Logaritma mengubah perkalian menjadi penjumlahan, sehingga rasio tetap menjadi beda tetap.
\kbahas Suku ke-$k$: $\log(ar^{k-1})=\log a+(k-1)\log r$. Ini berbentuk $A+(k-1)B$ dengan $A=\log a$, $B=\log r$ konstan, jadi aritmetika dengan \textbf{beda $=\log r$}.

\knomor{C.2}
\ksoal Jika ${}^{2}\!\log x$, ${}^{2}\!\log(x+2)$, dan ${}^{2}\!\log(x+6)$ membentuk barisan aritmetika, tentukan $x$.
\kanalisis Syarat aritmetika: suku tengah = rata-rata dua suku tepi, yakni $2U_2=U_1+U_3$.
\kbahas $2\,{}^{2}\!\log(x+2)={}^{2}\!\log x+{}^{2}\!\log(x+6)\Rightarrow {}^{2}\!\log(x+2)^2={}^{2}\!\log\big(x(x+6)\big)$. Maka $(x+2)^2=x(x+6)\Rightarrow x^2+4x+4=x^2+6x\Rightarrow 4=2x\Rightarrow x=2.$

\knomor{C.3}
\ksoal Diketahui $x^{\log x}=100x$ (log basis 10). Tentukan semua nilai $x$.
\kanalisis Ambil $\log_{10}$ kedua ruas dan misalkan $u=\log x$ agar menjadi kuadrat.
\kbahas $\log\big(x^{\log x}\big)=\log(100x)\Rightarrow (\log x)^2=2+\log x$. Misal $u=\log x$: $u^2-u-2=0\Rightarrow(u-2)(u+1)=0\Rightarrow u=2$ atau $u=-1$. Maka $x=100$ atau $x=0{,}1$.

\knomor{C.4}
\ksoal Hasil kali $n$ suku pertama barisan geometri positif adalah $P$. Jika ${}^{2}\!\log U_1 + {}^{2}\!\log U_n = 10$, tentukan ${}^{2}\!\log P$ dalam $n$.
\kanalisis Hasil kali suku geometri bersifat simetris: $U_k\cdot U_{n+1-k}=U_1 U_n$ tetap. Gunakan sifat logaritma hasil kali.
\kbahas ${}^{2}\!\log P={}^{2}\!\log(U_1U_2\cdots U_n)=\sum_{k=1}^{n}{}^{2}\!\log U_k$. Karena ${}^{2}\!\log U_k$ aritmetika, jumlahnya $=\dfrac n2\big({}^{2}\!\log U_1+{}^{2}\!\log U_n\big)=\dfrac n2(10)=5n.$

\knomor{C.5}
\ksoal Tentukan jumlah ${}^{2}\!\log\tfrac{2}{1}+{}^{2}\!\log\tfrac{3}{2}+\cdots+{}^{2}\!\log\tfrac{64}{63}$.
\kanalisis Sifat penjumlahan logaritma mengubah jumlah menjadi logaritma hasil kali yang teleskopik.
\kbahas Jumlah $={}^{2}\!\log\left(\tfrac21\cdot\tfrac32\cdots\tfrac{64}{63}\right)={}^{2}\!\log\tfrac{64}{1}={}^{2}\!\log 64=6.$

\kbagian{Bagian D}

\knomor{D.1}
\ksoal \textbf{(Eksponen + Fungsi).} Diberikan $f(x)=\dfrac{a^x}{a^x+\sqrt a}$ dengan $a>0$. Buktikan $f(x)+f(1-x)=1$, lalu hitung $\sum_{k=1}^{2026} f\!\left(\tfrac{k}{2027}\right)$.
\kanalisis Hitung $f(1-x)$ dan sederhanakan; manfaatkan simetri untuk menjumlahkan secara berpasangan.
\kbahas $f(1-x)=\dfrac{a^{1-x}}{a^{1-x}+\sqrt a}=\dfrac{a}{a+\sqrt a\,a^{x}}=\dfrac{\sqrt a}{\sqrt a+a^{x}}$ (bagi pembilang dan penyebut dengan $a^{x}\cdot$, lalu dengan $\sqrt a$). Maka $f(x)+f(1-x)=\dfrac{a^x}{a^x+\sqrt a}+\dfrac{\sqrt a}{a^x+\sqrt a}=1.$ Untuk jumlah, pasangkan $k$ dengan $2027-k$: tiap pasang bernilai $1$. Ada $2026$ suku $=1013$ pasang (tanpa suku tengah karena $2027$ ganjil), sehingga jumlah $=1013.$

\knomor{D.2}
\ksoal \textbf{(Logaritma + Peluang).} Entropi $H=-\sum p_i\,{}^{2}\!\log p_i$. Hitung $H$ untuk pelemparan dua koin adil, dan jelaskan maknanya.
\kanalisis Dua koin punya $4$ keluaran sama mungkin, masing-masing $p=\tfrac14$.
\kbahas $H=-\sum_{i=1}^{4}\tfrac14\,{}^{2}\!\log\tfrac14=-4\cdot\tfrac14\cdot(-2)=2$ bit. Maknanya: diperlukan rata-rata $2$ bit informasi untuk menyatakan hasil percobaan (sesuai $2^2=4$ kemungkinan).

\knomor{D.3}
\ksoal \textbf{(Eksponen + Kalkulus dasar).} Fungsi peluruhan $N(t)=N_0 e^{-kt}$ dengan waktu paruh $T$. Nyatakan $k$ dalam $T$, lalu tentukan kapan tersisa $\tfrac18 N_0$.
\kanalisis Waktu paruh berarti $N(T)=\tfrac12N_0$; gunakan logaritma natural.
\kbahas $\tfrac12N_0=N_0e^{-kT}\Rightarrow e^{-kT}=\tfrac12\Rightarrow kT=\ln2\Rightarrow k=\dfrac{\ln2}{T}.$ Untuk $N=\tfrac18N_0$: $e^{-kt}=\tfrac18\Rightarrow kt=\ln8=3\ln2\Rightarrow t=\dfrac{3\ln2}{k}=3T.$

\knomor{D.4}
\ksoal \textbf{(Logaritma + Statistika).} Jika $x_1,\dots,x_n$ punya rata-rata geometrik $G$, buktikan bahwa rata-rata aritmetika dari $\log x_i$ sama dengan $\log G$.
\kanalisis Definisi rata-rata geometrik $G=\sqrt[n]{x_1\cdots x_n}$; ambil logaritma.
\kbahas $\log G=\log\Big(\prod x_i\Big)^{1/n}=\dfrac1n\sum_{i=1}^{n}\log x_i=\overline{\log x}.$ Jadi rata-rata aritmetika dari $\log x_i$ tepat sama dengan $\log G$.

\knomor{D.5}
\ksoal \textbf{(Eksponen + Pembuktian).} Buktikan dengan induksi bahwa $2^n > n^2$ untuk setiap bilangan asli $n\ge 5$.
\kanalisis Induksi matematika: cek basis $n=5$, lalu langkah induktif memakai $2^{k+1}=2\cdot2^k$.
\kbahas \emph{Basis:} $n=5$: $2^5=32>25=5^2$ benar. \emph{Langkah:} andaikan $2^k>k^2$ untuk suatu $k\ge5$. Maka $2^{k+1}=2\cdot2^k>2k^2$. Cukup tunjukkan $2k^2\ge(k+1)^2$, yakni $k^2-2k-1\ge0$, benar untuk $k\ge5$ (sebab $k^2\ge2k+1$). Jadi $2^{k+1}>(k+1)^2$. Berdasarkan induksi, pernyataan benar untuk semua $n\ge5$. $\blacksquare$

\kbagian{Challenge Corner}

\knomor{Challenge 1}
\ksoal Tentukan banyaknya digit dari $2^{100}$.
\kanalisis Banyak digit bilangan $N$ adalah $\lfloor \log_{10} N\rfloor+1$.
\kbahas $\log_{10}2^{100}=100\log_{10}2\approx100(0{,}30103)=30{,}103$. Banyak digit $=\lfloor30{,}103\rfloor+1=31.$

\knomor{Challenge 2}
\ksoal Selesaikan $x^{x^{x^{\cdot^{\cdot^{\cdot}}}}}=2$ (menara pangkat tak hingga). Untuk nilai apa menara ini konvergen?
\kanalisis Jika menara $=y$, maka $x^y=y$; substitusikan $y=2$. Periksa syarat konvergensi $e^{-e}\le x\le e^{1/e}$.
\kbahas Misal nilai menara $y=2$. Karena eksponen teratas juga bernilai $y$, berlaku $x^{y}=y\Rightarrow x^2=2\Rightarrow x=\sqrt2.$ Menara $h(x)=x^{x^{\cdots}}$ konvergen untuk $e^{-e}\le x\le e^{1/e}\approx1{,}4447$; karena $\sqrt2\approx1{,}414<1{,}4447$, nilai ini sah.

\knomor{Challenge 3}
\ksoal Buktikan bahwa ${}^{2}\!\log 3$ adalah bilangan irasional.
\kanalisis Bukti kontradiksi: andaikan rasional, lalu gunakan ketunggalan faktorisasi (paritas).
\kbahas Andaikan ${}^{2}\!\log3=\tfrac pq$ dengan $p,q$ bilangan asli. Maka $2^{p/q}=3\Rightarrow 2^p=3^q$. Ruas kiri genap (kelipatan $2$) sedangkan ruas kanan ganjil (pangkat $3$) untuk $p,q\ge1$ --- kontradiksi. Jadi ${}^{2}\!\log3$ tidak mungkin rasional, sehingga irasional. $\blacksquare$

% ============================================================
\chapter{Pembahasan Bab 2: Barisan dan Deret}
% ============================================================

\kbagian{Bagian A}

\knomor{A.1}
\ksoal Tentukan $U_{10}$ dari $4,7,10,\dots$.
\kanalisis Barisan aritmetika dengan $a=4$, beda $b=3$; pakai $U_n=a+(n-1)b$.
\kbahas $U_{10}=4+9(3)=31.$

\knomor{A.2}
\ksoal Hitung $S_{12}$ dari $1+3+5+\cdots$.
\kanalisis Deret aritmetika $a=1$, $b=2$. (Jumlah $n$ bilangan ganjil pertama $=n^2$.)
\kbahas $S_{12}=\tfrac{12}{2}\big(2(1)+11(2)\big)=6(24)=144.$ (Sama dengan $12^2=144$.)

\knomor{A.3}
\ksoal Tentukan rasio dan $U_5$ dari $2,6,18,\dots$.
\kanalisis Barisan geometri; $r=U_2/U_1$, lalu $U_n=ar^{n-1}$.
\kbahas $r=3$, $U_5=2\cdot3^{4}=2(81)=162.$

\knomor{A.4}
\ksoal Hitung $S_\infty$ dari $12+4+\tfrac43+\cdots$.
\kanalisis Geometri dengan $a=12$, $r=\tfrac13$ ($|r|<1$ sehingga konvergen).
\kbahas $S_\infty=\dfrac{12}{1-\frac13}=\dfrac{12}{2/3}=18.$

\knomor{A.5}
\ksoal Diketahui $a=5$, $b=-2$. Tentukan $U_8$.
\kanalisis Langsung gunakan $U_n=a+(n-1)b$.
\kbahas $U_8=5+7(-2)=5-14=-9.$

\knomor{A.6}
\ksoal Suku ke berapa dari $3,7,11,\dots$ yang bernilai $99$?
\kanalisis $a=3$, $b=4$; selesaikan $U_n=99$.
\kbahas $3+(n-1)4=99\Rightarrow 4(n-1)=96\Rightarrow n-1=24\Rightarrow n=25.$

\knomor{A.7}
\ksoal Tentukan jumlah 6 suku pertama dari $1,2,4,8,\dots$.
\kanalisis Geometri $a=1$, $r=2$; pakai $S_n=\dfrac{a(r^n-1)}{r-1}$.
\kbahas $S_6=\dfrac{1(2^6-1)}{2-1}=63.$

\knomor{A.8}
\ksoal Tiga suku berurutan aritmetika $x-2,\ x,\ x+2$ berjumlah 30. Tentukan $x$.
\kanalisis Jumlah tiga suku simetris $=3\times$ suku tengah.
\kbahas $(x-2)+x+(x+2)=3x=30\Rightarrow x=10.$

\knomor{A.9}
\ksoal Ubah $0{,}777\ldots$ menjadi pecahan.
\kanalisis Desimal berulang $=$ deret geometri tak hingga $a=0{,}7$, $r=0{,}1$.
\kbahas $S_\infty=\dfrac{0{,}7}{1-0{,}1}=\dfrac{0{,}7}{0{,}9}=\dfrac{7}{9}.$

\knomor{A.10}
\ksoal Tentukan beda jika $U_5=20$ dan $U_{12}=48$.
\kanalisis Selisih dua suku $=$ (selisih indeks)$\times b$.
\kbahas $U_{12}-U_5=7b=48-20=28\Rightarrow b=4.$

\kbagian{Bagian B}

\knomor{B.1}
\ksoal Jumlah $n$ suku pertama deret aritmetika adalah $S_n=3n^2-n$. Tentukan $U_{10}$ dan beda.
\kanalisis Gunakan $U_n=S_n-S_{n-1}$.
\kbahas $U_n=(3n^2-n)-\big(3(n-1)^2-(n-1)\big)=6n-4.$ Maka $U_{10}=56$ dan beda $b=6$ (cek $U_1=2=S_1$).

\knomor{B.2}
\ksoal Suku ke-3 barisan geometri adalah 18 dan suku ke-6 adalah 486. Tentukan $S_5$.
\kanalisis Bagi $U_6/U_3$ untuk memperoleh $r^3$, lalu cari $a$.
\kbahas $\dfrac{ar^5}{ar^2}=r^3=\dfrac{486}{18}=27\Rightarrow r=3$. Dari $ar^2=18\Rightarrow a=2$. $S_5=\dfrac{2(3^5-1)}{3-1}=\dfrac{2(242)}{2}=242.$

\knomor{B.3}
\ksoal Sebuah tali dipotong menjadi 6 bagian membentuk barisan geometri; terpendek 3 cm, terpanjang 96 cm. Tentukan panjang tali semula.
\kanalisis $a=3$, $U_6=ar^5=96$ memberi $r$; jumlahkan keenam potongan.
\kbahas $r^5=\dfrac{96}{3}=32\Rightarrow r=2$. $S_6=\dfrac{3(2^6-1)}{2-1}=3(63)=189$ cm.

\knomor{B.4}
\ksoal Tiga bilangan aritmetika berjumlah 36. Jika suku pertama dan ketiga ditambah 1 (tengah tetap), terbentuk geometri. Tentukan ketiga bilangan semula.
\kanalisis Misalkan $12-d,\ 12,\ 12+d$ (sebab jumlah $=3\cdot12$); terapkan syarat geometri $U_2^2=U_1U_3$ pada barisan baru.
\kbahas Barisan baru $(13-d),12,(13+d)$ geometri: $12^2=(13-d)(13+d)=169-d^2\Rightarrow d^2=25\Rightarrow d=5$. Bilangan semula: $7,12,17.$

\knomor{B.5}
\ksoal Deret geometri tak hingga berjumlah 12 dan suku pertamanya 8. Tentukan rasio dan suku ke-4.
\kanalisis Pakai $S_\infty=\dfrac{a}{1-r}$ untuk mencari $r$, lalu $U_4=ar^3$.
\kbahas $\dfrac{8}{1-r}=12\Rightarrow 1-r=\tfrac23\Rightarrow r=\tfrac13$. $U_4=8\left(\tfrac13\right)^3=\dfrac{8}{27}.$

\kbagian{Bagian C}

\knomor{C.1}
\ksoal Jika $a,b,c$ membentuk barisan geometri, buktikan $\log a,\log b,\log c$ membentuk barisan aritmetika.
\kanalisis Syarat geometri $b^2=ac$; ambil logaritma kedua ruas.
\kbahas Dari $\dfrac ba=\dfrac cb$ diperoleh $b^2=ac$. Maka $\log b^2=\log(ac)\Rightarrow 2\log b=\log a+\log c$, yaitu syarat aritmetika dengan suku tengah $\log b$. $\blacksquare$

\knomor{C.2}
\ksoal Tentukan jumlah $\displaystyle\sum_{k=1}^{n} k\cdot 2^{k}$.
\kanalisis Teknik "kali rasio lalu kurangkan" (mirip penurunan deret geometri).
\kbahas Misal $S=\sum_{k=1}^n k2^k$. Maka $2S=\sum k2^{k+1}=\sum_{k=2}^{n+1}(k-1)2^k$. Kurangkan: $2S-S=n2^{n+1}-\sum_{k=1}^{n}2^k=n2^{n+1}-(2^{n+1}-2)$. Jadi $S=(n-1)2^{n+1}+2.$ (Cek $n=2$: $1\cdot8+2=10=2+8$.)

\knomor{C.3}
\ksoal Jika $\log 2,\ \log(2^x-1),\ \log(2^x+3)$ membentuk barisan aritmetika, tentukan $x$.
\kanalisis Syarat aritmetika: $2\times$ suku tengah $=$ jumlah dua tepi; ubah ke perkalian argumen, lalu misalkan $y=2^x$.
\kbahas $2\log(2^x-1)=\log2+\log(2^x+3)\Rightarrow (2^x-1)^2=2(2^x+3)$. Misal $y=2^x$: $y^2-2y+1=2y+6\Rightarrow y^2-4y-5=0\Rightarrow(y-5)(y+1)=0$. Karena $y=2^x>0$, $y=5$, jadi $2^x=5\Rightarrow x={}^{2}\!\log 5\approx2{,}32.$

\knomor{C.4}
\ksoal Suku-suku barisan geometri positif memenuhi $U_1 U_2 \cdots U_9 = 3^{18}$. Tentukan $U_5$.
\kanalisis Pada hasil kali $9$ suku geometri, pasangan simetris $U_k U_{10-k}=U_5^2$, sehingga hasil kali $=U_5^9$.
\kbahas $U_1U_2\cdots U_9=U_5^{9}=3^{18}\Rightarrow U_5=3^{18/9}=3^2=9.$

\knomor{C.5}
\ksoal Deret geometri tak hingga memiliki jumlah $S$ dan jumlah kuadrat sukunya $Q$. Nyatakan rasio $r$ dalam $S$ dan $Q$.
\kanalisis $S=\dfrac{a}{1-r}$; barisan kuadrat juga geometri dengan rasio $r^2$, jumlah $Q=\dfrac{a^2}{1-r^2}$. Bandingkan $Q$ dan $S^2$.
\kbahas $\dfrac{Q}{S^2}=\dfrac{a^2/(1-r^2)}{a^2/(1-r)^2}=\dfrac{(1-r)^2}{(1-r)(1+r)}=\dfrac{1-r}{1+r}$. Selesaikan: $r=\dfrac{S^2-Q}{S^2+Q}.$

\kbagian{Bagian D}

\knomor{D.1}
\ksoal \textbf{(Deret + Induksi).} Buktikan $\displaystyle\sum_{k=1}^{n} k^2 = \dfrac{n(n+1)(2n+1)}{6}$ dengan induksi matematika.
\kanalisis Cek basis $n=1$, lalu tambahkan $(k+1)^2$ pada hipotesis.
\kbahas \emph{Basis:} $n=1$: ruas kiri $1$, ruas kanan $\tfrac{1\cdot2\cdot3}{6}=1$. \emph{Langkah:} andaikan benar untuk $n=k$. Maka
$\sum_{i=1}^{k+1}i^2=\dfrac{k(k+1)(2k+1)}{6}+(k+1)^2=\dfrac{(k+1)\big(k(2k+1)+6(k+1)\big)}{6}=\dfrac{(k+1)(2k^2+7k+6)}{6}=\dfrac{(k+1)(k+2)(2k+3)}{6},$
yang merupakan rumus untuk $n=k+1$. $\blacksquare$

\knomor{D.2}
\ksoal \textbf{(Deret geometri + limit).} Pada kurva Koch tiap iterasi panjang menjadi $\tfrac43$ kali. Tunjukkan panjang total menuju tak hingga sedangkan luas konvergen, dan jelaskan paradoksnya.
\kanalisis Panjang membentuk geometri rasio $\tfrac43>1$ (divergen); luas tambahan tiap iterasi membentuk geometri rasio $<1$ (konvergen).
\kbahas Panjang $L_n=\left(\tfrac43\right)^nL_0\to\infty$ karena $\tfrac43>1$. Luas yang ditambahkan tiap iterasi mengecil dengan rasio $\tfrac49<1$, sehingga total luas berupa deret geometri konvergen (berhingga). \emph{Paradoks:} sebuah kurva berkeliling tak hingga namun melingkupi luas berhingga --- ciri khas objek fraktal.

\knomor{D.3}
\ksoal \textbf{(Deret + Peluang).} Peluang berhenti pada lemparan ke-$n$ adalah $(\tfrac12)^n$. Buktikan jumlah seluruh peluang $=1$ dan hitung ekspektasi $\sum n(\tfrac12)^n$.
\kanalisis Jumlah peluang $=$ deret geometri; ekspektasi $=$ deret aritmetika-geometri.
\kbahas $\sum_{n=1}^{\infty}\left(\tfrac12\right)^n=\dfrac{1/2}{1-1/2}=1.$ Ekspektasi: gunakan $\sum_{n\ge1}nx^n=\dfrac{x}{(1-x)^2}$ dengan $x=\tfrac12$: $E=\dfrac{1/2}{(1/2)^2}=2.$ Rata-rata $2$ lemparan.

\knomor{D.4}
\ksoal \textbf{(Deret + Eksponen).} Tunjukkan $\displaystyle\sum_{n=0}^{\infty}\frac{1}{2^n}=2$ dan bandingkan dengan deret harmonik $\sum\tfrac1n$ yang divergen.
\kanalisis Yang pertama geometri ($r=\tfrac12<1$); yang kedua sukunya menyusut terlalu lambat.
\kbahas $\sum_{n=0}^{\infty}\left(\tfrac12\right)^n=\dfrac{1}{1-\frac12}=2.$ Sebaliknya, $\sum\tfrac1n$ divergen karena suku $\tfrac1n$ hanya mengecil seperti $1/n$ (tak cukup cepat) --- pengelompokan $\tfrac13+\tfrac14>\tfrac12$, $\tfrac15+\cdots+\tfrac18>\tfrac12$, dst., menambah $\tfrac12$ tanpa henti. Kuncinya: deret geometri menyusut \emph{eksponensial}, harmonik hanya \emph{polinomial}.

\knomor{D.5}
\ksoal \textbf{(Anuitas).} Pinjaman Rp12.000.000 dilunasi 12 angsuran bulanan tetap, bunga majemuk $1\%$/bulan. Turunkan rumus angsuran $A$ dengan deret geometri, lalu hitung $A$.
\kanalisis Nilai sekarang dari $n$ angsuran $=$ deret geometri dengan rasio $(1+i)^{-1}$ yang harus sama dengan pinjaman $P$.
\kbahas $P=A\sum_{k=1}^{n}(1+i)^{-k}=A\cdot\dfrac{1-(1+i)^{-n}}{i}\Rightarrow A=\dfrac{P\,i}{1-(1+i)^{-n}}.$ Dengan $P=12.000.000$, $i=0{,}01$, $n=12$: $1{,}01^{-12}\approx0{,}88745$, sehingga $A=\dfrac{120.000}{0{,}11255}\approx \text{Rp}1.066.200.$

\kbagian{Challenge Corner}

\knomor{Challenge 1}
\ksoal Hitung $\displaystyle\sum_{k=1}^{100}\frac{1}{k(k+1)}$ tanpa menjumlah satu per satu.
\kanalisis Pecahan parsial $\dfrac{1}{k(k+1)}=\dfrac1k-\dfrac1{k+1}$ menghasilkan deret teleskopik.
\kbahas $\sum_{k=1}^{100}\left(\dfrac1k-\dfrac1{k+1}\right)=1-\dfrac{1}{101}=\dfrac{100}{101}.$

\knomor{Challenge 2}
\ksoal Bilangan $1$ ditulis, lalu tiap langkah setiap "1" diganti "10" dan "0" diganti "1". Selidiki banyak digit setelah $n$ langkah.
\kanalisis Aturan "1$\to$10, 0$\to$1" memetakan jumlah simbol; lacak banyak digit tiap langkah.
\kbahas Urutannya: $1\to10\to101\to10110\to10110101\to\cdots$ dengan banyak digit $1,2,3,5,8,13,\dots$ --- tepat \textbf{barisan Fibonacci}. Jika $d_n$ banyak digit setelah $n$ langkah, maka $d_n=F_{n+2}$ (dengan $F_1=F_2=1$), sebab tiap "1" menghasilkan 2 simbol dan tiap "0" menghasilkan 1, memberi rekursi $d_n=d_{n-1}+d_{n-2}$.

\knomor{Challenge 3}
\ksoal Buktikan bahwa $\displaystyle\sum_{k=1}^{n} F_k = F_{n+2}-1$, dengan $F_1=F_2=1$.
\kanalisis Gunakan $F_k=F_{k+2}-F_{k+1}$ agar deret menjadi teleskopik (atau induksi).
\kbahas Karena $F_k=F_{k+2}-F_{k+1}$, maka $\sum_{k=1}^{n}F_k=\sum_{k=1}^{n}(F_{k+2}-F_{k+1})=F_{n+2}-F_2=F_{n+2}-1.$ $\blacksquare$

% ============================================================
\chapter{Pembahasan Bab 3: Vektor}
% ============================================================

\kbagian{Bagian A}

\knomor{A.1}
\ksoal Diketahui $\vec a=(3,-2)$, $\vec b=(1,5)$. Tentukan $\vec a+\vec b$.
\kanalisis Jumlahkan komponen yang bersesuaian.
\kbahas $\vec a+\vec b=(3+1,\,-2+5)=(4,3).$

\knomor{A.2}
\ksoal Hitung $|\vec a|$ untuk $\vec a=(8,6)$.
\kanalisis Panjang vektor $=\sqrt{a_1^2+a_2^2}$ (Pythagoras).
\kbahas $|\vec a|=\sqrt{8^2+6^2}=\sqrt{100}=10.$

\knomor{A.3}
\ksoal Tentukan $3\vec a$ jika $\vec a=(-1,2)$.
\kanalisis Perkalian skalar mengalikan tiap komponen.
\kbahas $3\vec a=(-3,6).$

\knomor{A.4}
\ksoal Hitung $\vec a\cdot\vec b$ untuk $\vec a=(2,1)$, $\vec b=(3,4)$.
\kanalisis Hasil kali titik $=a_1b_1+a_2b_2$ (menghasilkan skalar).
\kbahas $\vec a\cdot\vec b=2(3)+1(4)=10.$

\knomor{A.5}
\ksoal Tentukan $\vec{AB}$ jika $A(1,2)$, $B(4,6)$.
\kanalisis $\vec{AB}=B-A$ (ujung dikurangi pangkal).
\kbahas $\vec{AB}=(4-1,\,6-2)=(3,4).$

\knomor{A.6}
\ksoal Tentukan vektor satuan searah $(0,5)$.
\kanalisis Bagi vektor dengan panjangnya: $\hat a=\vec a/|\vec a|$.
\kbahas $|(0,5)|=5$, jadi $\hat a=\tfrac15(0,5)=(0,1).$

\knomor{A.7}
\ksoal Apakah $(2,3)$ dan $(4,6)$ sejajar?
\kanalisis Dua vektor sejajar jika komponennya proporsional.
\kbahas $\tfrac42=\tfrac63=2$, jadi $(4,6)=2(2,3)$ --- \textbf{sejajar}.

\knomor{A.8}
\ksoal Hitung $\vec a-\vec b$ untuk $\vec a=(5,5)$, $\vec b=(2,1)$.
\kanalisis Kurangkan komponen bersesuaian.
\kbahas $\vec a-\vec b=(3,4).$

\knomor{A.9}
\ksoal Tentukan $k$ agar $(k,2)\perp(2,1)$.
\kanalisis Tegak lurus $\iff$ hasil kali titik $=0$.
\kbahas $2k+2=0\Rightarrow k=-1.$

\knomor{A.10}
\ksoal Hitung jarak $P(0,0)$ ke $Q(3,4)$.
\kanalisis Jarak $=|\vec{PQ}|$.
\kbahas $|\vec{PQ}|=\sqrt{3^2+4^2}=5.$

\kbagian{Bagian B}

\knomor{B.1}
\ksoal Titik $A(1,2)$, $B(4,3)$, $C(2,7)$. Tentukan jenis segitiga $ABC$.
\kanalisis Hitung panjang ketiga sisi; bandingkan kuadrat sisi terpanjang dengan jumlah kuadrat dua sisi lain.
\kbahas $|AB|=\sqrt{3^2+1^2}=\sqrt{10}$, $|BC|=\sqrt{(-2)^2+4^2}=\sqrt{20}$, $|AC|=\sqrt{1^2+5^2}=\sqrt{26}$. Ketiga sisi berbeda (\emph{sembarang}). Karena $26<10+20=30$, sudut terbesar lancip, jadi segitiga \textbf{lancip sembarang}.

\knomor{B.2}
\ksoal Diketahui $\vec a=(2,3)$, $\vec b=(t,4)$ membentuk sudut $45^\circ$. Tentukan $t$.
\kanalisis Pakai $\cos45^\circ=\dfrac{\vec a\cdot\vec b}{|\vec a||\vec b|}=\tfrac{1}{\sqrt2}$, lalu kuadratkan.
\kbahas $\dfrac{2t+12}{\sqrt{13}\sqrt{t^2+16}}=\tfrac{1}{\sqrt2}\Rightarrow 2(2t+12)^2=13(t^2+16)\Rightarrow 5t^2-96t-80=0\Rightarrow t=20$ atau $t=-\tfrac45$. (Keduanya memberi hasil kali titik positif, jadi keduanya valid.)

\knomor{B.3}
\ksoal Tentukan titik $D$ agar $ABCD$ jajaran genjang, dengan $A(1,1)$, $B(4,2)$, $C(5,5)$.
\kanalisis Pada jajaran genjang $ABCD$ berlaku $\vec{AB}=\vec{DC}$, sehingga $D=A+\vec{CB}$ atau $D=C-\vec{AB}$.
\kbahas $\vec{AB}=(3,1)$. $D=C-\vec{AB}=(5-3,\,5-1)=(2,4).$ Cek $\vec{AD}=(1,3)=\vec{BC}$. $\checkmark$

\knomor{B.4}
\ksoal Vektor $\vec a=(3,4)$ diuraikan menjadi komponen sejajar dan tegak lurus terhadap $\vec b=(1,0)$. Tentukan keduanya.
\kanalisis Komponen sejajar $=$ proyeksi $\dfrac{\vec a\cdot\vec b}{|\vec b|^2}\vec b$; komponen tegak lurus $=\vec a$ dikurangi proyeksi.
\kbahas Proyeksi $=\dfrac{3}{1}(1,0)=(3,0)$ (sejajar). Tegak lurus $=\vec a-(3,0)=(0,4)$.

\knomor{B.5}
\ksoal Perahu menyeberang dengan kecepatan $(0,4)$ m/s, arus sungai $(3,0)$ m/s. Tentukan kecepatan resultan, arah, dan besarnya.
\kanalisis Kecepatan resultan $=$ jumlah vektor; besar $=$ panjangnya; arah dari $\tan\theta$.
\kbahas Resultan $=(3,4)$, besar $=\sqrt{3^2+4^2}=5$ m/s. Arah $\tan\theta=\tfrac43\Rightarrow\theta\approx53^\circ$ terhadap arah arus (sumbu-$x$); perahu hanyut $53^\circ$ dari garis tepi.

\kbagian{Bagian C}

\knomor{C.1}
\ksoal Buktikan aturan kosinus $c^2=a^2+b^2-2ab\cos C$ menggunakan vektor.
\kanalisis Nyatakan sisi seberang $C$ sebagai selisih dua vektor sisi, lalu kuadratkan panjangnya.
\kbahas Misalkan dua sisi dari titik $C$ adalah vektor $\vec u,\vec v$ dengan $|\vec u|=b$, $|\vec v|=a$, dan sudut $C$ di antaranya. Sisi $c=|\vec u-\vec v|$. Maka $c^2=(\vec u-\vec v)\cdot(\vec u-\vec v)=|\vec u|^2+|\vec v|^2-2\vec u\cdot\vec v=b^2+a^2-2ab\cos C.$ $\blacksquare$

\knomor{C.2}
\ksoal Diketahui $|\vec a|=|\vec b|=1$ dan $|\vec a+\vec b|=\sqrt3$. Tentukan sudut antara $\vec a$ dan $\vec b$.
\kanalisis Kuadratkan $|\vec a+\vec b|$ dan gunakan $\vec a\cdot\vec b=\cos\theta$.
\kbahas $|\vec a+\vec b|^2=1+1+2\cos\theta=3\Rightarrow\cos\theta=\tfrac12\Rightarrow\theta=60^\circ.$

\knomor{C.3}
\ksoal Titik $A$, $B$, $C$ pada lingkaran satuan dengan sudut pusat $0^\circ, 120^\circ, 240^\circ$. Buktikan $\vec{OA}+\vec{OB}+\vec{OC}=\vec 0$.
\kanalisis Tuliskan tiap vektor sebagai $(\cos,\sin)$ lalu jumlahkan komponen.
\kbahas $\vec{OA}=(1,0)$, $\vec{OB}=(-\tfrac12,\tfrac{\sqrt3}{2})$, $\vec{OC}=(-\tfrac12,-\tfrac{\sqrt3}{2})$. Jumlah $=(1-\tfrac12-\tfrac12,\ 0+\tfrac{\sqrt3}{2}-\tfrac{\sqrt3}{2})=(0,0)=\vec0.$ $\blacksquare$

\knomor{C.4}
\ksoal Diberikan segitiga dengan titik berat $G$. Buktikan $\vec{GA}+\vec{GB}+\vec{GC}=\vec0$.
\kanalisis Titik berat $G=\tfrac13(A+B+C)$; gunakan $\vec{GA}=A-G$.
\kbahas $\vec{GA}+\vec{GB}+\vec{GC}=(A+B+C)-3G=(A+B+C)-(A+B+C)=\vec0.$ $\blacksquare$

\knomor{C.5}
\ksoal Dua vektor satuan membentuk sudut $\theta$. Nyatakan $|\vec a-\vec b|$ dalam $\theta$ dan tafsirkan untuk $\theta=60^\circ$.
\kanalisis Kuadratkan; gunakan identitas $1-\cos\theta=2\sin^2\tfrac\theta2$.
\kbahas $|\vec a-\vec b|^2=2-2\cos\theta=4\sin^2\tfrac\theta2\Rightarrow|\vec a-\vec b|=2\sin\tfrac\theta2.$ Untuk $\theta=60^\circ$: $2\sin30^\circ=1$ (sisi segitiga sama sisi berjari-jari 1).

\kbagian{Bagian D}

\knomor{D.1}
\ksoal \textbf{(Vektor + Trigonometri).} Buktikan ketaksamaan Cauchy--Schwarz $|\vec a\cdot\vec b|\le|\vec a||\vec b|$ dan jelaskan kapan kesamaan tercapai.
\kanalisis Gunakan $\vec a\cdot\vec b=|\vec a||\vec b|\cos\theta$ dan $|\cos\theta|\le1$.
\kbahas $|\vec a\cdot\vec b|=|\vec a||\vec b|\,|\cos\theta|\le|\vec a||\vec b|$ sebab $|\cos\theta|\le1$. Kesamaan tercapai bila $|\cos\theta|=1$, yaitu $\theta=0^\circ$ atau $180^\circ$ --- saat $\vec a$ dan $\vec b$ \emph{kolinear} (sejajar). $\blacksquare$

\knomor{D.2}
\ksoal \textbf{(Vektor + Fungsi Kuadrat).} Tinjau $f(t)=|\vec a+t\vec b|^2$ sebagai fungsi kuadrat dalam $t$. Tentukan $t$ yang meminimumkan $f$, dan hubungannya dengan proyeksi.
\kanalisis Jabarkan menjadi kuadrat dalam $t$; titik puncak di $t=-\tfrac{B}{2A}$.
\kbahas $f(t)=|\vec b|^2t^2+2(\vec a\cdot\vec b)t+|\vec a|^2$. Minimum di $t^*=-\dfrac{\vec a\cdot\vec b}{|\vec b|^2}$. Nilai $|t^*|\,|\vec b|=\dfrac{|\vec a\cdot\vec b|}{|\vec b|}$ tepat panjang proyeksi $\vec a$ pada $\vec b$; jadi $\vec a+t^*\vec b$ adalah komponen $\vec a$ yang tegak lurus $\vec b$ (jarak minimum).

\knomor{D.3}
\ksoal \textbf{(Vektor + Statistika).} Tunjukkan bahwa "cosine similarity" dua vektor data sama dengan koefisien korelasi data yang telah dipusatkan.
\kanalisis Pusatkan data ($x_i-\bar x$); tulis korelasi Pearson sebagai hasil kali titik dibagi panjang.
\kbahas Untuk $\vec u=(x_i-\bar x)$, $\vec v=(y_i-\bar y)$: $\cos\theta=\dfrac{\vec u\cdot\vec v}{|\vec u||\vec v|}=\dfrac{\sum(x_i-\bar x)(y_i-\bar y)}{\sqrt{\sum(x_i-\bar x)^2}\sqrt{\sum(y_i-\bar y)^2}}=r.$ Inilah koefisien korelasi Pearson. $\blacksquare$

\knomor{D.4}
\ksoal \textbf{(Vektor + Geometri).} Buktikan diagonal belah ketupat saling tegak lurus dengan dot product.
\kanalisis Sisi-sisi $\vec a,\vec b$ dengan $|\vec a|=|\vec b|$; diagonalnya $\vec a+\vec b$ dan $\vec a-\vec b$.
\kbahas $(\vec a+\vec b)\cdot(\vec a-\vec b)=|\vec a|^2-|\vec b|^2=0$ karena $|\vec a|=|\vec b|$. Jadi kedua diagonal tegak lurus. $\blacksquare$

\knomor{D.5}
\ksoal \textbf{(Vektor 3D).} Hitung $\vec a\times\vec b$ untuk $\vec a=(1,0,0)$, $\vec b=(0,1,0)$ dan tafsirkan arah serta besarnya sebagai luas.
\kanalisis Cross product menghasilkan vektor tegak lurus kedua vektor; besarnya $=|\vec a||\vec b|\sin\theta$.
\kbahas $\vec a\times\vec b=(0\cdot0-0\cdot1,\ 0\cdot0-1\cdot0,\ 1\cdot1-0\cdot0)=(0,0,1).$ Besarnya $1=$ luas persegi satuan yang dibentuk $\vec a,\vec b$; arahnya sepanjang sumbu-$z$ (kaidah tangan kanan).

\kbagian{Challenge Corner}

\knomor{Challenge 1}
\ksoal Buktikan titik tengah sisi sebuah segi empat membentuk jajaran genjang (Teorema Varignon).
\kanalisis Nyatakan titik tengah sebagai rata-rata dua titik sudut, lalu bandingkan sisi-sisi segi empat tengah.
\kbahas Misal sudut $A,B,C,D$ dan titik tengah $P=\tfrac{A+B}2,\ Q=\tfrac{B+C}2,\ R=\tfrac{C+D}2,\ S=\tfrac{D+A}2$. Maka $\vec{PQ}=Q-P=\tfrac{C-A}2$ dan $\vec{SR}=R-S=\tfrac{C-A}2$. Karena $\vec{PQ}=\vec{SR}$, maka $PQRS$ jajaran genjang. $\blacksquare$

\knomor{Challenge 2}
\ksoal Tiga gaya sama besar bekerja pada satu titik dan setimbang ($\vec0$). Buktikan sudut antar gaya $120^\circ$.
\kanalisis Dari $\vec{F_1}+\vec{F_2}+\vec{F_3}=\vec0$ dengan $|\vec{F_i}|=F$, isolasi satu gaya lalu kuadratkan.
\kbahas $\vec{F_3}=-(\vec{F_1}+\vec{F_2})\Rightarrow F^2=|\vec{F_1}+\vec{F_2}|^2=F^2+F^2+2F^2\cos\theta_{12}\Rightarrow\cos\theta_{12}=-\tfrac12\Rightarrow\theta_{12}=120^\circ.$ Karena simetris, semua sudut antar gaya $=120^\circ$. $\blacksquare$

\knomor{Challenge 3}
\ksoal Tentukan vektor $\vec v$ di bidang yang memenuhi $\vec v\cdot(1,2)=5$ dan $|\vec v|$ minimum.
\kanalisis $\vec v\cdot\vec n=5$ adalah garis; jarak minimum dari titik asal dicapai oleh $\vec v$ yang \emph{sejajar} $\vec n=(1,2)$.
\kbahas Ambil $\vec v=k(1,2)$. Maka $\vec v\cdot(1,2)=5k=5\Rightarrow k=1$, sehingga $\vec v=(1,2)$ dengan $|\vec v|=\sqrt5$ (minimum). Setiap $\vec v$ lain pada garis itu lebih panjang.

% ============================================================
\chapter{Pembahasan Bab 4: Trigonometri Dasar}
% ============================================================

\kbagian{Bagian A}

\knomor{A.1}
\ksoal Hitung $\sin30^\circ+\cos60^\circ$.
\kanalisis Gunakan nilai sudut istimewa.
\kbahas $\tfrac12+\tfrac12=1.$

\knomor{A.2}
\ksoal Diketahui depan $=5$, miring $=13$. Tentukan $\cos\alpha$.
\kanalisis Cari sisi samping lewat Pythagoras, lalu $\cos=\tfrac{sa}{mi}$.
\kbahas Samping $=\sqrt{13^2-5^2}=\sqrt{144}=12$, jadi $\cos\alpha=\tfrac{12}{13}.$

\knomor{A.3}
\ksoal Hitung $\tan45^\circ$.
\kanalisis Nilai istimewa.
\kbahas $\tan45^\circ=1.$

\knomor{A.4}
\ksoal Tentukan $\sin90^\circ-\cos0^\circ$.
\kanalisis $\sin90^\circ=1$, $\cos0^\circ=1$.
\kbahas $1-1=0.$

\knomor{A.5}
\ksoal Jika $\tan\alpha=\tfrac34$ (lancip), tentukan $\sin\alpha$.
\kanalisis Segitiga $3$-$4$-$5$: depan $3$, samping $4$, miring $5$.
\kbahas $\sin\alpha=\tfrac{depan}{miring}=\tfrac35.$

\knomor{A.6}
\ksoal Hitung $\sin^2 45^\circ+\cos^2 45^\circ$.
\kanalisis Identitas pokok $\sin^2+\cos^2=1$.
\kbahas $=1.$

\knomor{A.7}
\ksoal Tinggi tiang 6 m, bayangan 6 m. Berapa sudut elevasi matahari?
\kanalisis $\tan(\text{elevasi})=\tfrac{tinggi}{bayangan}$.
\kbahas $\tan\theta=\tfrac66=1\Rightarrow\theta=45^\circ.$

\knomor{A.8}
\ksoal Hitung $\cos60^\circ\cdot\tan45^\circ$.
\kanalisis Substitusi nilai istimewa.
\kbahas $\tfrac12\cdot1=\tfrac12.$

\knomor{A.9}
\ksoal Tentukan $\sin150^\circ$.
\kanalisis $150^\circ$ kuadran II (sin positif), relasi $180^\circ-150^\circ=30^\circ$.
\kbahas $\sin150^\circ=\sin30^\circ=\tfrac12.$

\knomor{A.10}
\ksoal Sederhanakan $\dfrac{\sin\alpha}{\tan\alpha}$.
\kanalisis $\tan\alpha=\tfrac{\sin\alpha}{\cos\alpha}$.
\kbahas $\dfrac{\sin\alpha}{\sin\alpha/\cos\alpha}=\cos\alpha.$

\kbagian{Bagian B}

\knomor{B.1}
\ksoal Diketahui $\sin\alpha+\cos\alpha=\tfrac75$. Tentukan $\sin\alpha\cos\alpha$.
\kanalisis Kuadratkan; pakai $\sin^2+\cos^2=1$.
\kbahas $(\sin\alpha+\cos\alpha)^2=1+2\sin\alpha\cos\alpha=\tfrac{49}{25}\Rightarrow 2\sin\alpha\cos\alpha=\tfrac{24}{25}\Rightarrow\sin\alpha\cos\alpha=\tfrac{12}{25}.$

\knomor{B.2}
\ksoal Sederhanakan $\dfrac{1}{1+\cos\alpha}+\dfrac{1}{1-\cos\alpha}$ menjadi bentuk $\sin$.
\kanalisis Samakan penyebut; gunakan $1-\cos^2\alpha=\sin^2\alpha$.
\kbahas $\dfrac{(1-\cos\alpha)+(1+\cos\alpha)}{1-\cos^2\alpha}=\dfrac{2}{\sin^2\alpha}.$

\knomor{B.3}
\ksoal Dari atas mercusuar tinggi 40 m, sudut depresi dua kapal $30^\circ$ dan $45^\circ$ segaris. Tentukan jarak antar kapal.
\kanalisis Sudut depresi $=$ sudut elevasi dari kapal; jarak mendatar $=\dfrac{40}{\tan\theta}$.
\kbahas Kapal dekat: $d_1=\dfrac{40}{\tan45^\circ}=40$ m. Kapal jauh: $d_2=\dfrac{40}{\tan30^\circ}=40\sqrt3$ m. Jarak $=40\sqrt3-40=40(\sqrt3-1)\approx29{,}3$ m.

\knomor{B.4}
\ksoal Jika $\sin\alpha=\tfrac{2}{3}$ dan $\alpha$ tumpul, tentukan $\tan\alpha$.
\kanalisis Di kuadran II $\cos\alpha<0$; cari $\cos$ lewat identitas.
\kbahas $\cos\alpha=-\sqrt{1-\tfrac49}=-\tfrac{\sqrt5}{3}$. $\tan\alpha=\dfrac{2/3}{-\sqrt5/3}=-\dfrac{2}{\sqrt5}=-\dfrac{2\sqrt5}{5}.$

\knomor{B.5}
\ksoal Buktikan $\dfrac{\cos\alpha}{1-\sin\alpha}=\dfrac{1+\sin\alpha}{\cos\alpha}$.
\kanalisis Perkalian silang; pakai $1-\sin^2\alpha=\cos^2\alpha$.
\kbahas Silang: $\cos^2\alpha=(1-\sin\alpha)(1+\sin\alpha)=1-\sin^2\alpha=\cos^2\alpha$. Identitas benar. $\blacksquare$

\kbagian{Bagian C}

\knomor{C.1}
\ksoal Gunakan dot product untuk membuktikan $\cos(A-B)=\cos A\cos B+\sin A\sin B$.
\kanalisis Ambil dua vektor satuan bersudut $A$ dan $B$ terhadap sumbu-$x$; hitung dot product dua cara.
\kbahas Misal $\vec u=(\cos A,\sin A)$, $\vec v=(\cos B,\sin B)$. Maka $\vec u\cdot\vec v=\cos A\cos B+\sin A\sin B$. Di sisi lain $\vec u\cdot\vec v=|\vec u||\vec v|\cos(A-B)=\cos(A-B)$. Kedua bentuk sama, terbukti. $\blacksquare$

\knomor{C.2}
\ksoal Benda pada bidang miring $30^\circ$ dengan gaya berat 50 N. Uraikan gaya menjadi komponen sejajar dan tegak lurus bidang.
\kanalisis Berat (vertikal) diuraikan: komponen sejajar bidang $=W\sin\theta$, tegak lurus $=W\cos\theta$.
\kbahas Sejajar bidang $=50\sin30^\circ=25$ N (komponen yang menggerakkan benda). Tegak lurus bidang $=50\cos30^\circ=25\sqrt3\approx43{,}3$ N (ditahan bidang).

\knomor{C.3}
\ksoal Buktikan identitas $\sin3\alpha=3\sin\alpha-4\sin^3\alpha$.
\kanalisis Uraikan $\sin3\alpha=\sin(2\alpha+\alpha)$ dan gunakan rumus sudut rangkap.
\kbahas $\sin3\alpha=\sin2\alpha\cos\alpha+\cos2\alpha\sin\alpha=2\sin\alpha\cos^2\alpha+(1-2\sin^2\alpha)\sin\alpha=2\sin\alpha(1-\sin^2\alpha)+\sin\alpha-2\sin^3\alpha=3\sin\alpha-4\sin^3\alpha.$ $\blacksquare$

\knomor{C.4}
\ksoal Pada segitiga $ABC$, buktikan aturan sinus $\dfrac{a}{\sin A}=\dfrac{b}{\sin B}=\dfrac{c}{\sin C}$.
\kanalisis Tarik garis tinggi dari satu titik sudut; nyatakan tinggi dua cara.
\kbahas Garis tinggi dari $C$ ke $AB$ bernilai $h=b\sin A=a\sin B$, sehingga $\dfrac{a}{\sin A}=\dfrac{b}{\sin B}$. Dengan garis tinggi lain diperoleh $\dfrac{c}{\sin C}$ juga sama. $\blacksquare$

\knomor{C.5}
\ksoal Vektor $\vec a$ membentuk sudut $\alpha$ dengan sumbu-$x$. Nyatakan $\vec a$ dalam $|\vec a|,\cos\alpha,\sin\alpha$, lalu tentukan sudut jika $\vec a=(-1,\sqrt3)$.
\kanalisis Komponen vektor $=$ panjang $\times$ (kosinus, sinus) arah.
\kbahas $\vec a=|\vec a|(\cos\alpha,\sin\alpha)$. Untuk $\vec a=(-1,\sqrt3)$: $|\vec a|=\sqrt{1+3}=2$, $\cos\alpha=-\tfrac12$, $\sin\alpha=\tfrac{\sqrt3}{2}\Rightarrow\alpha=120^\circ.$

\kbagian{Bagian D}

\knomor{D.1}
\ksoal \textbf{(Trigonometri + Fungsi Kuadrat).} Tentukan nilai maksimum dan minimum $f(\alpha)=3\sin\alpha+4\cos\alpha$ dengan menuliskannya sebagai $R\sin(\alpha+\varphi)$.
\kanalisis $a\sin\alpha+b\cos\alpha=R\sin(\alpha+\varphi)$ dengan $R=\sqrt{a^2+b^2}$.
\kbahas $R=\sqrt{3^2+4^2}=5$, sehingga $f(\alpha)=5\sin(\alpha+\varphi)$. Karena $-1\le\sin(\cdot)\le1$, nilai maksimum $=5$ dan minimum $=-5$.

\knomor{D.2}
\ksoal \textbf{(Trigonometri + Vektor).} Buktikan luas segitiga $L=\tfrac12 ab\sin C$ dan kaitkan dengan besar cross product.
\kanalisis Luas $=\tfrac12\times$ alas $\times$ tinggi; tinggi $=b\sin C$.
\kbahas Dengan alas $a$ dan tinggi $b\sin C$: $L=\tfrac12 a(b\sin C)=\tfrac12 ab\sin C$. Karena $|\vec a\times\vec b|=ab\sin C$, maka $L=\tfrac12|\vec a\times\vec b|$ --- luas segitiga adalah setengah besar cross product dua vektor sisinya.

\knomor{D.3}
\ksoal \textbf{(Trigonometri + Deret).} Tunjukkan bahwa untuk sudut kecil (radian), $\sin\theta\approx\theta$. Gunakan untuk menaksir $\sin 1^\circ$.
\kanalisis Untuk $\theta$ kecil, busur $\approx$ tali busur; secara deret $\sin\theta=\theta-\tfrac{\theta^3}{6}+\cdots\approx\theta$.
\kbahas $1^\circ=\dfrac{\pi}{180}\approx0{,}01745$ rad, sehingga $\sin1^\circ\approx0{,}01745$ (nilai sebenarnya $0{,}017452$ --- selisih sangat kecil).

\knomor{D.4}
\ksoal \textbf{(Trigonometri + Identitas).} Tentukan semua $x\in[0^\circ,360^\circ)$ yang memenuhi $\sin^4x+\cos^4x=\tfrac{7}{8}$ dan hitung jumlah semua solusinya.
\kanalisis $\sin^4x+\cos^4x = (\sin^2x+\cos^2x)^2-2\sin^2x\cos^2x = 1-\tfrac12\sin^22x$.
\kbahas Substitusi: $1-\tfrac12\sin^22x=\tfrac78 \Rightarrow \sin^22x=\tfrac14 \Rightarrow \sin2x=\pm\tfrac12$.

Karena $x\in[0^\circ,360^\circ)$ maka $2x\in[0^\circ,720^\circ)$.

Dari $\sin2x=\tfrac12$: $2x=30^\circ,150^\circ,390^\circ,510^\circ$ sehingga $x=15^\circ,75^\circ,195^\circ,255^\circ$.

Dari $\sin2x=-\tfrac12$: $2x=210^\circ,330^\circ,570^\circ,690^\circ$ sehingga $x=105^\circ,165^\circ,285^\circ,345^\circ$.

Total 8 solusi; jumlah $= (15+345)+(75+285)+(105+255)+(165+195) = 4\times360^\circ = 1440^\circ$.

\knomor{D.5}
\ksoal \textbf{(Pemodelan).} Tinggi air pasang $h(t)=2+1{,}5\sin\!\left(\tfrac{\pi}{6}t\right)$ meter. Tentukan tinggi maksimum, periode, dan kapan pertama kali $h=3{,}5$ m.
\kanalisis Maksimum saat $\sin=1$; periode $=\tfrac{2\pi}{\text{koefisien }t}$.
\kbahas Maksimum $=2+1{,}5=3{,}5$ m. Periode $=\dfrac{2\pi}{\pi/6}=12$ jam. $h=3{,}5\Rightarrow\sin\tfrac{\pi}{6}t=1\Rightarrow\tfrac{\pi}{6}t=\tfrac\pi2\Rightarrow t=3$ jam.

\kbagian{Challenge Corner}

\knomor{Challenge 1}
\ksoal Buktikan $\tan 1^\circ\tan2^\circ\tan3^\circ\cdots\tan89^\circ=1$.
\kanalisis Pasangkan $\tan k^\circ$ dengan $\tan(90^\circ-k^\circ)=\cot k^\circ$.
\kbahas $\tan k^\circ\cdot\tan(90^\circ-k^\circ)=\tan k^\circ\cot k^\circ=1$. Pasangan $(1,89),(2,88),\dots,(44,46)$ masing-masing $=1$, dan suku tengah $\tan45^\circ=1$. Jadi hasil kali $=1$. $\blacksquare$

\knomor{Challenge 2}
\ksoal Hitung $\cos36^\circ$ dalam bentuk akar (kaitannya dengan rasio emas).
\kanalisis Gunakan $\cos36^\circ-\cos72^\circ=\tfrac12$ atau persamaan dari segi lima beraturan.
\kbahas Dengan identitas pada segi lima beraturan diperoleh $\cos36^\circ=\dfrac{1+\sqrt5}{4}\approx0{,}809$, yakni separuh rasio emas $\varphi=\dfrac{1+\sqrt5}{2}$.

\knomor{Challenge 3}
\ksoal Sebuah tangga melewati sudut koridor berbentuk L (lebar $a$ dan $b$). Tentukan panjang maksimum tangga sebagai masalah optimasi bersudut $\theta$.
\kanalisis Panjang tangga yang menyentuh kedua dinding pada sudut $\theta$ adalah $L(\theta)=\dfrac{a}{\sin\theta}+\dfrac{b}{\cos\theta}$; tangga terpanjang yang \emph{muat} adalah minimum fungsi ini.
\kbahas Minimumkan $L(\theta)=\dfrac{a}{\sin\theta}+\dfrac{b}{\cos\theta}$. Dengan turunan $L'(\theta)=0$ diperoleh $\tan^3\theta=\tfrac ab$, dan panjang maksimum tangga $L_{\min}=\left(a^{2/3}+b^{2/3}\right)^{3/2}.$

% ============================================================
\chapter{Pembahasan Bab 5: Sistem Persamaan dan Pertidaksamaan Linear}
% ============================================================

\kbagian{Bagian A}

\knomor{A.1}
\ksoal Selesaikan $\begin{cases} x+y=7\\ x-y=1\end{cases}$.
\kanalisis Jumlahkan kedua persamaan untuk mengeliminasi $y$.
\kbahas $(x+y)+(x-y)=8\Rightarrow 2x=8\Rightarrow x=4$, lalu $y=3$. $(x,y)=(4,3)$.

\knomor{A.2}
\ksoal Selesaikan $\begin{cases} 2x+y=8\\ x=3\end{cases}$.
\kanalisis Substitusi nilai $x$.
\kbahas $2(3)+y=8\Rightarrow y=2$. $(x,y)=(3,2)$.

\knomor{A.3}
\ksoal Tentukan $z$ jika $x=1,y=2$ dan $x+y+z=6$.
\kanalisis Substitusi langsung.
\kbahas $1+2+z=6\Rightarrow z=3.$

\knomor{A.4}
\ksoal Apakah $(0,0)$ memenuhi $x+y\le5$?
\kanalisis Uji titik ke pertidaksamaan.
\kbahas $0+0=0\le5$ benar, jadi \textbf{ya}.

\knomor{A.5}
\ksoal Selesaikan $\begin{cases} x+y+z=6\\ x=y\\ z=2\end{cases}$.
\kanalisis Substitusi $x=y$ dan $z=2$.
\kbahas $2x+2=6\Rightarrow x=2$, jadi $y=2$, $z=2$. $(2,2,2)$.

\knomor{A.6}
\ksoal Tentukan titik potong $x+y=4$ dan $x-y=0$.
\kanalisis $x=y$ dari persamaan kedua.
\kbahas $2x=4\Rightarrow x=2=y$. Titik $(2,2)$.

\knomor{A.7}
\ksoal Tuliskan pertidaksamaan untuk daerah di bawah garis $y=2x$.
\kanalisis "Di bawah" garis berarti $y$ lebih kecil dari nilai garis.
\kbahas $y<2x$ (atau $y\le2x$ jika garis disertakan).

\knomor{A.8}
\ksoal Selesaikan $3x=12$, lalu $x+y=7$.
\kanalisis Cari $x$ dulu, lalu $y$.
\kbahas $x=4\Rightarrow y=7-4=3.$

\knomor{A.9}
\ksoal Hitung $f=2x+y$ di titik $(3,1)$.
\kanalisis Substitusi koordinat.
\kbahas $f=2(3)+1=7.$

\knomor{A.10}
\ksoal Apakah garis $x+y=2$ dan $x+y=5$ berpotongan?
\kanalisis Dua garis dengan ruas kiri sama tetapi konstanta beda adalah sejajar.
\kbahas Tidak --- keduanya \textbf{sejajar} dan tak berpotongan.

\kbagian{Bagian B}

\knomor{B.1}
\ksoal Selesaikan $\begin{cases} x+y+z=12\\ 2x-y+3z=14\\ x+3y-z=4\end{cases}$.
\kanalisis Eliminasi $z$ dengan substitusi dari persamaan pertama.
\kbahas Dari (1) $z=12-x-y$. Substitusi ke (2): $-x-4y=-22\Rightarrow x+4y=22$. Dari (1)$+$(3): $2x+4y=16\Rightarrow x+2y=8$. Selisih: $2y=14\Rightarrow y=7$, $x=-6$, $z=11$. $(x,y,z)=(-6,7,11)$.

\knomor{B.2}
\ksoal Harga $3$ buku $+2$ pensil $+1$ penghapus $=$ Rp28rb; $1$ buku $+1$ pensil $+1$ penghapus $=$ Rp12rb; $2$ buku $+1$ pensil $+3$ penghapus $=$ Rp22rb. Tentukan harga tiap barang.
\kanalisis Misal buku $b$, pensil $p$, penghapus $e$ (ribuan); susun SPLTV lalu eliminasi.
\kbahas (1)$-$(2): $2b+p=16$. (3)$-2\cdot$(2): $-p+e=-2\Rightarrow e=p-2$. (2): $b+2p=14$. Dari $2b+p=16$ dan $b+2p=14$: $p=4$, $b=6$, $e=2$. Jadi \textbf{buku Rp6rb, pensil Rp4rb, penghapus Rp2rb}.

\knomor{B.3}
\ksoal Maksimumkan $f=4x+3y$ pada $2x+y\le10$, $x+3y\le15$, $x,y\ge0$.
\kanalisis Cari titik pojok daerah fisibel (termasuk perpotongan dua kendala), lalu uji $f$.
\kbahas Perpotongan $2x+y=10$ dan $x+3y=15$: $x=3$, $y=4$. Pojok: $(0,0)\to0$; $(5,0)\to20$; $(0,5)\to15$; $(3,4)\to24$. Maksimum $\boxed{24}$ di $(3,4)$.

\knomor{B.4}
\ksoal Tentukan nilai $a$ agar sistem $\begin{cases} x+y=2\\ ax+2y=4\end{cases}$ punya tak hingga penyelesaian.
\kanalisis Tak hingga penyelesaian $\iff$ kedua garis berimpit (koefisien proporsional).
\kbahas Kalikan (1) dengan 2: $2x+2y=4$. Agar identik dengan $ax+2y=4$, perlu $a=2.$

\knomor{B.5}
\ksoal Perusahaan membuat meja dan kursi: tiap meja butuh 4 jam, kursi 2 jam; tersedia 40 jam. Laba meja Rp80rb, kursi Rp30rb; kursi maksimal 15 buah. Maksimalkan laba.
\kanalisis Model: $4m+2k\le40$ (yaitu $2m+k\le20$), $k\le15$, $m,k\ge0$; maksimalkan $P=80m+30k$.
\kbahas Pojok daerah: $(0,0)\to0$; $(10,0)\to800$; $(2{,}5,15)\to650$; $(0,15)\to450$ (dalam ribuan). Maksimum laba $\textbf{Rp800.000}$ dengan membuat $10$ meja dan $0$ kursi.

\kbagian{Bagian C}

\knomor{C.1}
\ksoal Selesaikan sistem non-linear $\begin{cases} x+y+z=6\\ xy+yz+zx=11\\ xyz=6\end{cases}$.
\kanalisis Ketiga ruas adalah polinomial simetris dasar; $x,y,z$ adalah akar $t^3-e_1t^2+e_2t-e_3=0$.
\kbahas Persamaan akar: $t^3-6t^2+11t-6=0\Rightarrow(t-1)(t-2)(t-3)=0$. Jadi $\{x,y,z\}=\{1,2,3\}$ (dalam sembarang urutan).

\knomor{C.2}
\ksoal Tunjukkan bahwa sistem $A\vec x=\vec b$ ($2\times2$) punya solusi tunggal jika dan hanya jika vektor-vektor kolomnya bebas linear.
\kanalisis Solusi tunggal $\iff\det A\neq0\iff$ kolom tidak saling kelipatan.
\kbahas Untuk $A=\begin{pmatrix}a&b\\c&d\end{pmatrix}$, solusi tunggal $\iff\det A=ad-bc\neq0$. Determinan nol berarti kolom $(a,c)$ dan $(b,d)$ proporsional (bergantung linear), sehingga garis berimpit/sejajar (tak hingga/tanpa solusi). Jadi solusi tunggal $\iff$ kolom bebas linear. $\blacksquare$

\knomor{C.3}
\ksoal Tiga garis $x+y=2$, $2x-y=1$, $ax+y=b$ bertemu di satu titik. Tentukan hubungan $a$ dan $b$.
\kanalisis Cari titik potong dua garis pertama, lalu paksa garis ketiga melaluinya.
\kbahas $(x+y)+(2x-y)=3\Rightarrow 3x=3\Rightarrow x=1$, $y=1$. Garis ketiga melalui $(1,1)$: $a(1)+1=b\Rightarrow b=a+1.$

\knomor{C.4}
\ksoal Tentukan semua $(x,y,z)$ bilangan bulat positif dengan $x+y+z=10$ dan $x\le y\le z$.
\kanalisis Enumerasi partisi $10$ menjadi tiga bagian positif terurut.
\kbahas $(1,1,8),(1,2,7),(1,3,6),(1,4,5),(2,2,6),(2,3,5),(2,4,4),(3,3,4)$ --- ada \textbf{8} penyelesaian.

\knomor{C.5}
\ksoal Pada daerah $x+y\le6$, $x,y\ge0$, fungsi $f=2x+ky$ mencapai maksimum di dua titik pojok sekaligus. Tentukan $k$.
\kanalisis Maksimum di dua pojok $\iff$ garis $f=c$ sejajar dengan sisi yang menghubungkannya, yakni sisi $x+y=6$.
\kbahas Pojok daerah: $(0,0),(6,0),(0,6)$. Agar $f(6,0)=f(0,6)$: $12=6k\Rightarrow k=2.$ Saat $k=2$, $f=2(x+y)$ bernilai konstan $12$ di sepanjang sisi $x+y=6$.

\kbagian{Bagian D}

\knomor{D.1}
\ksoal \textbf{(Sistem + Matriks).} Selesaikan SPLTV kantin ($2x+y+z=32$, $x+2y+z=34$, $x+y+2z=30$) dengan aturan Cramer.
\kanalisis Hitung $\det A$ dan tiga determinan pengganti kolom.
\kbahas $\det A=\begin{vmatrix}2&1&1\\1&2&1\\1&1&2\end{vmatrix}=4$. $\det_x=32$, $\det_y=40$, $\det_z=24$. Maka $x=\tfrac{32}{4}=8$, $y=\tfrac{40}{4}=10$, $z=\tfrac{24}{4}=6$ --- identik dengan hasil eliminasi.

\knomor{D.2}
\ksoal \textbf{(Program linear + Pemodelan).} Sebuah pabrik membuat 2 pupuk dengan 3 kendala bahan. Rumuskan masalah, gambar daerah fisibel, dan temukan laba maksimum.
\kanalisis Contoh model: $x,y$ ton pupuk; kendala bahan $A,B,C$; maksimalkan $P$.
\kbahas Misal $x+2y\le8$, $3x+2y\le12$, $x,y\ge0$, laba $P=40x+30y$. Pojok: $(0,0)\to0$; $(4,0)\to160$; $(0,4)\to120$; perpotongan $x+2y=8$, $3x+2y=12$ di $(2,3)\to80+90=170$. Maksimum $\textbf{170}$ di $(2,3)$. (Nilai dapat disesuaikan dengan data nyata; prosedurnya tetap: daerah fisibel $\to$ uji titik pojok.)

\knomor{D.3}
\ksoal \textbf{(Sistem + Geometri).} Klasifikasikan semua kemungkinan posisi relatif tiga bidang di ruang dan kaitkan dengan banyaknya solusi.
\kanalisis Tiap bidang $=$ satu persamaan linear 3 variabel; tinjau perpotongannya.
\kbahas (i) Berpotongan di satu titik $\to$ \emph{solusi tunggal}. (ii) Berpotongan pada satu garis bersama $\to$ \emph{tak hingga solusi} (1 parameter). (iii) Ketiganya berimpit $\to$ tak hingga (2 parameter). (iv) Ada yang sejajar atau membentuk "prisma" tanpa titik bersama $\to$ \emph{tanpa solusi}.

\knomor{D.4}
\ksoal \textbf{(Sistem + Statistika).} Turunkan persamaan normal regresi linear $y=ax+b$ dengan metode kuadrat terkecil.
\kanalisis Minimumkan $S(a,b)=\sum(y_i-ax_i-b)^2$; turunan parsial terhadap $a$ dan $b$ disamakan nol.
\kbahas $\dfrac{\partial S}{\partial b}=0\Rightarrow a\sum x_i+bn=\sum y_i$. $\dfrac{\partial S}{\partial a}=0\Rightarrow a\sum x_i^2+b\sum x_i=\sum x_iy_i$. Inilah sistem dua persamaan linear dalam $a,b$ (persamaan normal).

\knomor{D.5}
\ksoal \textbf{(Sistem + Eksponen).} Selesaikan $\begin{cases} 2^x\cdot 3^y=12\\ 2^x\cdot 3^{-y}=\tfrac43\end{cases}$.
\kanalisis Bagi kedua persamaan untuk mengisolasi $3^{2y}$; atau ambil logaritma untuk menjadikan sistem linear.
\kbahas Bagi (1) oleh (2): $3^{2y}=\dfrac{12}{4/3}=9=3^2\Rightarrow 2y=2\Rightarrow y=1$. Substitusi: $2^x\cdot3=12\Rightarrow2^x=4\Rightarrow x=2$. $(x,y)=(2,1)$.

\kbagian{Challenge Corner}

\knomor{Challenge 1}
\ksoal Sebuah kebun memaksimalkan hasil dua tanaman dengan kendala lahan, air, dan pupuk. Rancang data masuk akal, modelkan, dan selesaikan.
\kanalisis Tentukan variabel ($x,y$ luas/unit), tiga kendala sumber daya, dan fungsi hasil.
\kbahas Contoh: $x,y$ petak. Lahan $x+y\le10$, air $2x+y\le16$, pupuk $x+3y\le18$, $x,y\ge0$; hasil $H=5x+4y$ (juta). Pojok dominan: perpotongan $x+y=10$ dan $2x+y=16$ di $(6,4)\to H=46$; cek $(0,6)\to24$, $(8,0)\to40$. Maksimum $\textbf{46 juta}$ di $(6,4)$ (memenuhi semua kendala). Prosedur sama untuk data lain.

\knomor{Challenge 2}
\ksoal Tentukan banyaknya solusi bilangan bulat non-negatif dari $x+y+z=20$.
\kanalisis Metode "bintang dan batang": membagi 20 unit ke 3 kotak $=\binom{20+2}{2}$.
\kbahas Banyak solusi $=\dbinom{22}{2}=\dfrac{22\cdot21}{2}=231.$

\knomor{Challenge 3}
\ksoal Selesaikan $\begin{cases} x+y+z=1\\ x^2+y^2+z^2=1\\ x^3+y^3+z^3=1\end{cases}$.
\kanalisis Gunakan polinomial simetris: $e_1=1$; dari $p_2$ cari $e_2$; dari $p_3$ (Newton) cari $e_3$.
\kbahas $e_1=1$. $p_2=e_1^2-2e_2=1-2e_2=1\Rightarrow e_2=0$. $p_3=e_1p_2-e_2p_1+3e_3=1-0+3e_3=1\Rightarrow e_3=0$. Persamaan akar $t^3-t^2=0\Rightarrow t^2(t-1)=0$, akar $0,0,1$. Jadi $(x,y,z)$ adalah permutasi dari $(1,0,0)$.

% ============================================================
\chapter{Pembahasan Bab 6: Fungsi Kuadrat}
% ============================================================

\kbagian{Bagian A}

\knomor{A.1}
\ksoal Tentukan akar $x^2-9=0$.
\kanalisis Selisih dua kuadrat.
\kbahas $x^2=9\Rightarrow x=\pm3.$

\knomor{A.2}
\ksoal Tentukan titik potong sumbu-$y$ dari $f(x)=x^2+2x-3$.
\kanalisis Sumbu-$y$ saat $x=0$.
\kbahas $f(0)=-3$, titik $(0,-3)$.

\knomor{A.3}
\ksoal Hitung diskriminan $x^2+4x+4=0$.
\kanalisis $D=b^2-4ac$.
\kbahas $D=16-16=0$ (akar kembar).

\knomor{A.4}
\ksoal Tentukan sumbu simetri $f(x)=x^2-2x$.
\kanalisis $x=-\tfrac{b}{2a}$.
\kbahas $x=-\dfrac{-2}{2(1)}=1.$

\knomor{A.5}
\ksoal Faktorkan $x^2+5x+6$.
\kanalisis Cari dua bilangan berjumlah 5, berkali 6.
\kbahas $(x+2)(x+3).$

\knomor{A.6}
\ksoal Tentukan nilai minimum $f(x)=x^2-1$.
\kanalisis Puncak di $x=0$ (parabola terbuka atas).
\kbahas Nilai minimum $f(0)=-1.$

\knomor{A.7}
\ksoal Apakah parabola $f(x)=-2x^2+x$ terbuka ke atas atau bawah?
\kanalisis Tanda $a$ menentukan arah.
\kbahas $a=-2<0\Rightarrow$ terbuka ke \textbf{bawah}.

\knomor{A.8}
\ksoal Tentukan akar $x^2-x-2=0$.
\kanalisis Faktorkan.
\kbahas $(x-2)(x+1)=0\Rightarrow x=2$ atau $x=-1.$

\knomor{A.9}
\ksoal Jika $x_1+x_2=4$, $x_1x_2=3$, susun persamaannya.
\kanalisis $x^2-(x_1+x_2)x+x_1x_2=0$.
\kbahas $x^2-4x+3=0.$

\knomor{A.10}
\ksoal Tentukan titik puncak $f(x)=(x-3)^2+2$.
\kanalisis Bentuk puncak $a(x-h)^2+k$ berpuncak $(h,k)$.
\kbahas Puncak $(3,2)$.

\kbagian{Bagian B}

\knomor{B.1}
\ksoal Tentukan $m$ agar $f(x)=x^2-(m+1)x+4$ menyinggung sumbu-$x$.
\kanalisis Menyinggung $\iff D=0$.
\kbahas $(m+1)^2-16=0\Rightarrow(m+1)^2=16\Rightarrow m+1=\pm4\Rightarrow m=3$ atau $m=-5.$

\knomor{B.2}
\ksoal Akar $x^2-6x+k=0$ memenuhi $x_1=2x_2$. Tentukan $k$ dan kedua akarnya.
\kanalisis Gunakan jumlah akar $x_1+x_2=6$ dengan $x_1=2x_2$.
\kbahas $2x_2+x_2=6\Rightarrow x_2=2$, $x_1=4$. $k=x_1x_2=8.$ Akar $4$ dan $2$.

\knomor{B.3}
\ksoal Tentukan HP dari $\dfrac{x^2-4}{x-1}\ge0$.
\kanalisis Tentukan pembuat nol pembilang ($x=\pm2$) dan penyebut ($x=1$, dikecualikan); uji tanda tiap interval.
\kbahas Faktor $\dfrac{(x-2)(x+2)}{x-1}$. Tanda $+$ pada $-2\le x<1$ dan $x\ge2$. HP $=[-2,1)\cup[2,\infty).$

\knomor{B.4}
\ksoal Fungsi kuadrat melalui $(0,3)$, $(1,0)$, $(3,0)$. Tentukan rumusnya.
\kanalisis Akar di $1$ dan $3$: $f(x)=a(x-1)(x-3)$; cari $a$ dari titik $(0,3)$.
\kbahas $f(0)=a(-1)(-3)=3a=3\Rightarrow a=1$. $f(x)=(x-1)(x-3)=x^2-4x+3.$

\knomor{B.5}
\ksoal Peluru ditembakkan, tingginya $h(t)=40t-5t^2$. Tentukan tinggi maksimum dan waktu kembali ke tanah.
\kanalisis Puncak di $t=-\tfrac{b}{2a}$; kembali ke tanah saat $h=0$.
\kbahas Puncak $t=\dfrac{-40}{2(-5)}=4$, $h(4)=160-80=80$ m. $h=0\Rightarrow t(40-5t)=0\Rightarrow t=8$ s.

\kbagian{Bagian C}

\knomor{C.1}
\ksoal Tentukan nilai maksimum $f(\theta)=\sin\theta\cos\theta+\cos^2\theta$.
\kanalisis Ubah ke sudut rangkap: $\sin\theta\cos\theta=\tfrac12\sin2\theta$, $\cos^2\theta=\tfrac{1+\cos2\theta}{2}$.
\kbahas $f=\tfrac12\sin2\theta+\tfrac12+\tfrac12\cos2\theta=\tfrac12(\sin2\theta+\cos2\theta)+\tfrac12$. Amplitudo $\sin2\theta+\cos2\theta$ adalah $\sqrt2$, jadi maksimum $f=\dfrac{\sqrt2}{2}+\dfrac12=\dfrac{\sqrt2+1}{2}\approx1{,}21.$

\knomor{C.2}
\ksoal Persamaan $x^2-2px+(p^2-1)=0$ punya kedua akar di antara $-2$ dan $4$. Tentukan rentang $p$.
\kanalisis Bentuk $(x-p)^2=1$ memberi akar $p\pm1$; syaratkan keduanya di $(-2,4)$.
\kbahas $x=p\pm1$. Syarat: $p-1>-2$ dan $p+1<4\Rightarrow p>-1$ dan $p<3$. Jadi $-1<p<3.$

\knomor{C.3}
\ksoal Diberikan $\vec a=(t,1)$ dan $\vec b=(1,t)$. Tentukan $t$ yang meminimalkan $|\vec a+\vec b|^2$.
\kanalisis $\vec a+\vec b=(t+1)(1,1)$; panjang kuadratnya fungsi kuadrat dalam $t$.
\kbahas $|\vec a+\vec b|^2=2(t+1)^2$, minimum saat $t=-1$ (bernilai $0$).

\knomor{C.4}
\ksoal Jika $x_1,x_2$ akar $x^2+bx+c=0$, susun persamaan baru berakar $x_1^2$ dan $x_2^2$.
\kanalisis Hitung jumlah dan hasil kali akar baru memakai Vieta.
\kbahas $x_1^2+x_2^2=(x_1+x_2)^2-2x_1x_2=b^2-2c$; $x_1^2x_2^2=(x_1x_2)^2=c^2$. Persamaan: $x^2-(b^2-2c)x+c^2=0.$

\knomor{C.5}
\ksoal Tentukan banyaknya pasangan bilangan bulat $(x,y)$ dengan $y=x^2$ dan $y\le 100$.
\kanalisis $x^2\le100$ membatasi $x$; tiap $x$ memberi satu $y$.
\kbahas $-10\le x\le10$ memberi $21$ nilai $x$, sehingga ada \textbf{21} pasangan.

\kbagian{Bagian D}

\knomor{D.1}
\ksoal \textbf{(Kuadrat + Vektor + Trigonometri).} Tunjukkan $f(t)=|\vec u+t\vec v|^2$ kuadrat dalam $t$; tentukan minimumnya dan tafsirkan sebagai jarak titik ke garis.
\kanalisis Jabarkan dot product.
\kbahas $f(t)=|\vec v|^2t^2+2(\vec u\cdot\vec v)t+|\vec u|^2$, minimum di $t^*=-\dfrac{\vec u\cdot\vec v}{|\vec v|^2}$. Nilai minimum $|\vec u|^2-\dfrac{(\vec u\cdot\vec v)^2}{|\vec v|^2}$ adalah kuadrat jarak tegak lurus titik (ujung $\vec u$) ke garis berarah $\vec v$.

\knomor{D.2}
\ksoal \textbf{(Kuadrat + Kalkulus).} Buktikan puncak parabola dari melengkapkan kuadrat sama dengan titik $f'(x)=0$.
\kanalisis Melengkapkan kuadrat memberi $h=-\tfrac{b}{2a}$; bandingkan dengan turunan.
\kbahas $f(x)=ax^2+bx+c=a\left(x+\tfrac{b}{2a}\right)^2-\tfrac{D}{4a}$, puncak di $x=-\tfrac{b}{2a}$. Turunan $f'(x)=2ax+b=0\Rightarrow x=-\tfrac{b}{2a}$ --- sama persis. $\blacksquare$

\knomor{D.3}
\ksoal \textbf{(Kuadrat + Statistika).} Tunjukkan bahwa nilai $c$ yang meminimalkan $\sum(x_i-c)^2$ adalah $\bar x$.
\kanalisis $g(c)$ kuadrat dalam $c$; cari puncak.
\kbahas $g(c)=\sum(x_i-c)^2=nc^2-2c\sum x_i+\sum x_i^2$, minimum di $c=\dfrac{\sum x_i}{n}=\bar x.$ $\blacksquare$

\knomor{D.4}
\ksoal \textbf{(Kuadrat + Peluang).} Nyatakan ragam dadu sebagai $E[X^2]-(E[X])^2$ dan hubungkan dengan minimum $g(c)=E[(X-c)^2]$.
\kanalisis $g(c)$ kuadrat dalam $c$, minimum di $c=E[X]$.
\kbahas $g(c)=E[X^2]-2cE[X]+c^2$, minimum di $c=E[X]$ dengan nilai $E[X^2]-(E[X])^2=\mathrm{Var}(X)$. Untuk dadu adil: $E[X]=3{,}5$, $E[X^2]=\tfrac{91}{6}$, $\mathrm{Var}=\tfrac{91}{6}-3{,}5^2=\tfrac{35}{12}\approx2{,}92.$

\knomor{D.5}
\ksoal \textbf{(Optimasi terkendala).} Dari semua persegi panjang berkeliling tetap $K$, buktikan persegi memberi luas maksimum; lalu untuk luas tetap $A$, buktikan persegi memberi keliling minimum.
\kanalisis Nyatakan luas/keliling sebagai fungsi satu sisi.
\kbahas Sisi $x$ dan $\tfrac K2-x$: luas $L(x)=x\left(\tfrac K2-x\right)$ maksimum di $x=\tfrac K4$ (persegi). Untuk luas tetap $A$ ($y=A/x$): keliling $P(x)=2\left(x+\tfrac Ax\right)$ minimum di $x=\sqrt A$ (AM--GM), yakni persegi. $\blacksquare$

\kbagian{Challenge Corner}

\knomor{Challenge 1}
\ksoal Tentukan semua $a$ agar grafik $y=x^2$ dan garis $y=ax+1$ berpotongan di dua titik berjarak $\sqrt{20}$.
\kanalisis Substitusi memberi $x^2-ax-1=0$; nyatakan jarak dua titik potong dengan Vieta.
\kbahas $x_1+x_2=a$, $x_1x_2=-1$. Jarak$^2=(\Delta x)^2\big(1+(x_1+x_2)^2\big)$ dengan $(\Delta x)^2=a^2+4$. Maka $(a^2+4)(a^2+1)=20\Rightarrow a^4+5a^2-16=0\Rightarrow a^2=\dfrac{-5+\sqrt{89}}{2}$, sehingga $a=\pm\sqrt{\dfrac{\sqrt{89}-5}{2}}\approx\pm1{,}49.$

\knomor{Challenge 2}
\ksoal Sebuah bilangan kuadrat sempurna empat digit berbentuk $\overline{aabb}$ (dua digit pertama sama, dua terakhir sama). Temukan bilangan itu.
\kanalisis $\overline{aabb}=1100a+11b=11(100a+b)$ harus habis dibagi $11^2$, sehingga $a+b$ kelipatan 11.
\kbahas Karena $100a+b\equiv a+b\pmod{11}$, perlu $a+b=11$. Maka $\overline{aabb}=121(9a+1)$; agar kuadrat, $9a+1$ harus kuadrat sempurna. Untuk $a=7$: $9(7)+1=64=8^2$. Jadi bilangan $=121\cdot64=7744=88^2.$

\knomor{Challenge 3}
\ksoal Buktikan bahwa di antara semua segitiga berkeliling tetap, segitiga sama sisi memiliki luas terbesar.
\kanalisis Gunakan rumus Heron $L=\sqrt{s(s-a)(s-b)(s-c)}$ dengan $s$ tetap, lalu AM--GM.
\kbahas Dengan $s=\tfrac{a+b+c}{2}$ tetap, $(s-a)+(s-b)+(s-c)=s$ tetap. Menurut AM--GM, hasil kali $(s-a)(s-b)(s-c)$ maksimum saat ketiganya sama, yakni $a=b=c$. Maka $L$ maksimum untuk segitiga sama sisi. $\blacksquare$

% ============================================================
\chapter{Pembahasan Bab 7: Statistika}
% ============================================================

\kbagian{Bagian A}

\knomor{A.1}
\ksoal Hitung mean dari $4,6,8,10$.
\kanalisis $\bar x=\tfrac{\sum x_i}{n}$.
\kbahas $\bar x=\dfrac{28}{4}=7.$

\knomor{A.2}
\ksoal Tentukan median dari $3,7,9,11,15$.
\kanalisis Data sudah terurut; ambil nilai tengah.
\kbahas Median $=9$ (data ke-3).

\knomor{A.3}
\ksoal Tentukan modus dari $2,2,3,4,4,4,5$.
\kanalisis Nilai paling sering muncul.
\kbahas Modus $=4$ (muncul 3 kali).

\knomor{A.4}
\ksoal Hitung jangkauan dari $5,8,12,20$.
\kanalisis $x_{\max}-x_{\min}$.
\kbahas $20-5=15.$

\knomor{A.5}
\ksoal Tentukan $Q_2$ dari $1,2,3,4,5,6$.
\kanalisis $Q_2$ = median; data genap, rata-rata dua tengah.
\kbahas $\dfrac{3+4}{2}=3{,}5.$

\knomor{A.6}
\ksoal Mean 4 bilangan $=10$; tentukan jumlahnya.
\kanalisis Jumlah $=n\cdot\bar x$.
\kbahas $4\times10=40.$

\knomor{A.7}
\ksoal Tentukan $Q_1$ dari $2,4,6,8,10,12$.
\kanalisis $Q_1$ = median paruh bawah $(2,4,6)$.
\kbahas $Q_1=4.$

\knomor{A.8}
\ksoal Apakah data $5,5,5,5$ punya simpangan baku? Berapa?
\kanalisis Semua data sama $\Rightarrow$ tak ada sebaran.
\kbahas Ya, $s=0.$

\knomor{A.9}
\ksoal Hitung mean dari frekuensi: nilai $1,2,3$ frekuensi $2,2,1$.
\kanalisis $\bar x=\tfrac{\sum f_ix_i}{\sum f_i}$.
\kbahas $\dfrac{1(2)+2(2)+3(1)}{5}=\dfrac{9}{5}=1{,}8.$

\knomor{A.10}
\ksoal Tentukan jangkauan antarkuartil data $2,4,6,8,10,12$.
\kanalisis JAK $=Q_3-Q_1$.
\kbahas $Q_1=4$, $Q_3=10$ (median paruh atas $8,10,12$), JAK $=6.$

\kbagian{Bagian B}

\knomor{B.1}
\ksoal Rata-rata 10 nilai $=70$. Setelah satu nilai diralat dari $50$ menjadi $80$, tentukan rata-rata baru.
\kanalisis Perubahan total $=80-50=30$; bagi dengan $n$.
\kbahas Total lama $=700$, total baru $=730$, $\bar x_{baru}=\dfrac{730}{10}=73.$

\knomor{B.2}
\ksoal Data $x, 2x, 3x, 4x, 5x$ memiliki simpangan baku $s$. Bila tiap data dikali 2, bagaimana $s$ berubah?
\kanalisis Mengalikan semua data dengan $k$ mengalikan simpangan baku dengan $|k|$.
\kbahas Karena tiap simpangan terhadap mean ikut dikali 2, maka simpangan baku menjadi $2s$.

\knomor{B.3}
\ksoal Tentukan median data berkelompok: interval $[10,20),[20,30),[30,40)$ dengan frekuensi $4,10,6$.
\kanalisis $n=20$; median di data ke-10. Gunakan $Q_2=L+\dfrac{n/2-F}{f}\cdot c$.
\kbahas Kelas median $[20,30)$ ($L=20$, $F=4$, $f=10$, $c=10$): $Q_2=20+\dfrac{10-4}{10}\cdot10=26.$

\knomor{B.4}
\ksoal Gabungan dua kelompok: I ($n=20$, $\bar x=6$) dan II ($n=30$, $\bar x=8$). Tentukan rata-rata gabungan.
\kanalisis Rata-rata gabungan $=\dfrac{n_1\bar x_1+n_2\bar x_2}{n_1+n_2}$.
\kbahas $\dfrac{20(6)+30(8)}{50}=\dfrac{360}{50}=7{,}2.$

\knomor{B.5}
\ksoal Diketahui $\sum x_i=100$, $\sum x_i^2=1100$, $n=10$. Hitung simpangan baku.
\kanalisis Gunakan $s^2=\dfrac{\sum x_i^2}{n}-\bar x^2$.
\kbahas $\bar x=10$, $s^2=\dfrac{1100}{10}-10^2=110-100=10\Rightarrow s=\sqrt{10}\approx3{,}16.$

\kbagian{Bagian C}

\knomor{C.1}
\ksoal Buktikan bahwa nilai $c$ yang meminimalkan $\sum(x_i-c)^2$ adalah $c=\bar x$.
\kanalisis $g(c)$ adalah fungsi kuadrat dalam $c$ yang terbuka ke atas.
\kbahas $g(c)=nc^2-2c\sum x_i+\sum x_i^2$, berpuncak di $c=\dfrac{\sum x_i}{n}=\bar x$, yaitu titik minimum. $\blacksquare$

\knomor{C.2}
\ksoal Tunjukkan identitas $s^2=\overline{x^2}-\bar x^2$.
\kanalisis Jabarkan $(x_i-\bar x)^2$ dalam definisi ragam.
\kbahas $s^2=\dfrac1n\sum(x_i-\bar x)^2=\dfrac1n\sum x_i^2-2\bar x\cdot\dfrac1n\sum x_i+\bar x^2=\overline{x^2}-2\bar x^2+\bar x^2=\overline{x^2}-\bar x^2.$ $\blacksquare$

\knomor{C.3}
\ksoal Lima bilangan asli berbeda memiliki mean $=$ median $=5$. Tentukan nilai maksimum bilangan terbesar.
\kanalisis Urutkan $x_1<\dots<x_5$ dengan $x_3=5$ dan jumlah $25$; minimalkan suku lain.
\kbahas Pilih $x_1=1,x_2=2,x_3=5,x_4=6$ (sekecil mungkin namun tetap berbeda dan terurut). Maka $x_5=25-1-2-5-6=11.$ Nilai maksimum $=11$.

\knomor{C.4}
\ksoal Data $a,b,c$ punya mean $m$. Buktikan $(a-m)+(b-m)+(c-m)=0$ dan jelaskan maknanya.
\kanalisis Jumlah simpangan terhadap mean selalu nol.
\kbahas $(a+b+c)-3m=3m-3m=0$. Maknanya: mean adalah \emph{titik keseimbangan} --- jumlah simpangan positif persis menutup jumlah simpangan negatif, dasar dari pengukuran sebaran.

\knomor{C.5}
\ksoal Jika menambahkan konstanta $k$ ke semua data, buktikan mean naik $k$ tetapi simpangan baku tetap.
\kanalisis Tinjau pengaruh pergeseran pada mean dan pada simpangan.
\kbahas Mean baru $=\dfrac1n\sum(x_i+k)=\bar x+k$. Simpangan tiap data $(x_i+k)-(\bar x+k)=x_i-\bar x$ tak berubah, sehingga $s$ tetap. $\blacksquare$

\kbagian{Bagian D}

\knomor{D.1}
\ksoal \textbf{(Statistika + Peluang + Logaritma).} Tunjukkan entropi $H=-\sum p_i\,{}^{2}\!\log p_i$ maksimum saat semua $p_i$ sama.
\kanalisis Fungsi $-p\log p$ cekung; gunakan ketaksamaan Jensen atau argumen simetri.
\kbahas Karena $H$ cekung terhadap $(p_1,\dots,p_n)$ dengan kendala $\sum p_i=1$, maksimum dicapai pada titik paling simetris $p_i=\tfrac1n$. Nilainya $H=-\sum\tfrac1n\,{}^{2}\!\log\tfrac1n={}^{2}\!\log n$. Ini berarti ketidakpastian terbesar saat semua keluaran sama mungkin. $\blacksquare$

\knomor{D.2}
\ksoal \textbf{(Statistika + Vektor).} Tunjukkan koefisien korelasi $=$ kosinus sudut antara dua vektor data terpusat.
\kanalisis Pusatkan data lalu hitung $\cos\theta=\dfrac{\vec u\cdot\vec v}{|\vec u||\vec v|}$.
\kbahas Dengan $\vec u=(x_i-\bar x)$, $\vec v=(y_i-\bar y)$: $r=\dfrac{\sum(x_i-\bar x)(y_i-\bar y)}{\sqrt{\sum(x_i-\bar x)^2}\sqrt{\sum(y_i-\bar y)^2}}=\dfrac{\vec u\cdot\vec v}{|\vec u||\vec v|}=\cos\theta.$ $\blacksquare$

\knomor{D.3}
\ksoal \textbf{(Statistika + Fungsi Kuadrat).} Turunkan garis regresi $y=ax+b$ kuadrat terkecil dan tunjukkan ia melalui $(\bar x,\bar y)$.
\kanalisis Minimalkan $S=\sum(y_i-ax_i-b)^2$; persamaan normal memberi $b$.
\kbahas $\dfrac{\partial S}{\partial b}=0$ memberi $\sum y_i=a\sum x_i+nb$, yaitu $\bar y=a\bar x+b$. Jadi titik $(\bar x,\bar y)$ memenuhi persamaan garis --- garis regresi selalu melaluinya. $\blacksquare$

\knomor{D.4}
\ksoal \textbf{(Statistika + Deret).} Jika data adalah barisan aritmetika $a, a+d,\dots,a+(n-1)d$, turunkan mean dan simpangan bakunya.
\kanalisis Mean barisan aritmetika $=$ rata-rata suku tepi; ragam $=d^2\times$ ragam $\{0,1,\dots,n-1\}$.
\kbahas Mean $=a+\dfrac{(n-1)d}{2}$. Ragam $=d^2\cdot\dfrac{n^2-1}{12}$, sehingga $s=d\sqrt{\dfrac{n^2-1}{12}}.$

\knomor{D.5}
\ksoal \textbf{(Literasi).} Dua grafik dari data sama dengan skala sumbu berbeda. Analisis bagaimana penyajian menyesatkan dan rumuskan kriteria penyajian jujur.
\kanalisis Tinjau pengaruh titik awal sumbu dan rentang skala terhadap persepsi.
\kbahas Sumbu-$y$ yang tidak dimulai dari nol atau yang dipotong membuat selisih kecil tampak dramatis; skala tak proporsional melebih-lebihkan tren. \emph{Kriteria jujur}: sumbu mulai dari nol (untuk diagram batang), skala seragam dan diberi label, rentang data utuh, serta sumber dan satuan dicantumkan.

\kbagian{Challenge Corner}

\knomor{Challenge 1}
\ksoal Mungkinkah menambah satu data baru \emph{menurunkan} mean tetapi \emph{menaikkan} median sekaligus? Beri contoh atau buktikan mustahil.
\kanalisis Mean dipengaruhi nilai pencilan, median oleh posisi tengah; keduanya dapat bergerak berlawanan.
\kbahas \textbf{Mungkin.} Contoh: data $\{1,2,3,100\}$ punya mean $26{,}5$, median $2{,}5$. Tambahkan $2{,}6$: data menjadi $\{1,2,2{,}6,3,100\}$ dengan mean $\dfrac{108{,}6}{5}=21{,}72$ (\emph{turun}) dan median $2{,}6$ (\emph{naik}). $\blacksquare$

\knomor{Challenge 2}
\ksoal Sepuluh bilangan bulat positif memiliki mean $10$ dan median $9$. Tentukan nilai maksimum bilangan terbesar.
\kanalisis Jumlah $=100$; median $=\dfrac{x_5+x_6}{2}=9$. Minimalkan sembilan suku pertama.
\kbahas $x_5+x_6=18$ dengan $x_5\le x_6$ memaksa $x_6\ge9$. Ambil $x_1=\cdots=x_4=1$, $x_5=x_6=x_7=x_8=x_9=9$. Jumlah sembilan suku $=4+45=49$, sehingga $x_{10}=100-49=51.$ Urutan $1,1,1,1,9,9,9,9,9,51$ memenuhi semua syarat. Maksimum $=51$.

\knomor{Challenge 3}
\ksoal Rancang survei mini di kelas (mis. waktu tidur), kumpulkan data, sajikan dengan box plot, dan tulis satu kesimpulan beserta keterbatasannya.
\kanalisis Tentukan pertanyaan, ukuran sampel, lalu hitung lima ukuran ringkasan untuk box plot.
\kbahas Contoh: kumpulkan waktu tidur 20 siswa, urutkan, hitung min, $Q_1$, median, $Q_3$, maks, lalu gambar box plot. Kesimpulan mengaitkan median dan sebaran (mis. "separuh siswa tidur antara $6$--$8$ jam"). \emph{Keterbatasan}: sampel kecil dan hanya satu kelas sehingga \textbf{tidak} dapat digeneralisasi ke semua siswa; data laporan diri rawan bias.

% ============================================================
\chapter{Pembahasan Bab 8: Peluang}
% ============================================================

\kbagian{Bagian A}

\knomor{A.1}
\ksoal Sebuah dadu dilempar. Peluang muncul angka 3?
\kanalisis $P=\tfrac{n(A)}{n(S)}$ dengan $n(S)=6$.
\kbahas $P=\tfrac16.$

\knomor{A.2}
\ksoal Peluang muncul Gambar pada satu koin?
\kanalisis Dua hasil sama mungkin.
\kbahas $P=\tfrac12.$

\knomor{A.3}
\ksoal Dari 1--10, peluang bilangan genap?
\kanalisis Genap: $\{2,4,6,8,10\}$.
\kbahas $P=\tfrac{5}{10}=\tfrac12.$

\knomor{A.4}
\ksoal Jika $P(A)=0{,}25$, tentukan $P(A^c)$.
\kanalisis Aturan komplemen.
\kbahas $P(A^c)=1-0{,}25=0{,}75.$

\knomor{A.5}
\ksoal Hitung $C(5,2)$.
\kanalisis $C(n,r)=\dfrac{n!}{r!(n-r)!}$.
\kbahas $\dfrac{5!}{2!\,3!}=10.$

\knomor{A.6}
\ksoal Hitung $4!$.
\kanalisis $4!=4\cdot3\cdot2\cdot1$.
\kbahas $24.$

\knomor{A.7}
\ksoal Dua koin: peluang dua-duanya Angka?
\kanalisis Kejadian bebas: kalikan peluang.
\kbahas $\tfrac12\cdot\tfrac12=\tfrac14.$

\knomor{A.8}
\ksoal Dari 6 merah, 4 putih, peluang terambil putih (satu bola)?
\kanalisis $\tfrac{\text{putih}}{\text{total}}$.
\kbahas $\tfrac{4}{10}=\tfrac25.$

\knomor{A.9}
\ksoal Hitung $P(4,2)$.
\kanalisis Permutasi $P(n,r)=\dfrac{n!}{(n-r)!}$.
\kbahas $\dfrac{4!}{2!}=12.$

\knomor{A.10}
\ksoal Peluang dadu muncul faktor dari 6?
\kanalisis Faktor 6: $\{1,2,3,6\}$.
\kbahas $P=\tfrac46=\tfrac23.$

\kbagian{Bagian B}

\knomor{B.1}
\ksoal Dua dadu dilempar. Peluang jumlah mata $\ge 10$?
\kanalisis Hitung pasangan dengan jumlah $10,11,12$ dari $36$ hasil.
\kbahas Jumlah 10: $3$ cara; 11: $2$ cara; 12: $1$ cara. Total $6$. $P=\tfrac{6}{36}=\tfrac16.$

\knomor{B.2}
\ksoal Dari 9 bola (5 merah, 4 biru) diambil 3 sekaligus. Peluang tepat 2 merah?
\kanalisis Pilih 2 merah dari 5 dan 1 biru dari 4; bagi dengan total cara.
\kbahas $P=\dfrac{C(5,2)\,C(4,1)}{C(9,3)}=\dfrac{10\cdot4}{84}=\dfrac{40}{84}=\dfrac{10}{21}.$

\knomor{B.3}
\ksoal Kotak berisi 3 bola merah, 2 putih. Diambil 2 berurutan tanpa pengembalian. Peluang keduanya merah?
\kanalisis Peluang bersyarat: kalikan peluang tiap pengambilan.
\kbahas $P=\tfrac35\cdot\tfrac24=\tfrac{6}{20}=\tfrac{3}{10}.$

\knomor{B.4}
\ksoal Tiga koin dilempar. Peluang muncul jumlah Gambar genap (0 atau 2)?
\kanalisis Hitung hasil dengan 0 atau 2 Gambar dari $8$ hasil.
\kbahas $0$ Gambar: $1$ cara; $2$ Gambar: $C(3,2)=3$ cara. Total $4$. $P=\tfrac48=\tfrac12.$

\knomor{B.5}
\ksoal Dari kata "STATISTIKA", berapa susunan berbeda yang dapat dibentuk?
\kanalisis Permutasi dengan unsur berulang: $S{=}2,T{=}3,A{=}2,I{=}2,K{=}1$ ($10$ huruf).
\kbahas $\dfrac{10!}{2!\,3!\,2!\,2!}=\dfrac{3.628.800}{48}=75.600.$

\kbagian{Bagian C}

\knomor{C.1}
\ksoal Dadu adil dilempar. Tentukan nilai harapan (mean) dan variansi mata yang muncul.
\kanalisis $E[X]=\sum x_iP(x_i)$; $\mathrm{Var}=E[X^2]-(E[X])^2$.
\kbahas $E[X]=\tfrac{1+2+\cdots+6}{6}=3{,}5$. $E[X^2]=\tfrac{91}{6}$. $\mathrm{Var}=\tfrac{91}{6}-3{,}5^2=\tfrac{35}{12}\approx2{,}92.$

\knomor{C.2}
\ksoal Peluang menang $0{,}4$, dimainkan 5 kali bebas. Peluang menang tepat 3 kali?
\kanalisis Distribusi binomial $P=C(n,k)p^k(1-p)^{n-k}$.
\kbahas $P=C(5,3)(0{,}4)^3(0{,}6)^2=10(0{,}064)(0{,}36)=0{,}2304.$

\knomor{C.3}
\ksoal Tiga orang menembak dengan peluang kena $0{,}6$, $0{,}7$, $0{,}8$ bebas. Peluang target terkena tepat satu orang?
\kanalisis Jumlahkan tiga skenario "tepat satu kena, dua lain meleset".
\kbahas $P=0{,}6(0{,}3)(0{,}2)+0{,}4(0{,}7)(0{,}2)+0{,}4(0{,}3)(0{,}8)=0{,}036+0{,}056+0{,}096=0{,}188.$

\knomor{C.4}
\ksoal Tunjukkan $\binom{n}{0}+\binom{n}{1}+\cdots+\binom{n}{n}=2^n$ dan kaitkan dengan banyaknya himpunan bagian.
\kanalisis Gunakan teorema binomial dengan $x=y=1$.
\kbahas $(1+1)^n=\sum_{k=0}^{n}\binom nk=2^n$. Maknanya: banyak himpunan bagian dari himpunan beranggota $n$ adalah $2^n$ (tiap unsur "masuk" atau "tidak"). $\blacksquare$

\knomor{C.5}
\ksoal Peubah acak bernilai $1$ dengan peluang $p$ dan $0$ dengan peluang $1-p$. Tentukan mean dan variansinya, lalu cari $p$ yang memaksimalkan variansi.
\kanalisis Hitung $E[X]$ dan $\mathrm{Var}$; maksimalkan fungsi kuadrat dalam $p$.
\kbahas $E[X]=p$, $E[X^2]=p$, $\mathrm{Var}=p-p^2=p(1-p)$. Maksimum di $p=\tfrac12$ dengan nilai $\tfrac14$.

\kbagian{Bagian D}

\knomor{D.1}
\ksoal \textbf{(Peluang + Statistika + Logaritma).} Hitung entropi $H=-\sum p_i\,{}^{2}\!\log p_i$ untuk dadu adil dan jelaskan maknanya.
\kanalisis Enam keluaran sama mungkin, $p_i=\tfrac16$.
\kbahas $H=-6\cdot\tfrac16\,{}^{2}\!\log\tfrac16={}^{2}\!\log6\approx2{,}585$ bit --- ukuran rata-rata informasi/ketidakpastian percobaan.

\knomor{D.2}
\ksoal \textbf{(Peluang bersyarat --- Bayes).} Tes penyakit: sensitivitas $99\%$, spesifisitas $95\%$, prevalensi $1\%$. Hitung peluang benar sakit jika hasil tes positif.
\kanalisis Teorema Bayes dengan hukum total peluang untuk $P(+)$.
\kbahas $P(+)=0{,}99(0{,}01)+0{,}05(0{,}99)=0{,}0594$. $P(S\mid+)=\dfrac{0{,}0099}{0{,}0594}\approx16{,}7\%.$ Mengejutkan karena pada prevalensi rendah, \emph{mayoritas hasil positif justru positif palsu}.

\knomor{D.3}
\ksoal \textbf{(Peluang + Deret).} Lempar koin sampai muncul Angka pertama. Hitung ekspektasi banyak lemparan dan tunjukkan jumlah seluruh peluang $=1$.
\kanalisis Peluang berhenti lemparan ke-$n$ adalah $(\tfrac12)^n$; deret geometri dan aritmetika-geometri.
\kbahas $\sum_{n\ge1}(\tfrac12)^n=1$. Ekspektasi $\sum_{n\ge1}n(\tfrac12)^n=\dfrac{1/2}{(1/2)^2}=2$ lemparan.

\knomor{D.4}
\ksoal \textbf{(Peluang + Fungsi Kuadrat).} Tunjukkan $g(c)=E[(X-c)^2]$ minimum saat $c=E[X]$, dan minimumnya variansi.
\kanalisis $g(c)$ kuadrat dalam $c$.
\kbahas $g(c)=E[X^2]-2cE[X]+c^2$, minimum di $c=E[X]$ dengan nilai $E[X^2]-(E[X])^2=\mathrm{Var}(X)$. $\blacksquare$

\knomor{D.5}
\ksoal \textbf{(Peluang geometris).} Dua orang datang acak antara pukul 12.00--13.00, masing-masing menunggu 15 menit. Tentukan peluang mereka bertemu.
\kanalisis Modelkan $(x,y)$ pada persegi $[0,60]^2$; bertemu jika $|x-y|\le15$.
\kbahas Luas daerah bertemu $=60^2-2\cdot\tfrac12(45)^2=3600-2025=1575$. Peluang $=\dfrac{1575}{3600}=\dfrac{7}{16}=0{,}4375.$

\kbagian{Challenge Corner}

\knomor{Challenge 1}
\ksoal Tentukan peluang minimal dua dari 23 orang berulang tahun di hari yang sama, dan jelaskan mengapa melebihi $50\%$.
\kanalisis Gunakan komplemen: peluang semua tanggal \emph{berbeda}.
\kbahas $P(\text{semua beda})=\dfrac{365}{365}\cdot\dfrac{364}{365}\cdots\dfrac{343}{365}\approx0{,}493$, sehingga $P(\ge2\text{ sama})\approx0{,}507>50\%$. Sebabnya: yang dibandingkan adalah \emph{pasangan}, dan ada $\binom{23}{2}=253$ pasangan --- jauh lebih banyak dari dugaan intuisi.

\knomor{Challenge 2}
\ksoal Masalah Monty Hall: setelah memilih satu dari 3 pintu, pembawa acara membuka satu pintu kosong lalu menawarkan pindah. Sebaiknya pindah? Buktikan.
\kanalisis Bandingkan peluang menang strategi "tetap" vs "pindah".
\kbahas Pilihan awal benar berpeluang $\tfrac13$ (tetap menang $\tfrac13$). Dengan peluang $\tfrac23$ pilihan awal salah; setelah satu pintu kosong dibuka, \emph{pindah} pasti menang. Jadi strategi \textbf{pindah} menang berpeluang $\tfrac23>\tfrac13$ --- sebaiknya pindah.

\knomor{Challenge 3}
\ksoal Semut berjalan acak pada titik segitiga $ABC$; tiap langkah pindah ke salah satu titik tetangga berpeluang sama. Tentukan peluang kembali ke titik awal setelah 3 langkah.
\kanalisis Misal $a_n$ peluang berada di titik awal; rekursi simetri $a_{n+1}=\dfrac{1-a_n}{2}$.
\kbahas $a_0=1,\ a_1=0,\ a_2=\tfrac12,\ a_3=\dfrac{1-\frac12}{2}=\dfrac14.$ (Dicek langsung: dari $8$ lintasan, tepat $2$ berakhir di titik awal.) Peluang $=\dfrac14.$

% ============================================================
%  PENANDA BAGIAN SELANJUTNYA (tinggal diisi)
% ============================================================
% \part*{\centering FASE F — KELAS 11 \& 12}
% \chapter{Komposisi Fungsi dan Fungsi Invers}
% \chapter{Lingkaran}
% \chapter{Statistika Lanjutan (Data Bivariat)}
%
% \part*{\centering MATEMATIKA TINGKAT LANJUT (FASE F+)}
% \chapter{Bilangan Kompleks}
% \chapter{Polinomial (Suku Banyak)}
% \chapter{Matriks}
% \chapter{Geometri Transformasi}
% \chapter{Fungsi dan Pemodelannya}
% \chapter{Trigonometri Lanjutan}
% \chapter{Irisan Kerucut}
% \chapter{Kalkulus Dasar}

% ============================================================
\chapter{Pembahasan Bab 9: Bunga Tunggal dan Bunga Majemuk}
% ============================================================

\kbagian{Bagian A}

\knomor{A.1}
\ksoal Rp~2.000.000 ditabung dengan bunga tunggal $10\%$/tahun selama 3 tahun. Nilai akhirnya?
\kanalisis Gunakan $M_n = M_0(1+ni)$.
\kbahas $M_3 = 2.000.000(1+3\times0{,}1)=2.000.000\times1{,}3=\text{Rp}\ 2.600.000.$

\knomor{A.2}
\ksoal Bunga total dari modal Rp~5.000.000, bunga tunggal $6\%$/tahun, selama 4 tahun.
\kanalisis Total bunga $= M_0 \cdot n \cdot i$.
\kbahas $B = 5.000.000\times4\times0{,}06 = \text{Rp}\ 1.200.000.$

\knomor{A.3}
\ksoal Modal Rp~4.000.000, bunga majemuk $5\%$/tahun, 2 tahun. Nilai akhir?
\kanalisis $M_n = M_0(1+i)^n$.
\kbahas $M_2 = 4.000.000\times(1{,}05)^2=4.000.000\times1{,}1025=\text{Rp}\ 4.410.000.$

\knomor{A.4}
\ksoal Suku bunga tunggal berapa agar Rp~8.000.000 menjadi Rp~9.600.000 dalam 2 tahun?
\kanalisis $M_n=M_0(1+ni)$; isolasi $i$.
\kbahas $9.600.000=8.000.000(1+2i)\Rightarrow 1{,}2=1+2i\Rightarrow i=0{,}1=10\%$ per tahun.

\knomor{A.5}
\ksoal Modal Rp~3.000.000 bunga tunggal $12\%$/tahun. Berapa tahun hingga Rp~4.080.000?
\kanalisis Isolasi $n$ dari $M_n=M_0(1+ni)$.
\kbahas $4.080.000=3.000.000(1+0{,}12n)\Rightarrow 1{,}36=1+0{,}12n\Rightarrow n=3$ tahun.

\knomor{A.6}
\ksoal Modal Rp~1.000.000, bunga majemuk $10\%$/tahun, 3 tahun.
\kanalisis $M_3=M_0(1{,}1)^3$.
\kbahas $M_3=1.000.000\times1{,}331=\text{Rp}\ 1.331.000.$

\knomor{A.7}
\ksoal Nilai akhir bunga tunggal membentuk barisan jenis apa?
\kanalisis $M_0,M_0(1+i),M_0(1+2i),\ldots$: beda konstan.
\kbahas Barisan \textbf{aritmetika} dengan beda $d=M_0\cdot i$.

\knomor{A.8}
\ksoal Modal Rp~6.000.000 didepositokan dengan bunga majemuk $12\%$/tahun per bulan. Suku bunga per bulan?
\kanalisis Bagi suku bunga tahunan dengan 12.
\kbahas $i_{\text{bulan}}=12\%/12=1\%$ per bulan.

\knomor{A.9}
\ksoal Modal awal agar dengan bunga tunggal $15\%$/tahun setelah 2 tahun menjadi Rp~11.500.000.
\kanalisis $11.500.000=M_0(1+2\times0{,}15)=1{,}3\,M_0$.
\kbahas $M_0=11.500.000/1{,}3\approx\text{Rp}\ 8.846.154.$

\knomor{A.10}
\ksoal Bandingkan bunga tunggal dan majemuk $10\%$/tahun, modal Rp~1.000.000, 1 tahun.
\kanalisis Untuk $n=1$, $(1+i)^1=1+i$, sehingga hasilnya sama.
\kbahas Keduanya menghasilkan Rp~1.100.000. Selisih muncul hanya untuk $n>1$.

\kbagian{Bagian B}

\knomor{B.1}
\ksoal Modal Rp~10.000.000, bunga majemuk $8\%$/tahun, 5 tahun. Hitung dan bandingkan dengan bunga tunggal. ($(1{,}08)^5\approx1{,}4693$)
\kanalisis Hitung keduanya lalu selisihkan.
\kbahas Majemuk: $10.000.000\times1{,}4693=\text{Rp}\ 14.693.000$. Tunggal: $10.000.000(1+5\times0{,}08)=\text{Rp}\ 14.000.000$. Selisih $=\text{Rp}\ 693.000$.

\knomor{B.2}
\ksoal Setelah 3 tahun dengan bunga majemuk $10\%$/tahun nilainya Rp~13.310.000. Berapa modal awalnya?
\kanalisis $M_3=M_0(1{,}1)^3=M_0\times1{,}331$.
\kbahas $M_0=13.310.000/1{,}331=\text{Rp}\ 10.000.000.$

\knomor{B.3}
\ksoal Investasi A: bunga tunggal $15\%$/tahun; B: bunga majemuk $12\%$/tahun; 5 tahun. Mana lebih untung? ($(1{,}12)^5\approx1{,}7623$)
\kanalisis Hitung faktor pertumbuhan keduanya.
\kbahas A: $1+5\times0{,}15=1{,}75$. B: $(1{,}12)^5\approx1{,}7623$. Investasi B lebih menguntungkan ($1{,}7623>1{,}75$).

\knomor{B.4}
\ksoal Bunga majemuk $6\%$/tahun. Berapa lama modal berlipat tiga?
\kanalisis $(1{,}06)^n=3\Rightarrow n=\log3/\log1{,}06$.
\kbahas $n=0{,}4771/0{,}0253\approx18{,}9\approx19$ tahun.

\knomor{B.5}
\ksoal Modal Rp~5.000.000, bunga majemuk $12\%$/tahun per triwulan, 1 tahun. ($(1{,}03)^4\approx1{,}1255$)
\kanalisis $i_{\text{triwulan}}=3\%$, $n=4$ triwulan.
\kbahas $M_4=5.000.000\times(1{,}03)^4=5.000.000\times1{,}1255=\text{Rp}\ 5.627.500.$

\kbagian{Bagian C}

\knomor{C.1}
\ksoal Rina menabung Rp~1.000.000 setiap awal tahun, bunga majemuk $10\%$/tahun, 3 tahun. Hitung nilai totalnya.
\kanalisis Tabungan tahun ke-1 bertumbuh 3 tahun, ke-2 bertumbuh 2 tahun, ke-3 bertumbuh 1 tahun.
\kbahas $S=1.000.000[(1{,}1)^3+(1{,}1)^2+(1{,}1)]=1.000.000[1{,}331+1{,}21+1{,}1]=\text{Rp}\ 3.641.000.$

\knomor{C.2}
\ksoal Tunjukkan $M_{n+1}/M_n$ konstan pada bunga majemuk dan kaitkan dengan barisan geometri.
\kanalisis $M_n=M_0(1+i)^n$.
\kbahas $\dfrac{M_{n+1}}{M_n}=\dfrac{M_0(1+i)^{n+1}}{M_0(1+i)^n}=(1+i)=\text{konstan}$. Ini persis definisi barisan geometri dengan rasio $r=1+i$. $\blacksquare$

\knomor{C.3}
\ksoal Selisih waktu penggandaan bank A ($6\%$) vs bank B ($8\%$).
\kanalisis $n_A=\log2/\log1{,}06$; $n_B=\log2/\log1{,}08$.
\kbahas $n_A=0{,}301/0{,}0253\approx11{,}9$ tahun; $n_B=0{,}301/0{,}0334\approx9{,}0$ tahun. Selisih $\approx 2{,}9$ tahun.

\knomor{C.4}
\ksoal Buktikan $(1+i)^n > 1+ni$ untuk $n\ge2$, $i>0$.
\kanalisis Ekspansi binomial $(1+i)^n=\sum_{k=0}^n\binom{n}{k}i^k$.
\kbahas $(1+i)^n=1+ni+\binom{n}{2}i^2+\cdots>1+ni$ karena semua suku tambahan $\binom{n}{k}i^k>0$ untuk $k\ge2$ dan $i>0$. $\blacksquare$

\knomor{C.5}
\ksoal $f(n)=M_0(1+i)^n-M_0(1+ni)$. Tunjukkan $f(0)=f(1)=0$, $f(n)>0$ untuk $n\ge2$, lalu hitung $f(3)$ untuk $M_0=1$, $i=0{,}1$.
\kanalisis Substitusi langsung.
\kbahas $f(0)=1-1=0$. $f(1)=(1+i)-(1+i)=0$. Untuk $n\ge2$: dari C.4, $(1{,}1)^n>1+0{,}1n$, jadi $f(n)>0$. $f(3)=(1{,}1)^3-(1+0{,}3)=1{,}331-1{,}3=0{,}031.$

\kbagian{Bagian D}

\knomor{D.1}
\ksoal Berapa modal sekarang (bunga majemuk $8\%$/tahun) agar 5 tahun mendatang nilainya Rp~14.693.000?
\kanalisis Present value: $M_0=M_n/(1+i)^n$.
\kbahas $M_0=14.693.000/(1{,}08)^5=14.693.000/1{,}4693=\text{Rp}\ 10.000.000.$ Maknanya: Rp~10.000.000 sekarang \emph{setara} dengan Rp~14.693.000 lima tahun mendatang pada tingkat bunga $8\%$.

\knomor{D.2}
\ksoal Inflasi $5\%$/tahun. Harga sekarang Rp~100.000. Harga 10 tahun mendatang dan tahun berlipat dua.
\kanalisis Inflasi majemuk: $H_n=H_0(1{,}05)^n$.
\kbahas $H_{10}=100.000\times1{,}6289=\text{Rp}\ 162.890$. Berlipat dua: $n=\log2/\log1{,}05=0{,}301/0{,}02119\approx14{,}2$ tahun.

\knomor{D.3}
\ksoal Cicilan Rp~2.000.000/akhir tahun selama 3 tahun, bunga $10\%$/tahun. Nilai sekarang total?
\kanalisis $PV=\sum_{k=1}^{3}\dfrac{C}{(1{,}1)^k}$.
\kbahas $PV=\dfrac{2.000.000}{1{,}1}+\dfrac{2.000.000}{1{,}21}+\dfrac{2.000.000}{1{,}331}=1.818.182+1.652.893+1.502.630\approx\text{Rp}\ 4.973.705.$

\knomor{D.4}
\ksoal Bunga nominal $12\%$/tahun diperhitungkan per tahun, triwulan, dan bulan. Bandingkan.
\kanalisis Hitung $(1+i/m)^m$ untuk $m=1,4,12$.
\kbahas Tahunan: $(1{,}12)^1=1{,}1200$. Triwulanan: $(1{,}03)^4\approx1{,}1255$. Bulanan: $(1{,}01)^{12}\approx1{,}1268$. Makin sering diperhitungkan, makin besar hasilnya.

\knomor{D.5}
\ksoal Bandingkan $e^{0{,}12}$ dengan per-bulan dan per-tahun untuk $r=12\%$, $t=1$ tahun.
\kanalisis Bunga kontinu sebagai limit frekuensi majemuk tak terhingga.
\kbahas Tahunan: $1{,}1200$. Bulanan: $1{,}1268$. Kontinu: $e^{0{,}12}\approx1{,}1275$. Makin tinggi frekuensi, mendekati $e^{rt}$; bedanya kecil secara praktis tetapi penting secara teoritis.

\kbagian{Challenge Corner}

\knomor{Challenge 1}
\ksoal Tunjukkan $\left(1+\frac{1}{n}\right)^n$ mendekati $e\approx2{,}718$ secara numerik untuk $n=1,2,4,12,365$.
\kanalisis Hitung tiap nilai dan amati konvergensinya.
\kbahas $n=1$: $2$; $n=2$: $2{,}25$; $n=4$: $2{,}4414$; $n=12$: $2{,}6130$; $n=365$: $2{,}7146$. Mendekati $e$ tetapi tidak pernah melampaui $3$. Implikasi: frekuensi majemuk yang lebih tinggi memberikan imbal hasil lebih tinggi, tetapi ada batas atas.

\knomor{Challenge 2}
\ksoal Penduduk tumbuh $3\%$/tahun, inflasi $5\%$/tahun. Kapan daya beli per kapita turun $20\%$?
\kanalisis Cari $n$ sehingga $(1{,}03/1{,}05)^n=0{,}8$, yaitu $(0{,}9810)^n=0{,}8$.
\kbahas $n=\log0{,}8/\log0{,}9810=(-0{,}0969)/(-0{,}00837)\approx11{,}6$ tahun. Artinya sekitar 12 tahun.

\knomor{Challenge 3}
\ksoal Pinjam Rp~10.000.000, bunga majemuk $2\%$/bulan, tidak cicil 6 bulan. Total bayar? Bandingkan bunga tunggal. ($(1{,}02)^6\approx1{,}1262$)
\kanalisis Majemuk: $M_6=M_0(1{,}02)^6$; tunggal: $M_6=M_0(1+6\times0{,}02)$.
\kbahas Majemuk: $10.000.000\times1{,}1262=\text{Rp}\ 11.262.000$. Tunggal: $10.000.000\times1{,}12=\text{Rp}\ 11.200.000$. Selisih Rp~62.000; terlihat kecil dalam 6 bulan, tetapi makin lama selisih ini membesar drastis.

\end{document}